% 本文件由 LaTeX 排版 Agent 根据有效 translation.json 自动生成。
% author=Tatian; work_id=1002; title=塔提安的四福音合参

\chapter{\Person{塔提安}的四福音合参}

% structure: p0001 source=h1 target=omitted reason=page-title-already-rendered-by-parent

\Emph{第1节}\newline{}\Emph{第2节}\newline{}\Emph{第3节}\newline{}\Emph{第4节}\newline{}\Emph{第5节}\newline{}\Emph{第6节}\newline{}\Emph{第7节}\newline{}\Emph{第8节}\newline{}\Emph{第9节}\newline{}\Emph{第10节}\newline{}\Emph{第11节}\newline{}\Emph{第12节}\newline{}\Emph{第13节}\newline{}\Emph{第14节}\newline{}\Emph{第15节}\newline{}\Emph{第16节}\newline{}\Emph{第17节}\newline{}\Emph{第18节}\newline{}\Emph{第19节}\newline{}\Emph{第20节}\newline{}\Emph{第21节}\newline{}\Emph{第22节}\newline{}\Emph{第23节}\newline{}\Emph{第24节}\newline{}\Emph{第25节}\newline{}\Emph{第26节}\newline{}\Emph{第27节}\newline{}\Emph{第28节}\newline{}\Emph{第29节}\newline{}\Emph{第30节}\newline{}\Emph{第31节}\newline{}\Emph{第32节}\newline{}\Emph{第33节}\newline{}\Emph{第34节}\newline{}\Emph{第35节}\newline{}\Emph{第36节}\newline{}\Emph{第37节}\newline{}\Emph{第38节}\newline{}\Emph{第39节}\newline{}\Emph{第40节}\newline{}\Emph{第41节}\newline{}\Emph{第42节}\newline{}\Emph{第43节}\newline{}\Emph{第44节}\newline{}\Emph{第45节}\newline{}\Emph{第46节}\newline{}\Emph{第47节}\newline{}\Emph{第48节}\newline{}\Emph{第49节}\newline{}\Emph{第50节}\newline{}\Emph{第51节}\newline{}\Emph{第52节}\newline{}\Emph{第53节}\newline{}\Emph{第54节}\newline{}\Emph{第55节}\par

\SourceRecord{Translated by Hope W. Hogg.；Ante-Nicene Fathers；Vol. 9.；Edited by Allan Menzies.；Buffalo, NY: Christian Literature Publishing Co.,；1896.；<http://www.newadvent.org/fathers/1002.htm>.}

% structure: p0001 source=h1 target=section reason=page-title

\section{四福音合参（第一部分）}

1 \BibleRef{若 1:1} 在起初已有圣言，圣言与天主同在，圣言就是天主 [2,3]。\BibleRef{若 1:2} 这道\Emph{原是}在起初与天主同在。\BibleRef{若 1:3} 万物是藉他的手造成的，4 凡受造的，没有一样不是藉着他\Emph{而造的}。\BibleRef{若 1:4} 在他内有生命，这生命 5 是人的光。\BibleRef{若 1:5} 光照在黑暗中，黑暗不能征服它。\par

6 \BibleRef{路 1:5} 在黑落德王的日子，有一位司祭，名叫匝加利亚，属于阿彼雅班次；他的妻子是亚郎的女儿，名叫依撒伯尔。\BibleRef{路 1:6} 两人在天主前都是义人，遵行主的一切诫命，在天主面前无可指摘。\BibleRef{路 1:7} 他们没有孩子，因为依撒伯尔不能生育，两人又都上了年纪。\BibleRef{路 1:8} 匝加利亚按班次在天主前尽司祭的职务时，\BibleRef{路 1:9} 照司祭的常例，抽中了签，得进主的圣殿烧香。\BibleRef{路 1:10} 烧香的时候，众百姓在外面祈祷。\BibleRef{路 1:11} 有上主的天使显现给匝加利亚，站在香坛右边。\BibleRef{路 1:12} 匝加利亚看见，惊惶失措，害怕起来。\BibleRef{路 1:13} 但天使对他说：\Quote{匝加利亚，不要害怕！因为你的祈祷已蒙应允，你的妻子依撒伯尔要给你生一个儿子，你要给他起名叫若翰。\BibleRef{路 1:14} 你必欢喜快乐，许多人也要因他的诞生而喜乐。\BibleRef{路 1:15} 他在主面前将是伟大的，淡酒浓酒都不喝，而且他还在母胎中就要充满圣神。\BibleRef{路 1:16} 他要使许多以色列子民归向上主、他们的天主。\BibleRef{路 1:17} 他要以厄里亚的精神和能力，走在主前面，使父亲的心转向儿子，使悖逆者转向义人的明智，为主准备一个完善的百姓。} \BibleRef{路 1:18} 匝加利亚对天使说：\Quote{我凭着什么能知道这事呢？因为我已经老了，我的妻子也上了年纪。} \BibleRef{路 1:19} 天使回答说：\Quote{我是加俾额尔，站在天主面前的，奉命来向你说话，报告你这个喜讯。\BibleRef{路 1:20} 看，你必成为哑巴，不能说话，直到这些事成就的那一天，因为你没有相信我的话；我的话到了时候必然应验。} \BibleRef{路 1:21} 百姓等候匝加利亚，都奇怪他滞留在圣殿内。\BibleRef{路 1:22} 及至他出来，竟不能对他们说话。他们就知道他在圣殿中见了异象；他向他们打手势，竟成了哑巴。\BibleRef{路 1:23} 他供职的日子满了，就回家去了。\par

25 \BibleRef{路 1:24} 这些日子以后，他的妻子依撒伯尔怀了孕；她隐藏了五个月，说：\BibleRef{路 1:25} \Quote{主在眷顾我的日子，这样待了我，除去了我在人前的羞辱。}\par

27 \BibleRef{路 1:26} 到了第六个月，天使加俾额尔奉天主差遣，往加里肋亚一座名叫纳匝肋的城去，\BibleRef{路 1:27} 到一位童贞女那里，她已许配给达味家族的一个名叫若瑟的男子，童贞女的名字叫玛利亚。\BibleRef{路 1:28} 天使进去向她问安说：\Quote{愿你平安！充满恩宠者，主与你同在！你在妇女中受赞颂。} \BibleRef{路 1:29} 她因这话惊惶不安，思虑这样的请安有什么意思。\BibleRef{路 1:30} 天使对她说：\Quote{玛利亚，不要害怕，因为你在天主前获得了宠幸。\BibleRef{路 1:31} 看，你将怀孕生子，并要给他起名叫耶稣。\BibleRef{路 1:32} 他将是伟大的，并被称为至高者的儿子，上主天主要把他祖先达味的御座赐给他。\BibleRef{路 1:33} 他要为王统治雅各伯家，直到永远；他的王权没有终结。} \BibleRef{路 1:34} 玛利亚对天使说：\Quote{这事怎能成就？因为我不认识男人。} \BibleRef{路 1:35} 天使回答她说：\Quote{圣神要临于你，至高者的能力要庇荫你，因此，那要诞生的圣者，将称为天主的儿子。\BibleRef{路 1:36} 且看，你的亲戚依撒伯尔，她虽在老年，却怀了男胎，本月已六个月了，她原是素称不生育的。\BibleRef{路 1:37} 因为在天主前没有不能的事。} \BibleRef{路 1:38} 玛利亚说：\Quote{看，上主的婢女，愿照你的话成就于我吧。} 天使就离开她去了。\par

40 \BibleRef{路 1:39} 玛利亚就在那几日起身，急忙往山区去，来到犹大的一座城，\BibleRef{路 1:40} 进了匝加利亚的家，问依撒伯尔安。\BibleRef{路 1:41} 依撒伯尔一听到玛利亚请安，胎儿就在她腹中欢跃。\BibleRef{路 1:42} 依撒伯尔充满了圣神，大声呼喊说：\Quote{你在妇女中受赞颂，你胎中的胎儿也是受赞颂的！\BibleRef{路 1:43} 我主的母亲到我这里来，这事哪里来的呢？\BibleRef{路 1:44} 看，你请安的声音一入我耳，胎儿就在我腹中欢喜踊跃。\BibleRef{路 1:45} 那信了由上主传于她的话必要完成的，是有福的。} \BibleRef{路 1:46} 玛利亚说：\par

我的灵魂颂扬上主，\newline{}48 \BibleRef{路 1:47} 我的心灵欢跃于天主，我的救主，\newline{}49 \BibleRef{路 1:48} 因为他垂顾了他婢女的卑微：\newline{}看，从今以后，\newline{}万代都要称我有福。\newline{}50 \BibleRef{路 1:49} 全能者为我行了大事，\newline{}他的名字是圣的。\newline{}51 \BibleRef{路 1:50} 他的仁慈世世代代归于敬畏他的人。\newline{}52 \BibleRef{路 1:51} 他挥动臂膀施展大能，\newline{}驱散那些心高气傲的人。\newline{}53 \BibleRef{路 1:52} 他从高位上推下权贵，\newline{}却举扬了卑微者。\newline{}54 \BibleRef{路 1:53} 使饥饿者饱享美物，\newline{}却使富者空手而去。\newline{}55 \BibleRef{路 1:54} 他扶助了他的仆人以色列，\newline{}纪念他的仁慈，\newline{}56 \BibleRef{路 1:55} 正如他对我们祖先所说的话，\newline{}对亚伯拉罕和他的后裔，直到永远。\par

57 \BibleRef{路 1:56} 玛利亚同依撒伯尔住了约三个月，就回家去了。\par

58, 59 \BibleRef{路 1:57} 依撒伯尔的产期到了，她生了一个儿子。\BibleRef{路 1:58} 她的邻居和亲戚听说上主大大施恩与她，59 就和她一同欢乐。\BibleRef{路 1:59} 到了第八天，他们来给孩子行割礼，并要叫他匝加利亚，\Emph{照他父亲的名字}称呼他。61 \BibleRef{路 1:60} 但他母亲回答说：\Quote{不，要叫他若翰。}62 \BibleRef{路 1:61} 他们便说：\Quote{你亲族中没有人叫这名字的。}63, 64 \BibleRef{路 1:62} 他们就向他父亲打手势，\Emph{问他}愿意叫他什么名字。\BibleRef{路 1:63} 他要了一块写字板，写道：\Quote{他的名字是若翰。}众人就惊讶。65 \BibleRef{路 1:64} 他的口和舌头立刻开了，就说话并赞美天主。\BibleRef{路 1:65} 所有邻居都害怕；这事传遍了犹大整个山区。67 \BibleRef{路 1:66} 凡听见的人都心里思想，说：\Quote{这孩子将会成为什么人物呢？}因为上主的手与他同在。\par

68 \BibleRef{路 1:67} 他的父亲匝加利亚充满了圣神，预言说：\par

69 \BibleRef{路 1:68} 上主，以色列的天主，应受赞美，\newline{}因他眷顾了他的子民，为他们施行了救恩；\newline{}70 \BibleRef{路 1:69} 并为我们兴起了救恩的角，\newline{}在他仆人达味的家中，\newline{}71 \BibleRef{路 1:70} 正如他借历代圣先知的口所说的，\newline{}72 \BibleRef{路 1:71} 拯救我们脱离敌人，\newline{}和一切恨我们的人手中。\newline{}73 \BibleRef{路 1:72} 他向我们的祖先施行仁慈，\newline{}纪念他的神圣盟约，\newline{}74 \BibleRef{路 1:73} 就是他向我们的祖先亚伯拉罕所起的誓，\newline{}75 \BibleRef{路 1:74} 要把我们从敌人手中救出，\newline{}使我们无所恐惧，在他面前侍奉，\newline{}76 \BibleRef{路 1:75} 以圣洁和正义度日。\newline{}77 \BibleRef{路 1:76} 至于你，孩子，你要称为至高者的先知，\newline{}你要走在主的前面，预备他的道路，\newline{}78 \BibleRef{路 1:77} 使他的百姓认识救恩，\newline{}以得赦免他们的罪过，\newline{}79 \BibleRef{路 1:78} 因我们天主的慈爱和怜悯，\newline{}他眷顾我们，使旭日从高天降临，\newline{}80 \BibleRef{路 1:79} 照亮坐在黑暗中、死影下的人，\newline{}引导我们的脚步走上平安的道路。\par

81 那孩子渐渐长大，心灵强健，住在旷野，直到他显现在以色列人面前的日子。\par

\SourceRecord{Translated by Hope W. Hogg.；Ante-Nicene Fathers；Vol. 9.；Edited by Allan Menzies.；Buffalo, NY: Christian Literature Publishing Co.,；1896.；<http://www.newadvent.org/fathers/100201.htm>.}

% structure: p0001 source=h1 target=section reason=page-title

\section{\WorkTitle{迪亚特萨隆}（第二部分）}

1 [阿拉伯文版第7页] \BibleRef{玛 1:18} 耶稣默西亚的诞生是这样的：当他的母亲许配给若瑟，在他们同居以前，2 她就因圣神怀孕了。\BibleRef{玛 1:19} 她的丈夫若瑟是个义人，不愿公开羞辱她，就打算秘密地休退她。3 \BibleRef{玛 1:20} 他正思虑这事的时候，上主的天使在梦中显现给他说：若瑟，达味之子，不要怕娶你的妻子玛利亚，因为在她内所怀的孕是从圣神来的。4 \BibleRef{玛 1:21} 她要生一个儿子，你要给他起名叫耶稣，因为他要把自己的民族从他们的罪恶中拯救出来。5 \BibleRef{玛 1:22} 这一切是为应验主藉先知所说的话：\par

6 \BibleRef{玛 1:23} 看，童贞女将怀孕生子，\newline{}人们将称他的名字为\OriginalTerm{厄玛奴耳}{Immanuel}，\par

7 那就是解释为：天主与我们同在。\BibleRef{玛 1:24} 若瑟从睡梦中醒来，就遵从上主的天使所吩咐的，娶了他的妻子；8 \BibleRef{玛 1:25} 但没有与她同房，直到她生了头胎儿子。\par

9 \BibleRef{路 2:1} 那时，凯撒奥古斯都出了一道旨意，叫天下的百姓都登记户口。10 \BibleRef{路 2:2} 这是季黎诺作叙利亚总督时的第一次登记。11, 12 \BibleRef{路 2:3} 于是众人各归本城去登记。\BibleRef{路 2:4} 若瑟也从加里肋亚的纳匝肋城上犹太去，到达味城名叫白冷（因为他本是达味家族和支派的人），13 \BibleRef{路 2:5} 与他的聘妻玛利亚一同去登记，那时她已怀孕。14 [阿拉伯文版第8页] \BibleRef{路 2:6} 15 他们在那里的时候，她分娩的日期到了。她生了头胎儿子，用襁褓裹起来，放在马槽里，因为在客栈里没有地方给他们住。\par

16 在那个地区有牧人露宿，夜间守更看顾羊群。17 \BibleRef{路 2:9} 忽然上主的天使来到他们身边，上主的光耀环照着他们，他们非常害怕。18 \BibleRef{路 2:10} 天使对他们说：不要害怕！因为我给你们带来大喜的信息，将临于万民；19 \BibleRef{路 2:11} 今天在达味城，为你们诞生了一位救主，就是主默西亚。20 \BibleRef{路 2:12} 这是给你们的标记：你们要找到一个婴儿，裹着襁褓，躺在马槽里。21 \BibleRef{路 2:13} 忽然有一大队天军与天使一同出现，赞美天主说：\par

22 \BibleRef{路 2:14} 在至高之处荣耀归于天主，\newline{}在地上平安，良善的意愿归于人。\par

23 \BibleRef{路 2:15} 众天使离开他们往天上去了，牧人们彼此说：我们到白冷去，看看所发生的事，就是上主所显示给我们的。24 \BibleRef{路 2:16} 他们急忙去了，找到了玛利亚和若瑟，还有那躺在马槽里的婴儿。25 \BibleRef{路 2:17} 他们看见以后，就把天使关于这孩子的话传开了。26 \BibleRef{路 2:18} 凡听见的人都诧异牧人们向他们所说的事。27 \BibleRef{路 2:19} 玛利亚却把这一切事铭记在心，反复思考。28 \BibleRef{路 2:20} 牧人们回去了，为所听见所看见的一切事，正如天使向他们所说的，就光荣赞美天主。\par

29 [阿拉伯文版第9页] \BibleRef{路 2:21} 满了八天，给孩子行割礼时，给他起名叫耶稣，这是他受孕前由天使所起的名字。\par

30 \BibleRef{路 2:22} 按梅瑟法律满了取洁的日子，他们便带孩子上耶路撒冷去献给上主，31 \BibleRef{路 2:23} （如主的法律上所记载的：凡开胎首生的男子，应称为主的圣物），32 \BibleRef{路 2:24} 并照主的法律上所说的，献上一对斑鸠或两只雏鸽。33 \BibleRef{路 2:25} 在耶路撒冷有一个人，名叫西默盎，这人正义虔诚，期待着以色列的安慰，而且圣神在他身上。34 \BibleRef{路 2:26} 他曾蒙圣神启示，在自己未死以前，必得看见上主的默西亚。35 \BibleRef{路 2:27} 他因圣神的感动，进入圣殿；正遇见耶稣的父母抱着孩子进来，要照法律的惯例为他行礼，36 \BibleRef{路 2:28} 他就接过孩子，抱在怀中，赞美天主说：\par

37 \BibleRef{路 2:29} 主啊！现在可以照你的话，\newline{}让你的仆人安然离世；\newline{}38 \BibleRef{路 2:30} 因为我亲眼看见了你的救恩，\newline{}39 \BibleRef{路 2:31} 即你为万民所预备的；\newline{}40 \BibleRef{路 2:32} 为启示异邦的光明，\newline{}为你百姓以色列的荣耀。\par

41 \BibleRef{路 2:33} 若瑟和他的母亲惊异关于耶稣所说的这些事。42 \BibleRef{路 2:34} 西默盎祝福他们，并对他的母亲玛利亚说：看，这孩子被立为以色列中许多人跌倒和兴起的原因，并成为反对的记号；43 \BibleRef{路 2:35} 至于你，一把利剑将刺透你的心灵，为叫许多人心中的思念显露出来。44 [阿拉伯文版第10页] \BibleRef{路 2:36} 又有一位女先知亚纳，是阿协尔支派法努尔的女儿，年纪已经老迈，从作童女出嫁，同丈夫住了七年，45 \BibleRef{路 2:37} 以后寡居，直到八十四岁。她总不离开圣殿，禁食祈祷，日夜事奉天主。46 \BibleRef{路 2:38} 那时她也上前来称谢上主，并向一切期待耶路撒冷得救的人讲论这孩子。47 \BibleRef{路 2:39} 他们按着上主的法律行完了一切，便返回加里肋亚，到自己的城纳匝肋去了。\par

\SourceRecord{Translated by Hope W. Hogg.；Ante-Nicene Fathers；Vol. 9.；Edited by Allan Menzies.；Buffalo, NY: Christian Literature Publishing Co.,；1896.；<http://www.newadvent.org/fathers/100202.htm>.}

% structure: p0001 source=h1 target=section reason=page-title

\section{\WorkTitle{迪亚泰萨隆（第三部分）}}

1, 2 \BibleRef{玛 2:1} 此后，贤士从东方来到耶路撒冷，\BibleRef{玛 2:2} 说：\BibleQuote{那生下来作犹太人之王的在哪里？我们在东方看见了他的星，特来朝拜他。} 3 \BibleRef{玛 2:3} 黑落德王听见了，就惊慌不安，全耶路撒冷也同他一起惊慌。4 \BibleRef{玛 2:4} 他便召集了众司祭长和民间的经师，问他们默西亚应当生在哪里。5 \BibleRef{玛 2:5} 他们说：\BibleQuote{在犹太的白冷，因为先知曾这样记载：}\par

6 \BibleRef{玛 2:6} \BibleQuote{你，犹大的白冷啊！\newline{}你在犹大的首领中决不是最小的，\newline{}因为将由你出来一位领袖，\newline{}他将牧养我的百姓以色列。}\par

7 \BibleRef{玛 2:7} 那时，黑落德暗暗地召见了贤士，仔细询问他们那颗星出现的时间。8 \BibleRef{玛 2:8} 然后打发他们往白冷去，说：\BibleQuote{你们去仔细寻访那婴孩，一旦找到了，就给我报信，好让我也去朝拜他。} 9 \BibleRef{玛 2:9} 他们听了王的话，就走了；看，他们在东方所见的那颗星，走在他们前面，直到婴孩所在的地方，就停在上面。10, 11 \BibleRef{玛 2:10} 他们一见那颗星，就极其高兴欢喜。\BibleRef{玛 2:11} 他们进了屋，看见婴孩和他的母亲玛利亚，就俯伏朝拜了他，打开自己的宝盒，奉献了礼物：黄金、乳香和没药。12 \BibleRef{玛 2:12} 他们在梦中得到指示，不要回到黑落德那里，就从另一条路返回了家乡。\par

13 \BibleRef{玛 2:13} 他们离去后，看，上主的天使在梦中显现给若瑟，说：\BibleQuote{起来，带着婴孩和他的母亲逃往埃及，住在那里，直到我再通知你；因为黑落德要搜寻这婴孩，把他杀死。} 14 \BibleRef{玛 2:14} 若瑟便起来，星夜带着婴孩和他的母亲，逃往埃及，15 \BibleRef{玛 2:15} 住在那里，直到黑落德去世。这就应验了上主借先知所说的话：\BibleQuote{我从埃及召回了我的儿子。} 16 \BibleRef{玛 2:16} 那时，黑落德见自己被贤士愚弄，就大怒，派人将白冷及其周围境内所有两岁及两岁以下的男婴全都杀死，照着他从贤士那里仔细查问的时间推算。17 \BibleRef{玛 2:17} 这就应验了耶肋米亚先知所说的话：\par

18 \BibleRef{玛 2:18} \BibleQuote{在辣玛听到了声音，\newline{}是痛哭和极大的哀号；\newline{}辣黑耳痛哭她的儿女，\newline{}不愿受安慰，因为他们不在了。}\par

19 \BibleRef{玛 2:19} 黑落德死后，看，上主的天使在埃及梦中显现给若瑟，20 \BibleRef{玛 2:20} 说：\BibleQuote{起来，带着婴孩和他的母亲，往以色列地去；因为那些谋害婴孩性命的人已经死了。} 21 \BibleRef{玛 2:21} 若瑟便起来，带着婴孩和他的母亲，来到以色列地。22 \BibleRef{玛 2:22} 但他听说阿尔赫劳继他父亲黑落德作了犹太王，就害怕到那里去；在梦中得到指示，便往加里肋亚境内去了，23 \BibleRef{玛 2:23} 并在一个名叫纳匝肋的城里住下：这就应验了先知所说的话：\BibleQuote{他将称为纳匝肋人。}\par

24 \BibleRef{路 2:40} 婴孩渐渐长大，在心灵上强壮起来，充满智慧；天主的恩宠常在他身上。\par

25 \BibleRef{路 2:41} 他的父母每年逾越节往耶路撒冷去。26 \BibleRef{路 2:42} 当他十二岁时，他们照节日的惯例上去了。27 \BibleRef{路 2:43} 过完了节日，他们回去的时候，孩童耶稣仍留在耶路撒冷，若瑟和他的母亲并不知道。28 \BibleRef{路 2:44} 他们以为他同同行的人在一起，走了一天的路程，就在亲戚和相识的人中寻找他，29 \BibleRef{路 2:45} 却没有找到，便返回耶路撒冷寻找他。30 \BibleRef{路 2:46} 过了三天，他们在圣殿里找到了他，坐在经师中，聆听他们，也询问他们\Emph{问题}；31 \BibleRef{路 2:47} 凡听见他的人，都惊奇他的智慧和应对。32 \BibleRef{路 2:48} 他们看见他，便大为惊异，他的母亲对他说：\BibleQuote{孩子，为什么你这样对待我们？看，你的父亲和我，一直痛苦地寻找你。} 33 \BibleRef{路 2:49} 耶稣对他们说：\BibleQuote{你们为什么寻找我？难道你们不知道我必须在我父亲家里吗？} 34 \BibleRef{路 2:50} 但他们不明白他对他们所说的话。35 \BibleRef{路 2:51} 他就同他们下去，来到纳匝肋，并服从他们；他的母亲把这一切话默存在心中。\par

36 \BibleRef{路 2:52} 耶稣在智慧和身量上，并在天主和人前的恩爱上，渐渐增长。\par

37 \BibleRef{路 3:1} 凯撒提庇留执政第十五年，般雀比拉多作犹太总督，黑落德作加里肋亚分封王，他的兄弟斐理伯作依突勒雅和特辣工尼雅地区的分封王，吕撒尼雅作阿彼肋乃分封王，38 \BibleRef{路 3:2} 亚纳斯和盖法作大司祭时，天主的话在旷野中传到了匝加利亚的儿子若翰身上。39 \BibleRef{路 3:3} 他来到约但河附近整个地区，宣讲悔改的洗礼，为得罪之赦。40 \BibleRef{玛 3:1} 他在犹太旷野宣讲说：41 \BibleRef{玛 3:2} \BibleQuote{你们悔改吧！因为天国临近了。} 42 \BibleRef{玛 3:3} 这人就是依撒意亚先知所预言的那位：\par

\BibleQuote{在旷野里有呼喊者的声音：\newline{}43 \BibleRef{路 3:4} 你们预备主的道路，\newline{}修直他的途径！\newline{}一切深谷要填满，\newline{}一切山岳丘陵要削平，\newline{}44 \BibleRef{路 3:5} 弯曲的要修直，崎岖的要成为坦途！}\par

45 \BibleRef{路 3:6} \BibleQuote{凡有血肉的，都要看见天主的救恩。}\par

46 \BibleRef{若 1:7} \BibleQuote{这个\Emph{人}来作证，为给光作证，为使 47 众人藉着他的媒介而相信。} \BibleRef{若 1:8} \BibleQuote{他不是那光，而是他 48 要为光作证。} \BibleRef{若 1:9} \BibleQuote{那是真理之光，光照 49 来到世界上的每一个人。} \BibleRef{若 1:10} \BibleQuote{他在世界，世界藉着他 50 造成，世界 51 却不认识他。} \BibleRef{若 1:11} \BibleQuote{他来到自己的人中间，自己的人却不接待他。} \BibleRef{若 1:12} \BibleQuote{凡接待他的，他就赐给他们权力，使他们 52 成为天主的儿女——那些信他名的人：} \BibleRef{若 1:13} \BibleQuote{他们不是从血统生的，53 不是从肉情生的，也不是从人意生的，而是从天主生的。} \BibleRef{若 1:14} \BibleQuote{圣言成了血肉，居住在我们中间；我们看见了他的光荣，正如 54 来自圣父的独生\Emph{子}的光荣，充满恩宠和真理。} \BibleRef{若 1:15} \BibleQuote{\Person{若翰}为他作证\EditorialNote{阿拉伯语文本，第14页}，呼喊说：\InnerQuote{这就是我曾说的：他是在我以后来的，并且 55 在我以前存在，因为他本来在我以前。}} \BibleRef{若 1:16} \BibleQuote{从他的满盈中，56 我们都领受了恩宠，而且恩上加恩。} \BibleRef{若 1:17} \BibleQuote{因为法律是藉着\Person{梅瑟}传授的，恩宠和真理却是由\Person{耶稣基督}而来的。}\par

\SourceRecord{Translated by Hope W. Hogg.；Ante-Nicene Fathers；Vol. 9.；Edited by Allan Menzies.；Buffalo, NY: Christian Literature Publishing Co.,；1896.；<http://www.newadvent.org/fathers/100203.htm>.}

% structure: p0001 source=h1 target=section reason=page-title

\section{\WorkTitle{四福音合参}（第四部分）}

1 \BibleRef{若 1:18} \Quote{从来没有人见过天主；只有那在父怀里的独\Emph{生子}，他是天主，他讲述了\Emph{他。}}\par

2 \BibleRef{若 1:19} \Quote{这就是若翰的见证：那时，犹太人从耶路撒冷派遣了司祭和肋未人到他那里，问他：\Quote{你是谁？}} 3 \BibleRef{若 1:20} \Quote{他承认了，并没有否认，他公开说：\Quote{我不是默西亚。}} 4 \BibleRef{若 1:21} \Quote{他们又问他说：\Quote{那么你是谁？你是厄里亚吗？}他说：\Quote{我不是。}\Quote{你是那位先知吗？}他回答说：\Quote{不是。}} 5 \BibleRef{若 1:22} \Quote{他们便对他说：\Quote{你到底是谁？好叫我们给那派遣我们来的人一个答复。你自己说你是谁？}} 6 \BibleRef{若 1:23} \Quote{他说：\Quote{我是那在旷野里呼喊者的声音：修直主的道路，正如依撒意亚先知所说的。}} 7 \BibleRef{若 1:24} \Quote{那些被派遣来的是法利塞人。} 8 \BibleRef{若 1:25} \Quote{他们又问他说：\Quote{你既不是默西亚，又不是厄里亚，也不是那位先知，那么你为什么施洗呢？}} 9 \BibleRef{若 1:26} \Quote{若翰回答他们说：\Quote{我用水施洗；但有一位站在你们中间，是你们所不认识的；} 10 \BibleRef{若 1:27} \Quote{就是那在我以后来的，我给他解鞋带也不配。}} 11 \BibleRef{若 1:28} \Quote{这事发生在约但河对岸的伯达尼，那里是若翰施洗的地方。}\par

12 \BibleRef{玛 3:4} \Quote{若翰穿着骆驼毛的衣服，腰间束皮带，吃的是蝗虫和野蜜。} 13 \BibleRef{玛 3:5} \Quote{那时，耶路撒冷、全犹太和约但河一带的人都出来到他那里去，} 14 \BibleRef{玛 3:6} \Quote{在约但河里受他的洗，承认自己的罪。} 15 \BibleRef{玛 3:7} \Quote{但他看见许多法利塞人和撒杜塞人也来受洗，就对他们说：\Quote{毒蛇的种类！谁指示你们逃避将来的忿怒呢？} 16 \BibleRef{玛 3:8} \Quote{那么，你们就结出与悔改相称的果实吧！} 17 \BibleRef{玛 3:9} \Quote{不要心里说：\InnerQuote{我们有亚巴郎为父。}我告诉你们：天主能从这些石头给亚巴郎兴起子孙来。} 18 \BibleRef{玛 3:10} \Quote{斧子已放在树根上；凡不结好果子的树，就砍下来，丢在火里。}} 19 \BibleRef{路 3:10} \Quote{群众问他说：\Quote{那么，我们该作什么呢？}} 20 \BibleRef{路 3:11} \Quote{他回答他们说：\Quote{有两件内衣的，要分给那没有的；有食物的，也要照样做。}} 21 \BibleRef{路 3:12} \Quote{税吏也来受洗，问他说：\Quote{师傅，我们该作什么呢？}} 22 \BibleRef{路 3:13} \Quote{他对他们说：\Quote{除了给你们规定的以外，不要多征取。}} 23 \BibleRef{路 3:14} \Quote{兵士也问他说：\Quote{我们该作什么呢？}他对他们说：\Quote{不要勒索人，也不要讹诈；要以你们的粮饷为满足。}}\par

24 那时，人民对若翰都怀着期待，心中思量：或许他就是默西亚。 25 \BibleRef{路 3:16} \Quote{若翰回答众人说：\Quote{我固然以水洗你们；但有一位比我更有能力的要来，我就是给他解鞋带也不配。他要以圣神和火洗你们。} 26 \Quote{他手里拿着簸箕，要扬净他的禾场，把麦子收进仓里，而把糠秕用不灭的火烧尽。}}\par

27 \BibleRef{路 3:18} \Quote{他还用许多别的话劝勉人民，向他们宣讲福音。}\par

28 \BibleRef{玛 3:13} \Quote{那时，耶稣从加里肋亚来到约但河若翰那里，要受他的洗。} 29 \BibleRef{路 3:23} \Quote{耶稣开始传道的时候，大约三十岁，人都以为他是若瑟的儿子。} 30 \BibleRef{若 1:29} \Quote{第二天，若翰看见耶稣向他走来，便说：\Quote{看哪，天主的羔羊，除免世罪者！} 31 \BibleRef{若 1:30} \Quote{这就是我曾说：\InnerQuote{有一个人在我以后来，却成了在我以前的，因为他原在我以前。}} 32 \BibleRef{若 1:31} \Quote{我本来不认识他；但为了使他显示于以色列，为此，我来用水施洗。}} 33 \BibleRef{玛 3:14} \Quote{若翰想要阻止他，说：\Quote{我需要受你的洗，你反而到我这里来吗？}} 34 \BibleRef{玛 3:15} \Quote{耶稣回答他说：\Quote{你暂且容许吧！因为我们应当这样履行全义。}于是若翰就容许了他。} 35 \BibleRef{路 3:21} \Quote{众百姓都受了洗，耶稣也受了洗。} 36 \BibleRef{玛 3:16} \Quote{耶稣刚受了洗，就从水里上来，忽然天为他开了，} 37 \BibleRef{路 3:22} \Quote{圣神以鸽子的形状降在他头上，} 38 \BibleRef{玛 3:17} \Quote{并有声音从天上说：\Quote{这是我的爱子，我所喜悦的。}} 39 \BibleRef{若 1:32} \Quote{若翰又作证说：\Quote{我看见圣神好像鸽子从天降下，住在他身上。} 40 \BibleRef{若 1:33} \Quote{我本来不认识他；但那派遣我来用水施洗的，对我说：\InnerQuote{你看见圣神降下，并停在他身上，他就是那要以圣神施洗的人。}} 41 \BibleRef{若 1:34} \Quote{我看见了，便作证：这就是天主子。}}\par

42 \BibleRef{路 4:1} \Quote{耶稣充满圣神，从约但河回来，} 43 \BibleRef{谷 1:12} \Quote{圣神立刻催他到旷野里去，} \BibleRef{谷 1:13} \Quote{他在那里四十天，受撒殚的试探；他与野兽在一起，} 44 \BibleRef{玛 4:2} \Quote{他四十昼夜禁食，后来就饿了。} \BibleRef{路 4:2} \Quote{那些日子，他什么也没有吃；日子满了，他就饿了。} 45 \Quote{试探者前来对他说：\Quote{你若是天主子，就吩咐这些石头变成饼。}} 46 \BibleRef{玛 4:4} \Quote{他回答说：\Quote{经上记载：\InnerQuote{人生活不只靠饼，而也靠天主口中所发的一切言语。}}} 47 \BibleRef{玛 4:5} \Quote{于是，魔鬼带他进了圣城，叫他站在圣殿顶上，} 48 \BibleRef{玛 4:6} \Quote{对他说：\Quote{你若是天主子，就跳下去，因为经上记载：\InnerQuote{他为你吩咐了自己的天使，他们要用手托着你，免得你的脚碰在石头上。}}}\par

\Quote{他为你吩咐了自己的天使，\newline{}他们要用手托着你，\newline{}免得你的脚碰在石头上。}\par

49 \BibleRef{玛 4:7} 耶稣对他说：\BibleQuote{经上也记载：你不可试探主你的 50 天主。} \BibleRef{路 4:5} 魔鬼就带他上了一座高山，在霎时间内把普世万国 【阿拉伯文，第18页】 及其荣华指给他看；\BibleRef{路 4:6} 魔鬼对他说：\BibleQuote{这一切权柄及其荣华，我都要给你，因为这一切原已交付给了 52 我，我愿意给谁就给谁。所以你若俯伏朝拜我，这一切全是你的。}\BibleRef{路 4:7}\par

\SourceRecord{Translated by Hope W. Hogg.；Ante-Nicene Fathers；Vol. 9.；Edited by Allan Menzies.；Buffalo, NY: Christian Literature Publishing Co.,；1896.；<http://www.newadvent.org/fathers/100204.htm>.}

% structure: p0001 source=h1 target=section reason=page-title

\section{《迪亚泰萨隆》（第五部分）}

1 \BibleRef{玛 4:10} \Person{耶稣}回答他说：\BibleQuote{去吧，撒旦！因为经上记载：\InnerQuote{你当朝拜主你的天主，唯独事奉他。}} 2 \BibleRef{路 4:13} 当魔鬼完成了他的一切试探，就暂时离开了他。 3 \BibleRef{玛 4:11} 看啊，天使们前来侍奉他。\par

4, 5 \BibleRef{若 1:35} 第二天，若翰站着，他的两个门徒也在那里； \BibleRef{若 1:36} 他看见耶稣走过，就说：\BibleQuote{看，天主的羔羊！} 6 \BibleRef{若 1:37} 他的两个门徒听见他说\Emph{这话}，就跟随了耶稣。 \BibleRef{若 1:38} 耶稣转过身来，看见他们跟着，就对他们说：\BibleQuote{你们找什么？} 他们对他说：\BibleQuote{老师，你住在哪里？} 8 \BibleRef{若 1:39} 他对他们说：\BibleQuote{你们来看看。} 他们就去看了他的住处，当天就和他住下了；那时大约是第十时辰。 9 \BibleRef{若 1:40} 那两个从若翰听见并跟随耶稣的人中，有一个是西满的兄弟安德肋。 10 \BibleRef{若 1:41} 他先找到自己的兄弟西满，对他说：\BibleQuote{我们找到了默西亚。} 11 \BibleRef{若 1:42} 就带他到耶稣跟前。耶稣注视着他说：\BibleQuote{你是若纳的儿子西满，你要称为刻法。}\par

12 \BibleRef{若 1:43} 第二天，耶稣决意往加里肋亚去，遇见斐理伯，就对他说：\BibleQuote{跟随我。} 13 \EditorialNote{阿拉伯文版，第19页} \BibleRef{若 1:44} 斐理伯是贝特赛达人，与安德肋和西满同城。 14 \BibleRef{若 1:45} 斐理伯找到纳塔乃耳，对他说：\BibleQuote{梅瑟在法律和先知书上所写的那位，我们找到了，就是若瑟的儿子，纳匝肋人耶稣。} 15 \BibleRef{若 1:46} 纳塔乃耳对他说：\BibleQuote{纳匝肋能出什么好事吗？} 斐理伯对他说：\BibleQuote{你来看看。} 16 \BibleRef{若 1:47} 耶稣看见纳塔乃耳向他走来，就指着他说：\BibleQuote{看，这确是一个以色列人，在他内毫无诡诈。} 17 \BibleRef{若 1:48} 纳塔乃耳对他说：\BibleQuote{你从那里认识我？} 耶稣回答说：\BibleQuote{斐理伯叫你以前，当你还在无花果树下时，我就看见了你。} 18 \BibleRef{若 1:49} 纳塔乃耳回答他说：\BibleQuote{老师，你是天主子，你是以色列的君王。} 19 \BibleRef{若 1:50} 耶稣对他说：\BibleQuote{因为我告诉你，我看见你在无花果树下，你就信了吗？你要看见比这更大的事。} 20 \BibleRef{若 1:51} 又对他说：\BibleQuote{我实实在在告诉你们：你们要看见天开了，天主的使者在人子身上上去下来。}\par

21 \BibleRef{路 4:14} 耶稣充满圣神的能力，回到加里肋亚。\par

22 \BibleRef{若 2:1} 第三天，在加里肋亚的加纳有婚宴，耶稣的母亲在那里； 23 \BibleRef{若 2:2} 耶稣和他的门徒也被请去赴宴。 24 \BibleRef{若 2:3} 酒缺了，耶稣的母亲对他说：\BibleQuote{他们没有酒了。} 25 \BibleRef{若 2:4} 耶稣对她说：\BibleQuote{妇人，这与我和你有什么关系？我的时刻还没有到。} 26 \BibleRef{若 2:5} 他的母亲对仆人说：\BibleQuote{他吩咐你们什么，你们就做什么。} 27 \BibleRef{若 2:6} 那里有六口石缸，按犹太人的取洁礼安放，每口可盛两三桶水。 28 \EditorialNote{阿拉伯文版，第20页} \BibleRef{若 2:7} 耶稣对他们说：\BibleQuote{把缸灌满水。} 他们就灌满了，直到缸沿。 29 \BibleRef{若 2:8} 耶稣说：\BibleQuote{现在舀出来，送给司席。} 他们就\Emph{这样}做了。 30 \BibleRef{若 2:9} 司席尝了那水变成的酒，并不知道是那里来的（舀水的仆人却知道），司席便叫新郎来， 31 \BibleRef{若 2:10} 对他说：\BibleQuote{众人都是先摆上好酒，等客人喝醉了，才摆上次等的；你却把好酒留到如今。} 32 \BibleRef{若 2:11} 这是耶稣在加里肋亚的加纳所行的第一个神迹，显示了他的光荣；他的门徒就信了他。 33 \BibleRef{路 4:14} 他的名声传遍了周围各地。 34 \BibleRef{路 4:15} 他在他们的会堂里施教，受到众人的称扬。 35 \BibleRef{路 4:16} 他来到了纳匝肋，就是他长大的地方；按照他的习惯，在安息日进了会堂，站起来要诵读。 36 \BibleRef{路 4:17} 有人把依撒意亚先知书递给他。他展开书卷，找到一处，上边写着：\par

37 \BibleRef{路 4:18} \BibleQuote{主的圣神在我身上，\newline{}因为他膏抹了我，向穷人传报喜讯；\newline{}他派遣我医治破碎的心灵，\newline{}向作恶者宣告赦免，向盲者宣告复明，\newline{}使破碎者得蒙赦免，\newline{}38 \BibleRef{路 4:19} 宣告上主恩慈之年。}\par

39 \BibleRef{路 4:20} 他把书卷卷起来，交给仆役，就坐下；会堂里众人都注视着他。 40 \BibleRef{路 4:21} 他便开始对他们说：\BibleQuote{今天，你们耳中所听的这段圣经，已经应验了。} 41 \BibleRef{路 4:22} 众人都为他作证，并惊奇他口中所出的恩宠之言。\par

42 \EditorialNote{阿拉伯文版，第21页} \BibleRef{玛 4:17} 从那时起，耶稣开始宣讲天国的福音，说：\BibleQuote{你们悔改吧！信从福音。} 43 \BibleRef{谷 1:15} \BibleQuote{时期已满，天国临近了。}\par

44 \BibleRef{玛 4:18} 耶稣沿着加里肋亚海行走时，看见两个兄弟：称为刻法的西满和他的兄弟安德肋，在海里撒网，他们原是渔夫。 45 \BibleRef{玛 4:19} 耶稣对他们说：\BibleQuote{来跟随我，我要使你们成为捕人的渔夫。} 46 \BibleRef{玛 4:20} 他们立刻撇下网，跟随了他。 47 \BibleRef{玛 4:21} 从那里再往前行，看见了另外两个兄弟：载伯德的儿子雅各伯和他的兄弟若望，在船上同他们的父亲载伯德修理渔网；耶稣就召叫了他们。 48 \BibleRef{玛 4:22} 他们立刻舍下船和父亲，跟随了他。\par

49 \BibleRef{路 5:1} 当群众聚集到他跟前，要听天主的话，他50正站在\Place{革乃撒勒湖}边，\BibleRef{路 5:2} 他看见两只船停在湖边，两个渔夫从船上出来正在洗51网。\BibleRef{路 5:3} 其中一只船属于\Person{西满·刻法}。耶稣上去，坐在船上，吩咐他们稍微离开陆地往52水里去。他就坐下，从船上教训群众。\BibleRef{路 5:4} 他讲完了，就对西满说：\Quote{划到深处，撒你们的53网捕鱼。}\BibleRef{路 5:5} 西满回答说：\Quote{老师，我们整夜54劳力，并没有打着什么；但依你的话，我要下网。}\BibleRef{路 5:6} 他们这样做了，就圈住了许多鱼，他们的网险些55裂开。\BibleRef{路 5:7} 他们就招呼另一只船上的同伴来协助。他们来了，把两只船都装满了，以致船要下沉。\par

\SourceRecord{Translated by Hope W. Hogg.；Ante-Nicene Fathers；Vol. 9.；Edited by Allan Menzies.；Buffalo, NY: Christian Literature Publishing Co.,；1896.；<http://www.newadvent.org/fathers/100205.htm>.}

% structure: p0001 source=h1 target=section reason=page-title

\section{\WorkTitle{四福音合参（第六部分）}}

1 [Arabic, p. 22] \BibleRef{路 5:8} \BibleQuote{但当西满刻法看见\Emph{这事}，就俯伏在耶稣脚前，对他说：\InnerQuote{主，我请求你离开我，因为我是个 2 有罪的人。}} \BibleRef{路 5:9} \BibleQuote{惊异攫住了他及所有与他在一起的人，3 因他们所捕到的鱼。} \BibleRef{路 5:10} \BibleQuote{同样，载伯德的儿子雅各伯和若望——西满的伙伴——也被占据。耶稣对西满说：\InnerQuote{4 不要怕；从今以后你将成为得人的渔夫。}} \BibleRef{路 5:11} \BibleQuote{他们把船划到岸上，舍弃一切，跟随了他。}\par

5 \BibleRef{若 3:22} \BibleQuote{此后，耶稣和门徒来到犹太地；他在那里与他们周游 6 并授洗。} \BibleRef{若 3:23} \BibleQuote{若翰也在艾农——靠近撒冷——授洗，因为那里水多；他们前来受洗。} 7, 8 \BibleRef{若 3:24} \BibleQuote{那时若翰尚未被下狱。} \BibleRef{若 3:25} \BibleQuote{于是若翰的一个门徒和一个犹太人关于洁净发生了争辩 9。} \BibleRef{若 3:26} \BibleQuote{他们来到若翰面前，对他说：\InnerQuote{老师，先前与你在约但河对岸、你曾为他作证的那位，看，他也授洗，且众人都往他那里去。}} 10 \BibleRef{若 3:27} \BibleQuote{若翰回答说：\InnerQuote{人不能领受什么，除非 11 是天堂赐予他的。}} \BibleRef{若 3:28} \BibleQuote{你们自己为我作证：我曾说：\InnerQuote{我不是默西亚，我只是在他前面奉差遣的一位。}} \BibleRef{若 3:29} \BibleQuote{娶新娘的是新郎；新郎的朋友站着听新郎的声音，就因新郎的声音而大大喜乐。如今，看，13, 14 [Arabic, p. 23] 我的喜乐得以满足。} \BibleRef{若 3:30} \BibleQuote{他必兴旺，我必衰微。} \BibleRef{若 3:31} \BibleQuote{从天堂来的那一位，超越一切；那出于地的，属于地，也谈论地；从天堂来的那一位，15 超越万有。} \BibleRef{若 3:32} \BibleQuote{他见证他所见所闻的，却没有人 16 接受他的见证。} \BibleRef{若 3:33} \BibleQuote{那接受他见证的人，就证实了天主是 17 真实的。} \BibleRef{若 3:34} \BibleQuote{天主所派遣的那一位，讲论天主的话；天主把圣神 18 无限量地赐予他。} \BibleRef{若 3:35} \BibleQuote{圣父爱圣子，并将万有交在 19 他手中。} \BibleRef{若 3:36} \BibleQuote{信圣子的人，得永生；不服从圣子的人，必不得见生命，天主的义怒常在他身上。}\par

20 \BibleRef{若 4:1} \BibleQuote{耶稣得知法利塞人听见他收了许多门徒，21 并施洗比若翰更多} \BibleRef{若 4:2} \BibleQuote{（其实耶稣本人不施洗，而是他的门徒施洗）} \BibleRef{若 4:3} \BibleQuote{，就\Emph{这样}离开了犹太。}\par

23 \BibleRef{路 3:19} \BibleQuote{分封侯黑落德，因受若翰指责——为兄弟斐理伯的妻子黑落狄雅，以及他所行的一切恶事——24 又加上了这一件：} \BibleRef{路 3:20} \BibleQuote{把若翰关在狱中。}\par

25 \BibleRef{玛 4:12} \BibleQuote{耶稣听见若翰被交付，就退到加里肋亚去。} 26 \BibleRef{若 4:46} \BibleQuote{他又来到加纳，就是他从前变水为酒的地方。那里 27 在葛法翁有一个王臣，他的儿子病了。} \BibleRef{若 4:47} \BibleQuote{这个\Emph{人}听见耶稣从犹太来到加里肋亚，就去见他，恳求他下去医治他的儿子；因为他快要死了。} 28, 29 \BibleRef{若 4:48} \BibleQuote{耶稣对他说：\InnerQuote{除非你们看见神迹和奇事，你们总是不信。}} \BibleRef{若 4:49} \BibleQuote{那[Arabic, p. 24]王臣对他说：\InnerQuote{主，求你下去，免得孩子死了。}} 30 \BibleRef{若 4:50} \BibleQuote{耶稣对他说：\InnerQuote{去吧！你的儿子活了。}那人信了耶稣所说的话 31 便走了。} \BibleRef{若 4:51} \BibleQuote{他正下去的时候，仆人们迎见他 32 告诉他说：\InnerQuote{你的儿子活了。}} \BibleRef{若 4:52} \BibleQuote{他就问他们什么时候好转的。他们回答说：\InnerQuote{昨天第七时辰，热症就离开了 33 他。}} \BibleRef{若 4:53} \BibleQuote{父亲就知道那正是耶稣对他说：\InnerQuote{你的儿子活了。}的那个时辰。} 34 \BibleRef{若 4:54} \BibleQuote{于是他和全家都信了。这是耶稣从犹太回到加里肋亚后行的 35 第二个神迹。} \BibleRef{路 4:44} \BibleQuote{他就在加里肋亚的各会堂中宣讲。}\par

36 \BibleRef{玛 4:13} \BibleQuote{他离开纳匝肋，来到葛法翁——海滨地区，在则步隆和纳斐塔里境内 37 居住：} \BibleRef{玛 4:14} \BibleQuote{这应验了依撒意亚先知所说的话：}\par

38 \BibleRef{玛 4:15} \BibleQuote{则步隆地，纳斐塔里地，\newline{}沿海之路，约但河那边，\newline{}外邦人的加里肋亚：} \newline{}39 \BibleRef{玛 4:16} \BibleQuote{那坐在黑暗中的百姓，\newline{}看见了皓光；\newline{}那些坐在死亡阴影之地的人，\newline{}有光为他们出现。}\par

40 \BibleRef{路 4:31} \BibleQuote{他在安息日教导他们。} \BibleRef{路 4:32} \BibleQuote{他们对他的教训感到惊奇：41 因为他的话具有权威。} \BibleRef{路 4:33} \BibleQuote{会堂里 42 有一个人被不洁的魔鬼附身，他大声喊叫说：} \BibleRef{路 4:34} \BibleQuote{\InnerQuote{任凭我们吧！纳匝肋的耶稣，我们与你有什么相干？你来毁灭我们 43 吗？我知道你是谁，你是天主的圣者。}} \BibleRef{路 4:35} \BibleQuote{耶稣斥责他说：\InnerQuote{闭口，从他身上出来！}魔鬼把他摔在 44 众人中间，就从他身上出来了，没有伤害他。} \BibleRef{路 4:36} \BibleQuote{众人大为惊骇[Arabic, p. 25]，彼此谈论说：\InnerQuote{这是怎么回事？他以权柄和能力命令不洁的魔鬼，他们竟 45 出来了？}} \BibleRef{路 4:37} \BibleQuote{于是他的名声传遍了周围各地。}\par

46 \BibleRef{路 4:38} \BibleQuote{耶稣出了会堂，看见一个名叫玛窦的人坐在税关那里，就对他说：\InnerQuote{跟随我。}他就起来跟随了他。}\par

47, 48 \BibleRef{谷 1:29} 耶稣同雅各伯和若望，来到西满和安德肋的家里。\BibleRef{路 4:38} 西满的岳母正发高烧，他们就为她49求耶稣。\BibleRef{路 4:39} 耶稣站在她上面，斥责热症，热就退了，她立刻50起来服事他们。\BibleRef{玛 8:16} 到了晚上，人们把许多附魔的51带到他跟前，他就用话语把他们身上的魔鬼赶了出去。\BibleRef{路 4:40} 凡有病人的，他们的病症各不相同\Emph{而且}恶性，都带到他那里；他就按手52在他们每人身上，治好了他们。\BibleRef{玛 8:17} 这是为应验先知依撒意亚53所说的话：\BibleQuote{他承担了我们的痛苦，担负了我们的疾病。}\BibleRef{谷 1:33} 全城54的人都聚集在耶稣的门口。\BibleRef{路 4:41} 他又从许多人身上赶出魔鬼，魔鬼喊叫说：\Quote{你是\OriginalTerm{默西亚}{Messiah}，\OriginalTerm{天主子}{Son of God}；}耶稣却斥责他们，不许魔鬼说话，因为他们知道他是\OriginalTerm{主默西亚}{Lord Messiah}。\par

\SourceRecord{Translated by Hope W. Hogg.；Ante-Nicene Fathers；Vol. 9.；Edited by Allan Menzies.；Buffalo, NY: Christian Literature Publishing Co.,；1896.；<http://www.newadvent.org/fathers/100206.htm>.}

% structure: p0001 source=h1 target=section reason=page-title

\section{\WorkTitle{迪亚提沙龙}（第七部分）}

1 [阿拉伯语第26页] \BibleRef{谷 1:35} 那天清晨，他很早就出去，到了一处 2 旷野地方，在那里祈祷。\BibleRef{谷 1:36} 西满和那些与他 3 在一起的人寻找他。\BibleRef{谷 1:37} 他们找到了他，就对他说：\Quote{众人都找你。}\BibleRef{谷 1:38} 他对他们说：\Quote{我们到邻近的村庄和城镇去吧，好叫我也 5 在那里宣讲，因为我正是为此而来的。}\BibleRef{路 4:42} 群众寻找他，来到他那里，挽留他，不要他 6 离开他们。\BibleRef{路 4:43} 但耶稣对他们说：\Quote{我也必须在别的城宣讲天 7 主的国，因为我正是为此被派遣的。}\BibleRef{玛 9:35} 耶稣走遍各城各村，在他们的会堂里教导，宣讲天国的福音，医治各种疾病和各种灾殃，8 \BibleRef{谷 1:39} 并驱逐魔鬼。\BibleRef{路 4:14} 他的名声传开了，\BibleRef{路 4:15} 他在 9 各处教导，受到众人的赞扬。\BibleRef{谷 2:14} 当他经过时，看见肋未——阿尔斐的儿子坐在税关上，就对他说：\Quote{跟随 10 我！}他就起来跟随了耶稣。\BibleRef{玛 4:24} 他的名声传遍了叙利亚全境；人们把所有患各种疾病、受各种痛苦、附魔的、癫痫的、瘫痪的人都带到他跟前，他就治好了他们。\par

11, 12 \BibleRef{谷 2:1} 过了几天，耶稣又进了葛法翁。\BibleRef{谷 2:2} 众人听说他在屋里，就聚集了许多人，以致连门前 13 [阿拉伯语第27页] 都容不下；他就给他们讲论天主的道。\BibleRef{路 5:17} 在那里有法利塞人和法律教师坐着，他们是从加里肋亚各村庄、犹太和耶路撒冷来的；主的能力临在，以医治他们。\BibleRef{路 5:18} 有人用床抬着一个瘫子。15 他们想法要把他抬进来，放在他面前。\BibleRef{路 5:19} 但因人多，找不到方法抬进去，就上了房顶，从屋瓦中间把他连床缒到耶稣面前。16 \BibleRef{路 5:20} 耶稣看见他们的信德，就对瘫子说：\Quote{孩子，你的罪赦了 17。}\BibleRef{路 5:21} 经师和法利塞人开始心里议论说：\Quote{这人是谁？竟说亵渎的话？除了天主以外，谁能赦罪呢？}18 \BibleRef{谷 2:8} 耶稣凭神知道了他们心里这样议论，就对他们 19 说：\Quote{你们心里为什么这样议论？\BibleRef{谷 2:9} 哪一样更容易呢？是对瘫子说\InnerQuote{你的罪赦了}，或是说 \Emph{那}\InnerQuote{起来，拿起你的床行走}呢？\BibleRef{谷 2:10} 为使你们知道人子在地上 21 有权柄赦罪——（他对瘫子说）——\BibleRef{谷 2:11} 我吩咐你：22\InnerQuote{起来，拿起你的床，回家去吧！}}\BibleRef{谷 2:12} 那人立刻起来，拿起床，当着众人面前出去了。\BibleRef{路 5:25} 他回家去，赞美天主。23 \BibleRef{玛 9:8} 群众看见，就害怕了；\BibleRef{路 5:26} 惊骇占据了他们 24，他们赞美天主，因为赐给了人这样的权柄。\BibleRef{路 5:26} 他们说：\Quote{今天我们看见了奇事，这样的事我们从未见过。}\par

25 [阿拉伯语第28页] \BibleRef{路 5:27} 此后，耶稣出去，看见一个名叫肋未的税吏，坐 26 在税吏中间，就对他说：\Quote{跟随我。}\BibleRef{路 5:28} 他便舍弃 27 一切，起来跟随了他。\BibleRef{路 5:29} 肋未在自己家中为他摆设了盛大的宴席。有许多税吏和其他人同他们一起坐席。28 \BibleRef{路 5:30} 经师和法利塞人私下议论，对他的门徒说：\Quote{你们为什么 29 同税吏和罪人一起吃喝？}\BibleRef{路 5:31} 耶稣回答他们说：\Quote{医生不是为健康的人，而是为身患 30、31 重病的人。\BibleRef{路 5:32} 我来不是召叫义人，而是召叫罪人悔改。}\BibleRef{路 5:33} 他们对他说：\Quote{为什么若翰的门徒常禁食祈祷，法利塞人 32 的门徒也这样，而你的门徒却吃喝呢？}\BibleRef{路 5:34} 耶稣对他们说：\Quote{你们不能使 33 婚礼之子在婚宴中禁食，因为新郎还同他们在一起。\BibleRef{路 5:35} 但日子将要来到，当新郎从他们中被带走时，他们将在那些 34 日子禁食。}\BibleRef{路 5:36} 他又对他们讲了一个比喻：\BibleRef{谷 2:21} 没有人把新布缝在旧衣服上，免得新布的补丁撕破旧衣服，造成更大 35 的破口。\BibleRef{谷 2:22} 也没有人把新酒装在旧皮囊里，免得酒胀破皮囊，皮囊损坏，酒也漏了；而是把 36 新酒装在新皮囊里，两者都得保全。\BibleRef{路 5:38} 没有人喝了陈酒，立刻想喝新酒，因为他说：\Quote{陈酒更好。}\par

37 \BibleRef{玛 12:1} 当耶稣在安息日走过麦田时，他的门徒\EditorialNote{阿拉伯文：第29页}饿了。他们用手搓麦穗，并且 38 吃了起来。\BibleRef{玛 12:2} 但有些法利塞人看见他们，\BibleRef{谷 2:24} 就对他说：\BibleQuote{看，39 为什么你的门徒在安息日作那不合规矩的事？}\BibleRef{谷 2:25} 但耶稣对他们说：\BibleQuote{你们没有读过，在古时达味和跟随他的人，在急需和 40 饥饿时作了什么？\BibleRef{谷 2:26}他进了天主的殿，亚比亚塔作大司祭时，吃了供桌上的饼，这饼除了司铎外任何人都不可吃，并且也给了跟随他的人吗？}41 \BibleRef{谷 2:27} 他又对他们说：\BibleQuote{安息日是为人设立的，人并不是 42 为安息日受造的。}\BibleRef{玛 12:5} \BibleQuote{或者你们没有在法律上读到，司铎在 43 圣殿内违犯安息日，\Emph{然而}他们仍是无罪的？}\BibleRef{玛 12:6} \BibleQuote{我告诉你们，44 这里有比圣殿更大的。}\BibleRef{玛 12:7} \BibleQuote{如果你们明白\Emph{这事}：我喜欢仁爱，45 而非祭献，你们就不会定那无罪之人的罪了。}\BibleRef{玛 12:8} \BibleQuote{46 安息日的主是人子。}\BibleRef{谷 3:21} 他的亲属听见了，就出来要拉住他，因为他们说：\Quote{他疯了。}\par

47 \BibleRef{路 6:6} 下一个安息日，他进了会堂教导。48 \BibleRef{路 6:7} 那里有一个人，右手枯干。经师和法利塞人窥探他，是否在安息日治病，49 好找借口控告他。\BibleRef{路 6:8} 但他知道他们的念头，就对那枯干手的人说：\BibleQuote{起来，站在会堂中间。}50 \BibleRef{路 6:9} 他起来站着。耶稣对他们说：\BibleQuote{我问你们，在安息日行善或行恶，救命或是 51 害命，哪样是合法的？}\EditorialNote{阿拉伯文：第30页} \BibleRef{谷 3:4} 他们却默不作声。\BibleRef{谷 3:5} 耶稣怒目环视他们，因他们心硬而忧愁。他对那人说：\BibleQuote{伸出手来。}他就一伸，手便复原了。52 \BibleRef{玛 12:11} 耶稣对他们说：\BibleQuote{你们中间谁有一只羊，若在安息日掉进坑里，不抓住拉上来呢？53 \BibleRef{玛 12:12} 人比羊贵重得多了！所以，在安息日行善是合法的。}\par

\SourceRecord{Translated by Hope W. Hogg.；Ante-Nicene Fathers；Vol. 9.；Edited by Allan Menzies.；Buffalo, NY: Christian Literature Publishing Co.,；1896.；<http://www.newadvent.org/fathers/100207.htm>.}

% structure: p0001 source=h1 target=section reason=page-title

\section{四福音合参（第八部分）}

\BibleRef{玛 12:14}法利塞人出去，一同商议怎样陷害耶稣。\BibleRef{玛 12:15}耶稣知道了，就离开那里。有许多人跟随他；他治好了他们众人，\BibleRef{玛 12:16}又嘱咐他们不要把他显露出来，\BibleRef{玛 12:17}这是要应验先知依撒意亚的话，说：\par

\BibleQuote{ \BibleRef{玛 12:18}看哪，我的仆人，我所喜悦的；\newline{}我所喜爱的，我心灵所喜爱的；\newline{}我已将我的神赐给他，\newline{}他必将公理传给外邦。\newline{}\BibleRef{玛 12:19}他不争竞，也不喧嚷；\newline{}街上也无人听见他的声音。\newline{}\BibleRef{玛 12:20}压伤的芦苇，他不折断；\newline{}将残的灯火，他不吹灭； }\par

\BibleQuote{ \BibleRef{玛 12:21}外邦人都要仰望他的名。 }\par

\BibleRef{路 6:12}在那些日子，耶稣出去上山祈祷，整夜祈祷天主。\BibleRef{路 6:13}到了天亮，他叫门徒来。\BibleRef{谷 3:7}耶稣就往海边去；有许多人从加里肋亚跟随他，\BibleRef{谷 3:8}并从犹太、耶路撒冷、厄东、约但河外、提洛、漆冬、十城区来；听见他所做的事，就有许多人来到他那里。\BibleRef{谷 3:9}耶稣因为人多，就吩咐门徒为他预备一条小船，免得众人拥挤他。\BibleRef{谷 3:10}他治好了许多人，以致他们几乎跌倒在他身上\EditorialNote{阿拉伯文版第31页}，因为他们想要接近他。那些有灾病和不洁之灵的，\BibleRef{谷 3:11}一看见他，就俯伏在地，喊叫说：\Quote{你是天主子}。\BibleRef{谷 3:12}耶稣严厉地嘱咐他们不要把他显露出来。\BibleRef{路 6:18}那些被不洁之灵束缚的也被治好了。\BibleRef{路 6:19}众人都想要靠近他，因为有能力从他身上出来，治好了他们所有人。\par

\BibleRef{玛 5:1}耶稣看见群众，就上了山。\BibleRef{路 6:13}他叫门徒来，从他们中拣选了十二人，称他们为宗徒：\BibleRef{路 6:14}西满（他称为\OriginalTerm{刻法}{Cephas}），和他的兄弟安德肋，雅各伯和若望，斐理伯和巴尔多禄茂，\BibleRef{路 6:15}玛窦和多默，阿尔斐的儿子雅各伯，\OriginalTerm{热忱者}{Zealot}西满，\BibleRef{路 6:16}雅各伯的儿子犹达，和\OriginalTerm{依斯加略人}{Iscariot}犹达斯（他就是出卖他的那个人）。\BibleRef{路 6:17}耶稣和他们下来，站在平地上，有许多门徒和一大群百姓。\BibleRef{谷 3:14}他拣选了这十二人，使他们常与他同在，并派遣他们去宣讲，并赐他们权柄医治疾病、驱逐魔鬼。\par

\BibleRef{路 6:20}于是他举目看着他们，\BibleRef{玛 5:2}开口教训他们说：\par

\BibleQuote{ \BibleRef{玛 5:3}神贫的人是有福的，因为天国是他们的。 }\par

\BibleQuote{ \BibleRef{玛 5:4}哀恸的人是有福的，因为他们要受安慰。 }\par

\BibleQuote{ \BibleRef{玛 5:5}温良的人是有福的，因为他们要承受土地。 }\par

\BibleQuote{ \BibleRef{玛 5:6}饥渴慕义的人是有福的，因为他们要得饱饫。 }\par

\BibleQuote{ \BibleRef{玛 5:7}怜悯人的人是有福的，因为他们要受怜悯。 }\par

\EditorialNote{阿拉伯文版第32页}\BibleQuote{ \BibleRef{玛 5:8}心里洁净的人是有福的，因为他们要看见天主。 }\par

\BibleQuote{ \BibleRef{玛 5:9}缔造和平的人是有福的，因为他们要称为天主的子女。 }\par

\BibleQuote{ \BibleRef{玛 5:10}为义而受迫害的人是有福的，因为天国是他们的。 }\par

\BibleQuote{ \BibleRef{路 6:22}几时，人恼恨你们，并弃绝你们，迫害你们，毁谤你们，\BibleRef{玛 5:11}并且为了我的缘故，以一切恶言虚假地反对你们，你们是有福的。\BibleRef{玛 5:12}那时，你们欢喜踊跃吧！因为你们在天上的赏报是丰厚的，因为以前的先知也这样受过迫害。 }\par

\BibleQuote{ \BibleRef{路 6:24}但祸哉，你们富足的人！因为你们已得了你们的安慰。 }\par

\BibleQuote{ \BibleRef{路 6:25}祸哉，你们饱足的人！因为你们将要饥饿。 }\par

\BibleQuote{祸哉，你们现今欢笑的人！因为你们将要哀恸哭泣。 }\par

\BibleQuote{ \BibleRef{路 6:26}祸哉，你们，当众人都称赞你们的时候！因为他们的祖先也同样对待了假先知。 }\par

\BibleQuote{ \BibleRef{路 6:27}我对你们说：你们这些\Emph{你们}听的，\BibleRef{玛 5:13}你们是地上的盐。盐若失了味，怎能再咸呢？它毫无用处，只好扔在外面，任人践踏。\BibleRef{玛 5:14}你们是世界的光。城建在山上，是不能隐藏的。\BibleRef{玛 5:15}没有人点灯放在斗底下，而是放在灯台上，就照亮全家的人。\BibleRef{玛 5:16}照样，你们的光也当在人前照耀，好使他们看见你们的善行，光荣你们在天之父。\BibleRef{谷 4:22}没有隐藏的事不被显露的，也没有秘密不被知道的。\BibleRef{谷 4:23}有耳可听的，就应当听。 }\par

\BibleQuote{ \BibleRef{玛 5:17}不要以为我来是废除法律或先知；我来不是为废除，而是为成全。\BibleRef{玛 5:18}我实在告诉你们：直到天地过去，\EditorialNote{阿拉伯文版第33页}一撇或一画也不能从法律中过去，直到一切都\Emph{完成}了。\BibleRef{玛 5:19}无论谁废除这些诫命中最小的一条，并这样教训人，他在天国里要称为最小的；但谁若实行且教训人，他在天国里要称为大的。\BibleRef{玛 5:20}我告诉你们：除非你们的义德超过经师和法利塞人的义德，你们绝不能进天国。 }\par

50 \BibleRef{玛 5:21} 你们听过对古人说：\Quote{不可杀人}；凡杀人的应受裁判。\BibleRef{玛 5:22} 但我告诉你们：凡向自己弟兄发怒而无故的，就要受裁判；谁若对弟兄说：\Quote{你这个坏蛋}，就要受议会的裁判；谁若说：\Quote{你这个傻瓜}，就要受地狱之火。\BibleRef{玛 5:23} 所以，若你在祭台前献礼时，想起你的弟兄对你有什么怨恨，\BibleRef{玛 5:24} 就把礼物留在祭台前，先去与你的弟兄和解，然后再来献礼。\BibleRef{玛 5:25} 趁你还在路上，赶快与你的仇敌和解，\BibleRef{路 12:58} 并付出赎金，从他手中解救自己；\BibleRef{玛 5:25} 免得你的仇敌把你交给判官，判官又交给税吏，\BibleRef{玛 5:26} 你便被投入狱中。我实在告诉你：除非你还清最后的一分钱，你绝不能从那里出来。\par

57, 58 \BibleRef{玛 5:27} 你们听过说：\Quote{不可奸淫}。\BibleRef{玛 5:28} 但我现在告诉你们：凡注视妇女而贪恋她的，已经在他的心里与她奸淫了。\BibleRef{玛 5:29} 如果你的右眼使你跌倒，就把它剜出来扔掉；因为丧失一个肢体，比你全身投入\Emph{地狱}之火更好。\BibleRef{玛 5:30} 如果你的右手使你跌倒，就把它砍下来扔掉；因为丧失一个肢体，比你全身投入地狱更好。\BibleRef{玛 5:31} 又说过：谁若休妻，\Emph{应该}给她休书。\BibleRef{玛 5:32} 但我告诉你们：凡休妻的，除非因奸淫的缘故，就是使她犯奸淫；谁若娶被休的妇人，就是犯奸淫。\par

\SourceRecord{Translated by Hope W. Hogg.；Ante-Nicene Fathers；Vol. 9.；Edited by Allan Menzies.；Buffalo, NY: Christian Literature Publishing Co.,；1896.；<http://www.newadvent.org/fathers/100208.htm>.}

% structure: p0001 source=h1 target=section reason=page-title

\section{\WorkTitle{四福音合参（第9节）}}

1 \BibleRef{玛 5:33}你们也听见对古人说：不可说谎，但在你们的誓言中要向天主履行；2 \BibleRef{玛 5:34}我却对你们说：总不可发誓；不可指天，因为它是天主的宝座；3 \BibleRef{玛 5:35}不可指地，因为它是他脚下的脚凳；也不可指耶路撒冷，因为它是大王的城；4 \BibleRef{玛 5:36}也不可指你的头起誓，因为你不能使一根头发变白或变黑。5 \BibleRef{玛 5:37}你们的话该是就是是，非就是非；多余的话，便是出于那恶者。\par

6,7 \BibleRef{玛 5:38}你们听见有话说：以眼还眼，以牙还牙；\BibleRef{玛 5:39}我却对你们说：不要抵抗恶人；但谁打你的右脸，你把另一面也转给他。8 \BibleRef{玛 5:40}那愿意告你、拿你内衣的，连外衣也给他。9 \BibleRef{玛 5:41}谁强迫你走一里，你就和他走二里。10 \BibleRef{玛 5:42}有求你的，就给他；\BibleRef{路 6:30}有向你借贷的，不可拒绝。那拿你财物的，不要追讨。11 \BibleRef{路 6:31}你们愿意人怎样待你们，你们也要怎样待人。[Arabic, p. 35]\par

12,13 \BibleRef{玛 5:43}你们听见有话说：爱你的邻人，恨你的仇敌。\BibleRef{玛 5:44}我却对你们说：爱你们的仇敌，为那咒骂你们的祈祷，善待那恨你们的，为那用暴力拿你们并迫害你们的祈祷；14 \BibleRef{玛 5:45}好成为你们天父的儿女，因为他使太阳上升，光照善人，也光照恶人，降雨给义人，也给不义的人。15 \BibleRef{玛 5:46}你们若爱那爱你们的人，有什么报酬呢？\BibleRef{路 6:32}税吏和罪人也爱那爱他们的人。16 \BibleRef{路 6:33}你们若善待那善待你们的人，有什么优胜呢？因为罪人也这样行。17 \BibleRef{路 6:34}你们若借钱给那你们希望得回报的人，有什么优胜呢？因为罪人也借钱给罪人，期望从他们得到偿还。18 \BibleRef{路 6:35}但是要爱你们的仇敌，善待他们，借钱给人，不要断绝任何人的希望；这样你们的报酬就大，你们要成为至高者的儿女；因为他对待忘恩和邪恶的人是宽容的。19 \BibleRef{路 6:36}你们要仁慈，如同你们的父是仁慈的一样。20 \BibleRef{玛 5:47}你们若只问你们的弟兄安，你们多做了什么，\Emph{比他人？}这不也是税吏的行为吗？21 \BibleRef{玛 5:48}你们应当是成全的，如同你们的天父是成全的一样。\par

22 \BibleRef{玛 6:1}你们要谨慎你们的施舍；不要在人前做，为让他们看见你们；不然，你们在天父前就没有报酬。23 \BibleRef{玛 6:2}所以，当你施舍时，不要在你前面吹号，像伪善者在会堂和街上所做的那样，为受人的称赞。我实在告诉你们，他们已获得了他们的报酬。24 \BibleRef{玛 6:3}但是，你施舍时，不要让左手知道右手所做的，\BibleRef{玛 6:4}好使你的施舍隐藏起来；你父在暗中看见，必要公开报答你。[Arabic, p. 36]\par

26 \BibleRef{玛 6:5}你们祈祷时，不要像伪善者那样；他们喜爱站在会堂和街角祈祷，为叫人看见。我实在告诉你们，他们已获得了他们的报酬。27 \BibleRef{玛 6:6}但是你祈祷时，要进入你的内室，关上门，向你在暗中的父祈祷；你父在暗中看见，必要公开报答你。28 \BibleRef{玛 6:7}你们祈祷时，不要唠唠叨叨地，如同外邦人一样；因为他们以为话多了，就会蒙垂听。29 \BibleRef{玛 6:8}你们不要像他们一样；因为你们的父在你们求他之前，已知道你们需要什么。30 \BibleRef{路 11:1}有一个门徒对他说：主，请教我们祈祷，如同若翰教了他的门徒一样。31 \BibleRef{路 11:2}耶稣对他们说：你们要这样祈祷：32 \BibleRef{玛 6:9}我们在天的父！愿你的名被尊为圣。33,34 \BibleRef{玛 6:10}愿你的国来临；愿你的旨意\Emph{承行}，在天上，也在地上。35 \BibleRef{玛 6:11}赐给我们今天所需的食粮。36 \BibleRef{玛 6:12}宽免我们的罪债，如同我们宽免亏负我们的人一样。37 \BibleRef{玛 6:13}不要让我们陷于诱惑，但救我们脱离那恶者。因为国度、权柄、荣耀，都是你的，直到永远。\BibleRef{玛 6:14}你们若宽免别人的过犯，你们的天父也必宽免你们。38 \BibleRef{玛 6:15}但你们若不宽免别人，你们的父也必不宽免你们的过犯。[Arabic, p. 37]\par

39 \BibleRef{玛 6:16}你们禁食时，不要像伪善者那样愁容满面；因为他们把自己的脸弄得难看，为叫人看出他们禁食。我实在告诉你们，他们已获得了他们的报酬。40 \BibleRef{玛 6:17}但你禁食时，要洗你的脸，抹你的头；\BibleRef{玛 6:18}好叫你不叫人看出你禁食，但只叫你在暗中的父看见；你父在暗中看见，必要报答你。\par

42 \BibleRef{路 12:32} \BibleQuote{小小羊群，不要害怕！因为你们的圣父喜欢把天国赐给你们。} 43 \BibleRef{路 12:33} \BibleQuote{变卖你们所有的，施舍给人；为自己准备永不旧的钱囊。} 44 \BibleRef{玛 6:19} \BibleQuote{不要在地上为自己积存财宝，那里有虫蛀、锈蚀，也有贼挖洞偷窃：} 45 \BibleRef{玛 6:20} \BibleQuote{但要在天堂为自己积财宝，那里没有虫蛀、锈蚀，也没有贼挖洞偷窃：} 46 \BibleRef{玛 6:21} \BibleQuote{因为你们的财宝在哪里，你们的心也在那里。} 47 \BibleRef{玛 6:22} \BibleQuote{身体的灯就是眼睛。所以，如果你的眼睛现在健全，全身就都光明。} 48 \BibleRef{玛 6:23} \BibleQuote{但如果你的眼睛有了病，全身就都黑暗。那么，你身上的光明如果成了黑暗，那黑暗是多么可怕！} 49 \BibleRef{路 11:35} \BibleQuote{要留神，免得你身上的光成了黑暗。} 50 \BibleRef{路 11:36} \BibleQuote{因为如果你全身光明，没有一点黑暗，它就完全光明，如同灯用它的火焰照耀你。}\par

\SourceRecord{Translated by Hope W. Hogg.；Ante-Nicene Fathers；Vol. 9.；Edited by Allan Menzies.；Buffalo, NY: Christian Literature Publishing Co.,；1896.；<http://www.newadvent.org/fathers/100209.htm>.}

% structure: p0001 source=h1 target=section reason=page-title

\section{福音合参（第10节）}

1 [Arabic, p. 38] \BibleRef{玛 6:24} \BibleQuote{没有人能事奉两个主人；因为必须恨一个而爱另一个，尊敬一个而轻视那2另一个。} \BibleRef{玛 6:25} \BibleQuote{你们不能事奉天主而又事奉钱财。因此我告诉你们：不要为你们的生命忧虑吃什么、喝什么；也不要为你们的身体忧虑穿什么。生命不是贵于食物，身体3不是贵于衣服吗？} \BibleRef{玛 6:26} \BibleQuote{观察天空的飞鸟，它们不播种，也不收割，也不收藏在仓里，而你们的天父\Emph{仍然}养活它们。你们不比它们更贵重吗？} \BibleRef{玛 6:27} \BibleQuote{你们中谁能用思虑使自己的身量加一4肘呢？} \BibleRef{路 12:26} \BibleQuote{如果你们连极小的事\Emph{都不能}，为什么还为5,6其余的事忧虑呢？} \BibleRef{玛 6:28} \BibleQuote{观察田野的百合花，它们怎样生长；它们不劳作，也不纺织；} \BibleRef{玛 6:29} \BibleQuote{我告诉你们：撒罗满在他极荣耀的时候，所穿戴的也不如这花中的一朵。} \BibleRef{玛 6:30} \BibleQuote{如果天主这样装饰田野的草，它今天还在，明天就被扔进炉里，何况你们呢，小信德的人哪！} \BibleRef{玛 6:31} \BibleQuote{所以不要忧虑说：我们吃什么？或喝什么？或穿什么？} \BibleRef{路 12:29} \BibleQuote{也不要让你们的内心为这些事焦虑；} \BibleRef{玛 6:32} \BibleQuote{这一切都是世上的外邦人所寻求的；而你们的天父知道你们需要这一切。} \BibleRef{玛 6:33} \BibleQuote{你们先寻求天主的国和他的义德，} 12 [Arabic, p. 39] \BibleQuote{这一切都要加给你们。} \BibleRef{玛 6:34} \BibleQuote{所以不要为明天忧虑，因为明天自有明天的忧虑；一天的苦足够一天受的。}\par

13 \BibleRef{玛 7:1} \BibleQuote{你们不要判断人，免得你们受判断；} \BibleRef{路 6:37} \BibleQuote{不要定罪，免得你们被定罪；} \BibleRef{路 6:38} \BibleQuote{赦免，\Emph{也}必被赦免；释放，你们必被释放；给予，你们必\Emph{得到}；用好的、满的、堆得多的尺度，他们将倒在你们怀里。} \BibleRef{谷 4:24} \BibleQuote{你们用什么尺度量，也要用什么尺度量给你们。注意\Emph{你们所听的}：你们用什么尺度量，也要用什么尺度量给你们，并且还要加给你们。} \BibleRef{谷 4:25} \BibleQuote{我告诉那些听的人：凡有的，还要给他；凡没有的，连他自以为有的，也要从他夺去。}\par

17 \BibleRef{路 6:39} \BibleQuote{他又对他们讲了一个比喻：瞎子岂能领瞎子？两人岂不都要掉进坑里吗？} \BibleRef{路 6:40} \BibleQuote{门徒不能胜过他的师傅；凡学成的，将要像他的师傅一样。} \BibleRef{路 6:41} \BibleQuote{为什么你看见你弟兄眼中的木屑，却不理会你\Emph{自己}眼中的大梁呢？} \BibleRef{路 6:42} \BibleQuote{或者，你怎能对你的弟兄说：弟兄，让我取出你眼中的木屑；而你看不见自己眼中的大梁呢？伪君子！先取出你眼中的大梁，然后你才看得清楚，取出你弟兄眼中的木屑。}\par

21 \BibleRef{玛 7:6} \BibleQuote{不要把圣物给狗，也不要把你们的珍珠丢在猪前，免得它们用脚践踏，转过来咬伤你们。}\par

22 \BibleRef{路 11:5} \BibleQuote{他又对他们说：你们中间谁有一个朋友，半夜到他那里去，对他说：朋友，借给我三个饼；} \BibleRef{路 11:6} \BibleQuote{因为我的朋友从路上来到我这里，我没有什么可以摆上给他；} \BibleRef{路 11:7} [Arabic, p. 40] \BibleQuote{那朋友从里面回答说：不要搅扰我，门已经关上，孩子们也同我在床上，我不能起来给你吗？} \BibleRef{路 11:8} \BibleQuote{我告诉你们：他虽不因是朋友而起来给他，但因\Emph{他的}切求，必要起来，把他所需要的给他。} \BibleRef{路 11:9} \BibleQuote{我也告诉你们：你们求，就必\Emph{得到}；寻找，就必找到；敲门，就必给你们开门。} \BibleRef{路 11:10} \BibleQuote{因为凡求的，就得着；寻找的，就找到；敲门的，就给他开门。} \BibleRef{路 11:11} \BibleQuote{你们中间哪一个为父的，他的儿子求饼，他反而给他石头呢？或是求鱼，反将蛇当鱼给他呢？} \BibleRef{路 11:12} \BibleQuote{或是求鸡蛋，反给他蝎子呢？} \BibleRef{路 11:13} \BibleQuote{你们虽然不好，尚且知道把好东西给你们的儿女，何况你们在天上的父，岂不更将圣神赐给求他的人吗？} \BibleRef{玛 7:12} \BibleQuote{所以，凡你们愿意人给你们做的，你们也要照样给人做：这就是法律和先知。}\par

32 \BibleRef{玛 7:13} \BibleQuote{你们要进窄门；因为宽门和大路通向丧亡，进去的人也多。} \BibleRef{玛 7:14} \BibleQuote{那门多么窄，路多么狭，以致通向生命！找到它的人何其少。}\par

34 \BibleRef{玛 7:15} 你们要提防假先知，他们来到你们这里，外面披着羊皮，里面却是35残暴的狼。\BibleRef{玛 7:16} 但凭他们的果实你们会认出他们。\BibleRef{路 6:44} 因为每一棵树凭自己的果实就能被认出。无花果不从荆棘上摘取，葡萄也不从36蒺藜上采摘。\BibleRef{玛 7:17} 同样，凡是好树都结好果子，但坏树却结37 [Arabic, p. 41] 坏果子。\BibleRef{玛 7:18} 好树不能结坏果子，坏树38也\Emph{不能}结好果子。\BibleRef{路 6:45} 善人从他心里所存的善，就发出\Emph{善}来；恶人从他心里所存的恶，就发出\Emph{恶}来；因为从39心里所充满的，口里就说出来。\BibleRef{玛 7:19} 凡不结好果子的树就砍\Emph{下}来，丢在40, 41火里。\BibleRef{玛 7:20} 所以，凭着他们的果实你们会认出他们。\BibleRef{玛 7:21} 不是每一个对我说：主啊，主啊的人都能进入天国，而是实行42我在天之父旨意的人。\BibleRef{玛 7:22} 到那一天，许多人要对我说：主啊，主啊，我们不是因你的名字预言过，因你的名字驱过43魔鬼，因你的名字行过许多异能吗？\BibleRef{玛 7:23} 那时我必向他们声明：我从来不44认识你们；你们这些作恶的人，离开我吧！\BibleRef{路 6:47} 凡到45我这里来，听了我的话而实行的，我要指给你们看他像什么人：\BibleRef{路 6:48} 他像一个人建造房屋，挖深，把46基础安在磐石上；\BibleRef{玛 7:25} 雨淋，水冲，风吹，冲击那房屋，它却不倒塌，因为它的基础是建在47磐石上。\BibleRef{玛 7:26} 凡听了这些话而不实行的，就像48一个愚昧人把房屋建在沙土上，没有基础；\BibleRef{玛 7:27} 雨淋，水冲，风吹，冲击那房屋，它就倒塌了，并且倒塌得很厉害。\par

\SourceRecord{Translated by Hope W. Hogg.；Ante-Nicene Fathers；Vol. 9.；Edited by Allan Menzies.；Buffalo, NY: Christian Literature Publishing Co.,；1896.；<http://www.newadvent.org/fathers/100210.htm>.}

% structure: p0001 source=h1 target=section reason=page-title

\section{《四福音合参》（第十一章）}

1 \BibleRef{玛 7:28} 耶稣讲完了这些话，众人都希奇他的教训；\BibleRef{玛 7:29} 因为他教训他们，正像有权柄的人，不像他们的经师和法利塞人。\par

3 \BibleRef{玛 8:1} 耶稣从山上下来，有许多群众跟着他。\par

4 耶稣进了葛法翁，有一个百夫长的仆人患病沉重，快要死了，这仆人是他所器重的。\BibleRef{路 7:3} 他听见耶稣的事，就来见耶稣，同来的还有犹太人的几个长老；他恳求耶稣说：\Quote{主啊，我的仆人瘫痪了，躺在家里，非常痛苦。}\BibleRef{路 7:4} 长老们恳切地求耶稣说：\Quote{你给他行这事，是他配得的；\BibleRef{路 7:5} 因为他爱护我们的民族，并且给我们建造了会堂。}\BibleRef{玛 8:7} 耶稣对他说：\Quote{我去治好他。}\BibleRef{玛 8:8} 百夫长回答说：\Quote{主啊，你到我舍下，我不敢当；只要你说一句话，我的仆人就必好了。\BibleRef{路 7:8} 因为我自己也是服从上级的人，也有兵在我以下；我对这个说：\InnerQuote{你去！}他就去；对那个说：\InnerQuote{你来！}他就来；对我的仆人说：\InnerQuote{你做这事！}他就做。}\BibleRef{路 7:9} 耶稣听见这话，就希奇他，转身对跟随的群众说：\BibleRef{玛 8:10} \Quote{我实在告诉你们：在以色列中，我没有见过这么大的信德。\BibleRef{玛 8:11} 我告诉你们：将有许多人从东方和西方来，与亚巴郎、依撒格、雅各伯一同坐席在天国里；\BibleRef{玛 8:12} 但本国的子民却被驱逐到外面的黑暗中，在那里要有哀号和切齿。}\BibleRef{玛 8:13} 耶稣对百夫长说：\Quote{你回去吧！照你所信的，给你成就。}\BibleRef{路 7:10} 就在那一刻，他的仆人好了。百夫长回到家里，看见那个患病的仆人已经痊愈了。\par

17 \BibleRef{路 7:11} 第二天，耶稣往一座名叫纳因的城去，他的门徒和一大群人跟着他。\BibleRef{路 7:12} 将近城门时，看，正抬出一个死人来，他是母亲独生的儿子；母亲又是寡妇；城中有许多人陪着她。\BibleRef{路 7:13} 主看见那寡妇，就对她动了怜悯的心，向她说：\Quote{不要哭！}\BibleRef{路 7:14} 于是上前按住棺材，抬棺材的人就站住了。他说：\Quote{青年人，我吩咐你，起来！}\BibleRef{路 7:15} 那死人就坐起来，并开口说话；耶稣便把他交给他的母亲。\BibleRef{路 7:16} 众人都害怕起来，光荣天主说：\Quote{在我们中间兴起了一位大先知！}又说：\Quote{天主眷顾了他的子民！}\BibleRef{路 7:17} 于是，关于耶稣的这番话，传遍了犹太和周围各地。\par

24 \BibleRef{玛 8:18} 耶稣看见许多群众围着自己，就吩咐往对岸去。他们正走的时候，有一个经师来对他说：\Quote{老师，你不论往哪里去，我要跟随你。}\BibleRef{玛 8:20} 耶稣给他说：\Quote{狐狸有穴，天上的飞鸟有巢，但是人子却没有枕头的地方。}\BibleRef{路 9:59} 又对另一个人说：\Quote{你跟随我吧！}那人说：\Quote{主，请许我先去埋葬我的父亲。}\BibleRef{路 9:60} 耶稣给他说：\Quote{任凭死人去埋葬自己的死人吧！至于你，你要去宣扬天主的国。}\BibleRef{路 9:61} 又有一个人说：\Quote{主！我要跟随你；但是请许我先告别我的家人。}\BibleRef{路 9:62} 耶稣对他说：\Quote{手扶着犁而往后看的，不适于天主的国。}\par

31 那天傍晚，耶稣对他们说：\Quote{我们渡到湖对岸去吧！}他就离开了群众。耶稣上了船，坐在船上，他的门徒跟着他，还有别的船同行。\BibleRef{路 8:23} 湖上忽然起了暴风，波浪打入船内，甚至船已满了水，快要沉没。\BibleRef{谷 4:38} 耶稣却在船尾，靠着枕头睡觉。\BibleRef{玛 8:25} 门徒前来叫醒他说：\Quote{主啊！救我们，我们丧命啦！}\BibleRef{路 8:24} 耶稣起来，斥责了风和浪，对海说：\Quote{不要作声，平静了吧！}\BibleRef{谷 4:39} 风就停了，大大地平静了。\BibleRef{谷 4:40} 耶稣对他们说：\Quote{为什么这样胆怯？你们怎么没有信德呢？}\BibleRef{路 8:25} 他们就非常惧怕，彼此说：\Quote{这人到底是谁？他竟命令风和浪，连海也听从了他！}\par

38 \BibleRef{路 8:26} 他们航行到加达辣人的地区，就是加里肋亚的对面。39 \BibleRef{路 8:27} 当他从船上出来到陆地，有一个从坟墓出来的人遇见他，那人长久被魔鬼附着，不穿衣服，不住房子，只住在坟墓中。\BibleRef{谷 5:3} 没有人能用链条捆住他，\BibleRef{谷 5:4} 因为每当被铁链和脚镣捆住，他就挣断铁链，弄碎脚镣；\BibleRef{路 8:29} 他被魔鬼抓到旷野，没有人能制伏他；昼夜常在坟墓和山中，\BibleRef{玛 8:28} 没有人能从那条路经过；\BibleRef{谷 5:5} 他喊叫，用石头伤害自己。\BibleRef{谷 5:6} 当他远远看见\Person{耶稣}，就跑来朝拜他，大声喊叫说：\Quote{至高天主的儿子\Person{耶稣}，我与你有什么相干？我因天主恳求你，不要折磨我。}\Person{耶稣}命令不洁之神从那人身上出来：那人自被魔鬼占据以来，已经\Emph{受折磨}了很长时间。\BibleRef{路 8:30} \Person{耶稣}问他：\Quote{你叫什么名字？}他回答说：\Quote{\OriginalTerm{军旅}{Legion}；}因为有许多魔鬼进入他里面。\BibleRef{路 8:31} 他们恳求\Person{耶稣}不要命令他们到深渊里去。\BibleRef{路 8:32} 那里有一大群猪在山上放牧，那些魔鬼恳求\Person{耶稣}准许他们进入猪群；\Person{耶稣}准许了他们。\BibleRef{路 8:33} 魔鬼就从那人身上出来，进入猪群。\BibleRef{谷 5:13} 那群猪冲向悬崖，掉进海里，约有二千，都淹死在水里。\BibleRef{路 8:34} 放猪的人看见所发生的事，就逃跑，在城里和乡下传报。\BibleRef{路 8:35} 众人出来要看所发生的事；他们来到\Person{耶稣}跟前，看见那个魔鬼已离开的人，穿着衣服，神志清醒，坐在\Person{耶稣}脚前，就害怕了。\BibleRef{路 8:36} 他们报告了所看见的事，以及那被鬼附的人如何得了痊愈，\BibleRef{谷 5:16} 也报告了关于猪的事。\par

\SourceRecord{Translated by Hope W. Hogg.；Ante-Nicene Fathers；Vol. 9.；Edited by Allan Menzies.；Buffalo, NY: Christian Literature Publishing Co.,；1896.；<http://www.newadvent.org/fathers/100211.htm>.}

% structure: p0001 source=h1 target=section reason=page-title

\section{\WorkTitle{戴亚泰撒隆（第十二段）}}

1 \BibleRef{路 8:37} 格拉森一带的众人全都恳求祂离开他们，因为极大的恐惧抓住了他们。\par

2, 3 \BibleRef{玛 9:1} 耶稣上了船，渡过去，来到自己的城。\BibleRef{路 8:38} 那鬼已出去的人恳求与祂同在；4 但耶稣打发他走，对他说：\BibleRef{路 8:39} \Quote{你回家去，宣扬天主为你做了何等大事。}5 \BibleRef{谷 5:20} 他就去了，在\Place{低加波利}开始传扬耶稣为他所做的事；众人都惊讶不已。\par

6 耶稣乘船渡到对岸时，一大群人迎接祂；他们都等候祂。7 \BibleRef{路 8:41} 有一个人名叫\Person{雅依洛}，是会堂主管，俯伏在耶稣脚前，8 \BibleRef{谷 5:23} 再三恳求祂，说：\Quote{我有一个独生女儿，快要死了；\BibleRef{玛 9:18} 求祢来按手在她身上，她就必活了。}9 \BibleRef{玛 9:19} 耶稣就起身，祂的门徒也跟着祂。10 \BibleRef{谷 5:24} 有一大群人跟随祂，拥挤着祂。\par

11, 12 \BibleRef{谷 5:25} 有一个妇人，患血漏十二年，\BibleRef{谷 5:26} 在多位医生手里受了许多苦，花尽了她所有的，不但毫无益处，13 反而更重了。\BibleRef{谷 5:27} 她听见耶稣的事，就从人群中从后面来，摸了祂的衣裳；14 \EditorialNote{阿拉伯语原稿第47页} \BibleRef{谷 5:28} 她心里想：\Quote{我只要摸到祂的衣裳，就必痊愈。}15 \BibleRef{谷 5:29} 她的血源立刻干了；她在身体上觉着已经治好了这灾病。16 \BibleRef{谷 5:30} 耶稣立时心里觉得有能力从自己身上出去；就转过身来对众人说：\Quote{谁摸了我的衣裳？}17 \BibleRef{路 8:45} 众人都否认，\Person{西满·则法}和同他在一起的人对祂说：\Quote{老师，群众拥挤着祢，祢还说\InnerQuote{谁摸了我}？}18 \BibleRef{路 8:46} 祂说：\Quote{有人摸了我，因为我知道有能力从我身上出去了。}19 \BibleRef{路 8:47} 那妇人见自己藏不住，就战战兢兢地前来（因为她知道在自己身上所发生的事），俯伏在祂面前，当着众百姓，把她摸祂的缘故，20 以及怎样立刻痊愈，都述说出来。21 耶稣对她说：\Quote{女儿，放心吧！你的信德救了你；平平安安地去吧，并愿你从这灾病中痊愈。}\par

22 \BibleRef{路 8:49} 祂还说话的时候，有人从会堂主管家里来，对他说：\Quote{你的女儿死了，不要劳烦老师了。}23 \BibleRef{路 8:50} 耶稣听见了，就对那女孩的父亲说：\Quote{不要怕；只要信，她就必得救。}24 \BibleRef{谷 5:37} 祂不许任何人跟着祂，除了\Person{西满·则法}、\Person{雅各伯}和\Person{雅各伯}的兄弟\Person{若望}。25 \BibleRef{谷 5:38} 他们来到会堂主管的家，看见他们慌乱，哭泣哀号。26 \BibleRef{谷 5:39} 耶稣进去，对他们说：\Quote{为什么慌乱哭泣呢？27 \EditorialNote{阿拉伯语原稿第48页} 女孩子没有死，是睡着了。}\BibleRef{路 8:53} 他们讥笑祂，因为他们知道她已经死了。28 \BibleRef{谷 5:40} 祂把众人都赶出去，带着女孩的父母和\Person{西满}、\Person{雅各伯}、\Person{若望}，29 进了女孩所在的地方。\BibleRef{谷 5:41} 祂拉着女孩的手，对她说：\Quote{女孩，起来。}\BibleRef{路 8:55} 她的灵魂回来了，立刻站起来行走；30 \BibleRef{谷 5:42} 她约十二岁。\BibleRef{路 8:55} 耶稣吩咐给她吃东西。31 \BibleRef{路 8:56} 她的父母极其惊奇；耶稣嘱咐他们，不要告诉别人所发生的事。32 \BibleRef{玛 9:26} 这风声传遍了那一带地方。\par

33 \BibleRef{玛 9:27} 耶稣从那里继续前行，有两个瞎子跟着祂，喊着说：\Quote{\Emph{祢}\Person{达味}之子，可怜我们吧！}34 \BibleRef{玛 9:28} 耶稣进了屋子，那两个瞎子来到祂跟前；耶稣对他们说：\Quote{你们信我能做这事吗？}他们对祂说：\Quote{是，主。}35 \BibleRef{玛 9:29} 耶稣就摸他们的眼睛，说：\Quote{照你们的信德，给你们成就吧。}36 \BibleRef{玛 9:30} 他们的眼睛立刻开了。耶稣严严地嘱咐他们说：\Quote{你们要当心，不可让人知道。}37 \BibleRef{玛 9:31} 但他们出去，把祂的名声传遍了那一带地方。\par

38 \BibleRef{玛 9:32} 耶稣出去的时候，有人带着一个被鬼附着的哑巴来见祂。39 \BibleRef{玛 9:33} 鬼被赶出去，哑巴就说出话来。群众都惊奇说：\Quote{在以色列从来没有见过这样的事。}\par

40 \BibleRef{玛 9:35} 耶稣周游各城和\Emph{在}各乡村，在会堂里教导，宣讲天国的福音，医治各种疾病41 [Arabic, p. 49] 和各种疼痛；有许多人跟随他。\BibleRef{玛 9:36} 耶稣看见群众，就怜悯他们，因为他们困苦流离，像羊42没有牧人一样。他又叫齐十二门徒，赐给他们权柄和能力，43制伏一切魔鬼和疾病；\BibleRef{路 9:2} 并派遣他们两个两个出去，使他们44传讲天主的国，并\Emph{去}医治病人。\BibleRef{玛 10:5} 耶稣嘱咐他们说：外邦人的路你们不要走，撒玛黎雅人的城你们不要进。45,46 \BibleRef{玛 10:6} 你们宁往以色列家迷失的羊那里去。\BibleRef{玛 10:7} 你们去，要宣讲说：天国临近了。\BibleRef{玛 10:8} 医治病人，洁净癞病人，驱逐魔鬼：你们白白地获得，白白地48,49给予。\BibleRef{玛 10:9} 腰袋里不要带金银铜钱；在路上不要带任何东西，除了手杖以外；不要带口袋，不要带食物，不要有两件内衣，50不要穿鞋，也不要带手杖，但要穿凉鞋；因为工人应当得食物。51 \BibleRef{玛 10:10} 你们无论进哪一城哪一村，要打听那里谁是当得起的人，就住在那里，直到52,53你们离去。\BibleRef{玛 10:11} 你们进那家时，要向那家请安；那家若当得起，你们的平安就临到那家；若当不起，你们的54平安仍归你们。凡不接待你们，不听你们话的人，你们离开那家或那村时，要把脚上的尘土55 [Arabic, p. 50] 跺下去，作为反对他们的证据。\BibleRef{玛 10:15} 我实在告诉你们：在审判的日子，索多玛和蛾摩拉地方所受的，比那城还容易受。\par

\SourceRecord{Translated by Hope W. Hogg.；Ante-Nicene Fathers；Vol. 9.；Edited by Allan Menzies.；Buffalo, NY: Christian Literature Publishing Co.,；1896.；<http://www.newadvent.org/fathers/100212.htm>.}

% structure: p0001 source=h1 target=section reason=page-title

\section{《四福音合参》（第十三篇）}

1 \BibleRef{玛 10:16} \BibleQuote{我派遣你们如同羔羊进入狼群：你们要机警如同蛇，纯朴如同鸽子。} 2 \BibleRef{玛 10:17} \BibleQuote{你们要提防世人，因为他们要把你们交给议会，并在他们的会堂里鞭打你们；} 3 \BibleRef{玛 10:18} \BibleQuote{并且要为我的缘故，被带到总督和君王面前，对他们和外邦人作证据。} 4 \BibleRef{玛 10:19} \BibleQuote{当他们把你们交出时，你们不要忧虑，也不要预先考虑说什么；但在那时刻，必赐给你们所当说的话。} 5 \BibleRef{玛 10:20} \BibleQuote{因为不是你们说话，而是你们父的圣神在你们内说话。} 6 \BibleRef{玛 10:21} \BibleQuote{弟兄要把弟兄，父亲要把儿子，交到死地；儿女要起来反对父母，并把他们处死。} 7 \BibleRef{玛 10:22} \BibleQuote{你们要为我的名字被众人恼恨；唯独坚持到底的，才可得救。} 8 \BibleRef{玛 10:23} \BibleQuote{当这城迫害你们时，你们就逃到另一城去。我实在告诉你们：你们还未走遍以色列的城邑，人子就来了。}\par

9,10 \BibleRef{玛 10:24} \BibleQuote{门徒不能高过老师，仆人也无法高过主人。} \BibleRef{玛 10:25} \BibleQuote{门徒能如老师，仆人能如主人，也就够了。如果人们称家主为贝耳则步，何况他的家人呢？} 11 \BibleRef{玛 10:26} \BibleQuote{所以你们不要害怕他们；因为没有什么掩盖的，将来不被揭露；也没有什么隐藏的，将来不被知道和宣扬。} 12 \BibleQuote{我在黑暗中对你们所说的，你们要在光中说出来；你们私下在耳房中所听到的，要在屋顶上宣扬出去。} 13 \BibleQuote{我现在告诉你们，我亲爱的，不要害怕那些杀害肉身，而不能杀害灵魂的；我要告诉你们应害怕谁：害怕那能使灵魂和肉身都陷入地狱中的。} 14 \BibleQuote{是的，我告诉你们，你们应当害怕这一位。} 15 \BibleRef{玛 10:29} \BibleQuote{两只麻雀不是卖一个铜钱吗？但若没有你们父的许可，一只也不会掉在地上。} 16 \BibleRef{玛 10:30} \BibleQuote{至于你们，连你们的头发也都数过了。} 17,18 \BibleRef{玛 10:31} \BibleQuote{所以不要害怕；你们比许多麻雀还贵重。} \BibleRef{玛 10:32} \BibleQuote{凡在人面前承认我的，我在我天上的父面前也必承认他；} 19 \BibleRef{玛 10:33} \BibleQuote{但谁在人面前否认我，我在我天上的父面前也必否认他。}\par

20 \BibleRef{路 12:51} \BibleQuote{你们以为我来是给地上送和平吗？不，我告诉你们，而是送分裂。} 21 \BibleRef{路 12:52} \BibleQuote{从此以后，一家五口将要分裂，三个反对两个，两个反对三个。} 22 \BibleRef{路 12:53} \BibleQuote{父亲反对儿子，儿子反对父亲；母亲反对女儿，女儿反对母亲；婆婆反对儿媳，儿媳反对婆婆。} 23 \BibleRef{玛 10:36} \BibleQuote{人的仇敌就是自己家里的人。} 24 \BibleRef{玛 10:37} \BibleQuote{谁爱父亲或母亲超过我，不配是我的；谁爱儿子或女儿超过对我的爱，不配是我的。} 25 \BibleRef{玛 10:38} \BibleQuote{谁不背起自己的十字架跟随我，不配是我的。} 26 \BibleRef{玛 10:39} \BibleQuote{谁找到自己的性命，必要丧失性命；谁为我的缘故丧失自己的性命，必要找到性命。}\par

27 \BibleRef{玛 10:40} \BibleQuote{谁接纳你们，就是接纳我；谁接纳我，就是接纳那派遣我的。} 28 \BibleRef{玛 10:41} \BibleQuote{谁因先知的名义接纳先知，必得先知的赏报；谁因义人的名义接纳义人，必得义人的赏报。} 29 \BibleQuote{谁若只给这些小子中的一个一杯凉水喝，因他是门徒，我实在告诉你们，他决失不了他的赏报。}\par

30 \BibleRef{玛 11:1} \BibleQuote{耶稣嘱咐完他的十二门徒，就从那里离开，到他们的城里去施教和宣讲。} 31 \BibleRef{路 10:38} \BibleQuote{他们走路的时候，耶稣进了一个村庄。有一个名叫玛尔大的女人，接待他到自己家中。} 32 \BibleRef{路 10:39} \BibleQuote{她有一个妹妹，名叫玛利亚，坐在主的脚前听他讲话。} 33 \BibleRef{路 10:40} \BibleQuote{玛尔大为伺候多事，心神烦乱，就上前来说：主！我的妹妹丢下我一个人伺候，你不介意吗？请叫她来帮助我。} 34 \BibleRef{路 10:41} \BibleQuote{主回答她说：玛尔大，玛尔大！你为许多事操心忙碌，} 35 \BibleRef{路 10:42} \BibleQuote{其实需要的惟有一件。玛利亚选择了更好的一份，是不能从她夺去的。}\par

36 \BibleRef{谷 6:12} \BibleQuote{宗徒们就出去宣讲，叫人悔改，} 37 \BibleRef{谷 6:13} \BibleQuote{并驱逐了许多魔鬼，给许多病人傅油，治好了他们。} 38,39 \BibleRef{路 7:18} \BibleQuote{若翰的门徒把这一切报告给若翰。} \BibleRef{路 7:19} \BibleQuote{若翰在狱中听到默西亚的作为，便叫了自己的两个门徒来，打发他们到主那里去说：你就是要来的那位，或者我们还要等候另一位？} 40 \BibleRef{路 7:20} \BibleQuote{他们来到耶稣跟前说：洗者若翰打发我们来问你：你就是要来的那位，或者我们还要等候另一位？} 41 \BibleRef{路 7:21} \BibleQuote{在那时刻，耶稣治好了许多患有疾病和恶灵折磨的人，又赐许多盲人看见。} 42 \BibleRef{路 7:22} \BibleQuote{耶稣回答他们说：你们去，把你们所见所闻的报告给若翰：盲人看见，瘸子行走，癞病人洁净，瞎子听见，死人复活，穷人得闻福音。} 43 \BibleRef{路 7:23} \BibleQuote{凡不因我而跌倒的，是有福的。}\par

44 若翰的门徒走了以后，耶稣就对群众说起若翰来：\BibleRef{路 7:24}\BibleQuote{你们出去到荒野里是为看什么呢？为看一枝被风摇曳的芦苇吗？}\BibleRef{路 7:25}\BibleQuote{你们出去到底是为看什么？为看一位穿细软衣服的人吗？看哪！那穿华丽衣服、生活奢侈的人是在王宫里。}\BibleRef{路 7:26}\BibleQuote{你们出去究竟是为看什么？为看一位先知吗？是的，我告诉你们，而且他比先知还大。}\BibleRef{路 7:27}\BibleQuote{就是关于他所记载的：}\par

\BibleQuote{看，我派遣我的使者在你面前，\newline{}他要在你前面预备你的道路。}\par

\SourceRecord{Translated by Hope W. Hogg.；Ante-Nicene Fathers；Vol. 9.；Edited by Allan Menzies.；Buffalo, NY: Christian Literature Publishing Co.,；1896.；<http://www.newadvent.org/fathers/100213.htm>.}

% structure: p0001 source=h1 target=section reason=page-title

\section{\WorkTitle{四福音合参（第十四部分）}}

1 我实在告诉你们，凡妇人所生的，没有一个大过\Person{若翰洗者}的；但在天国里最小的，也比他大。\BibleRef{玛 11:11}\par

2 \EditorialNote{阿拉伯语本第54页}\BibleRef{路 7:29}所有听见的百姓和税吏，都承认天主是公义的，因为他们曾受过若翰的洗礼。3 \BibleRef{路 7:30}但法利塞人和法学士废弃了天主在他们身上的旨意，因为他们没有受他的洗。4 \BibleRef{玛 11:12}从若翰洗者的时候直到如今，天国是以暴力夺取的。5 \BibleRef{路 16:16}法律和先知直到若翰为止；从此天主的福音被传扬，人人都奋力进入。6 \BibleRef{玛 11:12}他们努力夺取它。7 \BibleRef{玛 11:13}众先知和法律直到若翰都预言了。8 \BibleRef{玛 11:14}如果你们愿意接受，他就是那将要来的\Person{厄里亚}。9 \BibleRef{玛 11:15}有耳可听的，就应当听。10 \BibleRef{路 16:17}天地废去，比律法的一点一画落空还容易。11 \BibleRef{路 7:31}这样，我拿这世代的人比作什么呢？他们好像什么呢？12 \BibleRef{路 7:32}他们好像孩童坐在街市上，彼此呼叫说：\Quote{我们向你们吹笛，你们不跳舞；我们向你们举哀，你们不哭泣。}13 \BibleRef{路 7:33}若翰洗者来了，不吃不喝，你们说他有鬼；14 \BibleRef{路 7:34}人子来了，也吃也喝，你们说他是贪食好酒的人，是税吏和罪人的朋友。15 \BibleRef{路 7:35}但智慧之子总显智慧为是。他说这话的时候，他们来到屋里。又有多人聚集，甚至他们连饭也顾不上吃。16 \BibleRef{路 11:14}耶稣赶出一个叫人哑巴的鬼；鬼出去了，哑巴就说出话来。\EditorialNote{阿拉伯语本第55页}众人都希奇。17 \BibleRef{玛 12:24}法利塞人听见，就说：\Quote{这个人赶鬼，无非是靠着鬼王\OriginalTerm{贝耳则布}{Beelzebul}。}18 \BibleRef{路 11:16}又有人试探耶稣，向他求从天上来的神迹。19 \BibleRef{玛 12:25}耶稣知道他们的意念，就对他们说：\Quote{凡一国自相纷争，就成为荒场；凡一家或一城自相纷争，必站立不住。}20 \BibleRef{玛 12:26}若撒但赶逐撒但，就是自相纷争，他的国怎能站得住呢？他的结局将到。21 \BibleRef{玛 12:26}那时他的国怎能站立呢？因为你们说我是靠着贝耳则布赶鬼。22 \BibleRef{玛 12:27}我若靠着贝耳则布赶鬼，你们的子弟赶鬼，又靠着谁呢？这样，他们就要断定你们的是非。23 \BibleRef{玛 12:28}我若靠着天主的圣神赶鬼，这就是天主的国临到你们了。24 \BibleRef{玛 12:29}人怎能进壮士的家里，抢夺他的家具呢？除非先捆住那壮士，才可以洗劫他的家。25 \BibleRef{路 11:21}壮士披挂整齐，看守自己的住宅，他的家财是平安的。26 \BibleRef{路 11:22}但有一个比他更壮的来，胜过他，就夺去他所倚靠的盔甲兵器，又分了他的赃。27 \BibleRef{路 11:23}不与我相合的，就是敌我的；不同我收聚的，就是分散的。28 \BibleRef{谷 3:28}所以我告诉你们：一切的罪和亵渎的话都可得赦免；29 \BibleRef{谷 3:29}但亵渎圣神的，总不得赦免，乃要担当永远的罪。30 \BibleRef{谷 3:30}这话是因为他们说：他有污鬼。\EditorialNote{阿拉伯语本第56页}31 \BibleRef{玛 12:32}对你们说：凡说话干犯人子的，还可赦免；惟独说话干犯圣神的，今世来世总不得赦免。32 \BibleRef{玛 12:33}你们或以为树好，果子也好；树坏，果子也坏；因为看果子就可以知道树。33 \BibleRef{玛 12:34}毒蛇的种类！你们既是恶的，怎么能说出善来呢？因为心里所充满的，口里就说出来。34 \BibleRef{路 6:45}善人从他心里所存的善就发出善来；恶人从他心里所存的恶就发出恶来。35 \BibleRef{玛 12:36}我又告诉你们：凡人所说的闲话，当审判的日子，必要句句供出来。36 \BibleRef{玛 12:37}因为要凭你的话定你为义，也要凭你的话定你有罪。\par

37 \BibleRef{路 12:54}耶稣又对众人说：你们看见西边起了云彩，立刻说：要下大雨；果然这样。38 \BibleRef{路 12:55}起了南风，就说：将要燥热；也果然如此。39 \BibleRef{玛 16:2}晚上天发红，你们就说：天必要晴。40 \BibleRef{玛 16:3}早晨天发红，又发黑，你们就说：今日必有风雨。\EditorialNote{阿拉伯语本第57页}你们这些假善人，你们知道分辨天地的气色，怎么不知道分辨这时候呢？\par

41 \BibleRef{玛 12:22}当下有人将一个被鬼附着、又瞎又哑的人带到耶稣那里，耶稣就医治他，甚至那哑巴又能说话，又能看见。42 \BibleRef{玛 12:23}众人都惊奇，说：莫非这就是\Person{达味}的子孙吗？\par

43 \BibleRef{谷 6:30}宗徒们回到耶稣那里，将所行所作的一切事都告诉他。44 \BibleRef{谷 6:31}他就说：你们来，同我暗暗地到旷野地方去歇一歇。这是因为来往的人多，他们连吃饭的工夫也没有。\par

45 \BibleRef{路 7:36}此后，有一位法利塞人来見\Emph{他}，求他与他一同吃饭。他就进了那法利塞人的家里，46斜倚而坐。\BibleRef{路 7:37}在那城里有一个女人，她\Emph{是个}罪人；她听说他在那法利塞人家里坐席，就带了一玉瓶香膏，47站在他背后，\BibleRef{路 7:38}挨着他的脚哭，眼泪滴湿了他的脚，用她的头发擦干，又连连亲他的脚，48把香膏抹上。\BibleRef{路 7:39}那请他的法利塞人看见了\Emph{这事}，心里说：\Quote{这\Emph{人}若是先知，必知道摸他的是谁，是个怎样的女人，乃是个罪人。}\par

\SourceRecord{Translated by Hope W. Hogg.；Ante-Nicene Fathers；Vol. 9.；Edited by Allan Menzies.；Buffalo, NY: Christian Literature Publishing Co.,；1896.；<http://www.newadvent.org/fathers/100214.htm>.}

% structure: p0001 source=h1 target=section reason=page-title

\section{四福音合参（第十五部分）}

1 \BibleRef{路 7:40} \BibleQuote{耶稣回答对他说：\Quote{西满，我有一句话要同你说。}} 2 他对他说：\Quote{请\Emph{说下去}，我的夫子。} \BibleRef{路 7:41} \BibleQuote{耶稣对他说：\Quote{有一个债主，有两个债户，一个欠五百银币，另一个欠五十个银币。[阿拉伯语原文第58页]}} 3 \BibleRef{路 7:42} \BibleQuote{因为他们没有可偿还的，他就赦免了他们两个。那么，他们中哪一个更爱他呢？} 4 \BibleRef{路 7:43} \BibleQuote{西满回答说：\Quote{我想是那多得赦免的。}耶稣对他说：\Quote{你判断得对。}} 5 \BibleRef{路 7:44} \BibleQuote{于是他转向那妇人，对西满说：\Quote{你看见这\Emph{妇人}吗？我进了你的家，你没有给我水洗脚，但这\Emph{妇人}却用眼泪洗了我的脚，并用头发擦干。}} 6 \BibleRef{路 7:45} \BibleQuote{\Quote{你没有亲吻我，但这\Emph{妇人}自从进来，就不住地亲我的脚。}} 7 \BibleRef{路 7:46} \BibleQuote{\Quote{你没有用油抹我的头，但这\Emph{妇人}却用香膏抹了我的脚。}} 8 \BibleRef{路 7:47} \BibleQuote{\Quote{为此，我告诉你：她那许多的罪得了赦免，因为她爱得多；但那少得赦免的，爱得也少。}} 9,10 \BibleRef{路 7:48} \BibleQuote{于是他对那妇人说：\Quote{你的罪得了赦免。}} \BibleRef{路 7:49} \BibleQuote{同席的人心里说：\Quote{这人是谁，竟然也赦罪呢？}} 11 \BibleRef{路 7:50} \BibleQuote{耶稣对那妇人说：\Quote{你的信德救了你，平安地去吧！}}\par

12 \BibleRef{若 2:23} \BibleQuote{当他在耶路撒冷过逾越节时，许多人看见他所行的神迹，就信从了他。} 13,14 \BibleRef{若 2:24} \BibleQuote{但耶稣却不信任他们，} \BibleRef{若 2:25} \BibleQuote{因为他认识众人，不需要谁为别人作证，因为他知道人心里有什么。}\par

15 \BibleRef{路 10:1} \BibleQuote{此后，主又选定了另外七十二人，派遣他们两个两个地在他前面，往他将要去的各城各地。} 16 \BibleRef{路 10:2} \BibleQuote{他对他们说：\Quote{庄稼多，工人少；所以你们应当求庄稼的主人，派遣工人来收他的庄稼。}} 17 \BibleRef{路 10:3} \BibleQuote{你们去吧！[阿拉伯语原文第59页]看，我派遣你们如同羔羊进入狼群。} 18 \BibleRef{路 10:4} \BibleQuote{不要带钱囊，不要带口袋，也不要带鞋；路上也不要问候人。} 19 \BibleRef{路 10:5} \BibleQuote{无论进了哪一家，先要说：愿这一家平安。} 20 \BibleRef{路 10:6} \BibleQuote{那里若有和平之子，你们的平安就临到他；不然，你们的平安仍归于你们。} 21 \BibleRef{路 10:7} \BibleQuote{你们要住在那家，吃喝他们所供给的，因为工人应得他的工资。不要从这家搬到那家。} 22 \BibleRef{路 10:8} \BibleQuote{无论进了哪座城，他们若接待你们，就吃摆给你们的；} 23 \BibleRef{路 10:9} \BibleQuote{并治好其中的病人，对他们说：天主的国临近你们了。} 24 \BibleRef{路 10:10} \BibleQuote{但无论进了哪座城，他们若不接待你们，就到大街上去说：} 25 \BibleRef{路 10:11} \BibleQuote{连你们城中粘在我们脚上的尘土，我们也当众抖掉，但你们应当知道：天主的国临近了。} 26 \BibleRef{路 10:12} \BibleQuote{我告诉你们：在审判之日，索多玛所受的惩罚比那城还容易忍受。} 27 \BibleRef{玛 11:20} \BibleQuote{那时，耶稣就开始斥责那些有许多异能的城邑，因为它们没有悔改：} 28 \BibleRef{玛 11:21} \BibleQuote{\Quote{苛辣匝因，你有祸了！贝特赛达，你有祸了！因为在你们那里所行的异能，如果行在提洛和漆冬，它们早已披上苦衣，头上撒灰做补赎了。}} 29 \BibleRef{玛 11:22} \BibleQuote{但是我告诉你们：在审判之日，提洛和漆冬所受的惩罚比你们还容易忍受。} 30 \BibleRef{玛 11:23} \BibleQuote{葛法翁！你本已升到天上，却要下到阴府；因为在你那里所行的奇事，如果行在索多玛，它必会存留到今天。} 31 \BibleRef{玛 11:24} \BibleQuote{但是我告诉你们：在审判之日，索多玛地域所受的惩罚比你们还容易忍受。}\par

32 [阿拉伯语原文第60页] \BibleRef{路 10:16} \BibleQuote{他又对宗徒说：\Quote{听从你们的，就是听从我；拒绝你们的，就是拒绝我；拒绝我的，就是拒绝那派遣我的。}}\par

33 \BibleRef{路 10:17} \BibleQuote{那七十二人欢欢喜喜地回来，说：\Quote{主！因着你的名字，连恶魔都屈服于我们。}} 34 \BibleRef{路 10:18} \BibleQuote{耶稣对他们说：\Quote{我看见撒殚像闪电一样从天坠落。}} 35 \BibleRef{路 10:19} \BibleQuote{我已经给你们权柄去践踏蛇和蝎子，并胜过仇敌的一切势力；决没有什么能伤害你们。} 36 \BibleRef{路 10:20} \BibleQuote{但是，不要因为恶魔屈服于你们而欢喜，而要为你们的名字已登记在天上而欢喜。}\par

37 \BibleRef{路 10:21} \BibleQuote{就在那时，耶稣因圣神而欢欣说：\Quote{父啊！天地的主宰，我称谢你，因为你将这些事瞒住了明智和通达的人，而启示给了小孩子；是的，父啊！因为你本来喜欢这样做。}} 38 \BibleRef{路 10:22} \BibleQuote{耶稣转过身来对门徒说：\Quote{我父将一切都交给了我；除了父，没有人知道子是谁；除了子和子所愿意启示的人，没有人知道父是谁。}} 39 \BibleRef{玛 11:28} \BibleQuote{\Quote{凡劳苦和负重担的，你们都到我这里来，我要使你们安息。}} 40 \BibleRef{玛 11:29} \BibleQuote{你们背起我的轭，跟我学吧！因为我是良善心谦的：这样你们必要找到你们灵魂的安息。} 41 \BibleRef{玛 11:30} \BibleQuote{因为我的轭是甘饴的，我的担子是轻省的。}\par

42 \BibleRef{路 14:25} 当许多群众与他同行时，他转过身来，对他们说：43 \BibleRef{路 14:26} 凡到我这里来，若不恨自己的父亲、母亲、弟兄、姊妹、妻子、儿女，甚至自己的性命，就不能做我的门徒。44 （阿拉伯文版第61页）\BibleRef{路 14:27} 凡不背着自己的十字架跟随我的，也不能做我的门徒。45 \BibleRef{路 14:28} 你们中间谁愿意建造一座塔，不先坐下计算花费，是否有\Emph{足够}完成它呢？46 \BibleRef{路 14:29} 免得安了地基，却不能完工，所有看见的人就讥笑他，47 \BibleRef{路 14:30} 说：\InnerQuote{这个人开始建造，却不能完工。}48 \BibleRef{路 14:31} 或者一个王去与另一个王打仗，岂不先坐下酌量，能否用一万兵去迎那领两万兵来攻打他的呢？49 \BibleRef{路 14:32} 若是不能，就趁敌人还远的时候，派使者去求和。50 \BibleRef{路 14:33} 这样，你们每个人，凡愿意做我门徒的，就当放下一切所有的；否则不能做我的门徒。\par

\SourceRecord{Translated by Hope W. Hogg.；Ante-Nicene Fathers；Vol. 9.；Edited by Allan Menzies.；Buffalo, NY: Christian Literature Publishing Co.,；1896.；<http://www.newadvent.org/fathers/100215.htm>.}

% structure: p0001 source=h1 target=section reason=page-title

\section{\WorkTitle{迪亚特萨隆}（第十六段）}

\BibleQuote{1 \BibleRef{玛 12:38}那时，有几个经师和法利塞人应声，为要试探祂，2 就说：\InnerQuote{老师，我们愿意从你看到一个征兆。}\BibleRef{玛 12:39}祂回答说：\InnerQuote{这邪恶淫乱的世代寻求征兆；除了先知约纳的征兆外，必不给它征兆。3 \BibleRef{路 11:30}因为约纳怎样为尼尼微人成了征兆，人子也要照样为这世代成了征兆。4 \BibleRef{玛 12:40}约纳怎样在大鱼腹中三天三夜，人子也要照样在地心三天三夜。5 \BibleRef{路 11:31}南方女王在审判时，要与这世代的人一同起来，并定他们的罪：因为她从地极而来，为要听撒罗满的智慧；看，这里有一位比撒罗满更大。\EditorialNote{阿拉伯文，第62页}6 \BibleRef{玛 12:41}尼尼微人在审判时，要与这世代一同起来，并定它的罪：因为他们因约纳的宣讲而悔改了；看，这里有一位比约纳更大。7 \BibleRef{路 11:24}不洁之魔从人身上出去后，就走过无水之地，寻找安歇之处；若找不到，8 就说：\InnerQuote{我要回到我出来的那屋里去。}\BibleRef{路 11:25}如果它来到，看见它已经打扫干净，修饰整齐，9 \BibleRef{路 11:26}就去另外带七个比自己更恶的魔鬼来，进去住在那屋里；那人末后的景况，就比先前更坏了。10 \BibleRef{玛 12:45}这邪恶的世代也将如此。}}\par

\BibleQuote{11 \BibleRef{路 11:27}祂正说这些话时，人群中有一个妇人高声对祂说：\InnerQuote{怀过祢的胎和哺养祢的乳房是有福的。}12 \BibleRef{路 11:28}祂却回答说：\InnerQuote{那听了天主的话而遵守的人，才是有福的。}}\par

13 祂正对群众说话时，祂的母亲和弟兄来到，想与祂说话；14 但因群众众多而不能；15 他们站在外面，打发人去叫祂。有一个人对祂说：\InnerQuote{看，祢的母亲和祢的弟兄\Emph{正在}外面站着，想与祢说话。}16 \BibleRef{玛 12:48}祂却回答那对祂说话的人说：\InnerQuote{谁是我的母亲？谁是我的弟兄？}17 \BibleRef{玛 22:49}祂伸出手指著祂的门徒说：\InnerQuote{看，我的母亲！看，我的弟兄！}18 \BibleRef{玛 12:50}凡奉行我在天之父的旨意的人，就是我的兄弟、姊妹和母亲。\par

19 此后，\Person{耶稣}走遍各城各乡，宣讲\EditorialNote{阿拉伯文，第63页}并传扬天主的国，十二门徒与祂同行，20 \BibleRef{路 8:2}还有曾被污灵和疾病治愈了的妇女，就是那称为玛利亚玛达肋纳的，从她身上赶出了七个魔鬼，21 \BibleRef{路 8:3}以及黑落德的家宰雇撒的妻子约安纳，和苏撒纳，还有许多别的妇女，她们都用自己的财物供给他们。\par

\BibleQuote{22 \BibleRef{玛 13:1}此后，\Person{耶稣}从屋里出来，坐在海边。\BibleRef{玛 13:2}23 有许多群众聚集到祂跟前。当人群拥挤祂时，祂上船坐下；众人都站在岸上。24 \BibleRef{玛 13:3}祂就用比喻对他们讲了许多事，说：\InnerQuote{有一个撒种的出去撒种：25 撒种的时候，有的落在路旁；26 被人践踏了，飞鸟来把它吃了。\BibleRef{玛 13:5}有的落在石头地上：那里土不深，立刻发芽，因为土不深；27 \BibleRef{玛 13:6}太阳出来一晒，就枯萎了，因为没有根，就干枯了。28 有的落在荆棘中，荆棘一同长起来，把它窒息了；29 它就不结果实。有的落在好而\Emph{且}肥沃的地里，就生长起来，结了果实，有三十倍的，有六十倍的，有一百倍的。}30 \BibleRef{路 8:8}祂说了这些话，就大声说：\InnerQuote{有耳可听的，就应当听！}31 当只剩下他们的时候，门徒来问祂说：\InnerQuote{这比喻是什么意思？你为什么要用比喻对他们讲话？}32 祂\EditorialNote{阿拉伯文，第64页}回答他们说：\InnerQuote{天国的奥秘只给了你们知道，但对外人却是不给的。33 \BibleRef{玛 13:12}凡有的，还要给他，使他有余；凡没有的，连他所有的也要从他夺去。34 \BibleRef{玛 13:13}因此，我所以用比喻对他们讲，是因为他们看却看不见，听却听不见，也不明白。35 \BibleRef{玛 13:14}在他们身上正应验了依撒意亚的预言说：}}\par

\BibleQuote{你们听是听见，却不明白；\newline{}看是看见，却不晓得；\newline{}36 \BibleRef{玛 13:15}因为这百姓的心迟钝了，\newline{}耳朵发沉，\newline{}眼睛闭着；\newline{}免得他们眼睛看见，\newline{}耳朵听见，\newline{}心里明白，\newline{}回转过来，\newline{}我就医治他们。}\par

37, 38 \BibleRef{玛 13:16} 但你们，你们有福的眼睛，因为它们看见；你们的耳朵，因为它们听见。\BibleRef{路 10:23} 看见你们所看见之事的眼睛，39 是有福的。\BibleRef{玛 13:17} 我实在告诉你们：许多先知和义人渴望看见你们所看见的，却没有看见；渴望听见你们所 40 听见的，却没有听见。\BibleRef{谷 4:13} 你们既不明白这个比喻，又怎能明白一切比喻呢？41, 42 \BibleRef{玛 13:18} 请听撒种者的比喻。\BibleRef{谷 4:14} 那撒种的，撒的是天主的圣言。43 \BibleRef{玛 13:19} 凡听见天国的圣言，而不领悟的，那恶者就来，把撒在他心里的圣言夺去：这就是那撒在大道旁的。44 \BibleRef{玛 13:20} 那撒在石头地里的，是指那听了圣言，立刻欣然接受；45, 46 \EditorialNote{阿拉伯语原文第 65 页} \BibleRef{玛 13:21} 但在心里没有根，只是暂时的；一旦为圣言遭遇艰难或迫害，他 47 便立刻跌倒。\BibleRef{玛 13:22} 那撒在荆棘中的，是指那听了圣言；\BibleRef{谷 4:19} 而世俗的焦虑、财富的迷惑，以及其余的各种私欲进来，把圣言窒息了，它就成了没有果实的。\BibleRef{路 8:15} 那撒在好地里的，是指那以纯洁善良的心听了我的圣言，且领悟持守，在忍耐中结出果实，\BibleRef{玛 13:23} 并结出百倍、六十倍或三十倍的果实。\par

49 \BibleRef{谷 4:26} 他又说：天主的国好比一个人把种子撒在地里，50 黑夜白天或睡或起，\BibleRef{谷 4:27} 种子发芽生长，如何成长，他不知道。\BibleRef{谷 4:28} 土地自然生出果实：先发苗，后吐穗，最后穗上结成饱满的麦粒。\BibleRef{谷 4:29} 果实一成熟，他便立刻动起镰刀，因为收割的时期到了。\par

\SourceRecord{Translated by Hope W. Hogg.；Ante-Nicene Fathers；Vol. 9.；Edited by Allan Menzies.；Buffalo, NY: Christian Literature Publishing Co.,；1896.；<http://www.newadvent.org/fathers/100216.htm>.}

% structure: p0001 source=h1 target=section reason=page-title

\section{四福音合编（第十七部分）}

1 \BibleRef{玛 13:24} 祂又给他们设了另一个比喻，说：\Quote{天国好比一个人，将好种子撒在自己的田里；2 \BibleRef{玛 13:25} 但人们睡觉的时候，他的仇敌来了，将莠子撒在麦子中间，就走了。\BibleRef{玛 13:26} 当苗长起来，结了果实，莠子也显出来了。\BibleRef{玛 13:27} 家主的仆人来对他说：\InnerQuote{主啊，你不是将好\EditorialNote{阿拉伯文版第66页}种子撒在你的田里了吗？那里面怎么有莠子呢？}\BibleRef{玛 13:28} 他对他们说：\InnerQuote{仇敌做了这事。}仆人对他说：\InnerQuote{你要我们去把它拔出来吗？}\BibleRef{玛 13:29} 他说：\InnerQuote{不必，恐怕你们拔莠子时，连麦子也一起拔出来。\BibleRef{玛 13:30} 让它们一同生长，直到收割的时候；在收割的时候，我要对收割者说：先把莠子拔出来，捆成捆，用火烧掉；把麦子收入我的仓里。}}\par

8,9 \BibleRef{玛 13:31} 祂又给他们设了另一个比喻，说：\Quote{\BibleRef{路 13:18} 天主的国好像什么？我要把它比作什么呢？在什么比喻里可以把它描述出来？它11好比一粒芥子，一个人拿来种在自己的田里；在所有撒在地里的东西中，它比地上所有撒下的东西都小；但当它长大时，比所有蔬菜都大，并长出大枝条，以致天上的飞鸟都能在它的枝条上做\Emph{它们的}窝。}\par

13,14 祂又给他们设了另一个比喻：\Quote{\BibleRef{路 13:20} 天主的国我要把它比作什么呢？\BibleRef{玛 13:33} 它好比酵母，一个妇人拿来，和在三斗面粉中，直到全部发酵。}\par

16 耶稣用比喻向群众宣讲了这一切，乃是按照他们所能听的。17 并且不用比喻，祂就不对他们讲什么；这是为应验主藉先知所说的话：\par

\Quote{我要开口说比喻；\newline{}我要说出从创立世界以来隐密的事。}\par

18 \BibleRef{谷 4:34} 但祂私下给祂的门徒解释了一切。\par

19 \BibleRef{玛 13:36} 那时，耶稣离开了群众，进了屋子。祂的门徒来到祂跟前，对祂说：\InnerQuote{请给我们解释那田间莠子\EditorialNote{阿拉伯文版第67页}的比喻。}\BibleRef{玛 13:37} 祂回答他们说：\Quote{那撒好种子的，就是21人子；\BibleRef{玛 13:38} 田地是世界；好种子是天国的22儿女；莠子是邪恶者的儿女；\BibleRef{玛 13:39} 那撒莠子的仇敌是撒殚；收割是世界的终结；收割者是天使。\BibleRef{玛 13:40} 正如莠子被拔出来用火烧掉，在世界的终结时也将如此。\BibleRef{玛 13:41} 人子要派遣祂的天使，从祂的国度里搜集一切使人跌倒的，\BibleRef{玛 13:42} 以及所有作恶的人，将他们抛入火窑里：在那里要有哀号和切齿。\BibleRef{玛 13:43} 那时，义人要在他们父的国里发光如同太阳。有耳能听的，就听吧。}\par

27 \BibleRef{玛 13:44} \Quote{天国好像藏在地里的宝贝；一个人找到了，把它藏起来，然后因欢喜，就去变卖他所有的一切，买了那块地。}\par

28 \BibleRef{玛 13:45} \Quote{天国又好像一个商人，\Emph{即是}一个寻找精美珍珠的商人；}\BibleRef{玛 13:46} \Quote{当他找到一颗价值连城的珍珠，就去变卖他所有的一切，买了它。}\par

30 \BibleRef{玛 13:47} \Quote{天国又好像撒在海里的网，网罗各种鱼；}\BibleRef{玛 13:48} \Quote{网满了，他们就拉上岸，坐下拣选；好的收在器具里，坏的丢在外面。}\BibleRef{玛 13:49} \Quote{世界的终结时也将如此：天使要出来，}\BibleRef{玛 13:50} \Quote{从义人中把恶人分开，将他们抛入火窑里；在那里要有哀号和切齿。}\par

34 \BibleRef{玛 13:51} 耶稣对他们说：\Quote{这一切\Emph{事}你们都明白了吗？}他们\EditorialNote{阿拉伯文版第68页}对祂说：\Quote{是的，主。}\BibleRef{玛 13:52} 祂对他们说：\Quote{因此，凡成为天国门徒的经师，就像一个家主，从他的宝库里拿出新的和旧的来。}\par

36, 37 \BibleRef{玛 13:53} 耶稣讲完了这一切比喻，就离开那里，来到自己的家乡；\BibleRef{玛 13:54} 并在他们的会堂里教训他们，以致他们困惑不已。38 \BibleRef{谷 6:2} 到了安息日，耶稣开始在会堂里教训人；许多39听见的人就惊奇，说：\Quote{\Emph{这人}从哪里得到这些事？}许多人嫉妒他，不理会他，却说：\Quote{这人所赐的智慧是什么？竟有这样的异能由他手中发生？}40 \BibleRef{玛 13:55} 这不是那木匠，那木匠的儿子吗？他的母亲不是叫玛利亚吗？41 他的弟兄们：雅各伯、若瑟、西满和犹大呢？\BibleRef{玛 13:56} 他的妹妹们，42 看，她们不都是在我们这里吗？\BibleRef{玛 13:57} \Emph{这人}从哪里得到这一切呢？他们就对他起了疑心。\BibleRef{路 4:23} 耶稣知道他们的心思，就对他们说：\Quote{你们必向我说这句俗语：医生，先医治你自己吧！还有，43 我们听见你在葛法翁所行的一切，也当在你自己的\Emph{家乡}行。}\BibleRef{路 4:24} 他又说：\Quote{我实在告诉你们，先知在自己的\Emph{家乡}是不被接纳的，44 在自己的弟兄中也不被接纳。}\BibleRef{谷 6:4} 因为先知除了在自己的\Emph{家乡}、45 \Emph{自己的}亲族和\Emph{自己的}本家外，是没有不受人轻视的。\BibleRef{路 4:25} 我实在告诉你们，在厄里亚先知的日子，天闭塞了46三年零六个月，遍地大饥荒，以色列中有许多寡妇；\BibleRef{路 4:26} 厄里亚[阿拉伯文，第69页]并没有被遣送到她们中任何一个那里，只到了漆冬的匝尔法特，到一个47寡妇那里。\BibleRef{路 4:27} 在厄里叟先知的时候，以色列中有许多癞病人；但他们中没有一个得洁净，只有纳阿曼纳巴泰人。48 \BibleRef{谷 6:5} 耶稣在那里因为他们的不信，不能行许多异能；49 只按手在几个病人身上，治好了\Emph{他们}。\BibleRef{谷 6:6} 他50诧异他们的不信。\BibleRef{路 4:28} 会堂里的人听见了，51 都怒气填胸；\BibleRef{路 4:29} 就起来，把他赶出城外，带到他们城所在的山崖上，52 要把他推下去；\BibleRef{路 4:30} 但他从他们中间走过，就去了。\par

53 \BibleRef{谷 6:6} 他走遍那\Emph{在}\Place{纳匝肋}周围的村庄，并在他们的会堂里教训人。\par

\SourceRecord{Translated by Hope W. Hogg.；Ante-Nicene Fathers；Vol. 9.；Edited by Allan Menzies.；Buffalo, NY: Christian Literature Publishing Co.,；1896.；<http://www.newadvent.org/fathers/100217.htm>.}

% structure: p0001 source=h1 target=section reason=page-title

\section{\WorkTitle{四福音合参（第18节）}}

1 那时，分封侯黑落德听见耶稣的名声，以及经他手所行的一切事，\BibleRef{谷 6:14} 就惊奇，因为他得到了关于他的确切消息。\BibleRef{路 9:7} 有人说洗者若翰从死人中复活了；\newline{}3 又有人说厄里亚显现了；也有人说耶肋米亚；\newline{}4 又有人说古先知中的一位复活了；另有人说他是像先知中的一位先知。\newline{}5 黑落德对他的仆人说：\Quote{这是洗者若翰，我斩了他的头；他从死人中复活了，所以大能的事从他而出。}\BibleRef{谷 6:17} 因为黑落德曾派人捉住若翰，把他捆锁在监狱里，为了他兄弟斐理伯的妻子黑落狄雅，因为他娶了她。\BibleRef{谷 6:18} 若翰对黑落德说：\Quote{你不可占有你兄弟的妻子。}\BibleRef{谷 6:19} 黑落狄雅就怀恨他，愿意杀他，却不能。\BibleRef{谷 6:20} 但黑落德敬畏若翰，知道他是个正义圣洁的人，就保护他；听他讲道，多次遵行，且乐意听从。\BibleRef{玛 14:5} 他也想杀他，但害怕百姓，因为他们认为若翰是先知。\BibleRef{谷 6:21} 有一天，到了好日子，黑落德在他的生日设宴招待他的大臣、军官和加里肋亚的首领。\BibleRef{谷 6:22} 黑落狄雅的女儿进来跳舞，使黑落德和同席的人欢喜。王对那女子说：\Quote{你无论求什么，我都给你。}\BibleRef{谷 6:23} 又对她发誓说：\Quote{凡你所求的，即使是我王国的一半，我也必给你。}\BibleRef{谷 6:24} 她出去对她母亲说：\Quote{我该求什么？}她母亲说：\Quote{洗者若翰的头。}\BibleRef{谷 6:25} 她便急忙进去，向王请求说：\Quote{我愿你现在就把洗者若翰的头放在盘子里给我。}\BibleRef{谷 6:26} 王非常忧愁；但因了誓言和同席的人，不愿拒绝她。\BibleRef{谷 6:27} 王立刻派了一个刽子手，吩咐把若翰的头拿来。刽子手就去，在监狱里斩了若翰的头，\BibleRef{谷 6:28} 放在盘子里，递给那女子；那女子给了她母亲。\BibleRef{谷 6:29} 若翰的门徒听见了，就来领去他的尸体，埋葬了。\BibleRef{玛 14:12} 他们来告诉耶稣所发生的事。\BibleRef{路 9:9} 为此黑落德说：\Quote{我斩了若翰，这人是谁？我竟听见关于他的这些事。}他就想见他。耶稣听见了，就离开那里，独自乘船到一处荒僻的地方，到提庇黎雅的加里肋亚海的另一边。\par

\BibleRef{谷 6:33} 许多人看见他们离去，认出了他们，便从各城步行，一起赶到那里，先到了他们那里。\BibleRef{若 6:2} 因为他们看见他在病人身上所行的奇迹。\BibleRef{若 6:3} 耶稣上了山，和门徒一起坐在那里。\BibleRef{若 6:4} 犹太人的逾越节近了。\BibleRef{若 6:5} 耶稣举目看见大批群众向他走来。\BibleRef{谷 6:34} 就怜悯他们，因为他们好像没有牧人的羊。\BibleRef{路 9:11} 他接待了他们，给他们讲论天主的国，并治好了那些需要医治的人。\par

27 \BibleRef{玛 14:15} 到了傍晚，他的门徒来到他跟前，对28他说：\BibleRef{谷 6:36} 这地方是荒野，时间已过；请遣散群众，让他们往周围的村庄和乡镇去，29为自己买饼吃，因为他们没有吃的东西。\BibleRef{玛 14:16} 但他对他们说：他们不必去；你们给他们吃的吧。\BibleRef{玛 14:17} 他们对他说：我们这里没有30 \Emph{足够}的饼。\BibleRef{若 6:5} 他对斐理伯说：我们从哪里买饼给这些人吃呢？31、32 \BibleRef{若 6:6} 他说这话是要试探他；他知道自己将要做什么。\BibleRef{若 6:7} 斐理伯回答他说［阿拉伯语原稿，第72页］：就是二百银币的饼也不够他们33每人取一小点之后。\BibleRef{若 6:8} 他的一个门徒，即西满伯多禄的兄弟安德肋，34对他说：\BibleRef{若 6:9} 这里有一个孩童，他有五个大麦饼和两条鱼；但这么多的人，这点东西算得了什么？难道你要我们去给众人买吃的东西吗？因为我们只有这五个饼和两条鱼。36那地方的草很多。耶稣对他们说：叫众人坐在草地上，37每五十人一组。\BibleRef{谷 6:40} 门徒就照做了。于是众人一组一组地坐下，有的一百，有的五十。\BibleRef{玛 14:18} 耶稣便对他们说：把那五个饼和两条鱼拿39到我这里来。\BibleRef{谷 6:41} 他们拿给他之后，耶稣拿起饼和鱼，望天祝福了，掰开，交给40门徒摆在他们面前；\BibleRef{玛 14:19} 门徒就摆在群众面前；41所有的人都吃了，并且吃饱了。他们吃饱以后，耶稣对门徒说：把剩下的碎块收集起来，免得糟蹋。\BibleRef{若 6:13} 他们就收集起来，装满了十二篮子，都是从那五个大麦饼和两条鱼所剩下42的碎块。\BibleRef{玛 14:21} 吃的人，除了妇女和孩童，约有五千。44［阿拉伯语原稿，第73页］\BibleRef{谷 6:45} 耶稣立刻催门徒上船，叫他们先往对岸的贝特赛达去，而他45\Emph{自己}却要去遣散群众。\BibleRef{若 6:14} 那些看见耶稣所行46神迹的人就说：这真是那位要到世界上来的先知。\BibleRef{若 6:15} 耶稣知道他们要来强迫他，立他为王，就独自退到山上祈祷去了。\par

47、48 \BibleRef{若 6:16} 到了黄昏将近，他的门徒下到海边，上了船，往葛法翁那边去。\BibleRef{若 6:17} 天黑下来，耶稣49还没有来到他们那里。\BibleRef{若 6:18} 海上因起了50大风，就翻腾起来。\BibleRef{玛 14:24} 那船已离岸很远，被波浪冲击，因为风不顺。\par

\SourceRecord{Translated by Hope W. Hogg.；Ante-Nicene Fathers；Vol. 9.；Edited by Allan Menzies.；Buffalo, NY: Christian Literature Publishing Co.,；1896.；<http://www.newadvent.org/fathers/100218.htm>.}

% structure: p0001 source=h1 target=section reason=page-title

\section{四福音合参（第十九段）}

1 \BibleRef{玛 14:25} 夜间四更时分，耶稣在海面上行走，往他们那里去。那时他们奋力划了约二十五至三十斯塔狄。2 \BibleRef{玛 14:26} 门徒看见他在海面上行走，就惊慌起来，说：\Quote{是个幻象！}并因恐惧而喊叫。3 \BibleRef{玛 14:27} 耶稣立刻对他们说：\Quote{放心！是我，不要害怕！}4 \BibleRef{玛 14:28} 刻法回答说：\Quote{主，如果是你，请吩咐我到你那里去，走在水面上。}5 \BibleRef{玛 14:29} 耶稣说：\Quote{来吧！}刻法就下了船，在水面上行走，要到耶稣那里去。6 \BibleRef{玛 14:30} 但他一见风势猛烈，就害怕起来，快要沉下去，便大声呼喊说：\Quote{主，救我！}7 \BibleRef{玛 14:31} 主立刻伸手拉住他，对他说：\Quote{小信德的人，你为什么怀疑？}8 \BibleRef{玛 14:32} 耶稣上了船，他和西满都上去，风立刻停了。9 \BibleRef{玛 14:33} 船上的人都前来朝拜他，说：\Quote{你真是天主子！}10 \BibleRef{若 6:21} 船立刻到达了他们所要去的地方。11 他们上了岸，大大惊奇，心里困惑。12 \BibleRef{谷 6:52} 因为他们没有领悟关于饼的事，心肠还是迟钝。\par

14 那地方的人一知道耶稣来了，就传遍全境，开始把病人用床抬到听说他在的地方。15 \BibleRef{谷 6:56} 凡他进入的村庄、城市或乡间，人们都把病人放在街市上，恳求他允许他们只摸他的衣边；凡摸着的都痊愈了，并得以存活。\par

16 \BibleRef{若 6:22} 第二天，站在海那边的群众看见那里没有别的船，只有门徒上的那条船，17 \BibleRef{若 6:23} 并且耶稣没有和门徒一起上船（但另有些从提庇黎雅来的船，靠近他们吃饼的地方——那饼是主祝福过的）。18 \BibleRef{若 6:24} 群众见耶稣和门徒都不在那里，就上了那些船，来到葛法翁，寻找耶稣。19 \BibleRef{若 6:25} 他们在海那边找到了他，就对他说：\Quote{老师，你什么时候到这里来了？}20 \BibleRef{若 6:26} 耶稣回答说：\Quote{我实实在在告诉你们：你们寻找我，不是因为见了神迹，而是因为吃了饼且吃饱了。21 \BibleRef{若 6:27} 不要为那可朽坏的食物劳碌，而要为那存留到永生的食物劳碌，这食物是人子要赐给你们的，因为父天主印证了他。}22 \BibleRef{若 6:28} 他们对他说：\Quote{我们该做什么，才算做天主的工作？}23 \BibleRef{若 6:29} 耶稣回答说：\Quote{天主的工作就是：你们信他所派遣的那位。}24 \BibleRef{若 6:30} 他们对他说：\Quote{你行什么神迹，好叫我们看见而信你？你做什么事？25 \BibleRef{若 6:31} 我们的祖先在旷野里吃过玛纳，如经上所载：\InnerQuote{他从天上赐给他们食粮吃。}}26 \BibleRef{若 6:32} 耶稣对他们说：\Quote{我实实在在告诉你们：不是梅瑟给了你们那从天上来的食粮，而是我父给你们赐下了那从天上来的真食粮。27 \BibleRef{若 6:33} 天主的食粮，就是那从天上降下，并赐给世界生命的。}28 \BibleRef{若 6:34} 他们对他说：\Quote{主，求你把这种食粮常赐给我们。}29 \BibleRef{若 6:35} 耶稣对他们说：\Quote{我就是生命的食粮；到我这里来的，永不会饿；信我的，永不会渴。30 \BibleRef{若 6:36} 但我曾对你们说过：你们看见了我，却不信。31 \BibleRef{若 6:37} 凡父赐给我的人，必到我这里来；到我这里来的，我决不把他遗弃在外面。32 \BibleRef{若 6:38} 因为我从天上降下来，不是为承行我的旨意，而是为承行派遣我者的旨意；33 \BibleRef{若 6:39} 派遣我者的旨意就是：凡他赐给我的，我什么都不失落，而在末日使他复活。34 \BibleRef{若 6:40} 这是我父的旨意：凡看见子并信他的人，必得永生，并且我将在末日使他复活。}\par

35 \BibleRef{若 6:41} 犹太人遂因为他说\Quote{我是从天上降下来的食粮}，而窃窃私议。36 \BibleRef{若 6:42} 他们说：\Quote{这不是若瑟的儿子耶稣吗？他的父亲和母亲我们岂不认识？那么他怎么说\InnerQuote{我从天上降下来}呢？}37 \BibleRef{若 6:43} 耶稣回答说：\Quote{你们不要彼此窃窃私议。38 \BibleRef{若 6:44} 除非派遣我的父吸引他，没有人能到我这里来；而我在末日要使他复活。39 \BibleRef{若 6:45} 先知书上记载：\InnerQuote{他们都要受天主的教导。}凡听了父的教导而学习的，必到我这里来。40 \BibleRef{若 6:46} 这不是说有人见过父，只有那从天主来的，才见过父。41 \BibleRef{若 6:47} 我实实在在告诉你们：信我的人，有永生。42 \BibleRef{若 6:48} 我就是生命的食粮。43 \BibleRef{若 6:49} 你们的祖先在旷野里吃过玛纳，却死了。44 \BibleRef{若 6:50} 这是从天上降下来的食粮，使人吃了不死。45 \BibleRef{若 6:51} 我是从天上降下来的活食粮；谁若吃了这食粮，将永远活着。我要赐给的食粮，就是我的肉，为世界的生命而赐的。}\par

46 \BibleRef{若 6:52} \BibleQuote{因此犹太人彼此争论说：\Quote{他怎能把他的肉给我们吃呢？}} \BibleRef{若 6:53} \BibleQuote{耶稣对他们说：\Quote{我实实在在告诉你们：你们若不吃人子的肉，不喝他的血，在你们内就没有生命。}} \BibleRef{若 6:54} \BibleQuote{谁吃我的肉，并喝我的血，必有永生，在末日我要使他复活。} \BibleRef{若 6:55} \BibleQuote{我的肉确实是食物，我的血确实是饮料。} \BibleRef{若 6:56} \BibleQuote{谁吃我的肉，并喝我的血，便住在我内，我也住在他内。} \BibleRef{若 6:57} \BibleQuote{就如生活的父派遣了我，我因父而生活；照样，那吃我的人，也要因我而生活。} \BibleRef{若 6:58} \BibleQuote{这是从天上降下来的食粮；不像你们的祖先吃过玛纳仍然死了；谁吃这食粮，将永远活着。} \BibleRef{若 6:59} 这些话是他在\Place{葛法翁}会堂教导时说的。 \BibleRef{若 6:60} \BibleQuote{他的门徒中有许多听了，便说：\Quote{这话生硬，谁能听得下去呢？}}\par

\SourceRecord{Translated by Hope W. Hogg.；Ante-Nicene Fathers；Vol. 9.；Edited by Allan Menzies.；Buffalo, NY: Christian Literature Publishing Co.,；1896.；<http://www.newadvent.org/fathers/100219.htm>.}

% structure: p0001 source=h1 target=section reason=page-title

\section{\WorkTitle{迪亚泰萨隆（第二十段）}}

\BibleRef{若 6:61} 耶稣心里知道门徒们因这话而窃窃私语，便对他们说：这话使你们反感吗？\BibleRef{若 6:62} 那么，如果你们看见人子升到他原先所在之处，又怎样呢？\BibleRef{若 6:63} 使人活的是圣神，肉体一无所益；我对你们所说的话，就是圣神，就是生命。\BibleRef{若 6:64} 但你们中间有些人不信。耶稣从起初就知道哪些人不信，谁要出卖他。\BibleRef{若 6:65} 于是他对他们说：因此我对你们说过，除非蒙父恩赐，没有人能到我这里来。\par

[阿拉伯文，第78页] \BibleRef{若 6:66} 从此，他许多门徒退去，不再与他同行。\BibleRef{若 6:67} 耶稣就对那十二人说：你们也愿意走吗？\BibleRef{若 6:68} 西满刻法回答说：主，我们去投奔谁呢？你有永生的话。\BibleRef{若 6:69} 我们信了，也知道你是默西亚，永生天主之子。\BibleRef{若 6:70} 耶稣对他们说：我不是拣选了你们十二个人吗？但你们中间有一个是魔鬼！\BibleRef{若 6:71} 他这话是指加略人西满的儿子犹大说的；因为他是十二人之一，想要出卖他。\par

\BibleRef{路 11:37} 耶稣说话的时候，有一个法利塞人请他吃饭；他便进去坐席。\BibleRef{路 11:38} 那法利塞人看见，诧异他饭前没有先洗手。\BibleRef{路 11:39} 主对他说：你们法利塞人洗净杯盘的外面，就自以为洁净了；但你们内心却充满不义和邪恶。\par

\BibleRef{路 11:40} 无知的人哪！那造外面的，不也造了里面吗？\BibleRef{路 11:41} 只要把你们所有的施舍给人，一切对你们就都洁净了。\par

\BibleRef{谷 7:1} 有法利塞人和从耶路撒冷来的经士聚集到耶稣跟前。\BibleRef{谷 7:2} 他们看见他的几个门徒用不洁的手，就是没有洗过的手吃饭，就指责他们。\BibleRef{谷 7:3} 因为所有的犹太人和法利塞人，如果不仔细洗手，就不吃饭；他们拘守先人的传授。\BibleRef{谷 7:4} 从市场回来，若不洗濯，也不吃饭；还有许多别的传授，如洗杯、量器、铜器、床榻等，他们都谨守。\BibleRef{谷 7:5} 经士和法利塞人问他说：你的门徒为什么不遵守先人的传授，而用不洁的手吃饭呢？\BibleRef{玛 15:3} 耶稣回答他们说：你们为什么为了你们的传授，而违犯天主的诫命呢？天主说：\BibleQuote{要孝敬你的父亲和你的母亲}；又说：\BibleQuote{咒骂父亲或母亲的，应处死刑。}\BibleRef{谷 7:11} 你们却说：人若对父亲或母亲说：我所能供养你的，已成了\InnerQuote{科尔班}，即奉献品，\BibleRef{谷 7:12} 你们就不许他再为父亲或母亲做什么；\BibleRef{谷 7:13} 这样你们为了你们所传授的，废弃了天主的话；并且你们还做了许多类似的事。\BibleRef{谷 7:8} 你们离弃了天主的诫命，而拘守人的传授。\BibleRef{谷 7:9} 你们真好！你们竟为遵守你们的传授，而废弃天主的诫命！\BibleRef{玛 15:7} 你们这些假善人！依撒意亚论你们预言得真对，他说：\par

\BibleRef{玛 15:8}\BibleQuote{这百姓用嘴唇尊敬我，\newline{}他们的心却远离我。\newline{}\BibleRef{玛 15:9} 他们恭敬我也是枉然，\newline{}因为他们把人的诫命当作道理传授。}\par

\BibleRef{谷 7:14}\BibleQuote{你们众人都要听我，且要明白：}\BibleRef{谷 7:15}\BibleQuote{从外面进入人内的，不能污秽人；但从人里面出来的，那才能污秽人。}\BibleRef{谷 7:16}\BibleQuote{凡有耳能听的，就听吧！}\BibleRef{玛 15:12}\BibleQuote{那时门徒前来对他说：你知道法利塞人听了这话，就愤慨吗？}\BibleRef{玛 15:13}\BibleQuote{他回答说：凡不是我天父所栽种的植物，必要连根拔除。}\BibleRef{玛 15:14}\BibleQuote{由他们吧！他们是盲人领盲人。如果盲人领盲人，两人必会掉进坑里。}\par

耶稣离开群众，进了屋子，西满刻法问他说：主，请给我们解释这个比喻。\BibleRef{谷 7:18}\BibleQuote{连你们也不明白吗？你们不知道：凡从外面进入人内的，不能污秽人，}\BibleRef{谷 7:19}\BibleQuote{因为它不进人心，只进肚腹，然后排泄出来——这是说一切食物都是洁净的？}\BibleRef{玛 15:18}\BibleQuote{但从口里出来的，是从心里出来的，那才污秽人。}\BibleRef{谷 7:21}\BibleQuote{因为从人心里出来的，有恶念、}\BibleRef{谷 7:22}\BibleQuote{邪淫、奸淫、偷盗、假见证、凶杀、不义、邪恶、诡诈、愚昧、嫉妒、毁谤、骄傲、愚妄：}\BibleRef{谷 7:23}\BibleQuote{这一切恶事，都是从里面出来，并且能污秽人。}\BibleRef{玛 15:20}\BibleQuote{至于不洗手吃饭，那并不污秽人。}\par

46 耶稣从那里出去，来到提洛和漆冬的边界。他进了一所房子，愿意没有人知道；\newline{}\BibleRef{谷 7:25} 但立刻有一个客纳罕妇人，她的女儿有48、49一个不洁之神，听到了他的事。\newline{}\BibleRef{谷 7:26} 那妇人是外邦人，来自叙利亚的厄梅撒。\newline{}\BibleRef{玛 15:22} 她从后面出来，喊着说：\BibleQuote{我主，\Emph{你}50达味之子，可怜我吧！我的女儿被撒殚折磨得苦。}\newline{}\BibleRef{玛 15:23} 耶稣[阿拉伯文，第81页]却一句话也不回答她。他的门徒前来求他说：\BibleQuote{请打发51她走吧：因为她在我们后面喊叫。}\newline{}\BibleRef{玛 15:24} 耶稣回答他们说：\BibleQuote{我被派遣52只是为了以色列家迷失的羊。}\newline{}\BibleRef{玛 15:25} 但那妇人前来朝拜他，说：\BibleQuote{主啊，帮助我，可怜我吧。}\par

53 \BibleRef{玛 15:26} 耶稣对她说：\BibleQuote{拿儿女的饼扔给小狗吃，是不合宜的。}\newline{}\BibleRef{玛 15:27} 但她回答说：\BibleQuote{是的，主；可是小狗也吃主人桌子上掉下来的碎屑，并且活着。}\newline{}\BibleRef{玛 15:28} 那时耶稣对她说：\BibleQuote{啊，妇人，56你的信德真大！照你所愿的，给你成就吧。}\newline{}\BibleRef{谷 7:29} \BibleQuote{那么你\Emph{走你的路}；57因这句话，魔鬼已从你女儿身上出去了。}\newline{}\BibleRef{玛 15:28} 就在那一刻，58她的女儿痊愈了。\newline{}\BibleRef{谷 7:30} 那妇人回到家里，见女儿躺在床上，魔鬼已经出去了。\par

\SourceRecord{Translated by Hope W. Hogg.；Ante-Nicene Fathers；Vol. 9.；Edited by Allan Menzies.；Buffalo, NY: Christian Literature Publishing Co.,；1896.；<http://www.newadvent.org/fathers/100220.htm>.}

% structure: p0001 source=h1 target=section reason=page-title

\section{戴亚忒隆（第21部分）}

1 \BibleRef{谷 7:31}耶稣又从提洛和漆冬的边境出来，来到加里肋亚海，靠近德卡波利斯边境。2 \BibleRef{谷 7:32}有人带了一个又聋又哑的人来见他，求他按手医治他。3 \BibleRef{谷 7:33}耶稣领他离开群众，独自走到一边，用手指蘸唾沫，并将\Emph{它们}插入他的耳朵，又摸他的舌头；4 \BibleRef{谷 7:34}望天叹息，对他说：\Quote{开了罢。}5 \BibleRef{谷 7:35}那时他的耳朵就开了，舌结也解了，说话也清楚了。6 \BibleRef{谷 7:36}耶稣切切嘱咐他们不要告诉人；但他越嘱咐，他们越\Emph{越发}传扬，并且极其惊讶，7 \BibleRef{谷 7:37}并[阿拉伯文版第82页]说：\Quote{这个\Emph{人}所作的一切都好：他使聋子听见，使哑巴说话。}\par

8 \BibleRef{若 4:4}他正经过撒玛黎雅地时，9 \BibleRef{若 4:5}来到一座撒玛黎雅城，名叫息哈尔，靠近雅各伯给他儿子若瑟的那块地。10 \BibleRef{若 4:6}那里有雅各伯泉。耶稣因行路疲乏，就坐在泉旁；那时大约是第六时辰。11 \BibleRef{若 4:7}有一个撒玛黎雅妇人来打水；耶稣对她说：\Quote{给我水喝，我要喝。}12 \BibleRef{若 4:8}原来他的门徒已进城买食物去了。13 \BibleRef{若 4:9}那撒玛黎雅妇人就对他说：\Quote{你既是犹太人，怎么向我一个撒玛黎雅妇人要水喝呢？}（原来犹太人和撒玛黎雅人是不往来的。）14 \BibleRef{若 4:10}耶稣回答她说：\Quote{若你知道天主的恩赐，并知道是谁对你说\InnerQuote{给我水喝}，你必早求他，他也必早给了你生命之水。}15 \BibleRef{若 4:11}那妇人对他说：\Quote{主，你没有打水器具，井又深，你从哪里得生命之水呢？16 \BibleRef{若 4:12}难道你比我们的祖先雅各伯还大吗？他给了我们这口井，他自己、他的儿子和他的牲畜都喝这井里的水。}17 \BibleRef{若 4:13}耶稣回答她说：\Quote{凡喝这水的，还要再渴；18 \BibleRef{若 4:14}但谁喝了我所赐的水，就永远不渴；我所赐的水要在人内成为涌到永生的水泉。}19 \BibleRef{若 4:15}那妇人对他说：\Quote{主，请给我这水，使我不再渴，也免得我到这里来打水。}20 \BibleRef{若 4:16}耶稣对她说：21 [阿拉伯文版第83页]\Quote{你去叫你丈夫，然后到这里来。}22 \BibleRef{若 4:17}那妇人回答说：\Quote{我没有丈夫。}23 \BibleRef{若 4:18}耶稣对她说：\Quote{你说\InnerQuote{我没有丈夫}说得不错；因为你曾经有过五个丈夫，你现在有的也不是你的丈夫；这话你说得真对。}24 \BibleRef{若 4:19}那妇人对他说：\Quote{主，我看出你是一位先知。25 \BibleRef{若 4:20}我们的祖先在这座山上朝拜，而你们却说，应该朝拜的地方是在耶路撒冷。}26 \BibleRef{若 4:21}耶稣对她说：\Quote{女人，你相信我罢！到了时候，你们朝拜圣父，既不在这座山上，也不在耶路撒冷。27 \BibleRef{若 4:22}你们朝拜你们所不认识的，我们朝拜我们所认识的，因为救恩出自犹太人。28 \BibleRef{若 4:23}然而时候要到，且现在就是，那些真正朝拜的人，将以心神和真理朝拜圣父；因为圣父正是寻找这样朝拜他的人。29 \BibleRef{若 4:24}天主是神，朝拜他的人，必须以心神和真理朝拜他。}30 \BibleRef{若 4:25}那妇人对他说：\Quote{我知道默西亚要来；\BibleRef{若 4:26}他一来了，必会告诉我们一切。}耶稣对她说：\Quote{我对你说话的，我就是。}\par

31 \BibleRef{若 4:27}正在说话的时候，他的门徒来了；他们惊异他为什么与一妇人谈话，但没有人问他说：\Quote{你找什么？}或：\Quote{你为什么同她说话？}32 \BibleRef{若 4:28}那妇人就留下水罐，往城里去，对众人说：33 \BibleRef{若 4:29}\Quote{你们来看！有一个人把我素来所行的一切事都给我说出来了。莫非这就是默西亚吗？}34 \BibleRef{若 4:30}众人就出城，往耶稣那里去。35 \BibleRef{若 4:31}同时，门徒请求耶稣说：\Quote{老师，请用饭。}36 \BibleRef{若 4:32}耶稣对他们说：\Quote{我有食物吃，是你们不知道的。}37 \BibleRef{若 4:33}门徒便彼此问说：\Quote{难道有人给他送来了吃的吗？}38 \BibleRef{若 4:34}耶稣对他们说：\Quote{我的食物就是承行派遣我者的旨意，完成他的工作。39 [阿拉伯文版第84页]\BibleRef{若 4:35}你们不是说：还有四个月才到收获期吗？看，我告诉你们：你们举目向田观看，庄稼已经发白，可以收割了。40 \BibleRef{若 4:36}收割的人已领到工钱，并为永生收集果实，使栽种的和收割的一同欢乐。41 \BibleRef{若 4:37}因为俗话说：\InnerQuote{这人栽种，那人收割}，这话是真的。42 \BibleRef{若 4:38}我派遣你们去收割你们没有劳苦的；别人劳苦了，你们却享受他们劳苦的成果。}\par

43 \BibleRef{若 4:39}那城里有许多撒玛黎雅人信从了耶稣，因为那妇人作证说：\Quote{他把我\Emph{曾经}所做的一切事都给我说出来了。}44 \BibleRef{若 4:40}于是，撒玛黎雅人来见他，请求他住在他们那里；他便在那里住了两天。45 \BibleRef{若 4:41}因着他的话，信从的人就更多了。46 \BibleRef{若 4:42}他们便对那妇人说：\Quote{现在我们信，不是因为你的话，而是因为我们亲自听见了，知道这真是世界的救世主——默西亚。}\par

47, 48 \BibleRef{若 4:43} 两天以后，耶稣从那里出来，往加里肋亚去。\BibleRef{若 4:44} 耶稣自己作证说：先知在自己的本乡是不受尊敬的。\BibleRef{若 4:45} 及至他来到加里肋亚，加里肋亚人接待了他。\par

\SourceRecord{Translated by Hope W. Hogg.；Ante-Nicene Fathers；Vol. 9.；Edited by Allan Menzies.；Buffalo, NY: Christian Literature Publishing Co.,；1896.；<http://www.newadvent.org/fathers/100221.htm>.}

% structure: p0001 source=h1 target=section reason=page-title

\section{\WorkTitle{达太福音（第二十二节）}}

1 \BibleRef{路 5:12} 当耶稣来到一个村庄时，有一个癞病人走近他，俯伏在他脚前，恳求他说：\Quote{你若愿意，就能洁净我。} 2 \BibleRef{谷 1:41} 耶稣动了怜悯的心，就伸手抚摸他，说：\Quote{我愿意，你洁净了吧。} 3 \BibleRef{谷 1:42} 癞病立刻离开了他，他就洁净了。 4 \BibleRef{谷 1:43} 耶稣严厉地嘱咐他，立即催他走， 5 \BibleRef{谷 1:44} 对他说：\Quote{你当谨慎，\Emph{不可}告诉任何人；但去叫司铎检验你，并为你的洁净献上梅瑟所规定的祭物，作为证据。} 6 \BibleRef{谷 1:45} 但他一出去，便大肆传扬，到处散布这事，以致耶稣不能再公开进入任何城，因为他的名声传得太广，他却留在外面荒野的地方。 7 \BibleRef{路 5:15} 有许多人从各处来到他那里，要听他的道，并希望他们的病得医治。 8 \BibleRef{路 5:16} 他却常退到荒野去祈祷。\par

9 此后，到了犹太人的节期，耶稣就上耶路撒冷去。\par

10 \BibleRef{若 5:2} 在耶路撒冷有一个沐浴的地方，希伯来语称为\Quote{慈悲之家}，周围有五个廊子。 11 \BibleRef{若 5:3} 在这些廊子里躺着许多病人，有瞎眼的、瘸腿的、瘫痪的，等候水动。 12 \BibleRef{若 5:4} 因为有天使按时下到池子里搅动那水，水动之后，谁先下去，无论患什么病，都会痊愈。 13 \BibleRef{若 5:5} 那里有一个人，患病已三十八年。 14 \BibleRef{若 5:6} 耶稣看见那\Emph{人}躺着，知道他病了许久，就问他说：\Quote{你愿意痊愈吗？} 15 \BibleRef{若 5:7} 那病人回答说：\Quote{主啊，水动的时候，没有人把我放到池子里；等我到了，别人已经先下去了。} 16 \BibleRef{若 5:8} 耶稣对他说：\Quote{起来，拿起你的床，行走吧！} 17 \BibleRef{若 5:9} 那人立刻痊愈了，就拿起床来行走。\par

18 \BibleRef{若 5:10} 那天正是安息日。犹太人看见那痊愈的人，就对他说：\Quote{今天是安息日，你不可拿床。} 19 \BibleRef{若 5:11} 他回答说：\Quote{使我痊愈的那位对我说：\InnerQuote{拿起你的床，行走吧。}} 20 \BibleRef{若 5:12} 他们便问他说：\Quote{对你说\InnerQuote{拿起床行走}的那人是谁？} 21 \BibleRef{若 5:13} 但那痊愈的人不知道是谁，因为耶稣已从那里离开，到了别处，因那地方人多拥挤。 22 \BibleRef{若 5:14} 过了两天，耶稣在圣殿里遇见他，对他说：\Quote{你已经痊愈了，不要再犯罪，免得遭遇更坏的事。} 23 \BibleRef{若 5:15} 那人就去告诉犹太人，是耶稣使他痊愈的。 24 \BibleRef{若 5:16} 因此犹太人迫害耶稣，想要杀他，因为他在安息日做这事。 25 \BibleRef{若 5:17} 耶稣对他们说：\Quote{我父一直工作到现在，我也工作。} 26 \BibleRef{若 5:18} 因此犹太人越发想要杀他，不仅因为他犯了安息日，而且因为他称天主为父，使自己与天主平等。 27 \BibleRef{若 5:19} 耶稣回答他们说：\Quote{我实实在在告诉你们：子不能凭自己做什么，只有看见父所做的，子才能做；父所做的，子也照样做。 28 \BibleRef{若 5:20} 父爱子，将自己所做的一切指示给他，并且还要将比这些更大的工作指示给他，使你们惊奇。 29 \BibleRef{若 5:21} 父怎样叫死人复活，赐给他们生命，子也照样随自己的意愿赐人生命。 30 \BibleRef{若 5:22} 父不审判任何人，却将审判权完全交给了子， 31 \BibleRef{若 5:23} 为叫众人尊敬子，如同尊敬父一样。不尊敬子的，就是不尊敬派遣他的父。 32 \BibleRef{若 5:24} 我实实在在告诉你们：那听我的话，并相信派遣我来的，就有永生，不受审判，而已出死入生。 33 \BibleRef{若 5:25} 我实实在在告诉你们：时候要到，而且现在就是，死人要听见天主子的声音，听见的就必存活。 34 \BibleRef{若 5:26} 父怎样在自己内有生命，也照样赐给子在自己内有生命， 35 \BibleRef{若 5:27} 并赐给他行审判的权柄，因为他是人子。 36 \BibleRef{若 5:28} 你们不要为这事惊奇；我是说，时候要到，凡在坟墓里的，都要听见他的声音， 37 \BibleRef{若 5:29} 就出来：行善的，复活进入生命；作恶\Emph{的事}的，复活进入审判。}\par

\BibleRef{若 5:30}\BibleQuote{38 我自己不能做什么；但我怎样听见，就怎样审判：我的审判39是公正的；我不寻求我\Emph{自己的}意愿，而寻求那派遣我来的那位的意愿。}\BibleRef{若 5:31}\BibleQuote{我为自己作证40，所以我的见证不真实。}\BibleRef{若 5:32}\BibleQuote{另一位为我作证41，并且我知道他为我所作的见证是真实的。}\BibleRef{若 5:33}\BibleQuote{你们曾派人42到若翰那里，他为真理作了见证。}\BibleRef{若 5:34}\BibleQuote{但我并不从人寻求43见证；我说这话，是要你们得救。}\BibleRef{若 5:35}\BibleQuote{他是盏灯，点着且44发光；你们情愿暂时在他的光中夸耀。}\BibleRef{若 5:36}\BibleQuote{但我有比若翰更大的见证：圣父交给我要我完成的工程45，就是我所行的这些工程，为我作证：证明圣父派遣了我。}\BibleRef{若 5:37}\BibleQuote{派遣我的圣父，他已为我作了见证。你们从未听过他的46声音，也从未见过他的容貌。}\BibleRef{若 5:38}\BibleQuote{并且他的话不在你们内，因为47你们不相信他所派遣的那一位。}\BibleRef{若 5:39}\BibleQuote{你们查考圣经，因你们以为48其中有永生；}\BibleRef{若 5:40}\BibleQuote{这些圣经为我作证；但你们却不愿意来到49, 50 [Arabic, p. 88] 我这里，以获得永生。}\BibleRef{若 5:41}\BibleQuote{我不求人的光荣。}\BibleRef{若 5:42}\BibleQuote{但我认识51你们，知道你们里面没有天主的爱。}\BibleRef{若 5:43}\BibleQuote{我因我圣父的名而来，你们却不接待我；但若另一个人凭自己的名而来，你们才要接待那\Emph{人}52。}\BibleRef{若 5:44}\BibleQuote{你们既然彼此接受光荣，也不求从唯一的天主而来的53光荣，怎能相信呢？}\BibleRef{若 5:45}\BibleQuote{难道你们以为我要在圣父面前控告你们吗？有一位控告你们的，就是你们所信赖的梅瑟54。}\BibleRef{若 5:46}\BibleQuote{假如你们相信梅瑟，你们也会相信我；因为他曾论及我写过。55}\BibleRef{若 5:47}\BibleQuote{如果你们不信他的著作，怎能相信我的话呢？}\par

\SourceRecord{Translated by Hope W. Hogg.；Ante-Nicene Fathers；Vol. 9.；Edited by Allan Menzies.；Buffalo, NY: Christian Literature Publishing Co.,；1896.；<http://www.newadvent.org/fathers/100222.htm>.}

% structure: p0001 source=h1 target=section reason=page-title

\section{\WorkTitle{四福音合参}（第23部分）}

1 \BibleRef{玛 15:29}耶稣从那里离开，来到加里肋亚海边，上了山，坐在那里。2 \BibleRef{玛 15:30}有许多群众到他那里，带着瘸子、瞎子、哑巴、残废的，和许多别的病人，都放在耶稣脚前；3 \BibleRef{若 4:45}因为他们见过他在耶路撒冷过节时所行的一切奇迹。4 \BibleRef{玛 15:30}他便治好了他们。\BibleRef{玛 15:31}群众看见哑巴的\Emph{人}说话，残废的\Emph{人}痊愈，瘸子的\Emph{人}行走，瞎子的\Emph{人}看见，都非常惊奇，并赞美以色列的天主。\par

5 \BibleRef{玛 15:32}耶稣叫了他的门徒来，对他们说：\Quote{我怜悯这群众，因为他们同我在一起已经三天，没有什么吃的；我不愿打发他们饿着回去，恐怕他们在路上晕倒。}6 \BibleRef{谷 8:3}他们中有些\EditorialNote{阿拉伯文抄本，第89页}是从远处来的。7 \BibleRef{玛 15:33}门徒对他说：\Quote{在这荒野里，我们从哪里得这么多的饼，使这群众吃饱呢？}8 \BibleRef{玛 15:34}耶稣问他们说：\Quote{你们有多少饼？}9 \BibleRef{玛 15:35}他们说：\Quote{七个，还有几条小鱼。}他便吩咐群众坐在地上；10 \BibleRef{玛 15:36}他拿起那七个饼和鱼，祝谢了，擘开，递给门徒，门徒再分给群众。11 \BibleRef{玛 15:37}众人都吃了，并且吃饱了；把剩下的碎块收集起来，装满了七篮子。12 \BibleRef{玛 15:38}吃的人，除了妇女和孩子，约有四千。13 \BibleRef{玛 15:39}耶稣遣散群众后，便上船，来到玛加丹境内。\par

14 法利塞人和撒杜塞人来到他那里，开始与他争辩，求他从天上显个记号给他们，试探他。15 耶稣心里叹息说：\Quote{这邪恶淫乱的世代寻求什么记号？它寻求记号，但必不给它记号，除了先知约纳的记号。}16 \BibleRef{谷 8:12}我实在告诉你们，这世代必得不到记号。17 \BibleRef{谷 8:13}他就离开他们，上船往对岸去了。\par

18 \BibleRef{谷 8:14}门徒忘了带饼，船上除了一个饼外，再没有别的。19 \BibleRef{谷 8:15}耶稣嘱咐他们说：\Quote{你们要谨慎，提防法利塞人和撒杜塞人的酵母，以及黑落德的酵母。}20 \BibleRef{玛 16:7}他们就彼此议论说：\Quote{这是因为我们没有带饼。}21 \BibleRef{玛 16:8}耶稣知道了，就说：\Quote{小信德的人哪！为什么你们彼此议论没有带饼的事呢？22 \BibleRef{谷 8:17}到现在你们还不明白，还不领悟吗？你们的心还是迟钝吗？23 \BibleRef{谷 8:18}你们有眼睛，却\Emph{看不见}吗？\EditorialNote{阿拉伯文抄本，第90页}你们有耳朵，却\Emph{听不见}吗？24 \BibleRef{谷 8:19}你们不记得吗？我擘开五个饼给五千人，你们收了多少篮子的碎块？}他们说：\Quote{十二篮。}25 \BibleRef{谷 8:20}他又问：\Quote{那七个饼给四千人，你们收了多少篮子的碎块？}他们说：\Quote{七篮。}26 他问他们说：\Quote{怎么还不明白呢？我不是对你们说，不是为了饼的事，而是要提防法利塞人和撒杜塞人的酵母吗？}27 \BibleRef{玛 16:12}他们这才领悟，他不是叫他们提防饼的酵母，而是提防法利塞人和撒杜塞人的教训，他把这个称为酵母。\par

28 \BibleRef{谷 8:22}他们来到贝特赛达。有人带了一个瞎子的\Emph{人}来，求耶稣抚摸他。29 \BibleRef{谷 8:23}耶稣拉着瞎子的手，领他到村外，吐唾沫在他的眼睛上，把手按在他身上，问他说：\Quote{你看见什么？}30 \BibleRef{谷 8:24}他抬头一看，说：\Quote{我看见人们，好像树木在行走。}31 \BibleRef{谷 8:25}耶稣又把手按在他的眼睛上，他的眼睛就开了，什么都看得清清楚楚了。32 \BibleRef{谷 8:26}耶稣打发他回家去，说：\Quote{不要进村里去，也不要告诉村里任何人。}\par

\BibleRef{谷 8:27} 耶稣出去，同他的门徒，往斐理伯的凯撒勒雅一带的村庄去。\BibleRef{玛 16:13} 在路上，只有他的门徒和他同走的时候，他问他的门徒说：\BibleQuote{人们说我是谁？人子\OriginalTerm{人子}{Son of man}是谁？}\BibleRef{玛 16:14} 他们回答说：\BibleQuote{有人说是洗者若翰；有人说是厄里亚；也有人说是耶肋米亚，或是先知中的一位。}\BibleRef{玛 16:15} 耶稣对他们说：\BibleQuote{你们说我是谁？}\BibleRef{玛 16:16} 西满·刻法\OriginalTerm{刻法}{Cephas}回答说：\EditorialNote{阿拉伯语，第91页}\BibleQuote{你是默西亚\OriginalTerm{默西亚}{Messiah}，永生天主之子。}\BibleRef{玛 16:17} 耶稣回答他说：\BibleQuote{约纳之子西满，你是有福的，因为不是血肉将\Emph{这事}启示给你，而是我在天之父。\BibleRef{玛 16:18} 我也告诉你，你是刻法，我要在这磐石上建立我的教会；阴府的门决不能战胜她。\BibleRef{玛 16:19} 我要把天国的钥匙交给你：凡你在地上束缚的，在天上也要被束缚；凡你在地上释放的，在天上也要被释放。}\BibleRef{玛 16:20} 于是他严格嘱咐门徒，不要对任何人说他就是默西亚。\BibleRef{玛 16:21} 从那时起，耶稣开始向门徒说明，他必须往耶路撒冷去，\BibleRef{谷 8:31} 要受许多苦，被长老、司祭长和经师弃绝，并且被杀，但第三天要复活。\BibleRef{谷 8:32} 耶稣明明地讲这话。\BibleRef{玛 16:22} 西满·刻法便拉他到一边，谏责他说：\BibleQuote{主，千万不可！这事绝不会临到你身上！}\BibleRef{谷 8:33} 耶稣转过身来，看着他的门徒，斥责西满说：\BibleRef{玛 16:23}\BibleQuote{撒殚\OriginalTerm{撒殚}{Satan}，退到我后面去！你是我的绊脚石，因为你所体会的，不是天主的事，而是人的事。}\par

耶稣召集群众和门徒，对他们说：\BibleQuote{谁若愿意跟随我，该弃绝自己，天天背着自己的十字架，跟随我。}\BibleRef{谷 8:35}\BibleQuote{因为谁若愿意救自己的性命，必要丧失性命；但谁若为我和福音的缘故，丧失自己的性命，必要救得性命。}\BibleRef{路 9:25}\BibleQuote{人纵然赚得了全世界，却赔上自己的灵魂，为他有什么益处？}\BibleRef{谷 8:37}\BibleQuote{或者，人还能拿什么\Emph{作为}自己灵魂的代价？}\EditorialNote{阿拉伯语，第92页}\BibleRef{谷 8:38}\BibleQuote{谁若在这淫乱罪恶的世代中，以我和我的话为耻，将来人子也要在他父的光荣中，同他的圣天使降来时，以他为耻。}\BibleRef{玛 16:27}\BibleQuote{因为人子将要在他父的光荣中同他的天使降来，那时他要按照每人的行为予以报应。}\par

\SourceRecord{Translated by Hope W. Hogg.；Ante-Nicene Fathers；Vol. 9.；Edited by Allan Menzies.；Buffalo, NY: Christian Literature Publishing Co.,；1896.；<http://www.newadvent.org/fathers/100223.htm>.}

% structure: p0001 source=h1 target=section reason=page-title

\section{四福音合参（第24章）}

1 祂对他们说：\Quote{\BibleQuote{我实在告诉你们，这里现在有一些站着的人，在未见到天主的国带着能力来临以前，决不会尝到死亡，} \BibleRef{谷 9:1} \BibleQuote{和那在自己的国中来临的人子。} \BibleRef{玛 16:28}}\par

2 过了六天，耶稣带着西满\OriginalTerm{刻法}{Cephas}、雅各伯和他的兄弟若望，3 领他们单独上了高山。\BibleRef{路 9:29}他们祈祷时，4 耶稣改变了，成为另一个人的样式；5 祂的面貌发光有如太阳，祂的衣裳洁白如雪，好像闪电的光，以致世上没有东西能这样漂白。\BibleRef{路 9:31}6 梅瑟和厄里亚显现，与耶稣谈话。\BibleRef{路 9:32}7 西满和同他在一起的人因困倦而沉睡；他们努力醒过来，看见祂的荣耀，以及那两位与祂站在一起的人。\BibleRef{路 9:33}8 当他们开始离开祂时，西满对耶稣说：我的老师，我们在这里真好；\BibleRef{玛 17:4}9 如果你愿意，我们在这里搭三个帐棚，一个为你，一个为梅瑟，一个为厄里亚。\BibleRef{路 9:33}10 由于占据他们的恐惧，他并不知道自己在说什么。\BibleRef{路 9:34}11 他正说这话时，一片明亮的云彩遮蔽了他们。\BibleRef{玛 17:5}12 有声音从云中说：这是我的爱子，我所拣选的；你们要听祂。\BibleRef{路 9:36}13 这声音听到后，只见耶稣单独一人。\BibleRef{玛 17:6}14 门徒们听见这声音，因恐惧而俯伏在地。\BibleRef{玛 17:7}15 耶稣前来抚摸他们说：起来，不要害怕。\BibleRef{玛 17:8}16 他们举目观看，只见耶稣如常。\par

17 \BibleRef{玛 17:9}当他们下山时，耶稣嘱咐他们说：不要将所见告诉任何人，直到人子从死人中复活。18 他们把这话存在心里，在那些日子没有向任何人提及所见之事。19 \BibleRef{谷 9:10}他们彼此思忖：祂对我们说的\Quote{当我从死人中复活时}这句话是什么意思？20 门徒问祂说：那么经师们说厄里亚必须先来，是什么意思？21 \BibleRef{谷 9:12}祂对他们说：厄里亚固然要先来，重整一切；正如经上论到人子所说，祂要受许多苦，并被轻视。22 \BibleRef{谷 9:13}但我告诉你们：厄里亚已经来了，他们却不认识他，反而任意对待了他，正如经上关于他所写的。23、24 \BibleRef{玛 17:12}同样，人子也要受他们的苦。\BibleRef{玛 17:13}门徒这才明白祂是指着洗者若翰说的。\par

25 \BibleRef{谷 9:14}在他们下山的那一天，有一大群人遇见祂，与祂的门徒站在一起，经师们正在和他们辩论。26 \BibleRef{谷 9:15}众人看见耶稣，都惊讶，急忙上前问候祂。27 \BibleRef{路 13:31}那一天，有几个法利塞人来对祂说：你出去，离开这里，因为黑落德要杀你。28 \BibleRef{路 13:32}耶稣对他们说：你们去告诉那只狐狸：看，我今天明天驱魔治病，第三天便成全了。29 \BibleRef{路 13:33}然而今天我明天必须谨慎，最后一天我要离去；因为先知在耶路撒冷之外丧命是不可能的。\par

30 之后，人群中有一个人来见祂，跪下说：我的主，求祢垂顾我的儿子；\BibleRef{路 9:38}他是我的独生\Emph{孩子}；不洁之神突然附上他。他得了癫痫，遭受痛苦。31 \BibleRef{谷 9:18}不洁之神一附上他，就打他；他口吐白沫，咬牙切齿，身体消瘦；多次把他扔进水里、火里，要毁灭他，且几乎不离开他，而使他受伤。32 \BibleRef{玛 17:16}我把他带到祢的门徒那里，他们却不能治好他。33 \BibleRef{玛 17:17}耶稣回答说：唉！无信而败坏的世代，我同你们在一起要到几时？我容忍你们要到几时？把孩子带到这里来。34 \BibleRef{谷 9:20}他们就把孩子带到祂跟前；不洁之神一见祂，立刻使他抽搐，倒在地上，打滚吐沫。35 \BibleRef{谷 9:21}耶稣问他的父亲：他这样有多久了？他回答说：从幼年直到如今。36 \BibleRef{谷 9:22}但主啊，求祢帮助我，怜悯我。37 \BibleRef{谷 9:23}耶稣对他说：你若能信！为信的人，一切都是可能的。38 \BibleRef{谷 9:24}孩子的父亲立刻喊着说：我信，我的主！求祢补足我的不信。39 \BibleRef{谷 9:25}耶稣看见群众都跑拢来，就斥责那不洁之神，说：你这聋哑的鬼，我命令你从他身上出来，再不要进去。40 \BibleRef{谷 9:26}那魔鬼大喊一声，使他剧烈地抽搐，就出来了；那孩子好像死了一样，以致许多人说：他死了。41 但耶稣握住他的手，扶他起来，把他交给他的父亲；那孩子从那时起就好了。42 众人都惊讶天主的伟大。\par

45 \BibleRef{谷 9:28} 耶稣进了屋子，门徒私下问他说：\Quote{为什么我们不能赶出他呢？}\BibleRef{玛 17:20} 耶稣对他们说：\Quote{因为你们的\OriginalTerm{不信}{unbelief}。我实在告诉你们：如果你们有像一粒芥菜子那样的\OriginalTerm{信德}{faith}，就是对这座山说：\InnerQuote{从这里移到那里！}它也会移开；并且没有一件事能难倒你们。\BibleRef{谷 9:29} 但这一类的，非用禁食和祈祷，是不能赶出去的。}\par

48 \BibleRef{谷 9:30} 他们从那里出去，经过加里肋亚；耶稣不愿意任何人知道，\BibleRef{谷 9:31} 因为他教训门徒，对他们说：\Quote{你们要把这话存在耳中和心里：人子将要被交在人手里，他们要杀害他；被杀以后，第三天必要复活。}\BibleRef{路 9:45} 但他们不明白这话，因为这话是隐藏的，不让他们领悟；他们又怕问这话的意思。\BibleRef{玛 17:23} 他们就非常忧愁。\par

\SourceRecord{Translated by Hope W. Hogg.；Ante-Nicene Fathers；Vol. 9.；Edited by Allan Menzies.；Buffalo, NY: Christian Literature Publishing Co.,；1896.；<http://www.newadvent.org/fathers/100224.htm>.}

% structure: p0001 source=h1 target=section reason=page-title

\section{\WorkTitle{迪亚特萨隆（第25节）}}

1 \BibleRef{路 9:46} 在那一天，这个想法出现在他的门徒们心中，他们想，他们中间究竟谁最大。2 \BibleRef{谷 9:33} 当他们来到\Place{葛法翁}，进了屋子，\Person{耶稣}对他们说：\Quote{你们在路上彼此议论了什么？}3 \BibleRef{谷 9:34} 他们就默不作声，因为他们曾议论那件事。\par

4 \BibleRef{玛 17:24} 当\Person{西满}出去时，那些收两个迪尔汗作税的人来到\Person{刻法}面前，对他说：\Quote{你的老师不纳两个迪尔汗吗？}5 他对他们说：\Quote{是。}\BibleRef{玛 17:25} 当\Person{刻法}进了屋子，\Person{耶稣}抢先对他说：\Quote{\Person{西满}，你以为怎样？地上的君王从谁那里征收关税和丁税？是从自己的儿子，还是从外人？}6 \BibleRef{玛 17:26} \Person{西满}对他说：\Quote{从外人。}\Person{耶稣}对他说：\Quote{那么儿子是自由的。}\Person{西满}对他说：\Quote{是。}\Person{耶稣}对他说：\Quote{也像外人一样给他们。}\BibleRef{玛 17:27} \BibleQuote{但为了避免触犯他们，你去海边钓鱼；先钓上来的鱼，打开它的口，你会找到一个斯塔特；拿那个，替我和你自己付了。}[Arabic, p. 97]\par

8 \BibleRef{玛 18:1} 就在那时，门徒们来到\Person{耶稣}面前，对他说：\Quote{你以为天国里谁是最大的？}9 \Person{耶稣}知道他们心里的意念，就叫了一个小孩子来，让他站在他们中间，又抱起他，10 对他们说：\BibleRef{玛 18:3} \BibleQuote{我实在告诉你们，你们若不回转，变成小孩子的样式，断不得进天国。}11 \BibleRef{路 9:48} \BibleQuote{凡因我的名接待这样一个孩子的，就是接待我；}\BibleRef{谷 9:37} \BibleQuote{凡接待我的，不是接待我，乃是接待那差我来的。}12 \BibleRef{路 9:48} \BibleQuote{你们中间最小的，他才是最大的。}13 \BibleRef{玛 18:6} \BibleQuote{凡使这信我的一个小子跌倒的，倒不如把大磨石拴在他的颈项上，沉在深海里。}\par

14 \BibleRef{路 9:49} \Person{若望}回答说：\Quote{老师，我们看见一个人奉你的名赶鬼；我们就阻止他，因为他不与我们一同跟从你。}15 \BibleRef{谷 9:39} \Person{耶稣}对他们说：\Quote{不要阻止他；因为没有人奉我的名行异能，反倒会急忙毁谤我。}16, 17 \BibleRef{路 9:50} \BibleQuote{不反对你们的，就是帮助你们的。祸哉，世界因试探而受祸！但祸哉，那试探由他手中而来的人！}[Arabic, p. 98] 18 \BibleRef{玛 18:8} \BibleQuote{如果你的一只手或一只脚使你跌倒，就砍下来丢掉；你残废或瘸腿进入永生，强如有两只手或两只脚，被丢进那永远燃烧的地狱的火里；}19, 20 \BibleRef{谷 9:44} \BibleQuote{那里的虫是不死的，火是不灭的。}21 \BibleRef{玛 18:9} \BibleQuote{如果你的一只眼使你跌倒，就把它剜出来丢掉；}\BibleRef{谷 9:47} \BibleQuote{你只有一只眼进入天主的国，强如有两只眼被丢进\OriginalTerm{革希纳}{Gehenna}的火里；}22, 23 \BibleRef{谷 9:48} \BibleQuote{那里的虫是不死的，火是不灭的。}24 \BibleRef{谷 9:49} \BibleQuote{每个人必被火盐腌，每样祭物必用盐腌。}25 \BibleRef{谷 9:50} \BibleQuote{盐本是好的；盐若失了味，可用什么叫它再咸呢？它既不适合土地，也不适合粪堆，只好丢在外面。有耳可听的，就应当听。}26 \BibleRef{谷 9:50} \BibleQuote{你们里头应当有盐，彼此和平相处。}\par

27 \BibleRef{谷 10:1} 他从那里起身，来到\Place{犹太}境界，\Place{约但}河外；有许多人聚集到他那里，他又照常教训他们。28 \BibleRef{谷 10:2} 有\Person{法利塞人}来问他说：\Quote{人休妻可以不可以？}意思是要试探他。29 \BibleRef{谷 10:3} 他回答说：\Quote{\Person{梅瑟}吩咐你们的是什么？}30 \BibleRef{谷 10:4} 他们说：\Quote{\Person{梅瑟}许人写了休书便可以休妻。}31 \BibleRef{谷 10:5} \Person{耶稣}对他们说：\BibleRef{玛 19:4} \BibleQuote{你们没有念过吗？那起初造人的，是造男造女，}32 并且说：\BibleRef{玛 19:5} \BibleQuote{因此，人要离开父母，与妻子连合，二人成为一体。}33 \BibleRef{玛 19:6} \BibleQuote{这样，他们不再是两个人，乃是一体的了。所以，天主所配合的，人不可分开。}34 \BibleRef{玛 19:7} \Person{法利塞人}对他说：\Quote{那么，\Person{梅瑟}为什么吩咐给妻子休书就可以休她呢？}35 \BibleRef{玛 19:8} \Person{耶稣}说：\Quote{\Person{梅瑟}因为你们心硬，所以许你们休妻，但起初并不是这样。}36 \BibleRef{玛 19:9} \BibleQuote{我告诉你们，凡休妻另娶的，若不是为淫乱的缘故，就是叫她犯奸淫了。}37 \BibleRef{谷 10:10} 门徒在屋子里又问他这事。38 \BibleRef{谷 10:11} 他对他们说：\Quote{凡休妻另娶的，就是叫她犯奸淫；}39 \BibleRef{谷 10:12} \BibleQuote{妻子若离弃丈夫另嫁，也是犯奸淫了。}\BibleRef{玛 19:9} \BibleQuote{凡娶这被休的妇人的，也是犯奸淫了。}40 \BibleRef{玛 19:10} 门徒对他说：\Quote{人和妻子既是这样，倒不如不娶。}41 \BibleRef{玛 19:11} 他对他们说：\Quote{这话不是人都能领受的，惟独赐给谁，谁才能领受。}42 \BibleRef{玛 19:12} 因为有\OriginalTerm{阉人}{eunuch}从母腹生来就是这样；有阉人是被人阉的；并有阉人为天国的缘故自阉的。\Quote{这话谁能领受，就可以领受。}[Arabic, p. 99]\par

43 \BibleRef{玛 19:13} 随后他们带了一些孩童到他那里，为叫他按手在他们身上并44祈祷：而门徒们却责备那些带孩童来的人。\Person{耶稣}看见了，这事使他感到难过；于是对他们说：\BibleQuote{让孩童\EditorialNote{阿拉伯文，第100页}到我这里来，不要阻止他们，因为那些像这样的人拥有\OriginalTerm{天主的国}{Kingdom of God}。} \BibleRef{谷 10:15} \BibleQuote{我实在告诉你们：谁若不像这个孩童一样接受天主的国，决不能进去。} \BibleRef{谷 10:16} 于是他抱起他们，按手在他们身上，祝福了他们。\par

\SourceRecord{Translated by Hope W. Hogg.；Ante-Nicene Fathers；Vol. 9.；Edited by Allan Menzies.；Buffalo, NY: Christian Literature Publishing Co.,；1896.；<http://www.newadvent.org/fathers/100225.htm>.}

% structure: p0001 source=h1 target=section reason=page-title

\section{四福音合参（第26节）}

1, 2 \BibleRef{路 15:1} 有税吏和罪人到耶稣跟前来听他的话。\BibleRef{路 15:2} 法利塞人和经师窃窃私议，说：\BibleQuote{这个人接待罪人，又同他们一起吃饭。}\BibleRef{路 15:3} 耶稣看见他们的议论，就对他们讲了这个比喻：\BibleRef{路 15:4} \BibleQuote{你们中间谁有一百只羊，失去了一只，而不把这九十九只丢在荒野，去寻觅那迷失的\Emph{一只}，直到找着呢？\BibleRef{玛 18:13} 我实在告诉你们：他一旦找着了，就为这一只羊欢喜，\Emph{胜过}为那没有迷失的九十九只；\BibleRef{路 15:5} 他把它放在肩上，带回家，并请朋友和邻居来，\BibleRef{路 15:6} 对他们说：\InnerQuote{与我一同欢乐吧！因为我那只迷失的羊已经找着了。}\BibleRef{玛 18:14} 同样，你们在天上的父也不愿意这些迷失的小子中有一个丧亡，而是寻找他们，使他们悔改。\BibleRef{路 15:7} 我告诉你们：同样，对于一个罪人悔改，在天上也要为他欢喜，\Emph{胜过}为那九十九个不需要悔改的\Emph{义人}。}\par

9 \BibleRef{路 15:8} 一个女人如果有十块\OriginalTerm{德拉克马}{drachma}，丢失了一块，岂不点上灯，打扫房屋，细心寻找，直到找着吗？\BibleRef{路 15:9} 找着了，就请朋友和邻居来，对他们说：\BibleQuote{与我一同欢乐吧！因为我丢失的那块德拉克马已经找着了。}\BibleRef{路 15:10} 我告诉你们：同样，在天上要有喜乐〔阿拉伯文抄本，页101〕，在天主的众天使面前，为一个悔改的罪人，\Emph{胜过}为那九十九个不需要悔改的\Emph{义人}。\par

12, 13 \BibleRef{路 15:11} 耶稣又对他们讲了一个比喻：\BibleRef{路 15:12} 一个人有两个儿子：小儿子对父亲说：\BibleQuote{父亲，请把我应得的家产分给我。}\BibleRef{路 15:13} 父亲就把家产分给了他们。过了不多几天，小儿子把所有的一切都收拾起来，往远方去了，在那里放荡不羁，挥霍了他的资财。\BibleRef{路 15:14} 当他耗尽了一切，那地方又发生了大饥荒。\BibleRef{路 15:15} 他在穷困中，去投靠当地的一个居民；那个\Emph{人}打发他到田里去放猪。\BibleRef{路 15:16} 他恨不得拿猪吃的\OriginalTerm{角豆}{carob}来填饱肚子，可是没有人给他。\BibleRef{路 15:17} 他醒悟过来，就说：\BibleQuote{我父亲有多少佣工，口粮丰足，我在这里却要饿死！\BibleRef{路 15:18} 我要起身到我父亲那里去，对他说：\InnerQuote{父亲！我得罪了天，也得罪了你。\BibleRef{路 15:19} 从今以后，我不配称为你的儿子，把我当作你的一个佣工吧。}}\BibleRef{路 15:20} 于是他起身往他父亲那里去。他离得还远的时候，父亲看见了他，就动了怜悯的心，跑上前去，扑到他的脖子上，亲吻他。\BibleRef{路 15:21} 儿子对他说：\BibleQuote{父亲！我得罪了天，也得罪了你，从今以后，我不配称为你的儿子。}\BibleRef{路 15:22} 父亲却吩咐仆人说：\BibleQuote{快把那上好的袍子拿出来，给他\Emph{穿上}；给他手上戴上戒指，脚上穿上鞋；把那只肥牛犊牵来宰了，我们应吃喝欢庆；\BibleRef{路 15:24} 因为我这儿子是死而复活〔阿拉伯文抄本，页102〕，是失而复得。}\BibleRef{路 15:25} 他们便开始欢庆。那时，大儿子正在田里。当他回来，离家不远时，听见奏乐跳舞的声音。\BibleRef{路 15:26} 便叫过一个仆人来，问这是什么事。\BibleRef{路 15:27} 仆人对他说：\BibleQuote{你弟弟回来了，你父亲因他平安无事地回来，便宰了那只肥牛犊。}\BibleRef{路 15:28} 他就生气，不肯进去；他父亲出来劝他。\BibleRef{路 15:29} 他回答父亲说：\BibleQuote{我服事你这么多年，从来没有违背过你的命令，而你从未给我一只山羊羔，让我与我的朋友一同欢乐；\BibleRef{路 15:30} 但你这个儿子，同娼妓耗尽了你的财产，一回来，你倒为他宰了肥牛犊。}\BibleRef{路 15:31} 父亲对他说：\BibleQuote{孩子，你常同我在一起，凡我所有的，都是你的；\BibleRef{路 15:32} 只应欢乐高兴，因为你这弟弟死而复活，是\Emph{曾经失去}，如今又找到的。}\par

\BibleQuote{34 \BibleRef{路 16:1}耶稣又对门徒讲了一个比喻：\Quote{有一个富人，他有一个管家；35 \BibleRef{路 16:2}有人在主人前告发他挥霍了主人的财物。主人便叫他来，对他说：\InnerQuote{我听说你这事做什么？把你管家的账目交出来，因为从今以后你不能再作管家了。}36 \BibleRef{路 16:3}管家心里想：\InnerQuote{主人撤去我管家的职务，我做什么才好呢？掘地，我没有气力；讨饭，我又害羞。}37 \BibleRef{路 16:4}我知道我要做什么，好叫人们在我被撤去管家职务之后，收留我到他家去。}38 \BibleRef{路 16:5}他便把主人的债户一个一个地叫来，对第一个说：\InnerQuote{你欠我主人多少？}39 \BibleRef{路 16:6}他说：\InnerQuote{一百桶油。}管家对他说：\InnerQuote{拿你的账单，坐下快写五十桶。}40 \BibleRef{路 16:7}然后对另一个说：\InnerQuote{你欠多少？}他说：\InnerQuote{一百石麦子。}管家对他说：\InnerQuote{拿你的账单，写八十石。}41 \EditorialNote{阿拉伯文，第103页} \BibleRef{路 16:8}主人就夸奖这个不义的管家，办事精明：因为世俗之子在应付世俗的事务上，比光明之子更为精明。}\Quote{42 \BibleRef{路 16:9}我也告诉你们：要用不义的钱财结交朋友，这样，当钱财用尽时，他们可以收留你们到永远的帐幕里。43 \BibleRef{路 16:10}在小事上忠信的，在大事上也忠信；在小事上不义的，在大事上也不义。44 \BibleRef{路 16:11}如果你们在不义的钱财上不忠信，谁还把真实的钱财托付给你们呢？45 \BibleRef{路 16:12}如果你们在别人的财物上不忠信，谁还把属于你们的交给你们呢？}\par

\SourceRecord{Translated by Hope W. Hogg.；Ante-Nicene Fathers；Vol. 9.；Edited by Allan Menzies.；Buffalo, NY: Christian Literature Publishing Co.,；1896.；<http://www.newadvent.org/fathers/100226.htm>.}

% structure: p0001 source=h1 target=section reason=page-title

\section{\WorkTitle{Diatessaron}（第27节）}

1 \BibleRef{玛 18:23} 因此，天国好比一个君王，要与他仆人结账。\newline{}2 \BibleRef{玛 18:24} 开始结账时，有人带了一个欠他一万塔冷通的人来。\newline{}3 \BibleRef{玛 18:25} 因他无力偿还，主人就下令把他和他的妻子儿女，并他所有的一切都卖掉来偿还。\newline{}4 \BibleRef{玛 18:26} 那仆人跪下来叩拜\Emph{他}，说：\Quote{主啊，宽容我吧！我必将一切还清。}\newline{}5 \BibleRef{玛 18:27} 那仆人的主人就动了怜悯的心，把他释放了，并且免了他的债。\newline{}6 \BibleRef{玛 18:28} 但那仆人出去，遇见了一个欠他一百德纳的同伴，就抓住他，掐住他的喉咙说：\Quote{还清你欠的债！}（阿拉伯语抄本第104页）\newline{}7 \BibleRef{玛 18:29} 那同伴俯伏在地求他说：\Quote{宽容我吧！我必要还你。}\newline{}8 \BibleRef{玛 18:30} 他却不肯，且把他下在狱里，直到还清欠债。\newline{}9 \BibleRef{玛 18:31} 其他同伴看见这事，非常悲伤，就去向主人报告了这一切。\newline{}10 \BibleRef{玛 18:32} 于是主人叫了他来，对他说：\Quote{\Emph{你}这恶仆！因为你求我，我就免了你全部的债；\newline{}11 \BibleRef{玛 18:33} 难道你不该也怜悯你的同伴，如同我怜悯你一样吗？}\newline{}12 \BibleRef{玛 18:34} 主人就大怒，把他交给刑役，直到还清所欠的一切。\newline{}13 \BibleRef{玛 18:35} 如果你们不从心里宽恕你们的弟兄，我的天父也必要这样对待你们。\newline{}14 \BibleRef{路 17:3} 你们要谨慎！如果你的弟兄犯了罪，你就规劝他；如果他后悔了，就宽恕他。\newline{}15 \BibleRef{路 17:4} 如果他一天七次得罪你，又七次转向你说：\Quote{我后悔了}，你也得宽恕他。\newline{}16 \BibleRef{玛 18:15} 如果你的弟兄得罪了你，去，只在你们两人之间规劝他；如果他听从你，你便赚得了你的弟兄。\newline{}17 \BibleRef{玛 18:16} 如果他不听，你就带一两个人去，这样凭两三个人的口供，一切事都可以成立。\newline{}18 \BibleRef{玛 18:17} 如果连他们也不听，你就告诉教会；如果连教会也不听，你就把他看作外邦人或税吏。\newline{}19 \BibleRef{玛 18:18} 我实在告诉你们：凡你们在地上束缚的，在天上也要被束缚；凡你们在地上释放的，在天上也要被释放。\newline{}20 \BibleRef{玛 18:19} 我也告诉你们：如果你们中有两个人在地上同心合意地祈求，我在天之父必成全他们。\newline{}21 \BibleRef{玛 18:20} 因为无论在哪里，有两三个人因我的名聚在一起，我就在他们中间。（阿拉伯语抄本第105页）\newline{}22 \BibleRef{玛 18:21} 那时，\Person{刻法}走近耶稣说：\Quote{主啊！若我的弟兄得罪我，我该宽恕他多少次？直到七次吗？}\newline{}23 \BibleRef{玛 18:22} 耶稣对他说：\Quote{我不对你说直到七次，而是到七十个七次。}\newline{}24 \BibleRef{路 12:47} 那知道主人的旨意，却不预备，或不按他的意思做的仆人，必受许多鞭打；\newline{}25 \BibleRef{路 12:48} 但那不知道的，做了该受鞭打的事，也必受轻打。凡多给谁，就向谁多取；多托付谁，就向谁多要。\newline{}26 \BibleRef{路 12:49} 我来是把火投在地上，我是多么希望它已经燃烧起来！\newline{}27 \BibleRef{路 12:50} 我有应受的圣洗，我是如何焦急，直到它完成！\newline{}28 \BibleRef{玛 18:10} 你们小心，不要轻视这些小子中的一个，因为他们信从我。我实在告诉你们：他们的天使在天上，常见我在天之父的面。\newline{}29 \BibleRef{玛 18:11} 人子来，是为拯救迷失的人。\par

30 \BibleRef{若 7:1} 此后，耶稣在加里肋亚行走；他不愿在犹太行走，因为犹太人想杀他。\newline{}31 \BibleRef{路 13:1} 那时有人来告诉他关于加里肋亚人的事，就是比拉多把他们的血掺和在他们的祭品中。\newline{}32 \BibleRef{路 13:2} 耶稣回答他们说：\Quote{你们以为这些加里肋亚人比其他加里肋亚人更有罪，才遭此祸吗？\newline{}33 \BibleRef{路 13:3} 不是的。我告诉你们，如果你们不悔改，你们也要同样灭亡。\newline{}34 \BibleRef{路 13:4} 或者，那些被史罗亚塔倒塌压死的十八个人，你们以为他们比住在耶路撒冷的其他居民更有罪吗？\newline{}35 \BibleRef{路 13:5} 不是的。我告诉你们，如果你们不全都悔改，你们也要像他们一样灭亡。}\par

36 \BibleRef{路 13:6} 他给他们讲了这个比喻：一个人有一棵无花果树栽在自己的葡萄园里；他来寻找果子，却找不到。\newline{}37 \BibleRef{路 13:7} 就对园丁说：\Quote{你看，我三年来在这棵无花果树上找果子，竟找不到；砍掉它，何必白占土地？}\newline{}38 \BibleRef{路 13:8} 园丁回答他说：\Quote{主人，再留它一年吧！让我在它周围挖松土壤，加上肥料；\newline{}39 \BibleRef{路 13:9} 将来如果结果子，就留着；不然的话，你就把它砍了。}\par

40 \BibleRef{路 13:10} 耶稣在安息日于一会堂中施教时，41 \BibleRef{路 13:11} 那里有一个妇人，被病魔缠身十八年；她佝偻着，完全不能直起腰来。42 \BibleRef{路 13:12} 耶稣看见她，便叫她过来，对她说：\Quote{妇人，你从疾病中解脱了。}43 \BibleRef{路 13:13} 遂按手在她身上；她立刻直起腰来，并光荣天主。44 \BibleRef{路 13:14} 会堂长因耶稣在安息日治病，便气忿忿地对群众说：\Quote{有六天应当工作；你们在这些日子里来求医，但不可在安息日。}45 \BibleRef{路 13:15} 耶稣却回答他说：\Quote{\Emph{你们}这些伪君子！你们每个人在安息日，难道不解开槽上的牛或驴，牵去饮水吗？46 \BibleRef{路 13:16} 这个妇人本是亚巴郎的女儿，被撒但捆绑了十八年，难道不该在安息日解开她的束缚吗？}47 \BibleRef{路 13:17} 他说这话时，所有反对他的人都感到惭愧；全体民众因他所行的一切奇事而欢喜。\par

\SourceRecord{Translated by Hope W. Hogg.；Ante-Nicene Fathers；Vol. 9.；Edited by Allan Menzies.；Buffalo, NY: Christian Literature Publishing Co.,；1896.；<http://www.newadvent.org/fathers/100227.htm>.}

% structure: p0001 source=h1 target=section reason=page-title

\section{戴亚塔撒隆（第二十八章）}

1, 2 [Arabic, p. 107] \BibleRef{若 7:2} 那时，犹太人的住棚节近了。\BibleRef{若 7:3} 于是耶稣的弟兄们对他说：\Quote{你离开这里，往犹太去罢，好叫你的门徒也看见你所行的事。\BibleRef{若 7:4} 因为没有人暗地里行事，却愿意显扬的。\BibleRef{若 7:5} 你若行这些事，就当将自己显明给世界看。} 因为连他的弟兄们也没有信从他。\BibleRef{若 7:6} 耶稣对他们说：\Quote{我的时候还没有到；你们的时候却常是方便的。\BibleRef{若 7:7} 世界不能恨你们，却是恨我，因为我指证它所做的事是恶的。\BibleRef{若 7:8} 你们上去过节罢，我现在不上去过这节，因为我的时候还没有满。} \BibleRef{若 7:9} 他说了这话，仍旧住在加利利。\par

9 但他的弟兄们上去过节以后，他就从加利利起程，来到犹太边境，到约但河外的地方；\BibleRef{玛 19:2} 有许多群众跟随着他，他就在那里治好了他们。\BibleRef{若 7:10} 但他出去，前往过节，不是公开的，而是隐藏着去的。\BibleRef{若 7:11} 犹太人就在节期寻找他，说：\Quote{那人在哪里？} \BibleRef{若 7:12} 群众中对他议论纷纷。有的说：\Quote{他是好人；} 有的说：\Quote{不，他迷惑众人。} \BibleRef{若 7:13} 但没有人公开讲论他，因为怕犹太人。\par

15 [Arabic, p. 108] \BibleRef{若 7:14} 但住棚节的日子过了一半，耶稣上到圣殿里去教训人。\BibleRef{若 7:15} 犹太人就希奇说：\Quote{这个人没有学过，怎么明白书呢？} \BibleRef{若 7:16} 耶稣回答他们说：\Quote{我的教训不是我自己的，乃是那差我来的。\BibleRef{若 7:17} 人若愿意遵行他的旨意，就必晓得这教训或是出于天主，或是我凭着自己说的。\BibleRef{若 7:18} 凭着自己说的，是寻求自己的荣耀；但那寻求差他来者之荣耀的，他是真实的，在他心里没有不义。\BibleRef{若 7:19} 摩西岂不是传给你们律法么？你们却没有一个人守律法。\BibleRef{若 7:20} 你们为什么想要杀我？} 群众回答说：\Quote{你是附了魔的，谁想要杀你？} \BibleRef{若 7:21} 耶稣回答他们说：\Quote{我做了一件事，你们都因此惊奇。\BibleRef{若 7:22} 摩西传给你们割礼（其实不是从摩西开始的，而是从祖先开始的），你们就在安息日给人行割礼。\BibleRef{若 7:23} 若有人在安息日受割礼，免得违背摩西的律法，那么，我在安息日使一个人全然痊愈，你们就向我发怒么？\BibleRef{若 7:24} 不要凭外貌断定是非，总要按公理断定是非。}\par

26 \BibleRef{若 7:25} 耶路撒冷有人问：\Quote{这不是他们想要杀的人么？\BibleRef{若 7:26} 你看，他明明地讲论，他们也不向他说什么。难道官长真知道这是基督么？\BibleRef{若 7:27} 然而这个人，我们知道他从哪里来；基督来的时候，没有人知道他从哪里来。} \BibleRef{若 7:28} 那时耶稣在圣殿里教训人，大声说：\Quote{你们也知道我，也知道我从哪里来；但我来并不是由于自己，但那差我来的是真实的，你们不认识他。\BibleRef{若 7:29} 我却认识他，因为我是从他来的，他也差了我来。} \BibleRef{若 7:30} 他们想要捉拿他，只是没有人下手，因为他的时候还没有到。\BibleRef{若 7:31} 但群众中有许多人信了他，说：\Quote{基督来的时候，他所行的神迹，难道会比这人所行的更多么？}\par

33 \BibleRef{路 12:13} 群众中有一个人对主说：\Quote{老师，请你吩咐我的兄弟同我分家业。} \BibleRef{路 12:14} 耶稣对他说：\Quote{你这个人，谁立我做你们的审判官和分家业的人呢？} \BibleRef{路 12:15} 他对门徒说：\Quote{你们要谨慎自守，免去一切的贪心，因为人的生命不在乎家道丰富。} \BibleRef{路 12:16} 就用比喻对他们说：\Quote{有一个财主，田产丰盛；\BibleRef{路 12:17} 自己心里思想说：\InnerQuote{我的出产没有地方收藏，怎么办呢？} \BibleRef{路 12:18} 又说：\InnerQuote{我要这么办：要把我的仓房拆了，另盖更大的，在那里好收藏我一切的粮食和财物。} \BibleRef{路 12:19} 然后要对我的灵魂说：\InnerQuote{灵魂哪，你有许多财物积存，可作多年的费用，只管安安逸逸的吃喝快乐罢。} \BibleRef{路 12:20} 天主却对他说：\InnerQuote{\Emph{你}这无知的人哪，今夜必要你的灵魂；你所预备的，要归谁呢？} \BibleRef{路 12:21} 凡为自己积财，在天主面前却不富足的，也是这样。}\par

42 \BibleRef{谷 10:17} 当耶稣在路上行走时，有一个官吏的年轻人来到他跟前，跪在他面前，问他说：\Quote{善师，43 我必须做什么才能获得永生？} \BibleRef{谷 10:18} 耶稣对他说：\Quote{你为什么称44我为善？除了那一位，\Emph{就是}天主，没有一个人是善的。} \BibleRef{谷 10:19} \Quote{你晓得诫命。45 你若愿意进入生命，就该遵守诫命。} 那年轻人[阿拉伯文，第110页]对他说：\Quote{哪一条诫命？} 耶稣对他说：46 \BibleRef{谷 10:19} \Quote{不可奸淫，不可偷盗，不可杀人，不可作假见证，不可伤害，孝敬你的父亲47和你的母亲，并要爱人如己。} 那年轻人对他48说：\Quote{这一切我自幼就遵守了，我还缺少什么？} \BibleRef{谷 10:21} 耶稣49注视着他，爱了他，对他说：\BibleRef{玛 19:21} \Quote{你若愿意是成全的，你所缺少的只有一件：去变卖你所有的一切，施舍给穷人，你必有宝藏在天上；然后背起你的50十字架，来跟随我。} 那年轻人听了这话，就愁眉不展，忧伤地走开了51，因为他很富有。耶稣看见他忧愁，就望着他的门徒，对他们说：\Quote{有钱财的人进入天主的国是多么难啊！}\par

\SourceRecord{Translated by Hope W. Hogg.；Ante-Nicene Fathers；Vol. 9.；Edited by Allan Menzies.；Buffalo, NY: Christian Literature Publishing Co.,；1896.；<http://www.newadvent.org/fathers/100228.htm>.}

% structure: p0001 source=h1 target=section reason=page-title

\section{四福音合参 第二十九篇}

1 \BibleQuote{我实在告诉你们：富人进天国是难的。}\BibleRef{玛 19:23} 2 \BibleQuote{我又告诉你们：骆驼穿过针眼比富人进天主的国更容易。}\BibleRef{玛 19:24} 3 \BibleQuote{门徒们对这些话感到惊奇。耶稣回答又对他们说：孩子们，倚靠钱财的人进天主的国是多么难啊！}\BibleRef{谷 10:24} 4 \BibleQuote{那些听的人更加惊奇，彼此激动地说：你们想，谁能得救呢？}\BibleRef{谷 10:26} 5 \BibleQuote{耶稣注视他们说：在人这是不可能的，但在天主却不然；因为天主能做一切事。}\BibleRef{谷 10:27}\EditorialNote{阿拉伯文原稿第111页} 6 \Person{西满刻法}对他说：\Quote{看，我们舍弃了一切跟随了你；那么，你认为我们会得到什么呢？} 7 \BibleQuote{耶稣对他们说：我实在告诉你们，你们这些跟随我的人，在新时代，当人子坐在他光荣的宝座上时，你们也要坐在十二个宝座上，审判以色列十二支派。}\BibleRef{玛 19:28} 8 \BibleQuote{我实在告诉你们，凡为了天主的国，或为了我、为了我的福音而舍弃房屋、兄弟、姊妹、父、母、妻子、儿女、田地的人，}\BibleRef{谷 10:29} 9 \BibleQuote{没有不在今世得到多倍的，并在来世承受永生的。}\BibleRef{路 18:30} 10 \BibleQuote{今世得房屋、兄弟、姊妹、母亲、儿女、田地，并受迫害；来世得\Emph{永远}的生命。}\BibleRef{谷 10:30} 11 \BibleQuote{有许多在先的要成为在后的，在后的要成为在先的。}\BibleRef{谷 10:31}\par

12 \BibleQuote{法利塞人听见这一切，因贪爱钱财而讥笑他。}\BibleRef{路 16:14} 13 \BibleQuote{耶稣知道他们的心，对他们说：你们在人前自充义人，但天主知道你们的心；因为人在高尚的，在天主面前却是可憎的。}\BibleRef{路 16:15}\par

14 \BibleQuote{有一个富人，身穿紫红袍和细麻衣，天天奢华宴乐。}\BibleRef{路 16:19} 15 \BibleQuote{有一个乞丐名叫拉匝禄，满身疮痍，被放在富人的门口，}\BibleRef{路 16:20} 16 \BibleQuote{他渴望用富人桌子上掉下来的碎屑充饥；但即使狗也来舔他的疮痍。}\BibleRef{路 16:21}\EditorialNote{阿拉伯文原稿第112页} 17 \BibleQuote{那乞丐死了，天使把他送到亚巴郎的怀里；富人也死了，被人埋葬了。}\BibleRef{路 16:22} 18 \BibleQuote{他在阴府受痛苦，举目远远望见亚巴郎和他怀中的拉匝禄，}\BibleRef{路 16:23} 19 \BibleQuote{便喊着说：父亲亚巴郎，可怜我吧！请打发拉匝禄用指头蘸点水来凉凉我的舌头，因为我在这火焰中非常痛苦。}\BibleRef{路 16:24} 20 \BibleQuote{亚巴郎说：孩子，你活着时享尽了福，而拉匝禄却受尽了苦；如今他在这里得安慰，而你受痛苦。}\BibleRef{路 16:25} 21 \BibleQuote{除此之外，在我们与你们之间有一道深渊，使从这里到你们那边去的人不能过去，从那边到我们这边来的人也不能过来。}\BibleRef{路 16:26} 22 \BibleQuote{他说：父亲，那么我求你打发拉匝禄到我父亲家去，}\BibleRef{路 16:27} 23 \BibleQuote{因为我有五个兄弟，让他去警戒他们，免得他们也来到这痛苦的地方。}\BibleRef{路 16:28} 24 \BibleQuote{亚巴郎说：他们有梅瑟和先知，让他们听从他们吧。}\BibleRef{路 16:29} 25 \BibleQuote{他说：不，父亲亚巴郎，倘若有人从死者中到他们那里去，他们必会悔改。}\BibleRef{路 16:30} 26 \BibleQuote{亚巴郎对他说：如果他们不听从梅瑟和先知，即使有人从死者中复活，他们也不会相信。}\BibleRef{路 16:31}\par

27 \BibleQuote{天国好像一个家主，清晨出去雇工人到他的葡萄园工作。}\BibleRef{玛 20:1} 28 \BibleQuote{他和工人议定一天一个德纳，就派他们到葡萄园去了。}\BibleRef{玛 20:2} 29 \BibleQuote{约在第三时辰，他又出去，看见另有人闲站在市场上。}\BibleRef{玛 20:3} 30 \BibleQuote{他对他们说：你们也去吧，我会给你们公道酬劳。}\BibleRef{玛 20:4}\EditorialNote{阿拉伯文原稿第113页} 31 \BibleQuote{他们就去了。约在第六和第九时辰，他又出去，照样做了，派他们去。}\BibleRef{玛 20:5} 32 \BibleQuote{约在第十一时辰，他出去，看见还有人闲站着，就问他们说：为什么你们整天闲站？}\BibleRef{玛 20:6} 33 \BibleQuote{他们回答说：因为没有人雇用我们。他对他们说：你们也去吧，会得到公道的报酬。}\BibleRef{玛 20:7} 34 \BibleQuote{到了晚上，葡萄园的主人对他的管家说：叫工人来，发给他们工资，从后来的开始，直到先来的。}\BibleRef{玛 20:8} 35 \BibleQuote{那些在第十一时辰来的人，每人领了一个德纳。}\BibleRef{玛 20:9} 36 \BibleQuote{先来的人以为会多领，但他们也每人领了一个德纳。}\BibleRef{玛 20:10} 37 \BibleQuote{他们领了\Emph{它}之后，就埋怨家主说：}\BibleRef{玛 20:11} 38 \BibleQuote{这些最后来的人只工作了一个时辰，而我们把一天的劳累和炎热都忍受了，你却使他们与我们平等。}\BibleRef{玛 20:12} 39 \BibleQuote{他回答其中一个说：朋友，我没有亏待你，\Emph{难道你}不是与我议定了一个德纳\Emph{吗}？}\BibleRef{玛 20:13} 40 \BibleQuote{拿你的走吧！我愿意给这个后来的如同给你一样。}\BibleRef{玛 20:14} 41 \BibleQuote{难道我不能用自己的财物做我所愿意的吗？还是因为我善良，你就眼红吗？}\BibleRef{玛 20:15} 42 \BibleQuote{这样，最后的将成为最先的，最先的将成为最后的。被召的人多，蒙选的人少。}\BibleRef{玛 20:16}\par

43 \BibleRef{路 14:1} 当耶稣在安息日进入一位法利塞首领的家吃饭时，他们正窥视他，要看他44,45做什么，\BibleRef{路 14:2} 在他面前有一个患水肿的人。\BibleRef{路 14:3} 耶稣回答并46对经师和法利塞人说：\Quote{在安息日治病，合乎律法吗？}\BibleRef{路 14:4} 但他们[阿拉伯文，第114页]却沉默不语。他就拉着那人，治好了他，叫他走了。47 \BibleRef{路 14:5} 又对他们说：\Quote{你们中间谁\Emph{有}儿子或牛在安息日掉进井里，不立刻拉他上来，给他48打水呢？}\BibleRef{路 14:6} 他们不能反驳他这句话。\par

\SourceRecord{Translated by Hope W. Hogg.；Ante-Nicene Fathers；Vol. 9.；Edited by Allan Menzies.；Buffalo, NY: Christian Literature Publishing Co.,；1896.；<http://www.newadvent.org/fathers/100229.htm>.}

% structure: p0001 source=h1 target=section reason=page-title

\section{\WorkTitle{迪亚特萨隆（第三十部分）}}

\BibleRef{路 14:7} 祂向那些被请的人讲了一个比喻，因为他们看见\newline{}他们选择坐席中最高处的座位：\BibleRef{路 14:8} \Quote{当有人请你赴宴时，不要先去坐在席上的首位，恐怕那里有比你更尊贵的人，\BibleRef{路 14:9} 那请你的人来对你说：\InnerQuote{请让座给这个人。}那时你就要羞愧地退到末位去。\BibleRef{路 14:10} 但你被请时，要去坐在末位，好让那请你的人来对你说：\InnerQuote{朋友，请上坐。}那时你在同席的人面前才有光彩。\BibleRef{路 14:11} 因为凡高举自己的，必被贬抑；凡贬抑自己的，必被高举。}\par

\BibleRef{路 14:12} 祂也对那请祂的人说：\Quote{当你设午宴或晚宴时，不要请你的朋友，也不要请你的弟兄，也不要请你的亲戚，也不要请富有的邻人，恐怕他们也要回请你，你就得了回报。\BibleRef{路 14:13} 但你设宴时，要请贫穷的、残废的、瘸腿的、瞎眼的：\BibleRef{路 14:14} 你就有福了，因为他们无力报答你；但在义人复活时，你必得报答。}\BibleRef{路 14:15} 同席的人中有一个听了这话，就对祂说：\Quote{将来能在天主的国里吃饭的，是有福的！}\par

\BibleRef{路 14:16} \Quote{耶稣又用比喻回答说：\EditorialNote{阿拉伯语原稿，第115页}天国好比一个国王，为自己的儿子摆设婚宴，并预备了盛大的筵席，请了许多人。\BibleRef{路 14:18} 他们却一致推辞。第一个对他说：\InnerQuote{我买了一块田地，必须前去看一看，请准我辞了。}\BibleRef{路 14:19} 另一个说：\InnerQuote{我买了五对牛，要去试试它们，请准我辞了。}\BibleRef{路 14:20} 另一个说：\InnerQuote{我才娶了妻，所以不能去。}\BibleRef{玛 22:4} 国王又打发别的仆人说：\InnerQuote{你们对被请的人说：我的宴席已经准备好了，我的公牛和肥畜已宰了，一切都齐备了，请来赴婚宴吧！}\BibleRef{玛 22:5} 他们却不理睬，有的往自己的田里去了，有的去做自己的生意去了；\BibleRef{玛 22:6} 其余的竟拿住他的仆人，凌辱后杀死了他们。\BibleRef{路 14:21} 那仆人回来，把这一切报告了主人。\BibleRef{玛 22:7} 国王就动了怒，打发自己的军队，消灭了那些凶手，焚毁了他们的城市。\BibleRef{玛 22:8} 然后对仆人说：\InnerQuote{婚宴已准备好了，但是被请的人都不配。\BibleRef{路 14:21} 你们快出去，到城中的大街小巷，领那些贫穷的、有病的、瘸腿的、瞎眼的进来。}\BibleRef{路 14:22} 仆人们照做了，回来说：\InnerQuote{主人，你吩咐的我们已经办了，可是还有空位子。}\BibleRef{路 14:23} 主人对仆人说：\InnerQuote{你们到路上和篱笆边去，勉强人进来，好坐满我的屋子。\BibleRef{路 14:24} 我告诉你们：先前被请的那些人，没有一个能尝到我的宴席。}\BibleRef{玛 22:10} 那些仆人出去，到大路上，凡遇到的，无论善恶都召集来了，婚宴上就满了坐席的人。\BibleRef{玛 22:11} 国王进来视察坐席的客人，看见那里有一个人没有穿婚宴礼服，\BibleRef{玛 22:12} 便对他说：\InnerQuote{朋友，你怎么不穿婚宴礼服就进来了？}他默不作声。\BibleRef{玛 22:13} 国王就对仆役说：\InnerQuote{捆起他的手和脚，把他丢到外面的黑暗中：那里要有哀号和切齿。}\BibleRef{玛 22:14} 因为被召的人多，被选的人少。}\EditorialNote{阿拉伯语原稿，第116页}\par

\BibleRef{若 5:1} 此后，犹太人的无酵节快到了，耶稣就出去往耶路撒冷去。在路上有十个癞病人遇见祂，远远地站着，\BibleRef{路 17:13} 高声说：\Quote{主耶稣，可怜我们罢！}\BibleRef{路 17:14} 祂看见他们，就对他们说：\Quote{你们去，把身体给司铎们看。}他们去的时候，就洁净了。\BibleRef{路 17:15} 其中一个看见自己洁净了，就回来大声赞美天主，\BibleRef{路 17:16} 俯伏在耶稣脚前感谢祂；这\Emph{人}是个撒玛黎雅人。\BibleRef{路 17:17} 耶稣回答说：\Quote{洁净了的不是十个人吗？那九个人在哪里呢？\BibleRef{路 17:18} 除了这个外邦人，没有一个回来归光荣于天主吗？}\BibleRef{路 17:19} 耶稣就对他说：\Quote{起来，去罢！你的信德救了你。}\par

40 \BibleRef{谷 10:32} 他们正上路往耶路撒冷去的时候，耶稣走在他们前面；他们感到惊奇，又害怕地跟随。他把十二门徒带到一边，41 开始私下告诉他们将要临到他身上的事。\BibleRef{路 18:31} 他对他们说（阿拉伯语，第117页）：\Quote{我们上耶路撒冷去，凡借先知所写关于人子的事都要应验42。\BibleRef{谷 10:33} 他要被交给司祭长和经师；他们要定他死罪，43 并把他交给外邦人；\BibleRef{谷 10:34} 他们要凌辱他，鞭打44他，向他脸上吐唾沫，羞辱他，钉死他，杀害他；\BibleRef{路 18:33} 并且在45第三天他要复活。} \BibleRef{路 18:34} 但他们一点也不明白这些话；这话对他们隐藏着，他们不能领悟所说的事。\par

46 \BibleRef{玛 20:20} 那时，载伯德两个儿子的母亲，同她的两个儿子前来，朝拜他，求他一件\Emph{特定}的事。\BibleRef{玛 20:21} 他对她说：47 你要什么？\BibleRef{谷 10:35} 雅各伯和若望，她的两个儿子，上前来对他说：老师，我们愿意你为我们做我们所求的事。48 \BibleRef{谷 10:36} 他对他们说：你们要我为你做什么？49 \BibleRef{谷 10:37} 他们对他说：赐我们在你的国和你的荣耀中，一个坐在你右边，一个坐在你左边。50 \BibleRef{谷 10:38} 耶稣对他们说：你们不知道你们所求的是什么。我所饮的杯，你们能饮吗？以我所受的51圣洗，你们能受吗？\BibleRef{谷 10:39} 他们对他说：我们能。耶稣对他们说：我所饮的杯，你们必要饮；以我所受的52圣洗，你们必要受；\BibleRef{谷 10:40} 但坐在我的右边或左边，不是我可以给的，而是\Emph{它是}为我父所预备\Emph{它}的那人。\par

\SourceRecord{Translated by Hope W. Hogg.；Ante-Nicene Fathers；Vol. 9.；Edited by Allan Menzies.；Buffalo, NY: Christian Literature Publishing Co.,；1896.；<http://www.newadvent.org/fathers/100230.htm>.}

% structure: p0001 source=h1 target=section reason=page-title

\section{四福音合参（第三十一部分）}

1 \BibleRef{谷 10:41} 十个人听见了，就向雅各伯和若望发怒。2 \BibleRef{谷 10:42} 耶稣叫了他们来，对他们说：\BibleQuote{你们知道，外邦人的统治者做主治理他们，他们的伟人也掌权管制他们。3 \BibleRef{谷 10:43} 但在你们中间不可这样：你们中间谁愿为大，就当作你们的仆役；4 \BibleRef{谷 10:44} 你们中间谁愿为首，就当作众人的奴仆。} 5 \BibleRef{玛 20:28} \Emph{就}如人子来不是受服事，而是服事人，并交出自己的生命，为大众作赎价。6 \BibleRef{路 13:22} 他说了这话，就继续在各村庄城市游行施教，往耶路撒冷去。7 \BibleRef{路 13:23} 有一个人问他：\BibleQuote{得救的人少吗？} 耶稣回答他们说：\BibleQuote{8 \BibleRef{路 13:24} 你们要努力从窄门进去：我告诉你们，将来有许多人想进去，却不能进去。9 \BibleRef{路 13:25} 当家主起来把门关上以后，你们站在外面敲门说：\InnerQuote{主啊，给我们开门吧！}他要回答你们说：\InnerQuote{我不认识你们，不知道你们是从哪里来的。}10 \BibleRef{路 13:26} 那时你们要说：\InnerQuote{我们在你面前吃过喝过，你也曾在我们的街市上教导过我们。}11 \BibleRef{路 13:27} 他要对你们说：\InnerQuote{我不认识你们，不知道你们是从哪里来的；你们这些作恶的人，离开我去吧！}12 \BibleRef{路 13:28} 你们要看见亚巴郎、依撒格、雅各伯和众先知在天主国里，自己却被赶在外面，在那里要有哀号和切齿。13 \BibleRef{路 13:29} 从东从西，从南从北，将有人来，在天主国里坐席。14 \BibleRef{路 13:30} 那时，有在后的将成为在先的，有在先的将成为在后的。}\par

15、16 \BibleRef{路 19:1} 耶稣进了耶里哥，正经过的时候，17 \BibleRef{路 19:2} 有一个人名叫匝凯，是个富有且作税吏长的。18 \BibleRef{路 19:3} 他想要看看耶稣是谁，但由于人多，不能看见，因为他身材矮小。\BibleRef{路 19:4} 于是他跑到前头，爬上一棵无花果树，要看耶稣，因为耶稣就要从那里经过。19 \BibleRef{路 19:5} 耶稣来到那地方，抬头一看，对他说：\BibleQuote{匝凯，快下来！今天我应当住在你家里。}20 \BibleRef{路 19:6} 他便赶快下来，欣喜地接待他。21 \BibleRef{路 19:7} 众人看见，都窃窃私议说：\BibleQuote{他竟到罪人家里去住宿。}22 \BibleRef{路 19:8} 匝凯站起来对主说：\BibleQuote{主啊，我现在把我一半的财产施舍给穷人；我若欺骗过谁，就还他四倍。}23 \BibleRef{路 19:9} 耶稣对他说：\BibleQuote{今天救恩到了这一家，因为他也是亚巴郎之子。24 \BibleRef{路 19:10} 因为人子来，是为寻找并拯救迷失了的人。}\par

25 耶稣从耶里哥出来的时候，同门徒一起，有许多群众跟着他。\BibleRef{路 18:35} 有一个瞎子坐在路旁讨饭。他名叫提买，是提买的儿子。他听见群众经过的声音，就问是什么事。\BibleRef{路 18:37} 有人告诉他：是纳匝肋人耶稣经过。他便喊叫说：\Quote{耶稣，达味之子，可怜我吧！} \BibleRef{路 18:39} 耶稣前头走的人斥责他，叫他不要出声；\BibleRef{谷 10:48} 但他越发喊叫说：\Quote{达味之子，可怜我吧！} \BibleRef{谷 10:49} 耶稣站住，吩咐把他带过来。他们叫那瞎子来，对他说：\Quote{放心！起来，他叫你哩！} \BibleRef{谷 10:50} 瞎子就丢下自己的外衣，跳起来，走到耶稣跟前。\BibleRef{谷 10:51} 耶稣问他说：\Quote{你愿意我给你做什么？} 瞎子说：\Quote{师傅！叫我能看见。} \BibleRef{玛 20:34} 耶稣动了怜悯的心，摸了摸他的眼睛，对他说：\Quote{你看见吧！你的信德救了你。} 他立刻看见了，就跟随耶稣，光荣天主；所有看见的人都赞美天主。\par

36 \BibleRef{路 19:11} 祂因临近耶路撒冷，又以为天主的国快要显现，就讲了一个比喻。他们以为37那时天主的国就要出现。\BibleRef{路 19:12} 祂对他们说：一个人，是一个贵族的儿子，往远方去，要得国，38然后回来。\BibleRef{路 19:13} 他叫了十个仆人来，给他们十个米纳，对他们39说：你们去做买卖，直到我回来的时候。\BibleRef{路 19:14} 但他本城的人却恨他，40打发使者随后去说：\Quote{我们不愿意这个\Emph{人}作我们的王。}\BibleRef{路 19:15} 他得了国回来，就吩咐叫那领了钱的仆人来，要知道他们各人41做买卖赚了多少。\BibleRef{路 19:16} 第一个前来，说：\Quote{主啊，你的米纳已经赚了42十个米纳。}\BibleRef{路 19:17} 王对他说：\Quote{你这善良忠信的仆人，你在小事上43显出了忠信，可以管理十座城。}\BibleRef{路 19:18} 第二个前来，44说：\Quote{主啊，你的那份赚了五份。}\BibleRef{路 19:19} 他也对他45说：\Quote{你也可以管理五座城。}\BibleRef{路 19:20} 又有一个前来，说：\Quote{主46啊，请看你的那份，我一直用手巾包着\BibleRef{路 19:21} 我害怕你，因为你是严厉的人，没有存放的也要取，没有给的也要47寻找，没有播种的也要收割。}\BibleRef{路 19:22} 主人对他说：\Quote{凭你的口定你的罪，你这恶懒的仆人，不可靠。你知道我是严厉的人，没有存放的也要取，没有48留下的也要收割。}\BibleRef{路 19:23} \Quote{为什么不把我的钱交给钱庄，49好让我回来时连本带利取回呢？}\BibleRef{路 19:24} 他对旁边站着的人说：\Quote{从他手里把那个米纳夺过来，给那有50、51 [Arabic, p. 121] 十个米纳的人。}\BibleRef{路 19:25} 他们对他说：\Quote{主啊，他已有十个米纳了。}\BibleRef{路 19:26} 他对他们说：\Quote{我告诉你们，凡有的，还要给他；那52没有的，连他所有的也要夺去。}\BibleRef{路 19:27} \Quote{至于那些反对我作王的敌人，把他们带到这里来，在我面前杀掉。}\par

\SourceRecord{Translated by Hope W. Hogg.；Ante-Nicene Fathers；Vol. 9.；Edited by Allan Menzies.；Buffalo, NY: Christian Literature Publishing Co.,；1896.；<http://www.newadvent.org/fathers/100231.htm>.}

% structure: p0001 source=h1 target=section reason=page-title

\section{四福音合参（第三十二节）}

1 当耶稣进入耶路撒冷时，他上了天主的圣殿，2 在那里看见有卖牛、羊、鸽子的。\BibleRef{玛 21:12} 又看见那些做买卖的人和兑换银钱的人坐着，\BibleRef{若 2:14} 他就用绳子做成鞭子，把众人连同羊和牛都赶出圣殿，又推倒兑换银钱的人的银钱，掀翻他们的桌子，3 以及卖鸽子的凳子。\BibleRef{玛 21:13} 他教训他们说：\Quote{经上不是记载着：\InnerQuote{我的殿宇将称为万民的祈祷之所}吗？你们竟把它变成了贼窝。} \BibleRef{若 2:16} 4 他又对卖鸽子的人说：\Quote{把这东西拿走，不要使我父的圣殿成为商场。} \BibleRef{谷 11:16} 5 也不许任何人携带器皿从圣殿里经过。\BibleRef{若 2:17} 6 他的门徒就想起经上所记载的：\BibleQuote{我为你的圣殿所发的热忱将把我吞灭。} \BibleRef{若 2:18} 7 \OriginalTerm{犹太人}{Jews}就问他说：\Quote{你给我们显什么神迹，证明你这样做有权柄呢？} \BibleRef{若 2:19} 8 耶稣回答他们说：\Quote{你们拆毁这\OriginalTerm{圣殿}{temple}，我三天内要把它重建起来。} \BibleRef{若 2:20} 9 犹太人说：\Quote{这圣殿四十六年才建成，你三天内就重建起来吗？} \BibleRef{若 2:21} 10 但耶稣所说的圣殿是指他的身体。他们拆毁它的时候，他三天内要把它重建起来。\BibleRef{若 2:22} 11 因此，当他从死人中复活以后，他的门徒就想起他说过这话；于是他们相信了圣经和耶稣所说的话。\par

12 \BibleRef{谷 12:41} 当耶稣坐在银库对面时，他看众人怎样投钱入库：有许多富\Emph{人}投了很多。13,14 \BibleRef{谷 12:42} 又有一个穷寡妇来了，投了两个\OriginalTerm{小钱}{mites}。\BibleRef{路 21:3} 耶稣叫门徒来，对他们说：\Quote{我实在告诉你们，这个穷寡妇投入库里的，比众人投的都多。\BibleRef{谷 12:44} 因为众人都是拿他们多余的来投；但这\Emph{妇人}却是在不足中，把她所有的一切都投进去了。}\par

16 \BibleRef{路 18:9} 他又对他们讲这个比喻，是关于那些相信自己17 是\OriginalTerm{义人}{righteous}，而轻视别人的人：18 \BibleRef{路 18:10} 有两个人上圣殿去祈祷：一个是\OriginalTerm{法利塞人}{Pharisee}，另一个是\OriginalTerm{税吏}{publican}。\BibleRef{路 18:11} 法利塞人站着，这样祈祷：\Quote{主啊，我感谢你，因为我不像别人，勒索、不义、奸淫，也不像这个税吏；19 \BibleRef{路 18:12} 我一周禁食两次，凡我所有的都献上十分之一。} \BibleRef{路 18:13} 那个税吏20 却远远地站着，连举目望天都不敢，只是21 捶着胸说：\Quote{主啊，可怜我这个罪人吧！} \BibleRef{路 18:14} 我告诉你们：这个人下去到他家里，比那法利塞人更算\OriginalTerm{成义}{justified}。因为凡高举自己的，必被贬抑；凡贬抑自己的，必被高举。\par

22 \EditorialNote{阿拉伯文页码第123页} 到了晚上，他离开众人，23 出城到\OriginalTerm{伯达尼}{Bethany}去，他和他的十二门徒一起住在那。\BibleRef{路 9:11} 众人因为知道那地方，都到他那，他接待了他们；凡需要医治的，他都治好了。24 \BibleRef{谷 11:12} 第二天早晨，他从伯达尼回来的时候，饿了。25 \BibleRef{谷 11:13} 远远地看见一棵\OriginalTerm{无花果树}{fig tree}，在路旁长得枝繁叶茂。他就前去，\Emph{期望}在树上能找到什么；到了跟前，除了叶子，什么也没有找着——26 因为还不是收无花果的季节——\BibleRef{谷 11:14} 他就对树说：\Quote{从今以后，永没有人吃你的果子。} 他的门徒也听见了。\par

27 他们来到耶路撒冷。那里有一个法利塞人，28 名叫尼苛德摩，是犹太人的首领。 \BibleRef{若 3:2} \BibleQuote{这人夜间来到耶稣那里，对他说：老师，我们知道你是由天主而来的师傅；因为你所行的这些神迹，若没有天主同在，无人能行。} 29 \BibleRef{若 3:3} \BibleQuote{耶稣回答说：我实实在在告诉你，人若不\Emph{第二次}出生，就不能见天主的国。} 30 \BibleRef{若 3:4} \BibleQuote{尼苛德摩说：人老了怎能出生呢？难道他能再进母腹而生吗？} 31 \BibleRef{若 3:5} \BibleQuote{耶稣回答说：我实实在在告诉你，人若不是由水和圣神而生，不能进天主的国。} 32 \BibleRef{若 3:6} \BibleQuote{从肉身生的就是肉身；从圣神生的就是神。} 33 \BibleRef{若 3:7} \BibleQuote{你不要惊奇，因我对你说，你们必须\Emph{第二次}出生。} 34 〔阿拉伯文版，第124页〕 \BibleRef{若 3:8} \BibleQuote{风随意向哪里吹，你听到它的响声，却不知道它从哪里来，往哪里去；凡由圣神而生的，也是这样。} 35 \BibleRef{若 3:9} \BibleQuote{尼苛德摩回答说：怎能这样呢？} 36 \BibleRef{若 3:10} \BibleQuote{耶稣回答说：你是以色列的老师，竟不知道这事吗？} 37 \BibleRef{若 3:11} \BibleQuote{我实实在在告诉你，我们所说的是我们所知道的，所见证的是我们所见过的；而你们却不接受我们的见证。} 38 \BibleRef{若 3:12} \BibleQuote{如果我对你们说地上的事，你们尚且不信；若我对你们说天上的事，你们怎能信呢？} 39 \BibleRef{若 3:13} \BibleQuote{没有人升过天，除了那从天上降下来的人子，他仍在天上。} 40 \BibleRef{若 3:14} \BibleQuote{梅瑟在旷野怎样举起了蛇，人子也照样被举起来，} 41 \BibleRef{若 3:15} \BibleQuote{为使凡信他的人不致丧亡，反而获得永生。} 42 \BibleRef{若 3:16} \BibleQuote{天主如此爱了世界，甚至赐下他的独生子，使凡信他的人不致丧亡，反而获得永生。} 43 \BibleRef{若 3:17} \BibleQuote{天主派遣子到世界上来，不是为审判世界，而是为叫世界藉他而得救。} 44 \BibleRef{若 3:18} \BibleQuote{信他的人不受审判；不信的人已受了审判，因为他没有信\Emph{天主独生子}的名字。} 45 \BibleRef{若 3:19} \BibleQuote{审判就在于此：光明来到了世界，世人却爱黑暗甚于光明，因为他们的行为是邪恶的。} 46 \BibleRef{若 3:20} \BibleQuote{凡作恶的，恨光明，也不来就光明，怕他的行为被暴露。} 47 \BibleRef{若 3:21} \BibleQuote{但履行真理的，却来就光明，为显示他的行为是在天主内完成的。}\par

\SourceRecord{Translated by Hope W. Hogg.；Ante-Nicene Fathers；Vol. 9.；Edited by Allan Menzies.；Buffalo, NY: Christian Literature Publishing Co.,；1896.；<http://www.newadvent.org/fathers/100232.htm>.}

% structure: p0001 source=h1 target=section reason=page-title

\section{\WorkTitle{四福音合参}（第33节）}

1 [Arabic, p. 125] \BibleRef{谷 11:19} 到了晚上，耶稣和门徒出城去了。\BibleRef{谷 11:20} 早晨他们经过的时候，门徒看见那棵无花果树连根都枯干了。\BibleRef{玛 21:20} 他们经过，说：\Quote{这无花果树怎么立刻枯乾了呢？}\BibleRef{谷 11:21} 伯多禄想起，就对他说：\Quote{老师，请看，你所咒诅的无花果树已经枯干了。}\BibleRef{谷 11:22} 耶稣回答他们说：\Quote{你们当有天主的信德。\BibleRef{谷 11:23} 我实在告诉你们，无论谁对这座山说：\InnerQuote{你挪开，投到海里去！}心里若不疑惑，只信他说的必成，就必给他成就。\BibleRef{玛 21:21} 你们若有信德，不疑惑，不但能行无花果树上所行的事，就是对这座山说：\InnerQuote{你挪开，投到海里去！}也必成就。\BibleRef{玛 21:22} 你们祷告，无论求什么，只要信，就必得着。}\BibleRef{路 17:5} 宗徒对主说：\Quote{请增加我们的信德。}\BibleRef{路 17:6} 主说：\Quote{如果你们有信德像一粒芥菜种，就是对这棵无花果树说：\InnerQuote{你要连根拔起，栽到海里去！}它也必服从你们。\BibleRef{路 17:7} 你们中间谁有仆人耕地或是放羊，从田里回来，就对他说：\InnerQuote{你快来坐下吃饭}呢？\BibleRef{路 17:8} 岂不是对他说：\InnerQuote{你给我预备晚饭，束上带子伺候我，等我吃喝完了，你才可以吃喝}吗？\BibleRef{路 17:9} 仆人照所吩咐的去作，主人还谢谢他吗？我想不会。\BibleRef{路 17:10} 这样，你们作完了一切所吩咐的，只当说：\InnerQuote{我们是无用的仆人，所作的本是我们应分作的。}}\par

15 \BibleRef{谷 11:24} 所以我告诉你们，凡你们祷告祈求的，无论是什么，只要信你们已经得着了，就必得着。[Arabic, p. 126] \BibleRef{谷 11:25} 你们站着祷告的时候，若想起有人得罪你们，就当饶恕他，好叫你们在天上的父也饶恕你们的过犯。\BibleRef{谷 11:26} 你们若不饶恕人，你们在天上的父也不饶恕你们的过犯。\par

18 \BibleRef{路 18:1} 耶稣又对他们讲一个比喻，是要他们常常祷告，不可灰心：\BibleRef{路 18:2} 说：\Quote{某城有一个审判官，不惧怕天主，也不尊重世人；\BibleRef{路 18:3} 那城里有个寡妇，常到他那里，说：\InnerQuote{求你帮我伸冤，治服我的对头。}\BibleRef{路 18:4} 他多日不准。后来心里说：\InnerQuote{我虽不惧怕天主，也不尊重世人，\BibleRef{路 18:5} 只因这寡妇烦扰我，我就帮她伸冤吧，免得她常来缠磨我。}\BibleRef{路 18:6} 主说：你们听这不义的审判官所说的话。\BibleRef{路 18:7} 天主的选民昼夜呼吁他，他纵然为他们忍了多时，岂不终久给他们伸冤吗？\BibleRef{路 18:8} 我告诉你们，要快快地给他们伸冤了。然而人子来的时候，遇得见世上有信德吗？}\par

26, 27 \BibleRef{谷 11:15} 他们又来到耶路撒冷。\BibleRef{路 20:1} 有一天，耶稣在殿里教训百姓，讲福音的时候，\BibleRef{路 20:2} 司祭长和经师并长老上前来，问他说：\Quote{你告诉我们，你仗着什么权柄作这些事？给你这权柄的是谁呢？}\BibleRef{谷 11:28} 耶稣说：\Quote{我也要问你们一句话，你们告诉我，\BibleRef{玛 21:25} 约翰的圣洗是从天上来的？还是从人间来的呢？}\BibleRef{谷 11:30} 你们告诉我。\BibleRef{玛 21:25} 他们就彼此商议说：\Quote{我们若说\InnerQuote{从天上来}，他必说：\InnerQuote{这样，你们为什么不信他呢？}\BibleRef{玛 21:26} 若说\InnerQuote{从人间来}，我们又怕百姓，因为他们都以约翰为真先知。}\BibleRef{谷 11:33} 于是回答耶稣说：\Quote{我们不知道。}\BibleRef{玛 21:28} 耶稣说：\Quote{我也不告诉你们我仗着什么权柄作这些事。你们想如何？一个人有两个儿子。他来到大儿子那里，说：\InnerQuote{我儿，你今天到葡萄园里去作工。}\BibleRef{玛 21:29} 他回答说：\InnerQuote{我不去。}后来自己懊悔，就去了。\BibleRef{玛 21:30} 又来到小儿子那里，也照样说。\BibleRef{玛 21:31} 他回答说：\InnerQuote{父亲，我去。}他却不去。你们想，这两个儿子中是哪一个遵行父命呢？}他们说：\Quote{大儿子。}耶稣说：\Quote{我实在告诉你们，税吏和娼妓倒比你们先进天主的国。\BibleRef{玛 21:32} 因为约翰遵着义路到你们这里来，你们却不信他；税吏和娼妓倒信他。你们看见了，后来还是不懊悔去信他。}[Arabic, p. 127]\par

40 \BibleRef{玛 21:33} 你们再听一个比喻：有一个家主，栽了一个葡萄园，周围围上篱笆，园里挖了一个榨酒池，盖了一座塔，41,42 租给佃户，就远行去了很久。\BibleRef{谷 12:34} 到了收果子的时节，他打发仆人到佃户那里，好叫他们43 把葡萄园的出产交给他。\BibleRef{谷 12:3} 那些佃户却打了那仆人，44 叫他空手回去。\BibleRef{谷 12:4} 他又打发另一个仆人去；他们45 用石头打他，伤了他，并且羞辱地打发他回去。\BibleRef{谷 12:5} 他又打发另一个去；他们就杀了他。他又打发许多别的仆人46 去。\BibleRef{玛 21:35} 佃户拿住他的仆人，打了一个，47 用石头砸了一个，杀了一个。\BibleRef{玛 21:36} 他又打发别的仆人，比前更多；48〔阿拉伯语，第128页〕他们也照样对待了他们。\BibleRef{路 20:13} 葡萄园的主人就说：我该怎么办呢？我要打发我的爱子去；或许他们见了就会49,50 惭愧。\BibleRef{谷 12:6} 最后他打发他的爱子去。\BibleRef{玛 21:38} 佃户看见儿子，就彼此说：这是继承人。51,52 \BibleRef{路 20:14} 他们说：我们杀了他，产业就归我们了。\BibleRef{玛 21:39} 于是他们抓住53 他，推出葡萄园外，杀了他。\BibleRef{玛 21:40} 那么，葡萄园的主人54 来的时候，要怎样处置那些佃户呢？\BibleRef{玛 21:41} 他们对他说：他要凶恶地消灭那些恶人，把葡萄园55 租给别的佃户，就是那按时交果子的。\BibleRef{玛 21:42} 耶稣对他们说：圣经上你们从来没有读过吗？\par

匠人弃而不用的石头，\newline{}\BibleRef{路 20:17} 反而成了屋角的基石：\newline{}56 \BibleRef{玛 21:42} 这是出于天主的作为，\newline{}在我们眼中，真是奇妙？\par

57 \BibleRef{玛 21:43} 因此我告诉你们，天主的国必从你们手中夺去，58 赐给一个能结果实的民族。\BibleRef{玛 21:44} 谁掉在这石头上，必被粉碎；但谁被它砸到，它必59 把他压碎成粉末。\BibleRef{玛 21:45} 司祭长和法利塞人听了他的比喻，60 明白他是指着他们说的。\BibleRef{玛 21:46} 他们就想要捉拿他；但惧怕群众，因为人们都拿他当作先知。\par

\SourceRecord{Translated by Hope W. Hogg.；Ante-Nicene Fathers；Vol. 9.；Edited by Allan Menzies.；Buffalo, NY: Christian Literature Publishing Co.,；1896.；<http://www.newadvent.org/fathers/100233.htm>.}

% structure: p0001 source=h1 target=section reason=page-title

\section{四福音合参（第34节）}

1 那时，法利塞人前去商议，如何就他的话语陷害他，2 将他交在审判官的手中，并交在掌权者的手中。\BibleRef{玛 22:16} 他们就打发自己的门徒，同黑落德党的人，到他跟前说：\Quote{\EditorialNote{阿拉伯文第129页}师傅，我们知道你说话真诚，并且本着公平教导天主的道，也不受任何人抬举：因为你行事不寻求人的眼见。3 \BibleRef{玛 22:17} 请告诉我们，你以为如何？向凯撒纳税，可不可以？我们该不该给？}\BibleRef{谷 12:15} 耶稣看透他们的奸诈，就对他们说：\Quote{你们为什么试探我，假善人？4 拿一个税银给我看。}\BibleRef{玛 22:20} 他们就拿一个银钱给他。耶稣对他们说：\Quote{这像和这名号是谁的？}5,6 他们回答说：\Quote{凯撒的。}\BibleRef{玛 22:21} 耶稣对他们说：\Quote{那么，凯撒的归给凯撒，天主的归给天主。}\BibleRef{路 20:26} 他们不能当民众从他的话语中找到把柄，就希奇他的回答，闭口不言。\par

9 \BibleRef{玛 22:23} 在那一天，撒杜塞人来到他跟前，说没有死人复活的事。10 \BibleRef{玛 22:24} 他们问他，说：\Quote{师傅，梅瑟对我们说过：若一个人死了，没有子女，他的弟弟应娶他的妻子，为哥哥立后。11 从前在我们这里有弟兄七人：第一个娶了妻，12 没有子女就死了；\BibleRef{路 20:30} 第二个娶了她的妻子，也没有子女就死了；13 \BibleRef{路 20:31} 第三个也娶了她；七个人都如此，没有留下子女就死了。14,15 \BibleRef{玛 22:27} 最后，那妇人也死了。\BibleRef{玛 22:28} 那么，复活的时候，这妇人该是哪一个人的妻子？因为他们都娶过她。}16 耶稣回答说：\Quote{你们错了，岂不是因为不明白圣经，也不晓得天主的能力吗？17 \BibleRef{路 20:34} 今世之子娶妻，女人出嫁；18 \BibleRef{路 20:35} 但那些得以进入那世界、并从死人中复活的人，\EditorialNote{阿拉伯文第130页}不娶不嫁，19 \BibleRef{路 20:36} 因为他们不能再死；他们像天使一样，是天主的儿女，因为他们成了复活之子。20 关于死人复活的事，你们没有读过梅瑟书中荆棘篇上天主对他所说的话吗？\InnerQuote{我是亚巴郎的天主，依撒格的天主，雅各伯的天主。}21 天主不是死人的天主，而是活人的天主：因为众人都为他而活着。你们大错了。}\par

22,23 \BibleRef{玛 22:33} 群众听见了，都惊奇他的教训。\BibleRef{路 20:39} 有几个经师回答说：\Quote{师傅，你说得好。}24 \BibleRef{玛 22:34} 其余的法利塞人见他使撒杜塞人在这事上闭口无言，就聚集起来同他争辩。\par

25 有一个经师，是那些熟悉法律的，见他回答卓越，就想试探他，问他说：\BibleRef{路 10:25}\Quote{我该做什么才能承受永生？}并且，\BibleRef{谷 12:28}\Quote{诫命中哪一条最大，在律法中居首位？}27 \BibleRef{谷 12:29} 耶稣对他说：\Quote{第一条最大的诫命是：\InnerQuote{以色列啊，请听！主我们的天主是唯一的主；28 你当全心、全灵、全意、全力爱上主你的天主。}29,30 \BibleRef{玛 22:38} 这是最大且首要的诫命。\BibleRef{谷 12:31} 第二条相似的是：\InnerQuote{你当爱邻人如同自己。}再没有比这两条更大的诫命。31 \BibleRef{玛 22:40} \EditorialNote{阿拉伯文第131页}全律法和先知都系于这两条诫命。}\EditorialNote{阿拉伯文第131页}32 \BibleRef{谷 12:32} 那经师对他说：\Quote{好极了，我主！你所说他只有一位，除他以外再没有别的，33 \BibleRef{谷 12:33} 并且人当全心、全意、全灵、全力爱他，并爱邻人如同自己，远胜于一切全燔祭和牺牲。}34 \BibleRef{谷 12:34} 耶稣见他回答得明智，就回答说：\Quote{你离天主的国不远了。}35,36 \BibleRef{路 10:28} \Quote{你说得对：你这样做，就必得生。}\BibleRef{路 10:29} 那人要想显明自己有理，就对耶稣说：\Quote{谁是我的邻人？}37 \BibleRef{路 10:30} 耶稣回答说：\Quote{有一个人从耶路撒冷下到耶里哥去，落在强盗手中；他们剥去他的衣服，把他打伤，丢下他半死不活就走了。38 \BibleRef{路 10:31} 恰巧有一个司铎从那条路下来，看见他，就从旁边过去了。39 \BibleRef{路 10:32} 同样，有一个肋未人来到那地方，看见他，也从旁边过去了。40 \BibleRef{路 10:33} 但有一个撒玛黎雅人，行路来到他那里，看见他，就动了怜悯的心，41 \BibleRef{路 10:34} 上前给他包扎伤口，倒上油和酒；把他扶上自己的牲口，带到客店里，细心照料他。42 \BibleRef{路 10:35} 第二天，他拿出两个银钱，交给店主说：\InnerQuote{请你照料他；若再多花费，我回来时必定还你。}43 \BibleRef{路 10:36} 你以为这三个人中，谁是那落在强盗手中的近人呢？}44 \BibleRef{路 10:37} 他说：\Quote{\EditorialNote{阿拉伯文第132页}是那怜悯他的。}45 \BibleRef{谷 12:34} 耶稣对他说：\Quote{你去，也照样做吧。}从此以后，没有人敢再问他什么。\par

46 \BibleRef{路 19:47}祂每日在殿里施教。但司祭长、经师和百姓的长老设法要除灭祂；\BibleRef{路 19:48}却\Emph{寻}不着办法如何处置祂，因为众百姓都侧耳听祂。48 \BibleRef{若 7:31}群众中有许多人信了祂，且说：\Quote{默西亚一来，祂所行的神迹难道会比这\Emph{人}所行的更多吗？}\BibleRef{若 7:32}法利塞人听见群众对祂这样议论，司祭长就派遣差役去捉拿祂。\BibleRef{若 7:33}耶稣遂对他们说：\Quote{我还有不多的时候与你们同在，然后我就要回到那派遣我者那里去。\BibleRef{若 7:34}你们要寻找我，却找不着；我所在的地方，你们也不能到。}\BibleRef{若 7:35}犹太人便彼此说：\Quote{这\Emph{人}要往哪里去，使我们\Emph{找}不着他呢？难道祂要往散居在\OriginalTerm{外邦人}{Gentiles}中的犹太人那里去，教训\OriginalTerm{外邦人}{Gentiles}么？}\BibleRef{若 7:36}犹太人便彼此说：\Quote{祂所说的\InnerQuote{你们要寻找我，却找不着；我所在的地方，你们也不能到}这句话，是什么意思？}\par

\SourceRecord{Translated by Hope W. Hogg.；Ante-Nicene Fathers；Vol. 9.；Edited by Allan Menzies.；Buffalo, NY: Christian Literature Publishing Co.,；1896.；<http://www.newadvent.org/fathers/100234.htm>.}

% structure: p0001 source=h1 target=section reason=page-title

\section{\WorkTitle{迪亚泰萨隆}（第35节）}

1 \BibleRef{若 7:37} \BibleQuote{在那节期的最后一天，就是最盛大的日子，耶稣站着高声说：} 2 \BibleRef{若 7:38} \BibleQuote{凡信我的，就如经上所说，从他的腹中要流出活水的江河来。} 3 \BibleRef{若 7:39} \BibleQuote{耶稣这话是指着信他之人要领受的圣神说的；那时圣神还没有赐下，因为耶稣尚未得着荣耀。} \EditorialNote{阿拉伯文版，第133页} 4 \BibleRef{若 7:40} \BibleQuote{众人听见这话，有的说：这真是那先知。} 5 \BibleRef{若 7:41} \BibleQuote{有的说：这是默西亚。但也有的说：岂能是默西亚从加利利出来吗？} 6 \BibleRef{若 7:42} \BibleQuote{经上岂不是说：默西亚是大卫的后裔，从大卫的故乡伯利恒出来吗？} 7 \BibleRef{若 7:43} \BibleQuote{于是众人因着他起了纷争。} 8 \BibleRef{若 7:44} \BibleQuote{其中有人想要捉拿他，只是没有人下手。}\par

9 \BibleRef{若 7:45} \BibleQuote{差役回到司祭长和法利塞人那里，他们对差役说：你们为什么没有带他来呢？} 10 \BibleRef{若 7:46} \BibleQuote{差役回答说：从来没有像他这样说话的。} 11 \BibleRef{若 7:47} \BibleQuote{法利塞人说：你们也受了迷惑吗？} 12 \BibleRef{若 7:48} \BibleQuote{官长或是法利塞人岂有信他的呢？} 13 \BibleRef{若 7:49} \BibleQuote{但这些不明白律法的百姓是被咒诅的。} 14 \BibleRef{若 7:50} \BibleQuote{内中有尼苛德摩，就是从前夜里去见耶稣的，} 15 \BibleRef{若 7:51} \BibleQuote{对他们说：不先听本人的口供，并知道他所作的事，难道我们的律法还定他的罪吗？} 16 \BibleRef{若 7:52} \BibleQuote{他们回答说：你也是从加利利出来的吗？你且去查考，就可知道加利利没有出过先知。}\par

17 \BibleRef{玛 22:41} \BibleQuote{法利塞人聚集的时候，耶稣问他们说：} 18 \BibleRef{玛 22:42} \BibleQuote{论到默西亚，你们的意见如何？他是谁的子孙呢？他们回答说：是大卫的子孙。} 19 \BibleRef{玛 22:43} \BibleQuote{耶稣说：这样，大卫被圣神感动，怎么还称他为主呢？他说：}\par

20 \BibleRef{玛 22:44} \BibleQuote{主对我主说：\newline{}你坐在我的右边，\newline{}等我把你的仇敌放在你的脚下。}\par

21 \BibleRef{玛 22:45} \BibleQuote{大卫既称他为主，他怎么又是大卫的子孙呢？} 22 \BibleRef{玛 22:46} \BibleQuote{他们没有一个人能回答他。从那日以后，也没有人敢再问他什么。}\par

23 \BibleRef{若 8:12} \BibleQuote{耶稣又对众人说：我是世界的光；跟从我的，就不在黑暗里走，必要得着生命的光。} 24 \BibleRef{若 8:13} \BibleQuote{法利塞人对他说：你是为自己作见证，你的见证不真。} \EditorialNote{阿拉伯文版，第134页} 25 \BibleRef{若 8:14} \BibleQuote{耶稣说：我虽然为自己作见证，我的见证还是真的；因我知道我从哪里来，往哪里去；你们却不知道我从哪里来，往哪里去。} 26 \BibleRef{若 8:15} \BibleQuote{你们是以外貌判断人，我却不判断人。} 27 \BibleRef{若 8:16} \BibleQuote{就是判断人，我的判断也是真的；因为不是我独自在这里，还有差我来的父与我同在。} 28 \BibleRef{若 8:17} \BibleQuote{你们的律法上也记着说：两个人的见证是真的。} 29 \BibleRef{若 8:18} \BibleQuote{我是为自己作见证，还有差我来的父也是为我作见证。} 30 \BibleRef{若 8:19} \BibleQuote{他们就问他说：你的父在哪里？耶稣回答说：你们不认识我，也不认识我的父；若是认识我，也就认识我的父。} 31 \BibleRef{若 8:20} \BibleQuote{这些话是耶稣在殿里的库房，教训人时所说的；也没有人拿他，因为他的时候还没有到。} 32 \BibleRef{若 8:21} \BibleQuote{耶稣又对他们说：我要去了，你们要找我，却找不到我，并且你们要死在罪中；我所去的地方，你们不能到。} 33 \BibleRef{若 8:22} \BibleQuote{犹太人说：他说\InnerQuote{我所去的地方，你们不能到}，难道他要自杀吗？} 34 \BibleRef{若 8:23} \BibleQuote{耶稣对他们说：你们是从下头来的，我是从上头来的；你们是属这世界的，我不是属这世界的。} 35 \BibleRef{若 8:24} \BibleQuote{所以我对你们说：你们要死在罪中；你们若不信我就是那一位，必要死在罪中。} 36 \BibleRef{若 8:25} \BibleQuote{他们就问他说：你是谁？耶稣对他们说：就是我从起初所告诉你们的。} 37 \BibleRef{若 8:26} \BibleQuote{我有许多事讲论你们，判断你们；但那差我来的是真的，我在他那里所听见的，我就传给世人。} 38 \BibleRef{若 8:27} \BibleQuote{他们不明白耶稣是指着父说的。} \EditorialNote{阿拉伯文版，第135页} 39 \BibleRef{若 8:28} \BibleQuote{所以耶稣说：你们举起人子以后，必知道我就是那一位，并且我没有一件事是凭着自己作的；我说这些话，乃是照着父所教训我的。} 40 \BibleRef{若 8:29} \BibleQuote{那差我来的是与我同在；他没有撇下我独自在这里，因为我常作他所喜悦的事。} 41 \BibleRef{若 8:30} \BibleQuote{耶稣说这话的时候，就有许多人信他。}\par

42 \BibleRef{若 8:31} 耶稣对那些信他的犹太人说：\Quote{如果你们存留在我话里，你们就真是我的门徒；\BibleRef{若 8:32} 并且你们将认识真理，真理必使你们自由。} 44 \BibleRef{若 8:33} 他们对他说：\Quote{我们是亚巴郎的后裔，从未做过任何人的奴隶45在奴役的方式中：你怎么说：你们将做自由之子？} \BibleRef{若 8:34} 耶稣对他们说：\Quote{我实实在在告诉你们，凡犯罪的，就是罪的奴隶46。\BibleRef{若 8:35} 奴隶不会永远住在家里，但儿子47,48直到永远。\BibleRef{若 8:37} 所以，如果圣子使你们自由，你们就真是自由之子。我知道你们是亚巴郎的后裔，但你们却想杀我，因为你们不能接纳我的话49。\BibleRef{若 8:38} 我所看见的在我父那里，我就说；而你们所看见的在你们父那里，你们就做50。} \BibleRef{若 8:39} 他们回答说：\Quote{我们的父亲是亚巴郎。} 耶稣对他们说：\Quote{如果你们是亚巴郎的子女，你们就会做亚巴郎所做的事。\BibleRef{若 8:40} 现在51，看，你们想杀我——这个给你们讲真理的人，这真理是我从天主那里听到的52；亚巴郎没有做过这样的事。\BibleRef{若 8:41} 你们却做你们父所做的事。} 他们对他说：\Quote{我们不是\Emph{出生}于淫乱，我们只有一个父，就是53天主。} \BibleRef{若 8:42} 耶稣对他们说：\Quote{如果天主是你们的父，你们就会爱我；因为我从天主那里出来而来到；我不是凭自己来的，而是他派遣了我54[阿拉伯文，第136页]。\BibleRef{若 8:43} 为什么你们不明白我的话呢？因为你们不能听我的话。\BibleRef{若 8:44} 你们出于父亲—魔鬼55，你们愿意随从你们父亲的欲望。他从起初就是杀人者，不站在真理中，因为真理不在他内。当他说谎时，是出于他自己，因为他是说谎者，也是谎言之父56。\BibleRef{若 8:45} 我虽然讲真理，你们57却不相信我。\BibleRef{若 8:46} 你们中谁能指证我有罪呢？既然我讲真理，你们58为什么不信我？\BibleRef{若 8:47} 出于天主的，必听天主的话；因此你们59不听，因为你们不出于天主。} \BibleRef{若 8:48} 犹太人回答说：\Quote{我们说你是个撒玛黎雅人，并附有魔鬼，60岂不正对吗？} \BibleRef{若 8:49} 耶稣对他们说：\Quote{我并没有附魔，我尊敬我的父，你们却61侮辱我。\BibleRef{若 8:50} 我不求自己的光荣；有一位为我求，并判断。}\par

\SourceRecord{Translated by Hope W. Hogg.；Ante-Nicene Fathers；Vol. 9.；Edited by Allan Menzies.；Buffalo, NY: Christian Literature Publishing Co.,；1896.；<http://www.newadvent.org/fathers/100235.htm>.}

% structure: p0001 source=h1 target=section reason=page-title

\section{\WorkTitle{Diatessaron}（第三十六段）}

1 \BibleRef{若 8:51}我实实在在告诉你们：谁若遵守我的话，就永远见不到死亡。2 \BibleRef{若 8:52}犹太人对他说：\Quote{现在我们知道你附了魔。亚巴郎死了，先知们也死了，你却说：谁若遵守我的话，就永远尝不到死味。3 \BibleRef{若 8:53}难道你比我们的父亲亚巴郎还大吗？他死了，先知们也死了。你把自己当作什么人呢？}4 \BibleRef{若 8:54}耶稣回答说：\Quote{我若光荣我自己，我的光荣算不了什么；光荣我的，是我的父，就是你们所称的\Emph{那一位}，\InnerQuote{我们的天主}；5 \BibleRef{若 8:55}你们\Emph{却}不认识他，我认识他；我若说不认识他，我就像你们一样是个说谎者；但我认识他，也遵守他的话。6 \BibleRef{若 8:56}你们的父亲亚巴郎曾欢欣期望看见我的日子，他看见了，就欢乐起来。}7 \BibleRef{若 8:57}犹太人对他说：\Quote{你还不到五十岁，就见过亚巴郎吗？}8 \BibleRef{若 8:58}耶稣对他们说：\Quote{我实实在在告诉你们：在亚巴郎出现以前，我就是。}9 \BibleRef{若 8:59}他们就拿起石头来要砸他；耶稣却隐藏了，从圣殿里出去。他穿过他们中间，走了。[阿拉伯文，第137页]\par

10 \BibleRef{若 9:1}耶稣前行时，看见了一个生来就瞎眼的人。11 \BibleRef{若 9:2}他的门徒就问他说：\Quote{老师，谁犯了罪？是这个\Emph{人}，还是他的父母，竟使他生来瞎眼呢？}12 \BibleRef{若 9:3}耶稣回答说：\Quote{不是他犯了罪，也不是他的父母，而是为叫天主的工作在他身上显扬出来。13 \BibleRef{若 9:4}趁着白天，我必须做那派遣我者的工作；黑夜要来，那时没有人能工作。14 \BibleRef{若 9:5}我在世界上的时候，我是世界的光。}15 \BibleRef{若 9:6}说完这话，便吐唾沫在地上，用唾沫和泥，把\Emph{它}抹在瞎子的眼上，16 对他说：\Quote{\BibleRef{若 9:7}你去\Place{西罗亚}池里洗一洗吧。}\BibleRef{若 9:8}他就去了，洗了，回来就看见了。他的邻人和那些素常见他讨饭的人就说：\Quote{这不是那曾坐着讨饭的人吗？}18 \BibleRef{若 9:9}有的说：\Quote{是他}，有的说：\Quote{不，只是像他。}他自己说：\Quote{就是我。}19,20 \BibleRef{若 9:10}他们问他说：\Quote{你的眼睛是怎样开的呢？}\BibleRef{若 9:11}他回答说：\Quote{名叫耶稣的那个人和了泥，抹在我的眼上，对我说：\InnerQuote{你去\Place{西罗亚}池里洗吧}；我去了，洗了，就看见了。}21 \BibleRef{若 9:12}他们问他说：\Quote{那个人在那里？}他说：\Quote{我不知道。}22,23 [阿拉伯文，第138页]\BibleRef{若 9:13}他们便将先前瞎眼的人带到法利塞人那里。\BibleRef{若 9:14}耶稣和泥开他眼睛的那天正是安息日。24 \BibleRef{若 9:15}法利塞人又问他是怎样看见的。他回答说：\Quote{他把泥抹在我的眼上，我洗了，就看见了。}25 \BibleRef{若 9:16}法利塞人中有的说：\Quote{这人不是从天主来的，因为他不守安息日。}有的却说：\Quote{一个\Emph{罪人}怎能行这样的奇迹呢？}他们中间便发生了分裂。26 \BibleRef{若 9:17}于是他们又对瞎子说：\Quote{你说那个人开了你的眼睛，你以为他是什么人？}他说：\Quote{我以为他是先知。}27 \BibleRef{若 9:18}犹太人不信他以前是瞎子而后来看见了，直到叫了那复明者的父母来，28 问他们说：\Quote{\BibleRef{若 9:19}这是你们的儿子吗？你们说\Emph{他}生来就是瞎眼的，怎么现在他看见了呢？}29 \BibleRef{若 9:20}他的父母回答说：\Quote{我们知道这是我们的儿子，也生来就是瞎眼；30 \BibleRef{若 9:21}但他如今怎么能看见，我们却不知道；或者谁开了他的眼睛，我们也不知道；他已经成年，你们问他吧，他自己会说的。}31 \BibleRef{若 9:22}他的父母这样说，是因为怕犹太人：原来犹太人已经议定，谁若认耶稣是默西亚，就把他赶出会堂。32 \BibleRef{若 9:23}为此，他的父母说\Quote{他已经成年，你们问他吧。}33 \BibleRef{若 9:24}于是他们又把那瞎眼的人叫来，对他说：\Quote{归光荣于天主吧！我们知道这人是个罪人。}34 \BibleRef{若 9:25}那人回答说：\Quote{他是不是罪人，我不知道；我只知道一件事：我原是瞎子，现在我看见了。}35 \BibleRef{若 9:26}他们又问他说：\Quote{他给你做了什么？怎样开了你的眼睛？}36 [阿拉伯文，第139页]\BibleRef{若 9:27}他回答说：\Quote{我已经告诉你们了，你们不听；为什么又要听呢？莫非你们也要做他的门徒吗？}37 \BibleRef{若 9:28}他们辱骂他说：\Quote{你是那个\Emph{人}的门徒；我们是\Person{梅瑟}的门徒。38 \BibleRef{若 9:29}我们知道天主曾对\Person{梅瑟}说过话；至于这个人，我们不知道他是从那里来的。}39 \BibleRef{若 9:30}那人回答说：\Quote{这真是奇怪！你们不知道他是从那里来的，他却开了我的眼睛。40 \BibleRef{若 9:31}我们知道天主不俯听罪人，只俯听那敬畏天主并遵行他旨意的人。41 \BibleRef{若 9:32}自古以来从未听说有人开了生来瞎眼的\Emph{人}的眼睛。42 \BibleRef{若 9:33}\Emph{这人}若不是从天主来的，他什么也不能做。}43 \BibleRef{若 9:34}他们回答说：\Quote{你整个生在罪中，竟敢教训我们吗？}便把他赶出去了。\par

44 \BibleRef{若 9:35}耶稣听说他们把他赶出去了，找到他，就向他说：\Quote{你信天主子吗？}45 \BibleRef{若 9:36}那复明的人回答说：\Quote{主，他是谁，好使我相信他呢？}46 \BibleRef{若 9:37}耶稣对他说：\Quote{你已经看见了他，现在和你说话的就是他。}47 \BibleRef{若 9:38}他说：\Quote{主，我信。}就俯伏朝拜了他。\par

\SourceRecord{Translated by Hope W. Hogg.；Ante-Nicene Fathers；Vol. 9.；Edited by Allan Menzies.；Buffalo, NY: Christian Literature Publishing Co.,；1896.；<http://www.newadvent.org/fathers/100236.htm>.}

% structure: p0001 source=h1 target=section reason=page-title

\section{\WorkTitle{四福音合参}（第37部分）}

\BibleRef{若 9:39}耶稣说：\Quote{我是为审判这个世界而来，叫那看不见的得以看见，那看得见的反而成为瞎子。}\BibleRef{若 9:40}同他在一起的几个法利塞人听见了这话，就对他说：\Quote{难道我们也是瞎子吗？}\BibleRef{若 9:41}耶稣对他们说：\Quote{如果你们是瞎子，就没有罪了；但你们却说：\InnerQuote{我们看得见}，因此你们的罪还在。}\par

[阿拉伯文，第140页]\BibleRef{若 10:1}\BibleQuote{我实实在在告诉你们：凡不从门进入羊栈，而从别处爬进去的，那\Emph{人}就是贼和强盗。}\BibleRef{若 10:2}\BibleQuote{但由门进入的，才是羊的牧人。}\BibleRef{若 10:3}\BibleQuote{看门的就给他开门，羊也听他的声音；他按名字呼唤自己的羊，并引领他们出来。}\BibleRef{若 10:4}\BibleQuote{他把羊放出来以后，就走在羊的前面，羊也跟随他，因为认得他的声音。}\BibleRef{若 10:5}\BibleQuote{羊决不跟随陌生人，反而逃避他，因为不认得陌生人的声音。}\BibleRef{若 10:6}耶稣给他们讲了这个比喻，他们却不明白他说的是什么。\par

\BibleRef{若 10:7}耶稣又对他们说：\BibleQuote{我实实在在告诉你们：我就是羊的门。}\BibleRef{若 10:8}\BibleQuote{凡在我以前来的，都是贼和强盗，羊却不听他们。}\BibleRef{若 10:9}\BibleQuote{我就是门；谁若经过我进来，必得安全，可以进，可以出，也能找到草场。}\BibleRef{若 10:10}\BibleQuote{盗贼来，无非为偷窃、杀害、毁灭；我来，却是为叫他们获得生命，且获得那\Emph{更好的}。}\BibleRef{若 10:11}\BibleQuote{我是善牧：善牧为羊舍命。}\BibleRef{若 10:12}\BibleQuote{佣工，不是牧人，羊也不是他自己的，一见狼来，便撇下羊逃跑，狼就抓住羊，把羊群驱散。}\BibleRef{若 10:13}\BibleQuote{佣工逃跑，因为他是佣工，对羊漠不关心。}\BibleRef{若 10:14}\BibleQuote{我是善牧，}\BibleRef{若 10:15}\BibleQuote{我认识我的羊，我的羊也认识我，正如父认识我，我也认识父一样；我并且为羊舍命。}\BibleRef{若 10:16}\BibleQuote{我还有别的羊，不属于这一栈，我也该把他们引回来，他们要听我的声音，这样，将只有一个羊群，一个牧人。}\BibleRef{若 10:17}\BibleQuote{父爱我，因为我舍弃我的性命，为再取回它来。}\BibleRef{若 10:18}\BibleQuote{没有人夺去我的性命，而是我甘心情愿舍弃的。我有权舍弃它，也有权再取回它来。这是我由我父所接受的命令。}\par

\BibleRef{若 10:19}因了这些话，犹太人中间又发生了分裂。\BibleRef{若 10:20}他们中有许多人说：\BibleQuote{他附了魔，疯了；为什么还听他的？}\BibleRef{若 10:21}另有些人说：\BibleQuote{这些话不是附魔的\Emph{人}所能说的；难道魔鬼能开瞎\Emph{人}的眼睛吗？}\par

\BibleRef{若 10:22}在耶路撒冷有重建节，正是冬天。\BibleRef{若 10:23}耶稣在圣殿里，在撒罗满廊下行走。\BibleRef{若 10:24}犹太人围起他来，对他说：\BibleQuote{你使我们的心神悬疑不定，要到几时呢？如果你是\OriginalTerm{默西亚}{Messiah}，就坦白告诉我们吧！}\BibleRef{若 10:25}耶稣回答他们说：\BibleQuote{我已经告诉你们，你们却不相信；我因我父的名所行的这些事，为我作证。}\BibleRef{若 10:26}\BibleQuote{但你们不信，因为你们不是我的羊，正如我已对你们说过的。}\BibleRef{若 10:27}\BibleQuote{我的羊听我的声音，我也认识他们，他们也跟随我；}\BibleRef{若 10:28}\BibleQuote{我赐给他们永生，他们决不会丧亡，谁也不能从我手中把他们夺去。}\BibleRef{若 10:29}\BibleQuote{我父，把\Emph{他们}赐给了我，他超越一切，所以谁也不能从我父手里把\Emph{他们}夺去。}\BibleRef{若 10:30}\BibleQuote{我与父原是一体。}\BibleRef{若 10:31}犹太人又拿起石头来，要砸死他。\BibleRef{若 10:32}耶稣对他们说：\BibleQuote{我从父那里显示出许多善事给你们看，你们为了哪一件要砸死我呢？}\BibleRef{若 10:33}犹太人回答他说：\BibleQuote{为了善事，我们不会砸死你；而是为了亵渎的话，因为你是一个人，却把自己当作天主。}\BibleRef{若 10:34}耶稣对他们说：\BibleQuote{你们的法律上不是记载着：\InnerQuote{我说过：你们是神}么？}\BibleRef{若 10:35}\BibleQuote{如果那些承受天主话的人，尚且被称为神——而圣经是不能废弃的——}\BibleRef{若 10:36}\BibleQuote{那么父所祝圣并派遣到世界上来的，你们却说：你说亵渎的话，因为我说了\InnerQuote{我是天主子}么？}\BibleRef{若 10:37}\BibleQuote{如果我不做我父的事，你们就不必信我；}\BibleRef{若 10:38}\BibleQuote{如果我做了，\Emph{即使}你们不信我，至少该信这些事，为使你们认识并相信：父在我内，我在父内。}\BibleRef{若 10:39}他们又想捉拿他，他却从他们手中走脱了。\par

\BibleRef{若 10:40}耶稣又到约旦河对岸，就是若翰先前施洗的地方，并住下了。\BibleRef{若 10:41}有许多人来到他那里说：\BibleQuote{若翰固然没有行过神迹，但若翰关于这人所说的一切，都是真的。}\BibleRef{若 10:42}许多人就在那里信了耶稣。\par

46 \BibleRef{若 11:1} 有一个患病的\Emph{人}，名叫\Person{拉匝禄}，住在伯达尼村，他是\Person{玛利亚}和\Person{玛尔大}的兄弟。47 \BibleRef{若 11:2} 这玛利亚就是用香膏傅抹耶稣的脚，并用头发擦\Emph{它们}的妇人；患病的拉匝禄是这\Emph{女人}的兄弟。48 \BibleRef{若 11:3} 她的姊妹打发人去找耶稣说：\BibleQuote{主，你所爱的人病了。}49 \BibleRef{若 11:4} 耶稣却说：\BibleQuote{这病不至于死，但为彰显天主的荣耀，好使天主子因此受到光荣。}50, 51 \BibleRef{若 11:5} 耶稣素爱\Person{玛尔大}、\Person{玛利亚}和\Person{拉匝禄}。\BibleRef{若 11:6} 他听说拉匝禄病了，仍在原地住了两天。52 \BibleRef{若 11:7} 此后，他对门徒说：\BibleQuote{我们再到犹太去吧。}53 \BibleRef{若 11:8} 门徒对他说：\BibleQuote{我们的老师 [Arabic, p. 143]，犹太人图谋砸死你，你还要到那里去吗？}54, 55 \BibleRef{若 11:9} 耶稣对他们说：\BibleQuote{白日不是有十二小时吗？人若在白日行走，不会跌倒，因为他看见世界的光。}\BibleRef{若 11:10} \BibleQuote{人若在夜间行走，就会跌倒，因为他没有光在身上。}56 \BibleRef{若 11:11} 耶稣说了这话，然后对他们说：\BibleQuote{我们的朋友拉匝禄睡着了，我要去唤醒他。}57 \BibleRef{若 11:12} 门徒对他说：\BibleQuote{主，他若睡着了，必定会好的。}58 \BibleRef{若 11:13} 耶稣是指他的死说的，门徒却以为他说的是安眠入睡。59 \BibleRef{若 11:14} 耶稣就明明地对他们说：\BibleQuote{拉匝禄死了。}60 \BibleRef{若 11:15} \BibleQuote{为了你们的缘故，我不在那里反而高兴，好叫你们相信；我们到他那里去吧。}61 \BibleRef{若 11:16} \Person{多默}——又称为塔玛，对同伴门徒说：\BibleQuote{我们也去，和他一起死吧。}\par

\SourceRecord{Translated by Hope W. Hogg.；Ante-Nicene Fathers；Vol. 9.；Edited by Allan Menzies.；Buffalo, NY: Christian Literature Publishing Co.,；1896.；<http://www.newadvent.org/fathers/100237.htm>.}

% structure: p0001 source=h1 target=section reason=page-title

\section{四福音合参 第三十八节}

\BibleQuote{1, 2 \BibleRef{若 11:17} \Person{耶稣}来到\Place{伯大尼}，发现他已经\Emph{已经}在坟墓里四天了。\BibleRef{若 11:18} \Place{伯大尼}靠近\Place{耶路撒冷}，相距约十五\OriginalTerm{斯塔狄}{furlongs}；3 \BibleRef{若 11:19} 许多\Person{犹太人}来到\Person{玛尔大}和\Person{玛利亚}那里，要因她们兄弟的事安慰她们的心。4 \BibleRef{若 11:20} \Person{玛尔大}听见\Person{耶稣}来了，就出去迎接他，\Person{玛利亚}却坐在家里。5 \BibleRef{若 11:21} \Person{玛尔大}于是对\Person{耶稣}说：\Quote{我主，你若早在这里，我的兄弟就不会死了。6 \BibleRef{若 11:22} 但如今我知道，你无论向\Person{天主}求什么，\Person{他}必赐给你。}7 \BibleRef{若 11:23} \Person{耶稣}对她说：\Quote{你的兄弟必要复活。}8 \BibleRef{若 11:24} \Person{玛尔大}对他说：\Quote{我知道他在末日复活时必要复活。}9 \BibleRef{若 11:25} \Person{耶稣}对她说：\Quote{我是复活，是生命：信我的人，即使死了，仍要活着；10 \BibleRef{若 11:26} \EditorialNote{阿拉伯文版第144页} 凡活着而信我的人，必永远不死。你信这话吗？}11 \BibleRef{若 11:27} 她对他说：\Quote{是，我主，我信你是\OriginalTerm{默西亚}{Messiah}，\Person{天主子}，要来到世界上的那一位。}12 \BibleRef{若 11:28} 她说了这话，就去暗暗地叫她妹妹\Person{玛利亚}，说：\Quote{\Person{师傅}来了，他在叫你。}13 \BibleRef{若 11:29} \Person{玛利亚}一听见，就急忙起来，到他那里去。14 \BibleRef{若 11:30} （那时\Person{耶稣}还没有进村庄，仍在\Person{玛尔大}迎接他的地方。）15 \BibleRef{若 11:31} 那些在屋里安慰她的\Person{犹太人}，见\Person{玛利亚}急忙起来出去，就跟随着她，以为她往坟墓那里去哭泣。16 \BibleRef{若 11:32} \Person{玛利亚}来到\Person{耶稣}所在的地方，看见他，就俯伏在他脚前，对他说：\Quote{我主，你若早在这里，我的兄弟就不会死了。}17 \BibleRef{若 11:33} \Person{耶稣}看见她哭泣，同她一起来的\Person{犹太人}也哭泣，就心里悲叹，难过起来，\BibleRef{若 11:34} 便说：\Quote{你们把他安放在哪里？}他们对他说：\Quote{主，请来看。}19 \BibleRef{若 11:35} \Person{耶稣}流泪了。20 \BibleRef{若 11:36} \Person{犹太人}便说：\Quote{看，他多么爱他！}21 \BibleRef{若 11:37} 但他们中有人说：\Quote{这个\Emph{人}，他开了那个\Emph{瞎子}的眼睛，难道不能也使这个\Emph{人}不死吗？}22 \BibleRef{若 11:38} \Person{耶稣}又心里悲叹，来到坟墓前。坟墓是个洞，上面有块石头挡着。23 \BibleRef{若 11:39} \Person{耶稣}说：\Quote{把石头\Emph{挪开}。}死人的姐姐\Person{玛尔大}对他说：\Quote{主，他已经臭了，因为已经四天了，他\Emph{死了}。}24 \BibleRef{若 11:40} \Person{耶稣}对她说：\Quote{我不是对你说过，你若信，就必看见\Person{天主}的荣耀吗？}25 \BibleRef{若 11:41} \EditorialNote{阿拉伯文版第145页} 他们便把石头挪开。\Person{耶稣}举目向上说：\Quote{父啊，我感谢你，因为你俯听了我。26 \BibleRef{若 11:42} 我本来知道你总俯听我，但我这样说，是为了四周站立的群众，叫他们信是你派遣了我。}27 \BibleRef{若 11:43} 说了这话，便大声喊说：\Quote{拉匝禄，出来！}28 \BibleRef{若 11:44} 那\Emph{死人}便出来了，手脚缠着布条，脸上包着汗巾。\Person{耶稣}对他们说：\Quote{解开他，让他去。}}\par

\BibleQuote{29 \BibleRef{若 11:45} 许多来到\Person{玛利亚}那里的\Person{犹太人}，看见\Person{耶稣}所做的事，就信了他。30 \BibleRef{若 11:46} 但他们中有些人去到\Person{法利塞人}那里，将\Person{耶稣}所做的事都告诉了他们。}\par

\BibleQuote{31 \BibleRef{若 11:47} \Person{司祭长}和\Person{法利塞人}聚集起来，说：\Quote{我们该怎么办？32 因为这人行了许多\OriginalTerm{神迹}{signs}。33 \BibleRef{若 11:48} 如果这样任凭他，众人都会信从他，\Person{罗马人}就会来夺去我们的土地和人民。}34 \BibleRef{若 11:49} 他们中有一个人，名叫\Person{盖法}，那年他是\Person{大司祭}，对他们说：\Quote{你们什么都不知道，35 \BibleRef{若 11:50} 也不考虑：对我们来说，一个人替人民死，免得全民族灭亡，这更有利。}36 \BibleRef{若 11:51} 他说这话不是出于自己，但因他是那年的大司祭，就\OriginalTerm{预言}{prophesied}\Person{耶稣}将要替人民死；37 \BibleRef{若 11:52} 不但替人民死，还要把\Person{天主}四散的子女聚集归一。38 \BibleRef{若 11:53} 从那一天起，他们就商议\Emph{怎样}杀害他。}\par

\BibleQuote{38 \EditorialNote{阿拉伯文版第146页}\BibleRef{若 11:54} \Person{耶稣}不再公开地在\Person{犹太人}中间往来，却离开那里，往近旷野的地方去，到了一座城，名叫\Place{厄弗辣因}，就和门徒在那里住下。39 \BibleRef{若 11:55} \Person{犹太人}的\OriginalTerm{逾越节}{Passover}近了，许多人从乡间上\Place{耶路撒冷}去，要在节前自洁。40 \BibleRef{若 11:56} 他们就寻找\Person{耶稣}，在\OriginalTerm{圣殿}{temple}里彼此谈论说：\Quote{你们以为怎样？他不来过节吗？}41 \BibleRef{若 11:57} \Person{司祭长}和\Person{法利塞人}早已发出命令：若有人知道他在哪里，就要报告，他们好捉拿他。}\par

\BibleQuote{42 \BibleRef{路 9:51} 当他被接升天的日子将到，他就决意向\Place{耶路撒冷}去，43 \BibleRef{路 9:52} 便打发使者在他前头走，他们去了，进入\Place{撒玛黎雅}的一个村庄，要为他准备。44 \BibleRef{路 9:53} 那里的人不接待他，因为他面向\Place{耶路撒冷}而去。45 \BibleRef{路 9:54} 他的门徒\Person{雅各伯}和\Person{若望}看见了\Emph{这事}，就说：\Quote{主，你愿意我们叫火从天降下，烧灭他们，像\Person{厄里亚}所作的吗？}46 \BibleRef{路 9:55} \Person{耶稣}转过身来，斥责他们，说：\Quote{你们不知道你们是属什么\OriginalTerm{神}{spirit}。47 \BibleRef{路 9:56} \OriginalTerm{人子}{Son of man}来不是要毁灭性命，而是要救性命。}于是他们往别的村庄去了。}\par

\SourceRecord{Translated by Hope W. Hogg.；Ante-Nicene Fathers；Vol. 9.；Edited by Allan Menzies.；Buffalo, NY: Christian Literature Publishing Co.,；1896.；<http://www.newadvent.org/fathers/100238.htm>.}

% structure: p0001 source=h1 target=section reason=page-title

\section{四福音合参（第39节）}

1 \BibleRef{若 12:1} 耶稣在逾越节前六天来到伯达尼，那里有拉匝禄，就是耶稣从死者中复活的那位。\BibleRef{若 12:2} 他们在那里为他设了筵席：3 玛尔大伺候，而拉匝禄也是同他一起坐席的人之一。4 \BibleRef{谷 14:3} 当耶稣在伯达尼癞病人西蒙家里的时候，\BibleRef{若 12:9} 许多犹太人听说耶稣在那里，就来了，不但是为耶稣，也是为要看拉匝禄\EditorialNote{阿拉伯文第147页}，就是他由死者中复活的那位。5、6 \BibleRef{若 12:10} 司祭长商议连拉匝禄也要杀掉；\BibleRef{若 12:11} 因为7 许多犹太人因他而离开，并信了耶稣。\BibleRef{若 12:3} 玛利亚拿了一斤极珍贵的纯哪哒香膏，\BibleRef{谷 14:3} 打开它，8 倒在耶稣的头上，那时他正斜卧着；\BibleRef{若 12:3} 她又抹了耶稣的脚，并用头发擦干；屋里便充满了香膏的香气。9、10 \BibleRef{若 12:4} 但加略人犹达斯，即那要出卖他的门徒之一说：\Quote{为什么不11 把这香膏卖三百银钱，施舍给穷人？}\BibleRef{若 12:6} 他说这话，并不是因为他关心穷人，而是因为他是个贼，且掌管钱囊，12 常取用其中所存。\BibleRef{谷 14:4} 其他门徒也心中不悦，说：\Quote{这香膏为什么13 浪费了呢？}\BibleRef{玛 26:9} 这香膏原可卖许多钱，施舍给穷人14它。\BibleRef{谷 14:5} 他们就对玛利亚生气。\BibleRef{玛 26:10} 耶稣知道了，就对他们说：\Quote{由她吧！为什么难为她？她在我身上做了一件善事：\BibleRef{若 12:7} 因为15 她是为我安葬之日而保存的。\BibleRef{若 12:8} 你们常有穷人在身边，16 若愿意，可以随时向他们行善：\BibleRef{谷 14:7} 但你们不常有我。\BibleRef{玛 26:12} 故此，她将这香膏倒在我身上，是为我的安葬而作的，17 预先膏抹了我的身体。我实在告诉你们：在全世界，无论这福音传到何处，她所行的，也将被传述，作为对她的纪念。}\par

18、19 \EditorialNote{阿拉伯文第148页} \BibleRef{路 19:28} 耶稣说了这话，就从容前行，要上耶路撒冷去。及至他到了贝特法革和伯达尼，靠近那叫橄榄山的山旁，耶稣派了两个门徒，对他们说：\Quote{你们21 往对面的村庄去，一进去，必看见一匹驴拴在那里，还有驴驹同在一处，是从来没有人骑过的；解开，牵来给我。若有人对你们说：\InnerQuote{为什么解它们？}你们要这样说：\InnerQuote{我们23 为主取它们；}他立刻就会让它们来。}这一切事的发生，是为应验先知所说的：\par

24 \BibleRef{玛 21:5} \BibleQuote{你们向熙雍女子说：\newline{}看，你的君王到你这里来，\newline{}温和地骑着一匹驴，\newline{}又骑着一匹驴驹，即驴的崽子。}\par

25 \BibleRef{若 12:16} 门徒当时不明白这事；但等到耶稣受光荣以后，他们才想起这些事是指着他写的，并且26 他们向他做了这事。两个门徒去了，所遇见的正如耶稣27所说的，他们就照耶稣所吩咐的做了。等他们解开它们时，28 主人对他们说：\Quote{为什么解它们？}\BibleRef{路 19:34} 他们回答说：\Quote{我们为主取它们。}他们就任凭他们牵去。他们牵了驴和驴驹，把衣服搭在驴驹上，耶稣就骑上。\BibleRef{玛 21:8} 大部分群众把衣服铺在路上，另有人从树上砍下树枝，铺在路上。\BibleRef{路 19:37} 当他临近橄榄山\EditorialNote{阿拉伯文第149页}的下坡时，众门徒因所见过的一切异能，都欢乐起来，大声赞美天主，32 说：\Quote{在至高之处赞美他！达味之子，赞美他！奉主名而来的当受赞美！33 \BibleRef{谷 11:10} 那即将来临的、我们祖先达味的王国，也是可赞美的！\BibleRef{路 19:38} 在天上有平安，在至高之处有赞美！}\par

34 \BibleRef{若 12:12} 有许多上来过节的人，听说耶稣要来耶路撒冷，35 就拿了棕榈枝，\BibleRef{若 12:13} 出去迎接他，呼喊说：\Quote{赞美他！奉主名来的，以色列的君王，当受赞美！}36 \BibleRef{路 19:39} 有几个法利塞人从群众中对他说：\Quote{老师，责斥你的门徒吧！}37 \BibleRef{路 19:40} 耶稣回答说：\Quote{我告诉你们：这些人若不作声，石头必要喊叫起来。}\par

38、39 \BibleRef{路 19:41} 耶稣临近的时候，看见城，就为她哭泣，说：\Quote{\BibleRef{路 19:42} 恨不得在这一天，你也知道关乎你平安的事！但如今这事40已由你眼前隐藏了。\BibleRef{路 19:43} 日子将到，你的仇敌要筑起壁垒，从四面围困你，41 \BibleRef{路 19:44} 并要消灭你和你城内的子女；在你内不留一块石头在另一块上，因为你没有认识你被眷顾的时候。}\par

42 \BibleRef{玛 21:10} 耶稣进了耶路撒冷，全城都惊动起来，他们说：\Quote{43 这人是谁？}\BibleRef{玛 21:11} 群众说：\Quote{这是从纳匝肋44的加里肋亚来的先知耶稣。}\BibleRef{若 12:17} 那些同他在一起的群众见证他叫45拉匝禄出坟墓，并使他从死者中复活。\BibleRef{若 12:18} 因此，有许多人出去迎接他，因为他们听说他行了这个奇迹。\par

\SourceRecord{Translated by Hope W. Hogg.；Ante-Nicene Fathers；Vol. 9.；Edited by Allan Menzies.；Buffalo, NY: Christian Literature Publishing Co.,；1896.；<http://www.newadvent.org/fathers/100239.htm>.}

% structure: p0001 source=h1 target=section reason=page-title

\section{\WorkTitle{迪亚特萨隆}（第四十节）}

1 [阿拉伯文，第150页] \BibleRef{玛 21:14} \Person{耶稣}进了圣殿，有人把瞎子和2瘸子带到祂跟前，祂就治好了他们。\BibleRef{玛 21:15} 但司祭长和法利塞人看见祂所做的奇事，又看见儿童在圣殿里喊叫3说：\Quote{愿\Person{达味}之子受赞美！}就很忿怒，\BibleRef{玛 21:16} 便对祂说：\Quote{你听见他们说什么吗？}\Person{耶稣}对他们说：\Quote{是的。你们从未读过：\InnerQuote{从婴孩和4吃奶者的口中，你预备了赞美}吗？}\BibleRef{若 12:19} 法利塞人彼此说：\Quote{瞧，你们没有看到吗？我们一无所成。看，全世界都跟从祂去了。}\par

5 \BibleRef{若 12:20} 在那些上来过节礼拜的人中，有些外邦人。6 \BibleRef{若 12:21} 他们来到\Place{加里肋亚}的\Place{贝特赛达}人\Person{斐理伯}那里，请求他说：\Quote{先生，我们愿意见\Person{耶稣}。}7 \BibleRef{若 12:22} \Person{斐理伯}去告诉\Person{安德肋}，\Person{安德肋}和\Person{斐理伯}便去告诉\Person{耶稣}。8 \BibleRef{若 12:23} \Person{耶稣}回答他们说：\Quote{人子受光荣的时刻临近了。9 \BibleRef{若 12:24} 我实实在在告诉你们：一粒麦子若不落在地里死去，仍旧是一粒；如果死了，就结出许多子实来。10 \BibleRef{若 12:25} 爱惜自己性命的，必要丧失性命；在现世憎恨自己性命的，必要保存性命进入永生。11 \BibleRef{若 12:26} 谁若事奉我，就该跟随我；我在哪里，我的仆人也必在那里；谁若事奉我，我父必尊重他。12 \BibleRef{若 12:27} 现在我心神烦乱，我可说什么呢？我父，救我脱离这个时辰罢？但正是为此，我才到了这个时辰。13 \BibleRef{若 12:28} 我父，光荣你的名罢！}于是有声音从天上传来：\Quote{我已光荣了\Emph{它}，还要再光荣\Emph{它}。}[阿拉伯文，第151页] 14 \BibleRef{若 12:29} 站在那里的群众听见了，就说：\Quote{这是打雷。}另有人说：\Quote{有位天使在对他说话。}15 \BibleRef{若 12:30} \Person{耶稣}回答他们说：\Quote{这声音不是为我，而是为你们。16 \BibleRef{若 12:31} 现在就是这世界的审判；这世界的元首现在要被赶出去。17 \BibleRef{若 12:32} 至于我，当我从地上被举起来时，要吸引众人归向我。}18 \BibleRef{若 12:33} 他说这话，是指明他将要怎样死去。19 \BibleRef{若 12:34} 群众回答他说：\Quote{我们从法律上听说：默西亚永远长存。你怎么说人子必须被举起来呢？这人子是谁？}20 \BibleRef{若 12:35} \Person{耶稣}对他们说：\Quote{光在你们中间还有不多的时候。你们趁着有光的时候行走，免得黑暗笼罩你们；因为在黑暗中行走的，不知道往哪里去。21 \BibleRef{若 12:36} 你们趁着有光的时候，应当信从光，好成为光明之子。}\par

22 \BibleRef{路 17:20} 法利塞人问\Person{耶稣}天主的国何时来临，祂回答他们说：\Quote{天主的国来临，并非可以观察的；23 \BibleRef{路 17:21} 人也不能说：看，在这里！或说：在那里！因为天主的国就在你们中间。}\par

24 \BibleRef{路 21:37} 祂白天在圣殿里施教，夜间出去，到那名叫\Place{橄榄山}的山上过夜。25 \BibleRef{路 21:38} 众百姓清早起来到圣殿里，到祂跟前听祂讲道。\par

26, 27 \BibleRef{玛 23:1} 那时，\Person{耶稣}对群众和祂的门徒们讲论说：\BibleRef{玛 23:2} \Quote{28 [阿拉伯文，第152页] 经师和法利塞人坐在\Person{梅瑟}的座位上。\BibleRef{玛 23:3} 凡他们所吩咐你们的，你们都要遵守并实行；但不要效法他们的行为，因为他们只说而不做。29 \BibleRef{玛 23:4} 他们把沉重难担的担子捆起来，放在人的肩上，自己却连一个指头也不肯动。30, 31 \BibleRef{玛 23:5} 他们所做的一切，都是为了引人注目。}\BibleRef{谷 12:37} 众百姓都喜欢听他讲论。\par

32 \BibleRef{谷 12:38} \Person{耶稣}在教训中向他们说：\Quote{你们要谨防经师，33他们喜欢穿长袍游行，\BibleRef{谷 12:39} 喜爱在街市上受人请安，在会堂里坐上座，在筵席上坐首席；34 \BibleRef{玛 23:5} 他们加宽经文匣，加长衣繸；35 \BibleRef{玛 23:7} \Emph{喜爱}人们称呼他们为老师。\BibleRef{谷 12:40} 他们吞没寡妇的家产，而以长久的祈祷作掩饰。这些人必要受更重的处罚。36 \BibleRef{玛 23:8} 至于你们，不要被称为老师；因为你们的老师只有一位，你们都是兄弟。37 \BibleRef{玛 23:9} 也不要称世上的人为父；因为你们的父只有一位，就是天上的父。38 \BibleRef{玛 23:10} 也不要被称为导师；因为你们的导师只有一位，\Emph{就是}默西亚。39, 40 \BibleRef{玛 23:11} 你们中最大的，该作你们的仆役。\BibleRef{玛 23:12} 凡高举自己的，必被贬抑；凡贬抑自己的，必被高举。}\par

41 \BibleRef{路 11:43} \Quote{祸哉，你们法利塞人！因为你们喜爱会堂里的上座，街市上的请安。}\par

42 \BibleRef{玛 23:14} \Quote{祸哉，你们经师和法利塞假善人！因为你们吞没寡妇的家产，而以长久的祈祷作掩饰。因此，你们必要受更重的处罚。}\par

43 \BibleRef{玛 23:13} \Quote{祸哉，你们经师和法利塞假善人！因为你们在人面前关闭了天国。}\par

44 [阿拉伯文，第153页] \BibleRef{路 11:52} \Quote{祸哉，你们精通法律的人！因为你们拿走了知识的钥匙。\BibleRef{玛 23:13} 你们自己不进去，那正要进去的，你们也加以阻止。}\par

45 \BibleRef{玛 23:15} \Quote{祸哉，你们经师和法利塞假善人！因为你们走遍海洋和陆地，为使一个人成为归依者；当他成了\Emph{这样}的人，你们却使他成为地狱之子，比你们加倍还坏。}\par

\BibleQuote{46 \BibleRef{玛 23:16}祸哉，你们这些瞎眼的向导！因为你们说：\InnerQuote{谁若指着圣所发誓，不算什么；但是谁若指着圣所中的金子发誓，47 就该还债。}\BibleRef{玛 23:17}你们这些瞎眼的愚人哪！哪个更大呢？是金子，还是那48使金子成圣的圣所？\BibleRef{玛 23:18}并且，\InnerQuote{谁若指着祭台发誓，不算什么；但是谁若指着上面的供物发誓，就该还债。}49 \BibleRef{玛 23:19}你们这些瞎眼的愚人哪！哪个更大呢？是供物，还是那50使供物成圣的祭台？\BibleRef{玛 23:20}所以，谁若指着祭台发誓，就是指着祭台和上面51的一切发誓；\BibleRef{玛 23:21}谁若指着圣所发誓，就是指着圣所和那52住在其中的天主发誓；\BibleRef{玛 23:22}谁若指着天发誓，就是指着天主的宝座和那坐在上面的天主发誓。}\par

\BibleQuote{53 \BibleRef{玛 23:23}祸哉，你们这些经师和法利塞人，假善人！因为你们将薄荷、茴香、莳萝、孜然和一切蔬菜献上十分之一，却忽略了法律上重要的事：正义、怜悯、信德和爱天主。这些是该做的，54而不可\Emph{忽略}那些。\BibleRef{玛 23:24}你们这些瞎眼的向导，滤出蚊子，却吞下骆驼。}\par

\BibleQuote{55 \BibleRef{玛 23:25}祸哉，你们这些经师和法利塞人，假善人！因为你们洗净杯盘的外面，里面却充满不义和邪恶。56 \BibleRef{玛 23:26}你们这些瞎眼的法利塞人，先洗净杯盘的里面，好使外面也成为洁净的。}\par

57 [阿拉伯文版，第154页] \BibleQuote{\BibleRef{玛 23:27}祸哉，你们这些经师和法利塞人，假善人！因为你们好像粉刷的坟墓，外面好看，里面却装满了死人的骨头和一切污秽。58 \BibleRef{玛 23:28}同样，你们外面在人前显出义人的样子，里面却充满邪恶和假善。}\par

\BibleQuote{59 \BibleRef{路 11:45}有一个经师回答他说：\InnerQuote{老师，你这样说，60 是在侮辱我们。} \BibleRef{路 11:46}他说：你们这些经师也有祸！因为你们把人难以负荷的重担放在他们身上，你们自己却连一个指头也不肯动一下。}\par

\BibleQuote{61 \BibleRef{玛 23:29}祸哉，你们这些经师和法利塞人，假善人！因为你们修建先知的坟墓，装饰义人的墓穴，62 并且说：\BibleRef{玛 23:30}\InnerQuote{如果我们生在我们祖先的时代，我们就不至于与他们同流合污，63 流先知的血了。}\BibleRef{玛 23:31}因此，你们64 自己证明自己是杀害先知者的子孙。\BibleRef{玛 23:32}你们65 也充满你们祖先的恶行。\BibleRef{玛 23:33}你们这些蛇，毒蛇的子孙，你们怎能逃脱离地狱的审判呢？}\par

\SourceRecord{Translated by Hope W. Hogg.；Ante-Nicene Fathers；Vol. 9.；Edited by Allan Menzies.；Buffalo, NY: Christian Literature Publishing Co.,；1896.；<http://www.newadvent.org/fathers/100240.htm>.}

% structure: p0001 source=h1 target=section reason=page-title

\section{\WorkTitle{四福音合参（第四十一节）}}

1 \BibleRef{玛 23:34} 因此，看哪，我、天主的智慧，正派遣先知、宗徒、智者和经师到你们这里来：你们要杀害并钉死其中一些，在会堂里鞭打其中一些，从这城逼迫到那城；\BibleRef{玛 23:35} 为使你们追偿地上所流一切义人的血，从义人亚伯尔的血，直到你们在圣所与祭台之间所杀的巴辣基亚的儿子则加黎雅的血。\BibleRef{玛 23:36} 我实在告诉你们：所有这些\Emph{事}都要归到这一代人身上。\par

\EditorialNote{阿拉伯文版第155页}4 \BibleRef{玛 23:37} 耶路撒冷，耶路撒冷，你杀害先知，用石头砸死那些派到你这里来的人！我多少次愿意聚集你的子女，如同母鸡将小鸡聚集在翅膀下，而你们却不愿意！\BibleRef{玛 23:38} 你们的房屋将荒凉地留给你们。\BibleRef{玛 23:39} 我实在告诉你们：从今以后，你们不得再见我，直到你们说：奉主名而来的，当受赞美。\par

7 \BibleRef{若 12:42} 连许多官长也信了他，但因了法利塞人的缘故，不敢公开承认\Emph{他}，免得被逐出会堂；\BibleRef{若 12:43} 因为他们喜爱人的称赞，胜过天主的称赞。\BibleRef{若 12:44} 耶稣呼喊说：信我的，不是信我，而是信那派遣我的；\BibleRef{若 12:45} 看见我的，就是看见那派遣我的。\BibleRef{若 12:46} 我来到世上为光，使凡信我的，不住在黑暗中。\BibleRef{若 12:47} 谁若听我的话而不遵守，我不审判他，因为我来不是为审判世界，而是为拯救世界。\BibleRef{若 12:48} 谁若轻视我，不接受我的话，自有审判他的：我所讲的话，在末日要审判他。\BibleRef{若 12:49} 我没有凭自己说话，而是派遣我的圣父给我规定了该说什么、该讲什么。\BibleRef{若 12:50} 我知道他的命令就是永生。所以，我所讲的，正如圣父对我所说的，我就\Emph{这样}讲。\par

16 \BibleRef{路 11:53} 耶稣从那里出来，经师和法利塞人开始恶意相向，对\Emph{他}发怒，并盘问他许多的事，\BibleRef{路 11:54} 窥伺从他的口中抓到什么，好能诬告他。\par

18 \EditorialNote{阿拉伯文版第156页}\BibleRef{路 12:1} 那时，成千累万的群众聚集在一起，几乎彼此践踏，耶稣开始对门徒说：你们要谨慎，防备法利塞人的酵母，即虚伪。\BibleRef{路 12:2} 因为没有隐藏的事，将来不被揭露的；也没有秘密的事，将来不被知道的。\BibleRef{路 12:3} 你们在暗处所说的，将来要在明处被听见；你们在内室贴近耳朵所说的，将来要在屋顶上宣扬出来。\par

21 \BibleRef{若 12:36} 耶稣说了这话，就离开他们隐藏了。22 \BibleRef{若 12:37} 他虽然在他们在前行了这些奇迹，他们仍不信他；23 \BibleRef{若 12:38} 这应验了依撒意亚先知的话，他说：\par

\Quote{吾主，有谁相信了我们的报道？\newline{}主的膀臂向谁显示了呢？}\par

24 \BibleRef{若 12:39} 因此，他们不能相信，因为依撒意亚又说过：\par

25 \BibleRef{若 12:40} \Quote{他们瞎了眼，硬了心，\newline{}免得他们眼睛看见，心里明白，\newline{}而回头，\newline{}使我治好他们。}\par

26 \BibleRef{若 12:41} 依撒意亚因为看见了他的光荣，就说了这话，是指着他说的。\par

27 \BibleRef{玛 24:1} 耶稣出了圣殿，28 有些门徒前来，把圣殿的建筑指给他看，29 论它的美丽和宏大，以及所砌石头的坚固和建筑的雅致，且装饰着贵重的石头和美丽的彩色。30 \BibleRef{玛 24:2} 耶稣回答他们说：你们看见这些伟大的建筑吗？我实在告诉你们：日子将到，这里没有一块石头留在另一块石头上，而不被拆毁。\par

31 \BibleRef{谷 14:1} 过两天就是逾越节和无酵节，司祭长和经师设法怎样用诡计捉拿耶稣，而把他杀害；32 \BibleRef{谷 14:2} 但他们说：不要在节庆期间，免得民众骚动。\par

\BibleRef{谷 13:3} 当耶稣坐在橄榄山上，对着圣殿的时候，他的门徒\Person{西满刻法}、\Person{雅各伯}、\Person{若望}和\Person{安德肋}，来到他跟前，私下对他说：\Quote{师傅，请告诉我们，什么时候要有这些事？你降临和世界末了有什么预兆？}耶稣回答他们说：\Quote{日子将到，你们渴望看见人子的一个日子，却不得看见。\BibleRef{玛 24:5} 你们要谨慎，免得有人迷惑你们。\BibleRef{路 21:8} 因为将来有好些人冒我的名来，说：\InnerQuote{我是默西亚}；又说：\InnerQuote{时候近了}，并且要迷惑许多人；你们不要跟从他们。当你们听见打仗和骚乱的消息时，要谨慎，\Emph{不要惊慌}：因为这些事\Emph{必须先发生}，但终结还没有立刻来到。\BibleRef{玛 24:7} 民要攻打民，国要攻打国；\BibleRef{路 21:11} 多处必有大地震、饥荒、瘟疫，又有可怕的异象和大神迹从天上显现，并有大的风暴。\BibleRef{玛 24:8} \Emph{这一切}是灾难的起头。\BibleRef{路 21:12} 但这一切事以前，他们要下手拿住你们，逼迫你们，把你们交给会堂，并且收在监里，又为我的名拉你们到君王和审判官面前。\BibleRef{路 21:13} 但这些事终必成为你们作见证的机会。\BibleRef{谷 13:10} 然而，福音必须先传给万民。\BibleRef{路 12:11} 人带你们到会堂、首领和有权柄的人面前，不要思虑怎样分诉，说什么话；\BibleRef{谷 13:11} 因为说话的不是你们，乃是圣神。\BibleRef{路 21:14} 所以你们应当立定心意，不要预先思想怎样分诉；\BibleRef{路 21:15} 因为我必赐你们口才智慧，是你们一切敌人所敌不住、驳不倒的。\BibleRef{玛 24:9} 那时，人要把你们陷在患难里，也要杀害你们；你们又要为我的名被万民恨恶。\BibleRef{玛 24:30} 那时，必有许多人被迷惑，也要彼此陷害、彼此恨恶。\BibleRef{路 21:16} 连你们的父母、弟兄、亲族、朋友也要把你们交官，你们也有被他们害死的。\BibleRef{路 21:18} 然而，你们连一根头发也必不损坏。\BibleRef{路 21:19} 你们常存忍耐，就必保全灵魂。\BibleRef{玛 24:11} 且有好些假先知起来，迷惑多人。\BibleRef{玛 24:12} 只因不法的事增多，许多人的爱德才渐渐冷淡了。\BibleRef{玛 24:13} 惟有忍耐到底的，必然得救。\BibleRef{玛 24:14} 这天国的福音要传遍天下，对万民作见证，然后末期才来到。}\par

\SourceRecord{Translated by Hope W. Hogg.；Ante-Nicene Fathers；Vol. 9.；Edited by Allan Menzies.；Buffalo, NY: Christian Literature Publishing Co.,；1896.；<http://www.newadvent.org/fathers/100241.htm>.}

% structure: p0001 source=h1 target=section reason=page-title

\section{迪亚泰萨隆（第四十二节）}

1 \BibleRef{路 21:20} \BibleQuote{但当你们看见耶路撒冷被军队围困时，就知道它的荒凉近了。} 2 \BibleRef{路 21:21} \BibleQuote{那时在\OriginalTerm{犹太地}{Judæa}的，要逃到山上；城内的要逃出去；乡间的不要进城。} 3 \BibleRef{路 21:22} \BibleQuote{因为这是报复的日子，为要应验所记载的一切。} 4 \BibleRef{玛 24:15} \BibleQuote{当你们看见\OriginalTerm{达尼尔}{Daniel}先知所说的可憎的荒凉之物立在圣地时——读者应当明白——} 5,6 \BibleRef{玛 24:16} \BibleQuote{那时在犹太的，要逃到山上；} \BibleRef{谷 13:15} \BibleQuote{在屋顶上的，不要下来，也不要进去拿家里的东西；} 7 \BibleRef{谷 13:16} \BibleQuote{在田里的，不要回头取衣服。}\EditorialNote{阿拉伯语版第159页} 8 \BibleRef{路 21:23} \BibleQuote{那些日子怀孕的和哺乳的有祸了！因为将有大灾难降在这地上，也有震怒临到这百姓。} 9 \BibleRef{路 21:24} \BibleQuote{他们要倒在刀下，被掳到各国；耶路撒冷要被外邦人践踏，直到外邦人的时期满了。}\par

10 \BibleRef{谷 13:21} \BibleQuote{那时，若有人对你们说：\OriginalTerm{默西亚}{Messiah}在这里；或说：他在那里；你们不要相信。} 11 \BibleRef{玛 24:24} \BibleQuote{因为将有\OriginalTerm{假默西亚}{false Messiahs}和假先知起来，行大奇事和神迹，为要迷惑\OriginalTerm{被选者}{elect}，如果可能的话。} 12 \BibleRef{谷 13:23} \BibleQuote{但你们要当心；因为我已预先将一切告诉了你们。} 13 \BibleRef{玛 24:26} \BibleQuote{所以，若他们对你们说：看，他在旷野里；不要出去；免得被捉拿。若他们说：看，他在内室里；不要相信。} 14 \BibleRef{玛 24:27} \BibleQuote{因为闪电从东方发出，直照到西方，\OriginalTerm{人子}{Son of man}的来临也是这样。} 15 \BibleRef{路 17:25} \BibleQuote{但他必须先受许多苦，并被这一代所弃绝。} 16 \BibleRef{玛 24:20} \BibleQuote{你们要祈祷，免得你们的逃亡在冬天或在安息日。} 17 \BibleRef{玛 24:21} \BibleQuote{因为那时将有大灾难，自从创世以来直到如今，从未有过，将来也必没有。} 18 \BibleRef{谷 13:20} \BibleQuote{若不是\OriginalTerm{主}{Lord}缩短了那些日子，凡有血肉的都没有得救的；但为了他所拣选的被选者，他缩短了那些日子。} 19 \BibleRef{路 21:25} \BibleQuote{在日月星辰上将有异兆；地上万国要困苦，因海涛的响声而惊惶失措。} 20 \BibleRef{路 21:26} \BibleQuote{人因恐惧和等待将要临到世上的事，都吓得魂不附体。}\EditorialNote{阿拉伯语版第160页} 在那些日子，那灾难之后，太阳要昏暗，月亮也不发光，星辰要从天上坠落，天上万象也要震动。 22 \BibleRef{玛 24:30} \BibleQuote{那时，人子的记号要显在天上；地上万民都要哀号，并要看见人子带着威能和大光荣，乘着天上的云彩降来。} 23 \BibleRef{玛 24:31} \BibleQuote{他要派遣天使，以洪亮的号角，从四方、从天这边到天那边，聚集他所拣选的被选者。} 24 \BibleRef{路 21:28} \BibleQuote{这些事开始发生时，你们应当挺起身来，抬起头来，因为你们的救恩近了。}\par

25 \BibleRef{玛 24:32} \BibleQuote{你们应从无花果树学个比喻：当它的枝条发嫩，长出叶子时，你们就知道夏天近了。} 26 \BibleRef{玛 24:33} \BibleQuote{同样，当你们看见这一切事开始发生时，就知道\OriginalTerm{天主的国}{kingdom of God}近了，已在门口。} 27 \BibleRef{玛 24:34} \BibleQuote{我实在告诉你们：这一代决不会过去，直到所有这些\Emph{事}都应验。} 28 \BibleRef{玛 24:35} \BibleQuote{天地要过去，但我的话决不会过去。}\par

29 \BibleRef{路 21:34} \BibleQuote{你们要谨慎，免得你们的心因纵欲、醉酒和生活的挂虑而沉甸甸，那日子就突然临到你们身上。} 30 \BibleRef{路 21:35} \BibleQuote{因为那日子如同罗网，要临到全地面上的一切居民。} 31 \BibleRef{路 21:36} \BibleQuote{你们要时时醒寤祈祷，为使你们能逃脱这一切要发生的事，并能站立在\OriginalTerm{人子}{Son of man}面前。}\EditorialNote{阿拉伯语版第161页} 32 \BibleRef{谷 13:32} \BibleQuote{至于那日子和那时刻，没有人知道，连天上的天使也不知道，\OriginalTerm{圣子}{the Son}也不知道，惟有\OriginalTerm{圣父}{the Father}知道。} 33 \BibleRef{谷 13:33} \BibleQuote{你们要当心，要醒寤祈祷，因为你们不知道那时间\Emph{将是什么时候}。} 34 \BibleRef{谷 13:34} \BibleQuote{\Emph{它犹如}一个人远行，离开自己的家，把权柄交给仆人，分配各人的工作，又嘱咐看门的要醒寤。} 35 \BibleRef{谷 13:35} \BibleQuote{所以你们要醒寤，因为你们不知道\OriginalTerm{家主}{lord of the house}何时回来：或晚上，或半夜，或鸡叫时，或清晨。} 36 \BibleRef{谷 13:36} \BibleQuote{免得他忽然来到，看见你们睡觉。} 37 \BibleRef{谷 13:37} \BibleQuote{我对你们说的话，也是对众人说的：你们要醒寤。}\par

38 \BibleRef{玛 24:37} \BibleQuote{因为就如诺厄的日子一样，人子的来临也将如此。} 39 \BibleRef{玛 24:38} \BibleQuote{他们在洪水之前，照常吃喝、嫁娶，和将妻子给男人，} 40 \BibleRef{玛 24:39} \BibleQuote{直到诺厄进入方舟的那天，他们没有察觉，直到洪水来了，把他们都冲走了；人子的来临也将这样。} 41 \BibleRef{路 17:28} \BibleQuote{又如罗特的日子一样；他们吃喝、买卖、栽种、建造，} 42 \BibleRef{路 17:29} \BibleQuote{到罗特从索多玛出去的那天，主从天上降下火与硫磺，摧毁了他们所有人} 43,44 \BibleRef{路 17:30} \BibleQuote{人子显现的日子也要这样。} \BibleRef{路 17:31} \BibleQuote{在那天，在屋顶上的，他的衣物在屋里，不要下来取它们；在田里的，也不要回头。} 45 \BibleRef{路 17:32} \BibleQuote{要记得罗特的妻子。} 46 \BibleRef{路 17:33} \BibleQuote{凡想保全性命的，必要丧失性命；凡丧失性命的，必要保全性命。} 47 \BibleRef{路 17:34} \BibleQuote{我实在告诉你们，那一夜，两个人同在一张床上，一个被提去，另一个被留下。} 48 \EditorialNote{阿拉伯文，第162页} \BibleRef{路 17:35} \BibleQuote{两个\Emph{女人}一同推磨，一个被提去，另一个被留下。} 49 \BibleRef{路 17:36} \BibleQuote{两个人在田里，一个被提去，另一个被留下。} 50 \BibleRef{路 17:37} \BibleQuote{他们回答说：主啊，往哪里去？他对他们说：哪里有尸体，鹰鹫就聚集在那里。} 51,52 \BibleRef{玛 24:42} \BibleQuote{所以现在要警醒，因为你们不知道你们的主在哪时刻来。} \BibleRef{玛 24:43} \BibleQuote{你们要知道：如果家主知道贼在哪个更时来，他必警醒，不容许他的房屋被挖穿。} 53 \BibleRef{玛 24:44} \BibleQuote{因此，你们也要准备，因为你们想不到的时候，人子就来了。}\par

\SourceRecord{Translated by Hope W. Hogg.；Ante-Nicene Fathers；Vol. 9.；Edited by Allan Menzies.；Buffalo, NY: Christian Literature Publishing Co.,；1896.；<http://www.newadvent.org/fathers/100242.htm>.}

% structure: p0001 source=h1 target=section reason=page-title

\section{\WorkTitle{四福音合编}（第43节）}

1 \BibleRef{路 12:41} \Person{西满·刻法}对祂说：\Quote{主，你讲这个比喻，\Emph{是对我们说的呢}，还是也对众人说的呢？}耶稣对他说：\Quote{你以为谁是那个忠信的管家，主人派他管理自己的家仆，按时分粮给他们呢？\BibleRef{玛 24:46} 那仆人在主人回来时看到他这样做，就有福了。我实在告诉你们，主人要派他管理一切所有的。但如果那个恶仆心里说：\Quote{我的主人必会迟来}，\BibleRef{玛 24:49} 就动手打他的同伴和婢女，\BibleRef{玛 24:50} 那仆人的主人要在想不到的日子，不知道的时刻回来，\BibleRef{玛 24:51} 就\EditorialNote{阿拉伯文版第163页}重重处罚他，使他与假善人同分，\BibleRef{路 12:46} 与不忠信的人在一起：\BibleRef{玛 24:51} 那里要有哀号和切齿。}\par

9 \BibleRef{玛 25:1} 那时，天国好比十个童女，她们拿着灯，出去迎接新郎和新娘。\BibleRef{玛 25:2} 其中五个是明智的，五个是糊涂的。\BibleRef{玛 25:3} 那些糊涂的拿着灯，却没有随身带油；\BibleRef{玛 25:4} 但明智的拿着灯，又在器皿里带了油。\BibleRef{玛 25:5} 新郎迟延的时候，她们都打盹睡着了。\BibleRef{玛 25:6} 半夜有人喊说：看，新郎来了，你们出来迎接他！\BibleRef{玛 25:7} 那些童女都起来，整理她们的灯。\BibleRef{玛 25:8} 糊涂的对明智的说：请分点油给我们，因为我们的灯灭了。\BibleRef{玛 25:9} 明智的回答说：怕不够我们和你们用的，不如你们到卖油的那里为自己买吧！\BibleRef{玛 25:10} 她们去买的时候，新郎到了；准备好的同他进去赴婚宴，门就关上了。\BibleRef{玛 25:11} 后来其余的童女也来了，说：主啊，主啊，给我们开门！\BibleRef{玛 25:12} 他回答说：我实在告诉你们，我不认识你们。\BibleRef{玛 25:13} 所以你们要警醒，因为你们不知道那日子，也不知道那时辰。\par

22 \BibleRef{玛 25:14} 又如一个人要远行，叫了自己的仆人，把财产托付给他们。\BibleRef{玛 25:15} 一个给了五个塔冷通，一个给了两个，一个给了一个；各按各的能力，然后他就起程了。\BibleRef{玛 25:16} \EditorialNote{阿拉伯文版第164页} 那领了五个的，立刻去做买卖，赚了另外五个。\BibleRef{玛 25:17} 同样，那领了两个的也赚了另外两个。\BibleRef{玛 25:18} 但那领了一个的，去掘开地，把主人的银子藏了。\BibleRef{玛 25:19} 过了很久，那些仆人的主人回来了，要跟他们算账。\BibleRef{玛 25:20} 那领了五个的上前来，又带了另外五个，说：主啊，你交给我五个塔冷通，请看，我赚了另外五个。\BibleRef{玛 25:21} 主人对他说：好！善良忠信的仆人，你在少许事上忠信，我要派你管理许多大事；进来享受你主人的欢乐吧！\BibleRef{玛 25:22} 那领了两个的也上前来说：主啊，你交给我两个塔冷通，请看，我赚了另外两个。\BibleRef{玛 25:23} 主人对他说：好！忠信的仆人，你在少许事上忠信，我要派你管理许多大事；进来享受你主人的欢乐吧！\BibleRef{玛 25:24} 那领了一个的也上前来说：主啊，我知道你是刻薄的人，在你没有播种的地方收割，没有散布的地方聚敛。\BibleRef{玛 25:25} 因此我害怕，就去把你的塔冷通藏在地里；请看，你的仍还给你。\BibleRef{玛 25:26} 主人回答他说：可恶懒惰的仆人！你既然知道我在没有播种的地方收割，没有散布的地方聚敛，\BibleRef{玛 25:27} 你就该把我的银子交给钱庄，及我回来时，可以连本带利收回。\BibleRef{玛 25:28} 把他这个塔冷通拿来，给那有十个的。\BibleRef{玛 25:29} \EditorialNote{阿拉伯文版第165页} 因为凡有的，还要给他，使他富足；凡没有的，连他所有的也要从他手中拿去。\BibleRef{玛 25:30} 把这无用的仆人丢在外面的黑暗里；那里要有哀号和切齿。\par

39, 40 \BibleRef{路 12:35} 你们腰里要束上带，灯也要点着；\BibleRef{路 12:36} 你们要像等候主人的人，当主人从婚宴回来时，他一来敲门，立刻给他开门。\BibleRef{路 12:37} 那些仆人，主人来时发现他们警醒，就有福了。我实在告诉你们，主人要束上带，请他们坐席，自己前来伺候他们。\BibleRef{路 12:38} 他二更或三更来，发现他们这样，那些仆人就有福了。\par

43 \BibleRef{玛 25:31}但当人子在自己的光荣中来临，所有他的圣天使与他一同来临的时候，44 那时他将坐在他光荣的宝座上；\BibleRef{玛 25:32}他将把万民聚集在他面前，分开他们，如同牧人分开45绵羊和山羊；\BibleRef{玛 25:33}把绵羊放在右边，山羊在46左边。\BibleRef{玛 25:34}那时君王要对那些在他右边的人说：你们这些蒙我父祝福的，来承受从创世以来为你们预备的国度吧！47 \BibleRef{玛 25:35}因为我饿了，你们给了我吃的；我渴了，你们给了我喝的；我48做客，你们收留了我；\BibleRef{玛 25:36}我赤身露体，你们给了我穿的；我49生病，你们看顾了我；我在监里，你们来探望了我。\BibleRef{玛 25:37}那时那些义人要回答说：主啊，我们什么时候见你饿了而给了你吃的？50或者渴了而给了你喝的？\BibleRef{玛 25:38}什么时候见你做客而收留51了你？或者赤身露体而给了你穿的？\BibleRef{玛 25:39}什么时候见你生病或坐监而来看顾52了你？\BibleRef{玛 25:40}君王要回答他们说：我实在告诉你们，\EditorialNote{阿拉伯文，第166页}你们对我这些最小兄弟中的一个所做的，就是对我做的。\BibleRef{玛 25:41}那时他也要对左边的人说：离开我，你们这些被咒骂的，54到那为魔鬼和他的使者所预备的永火里去吧！\BibleRef{玛 25:42}因为我饿了，你们没有55给我吃的；我渴了，你们没有给我喝的；\BibleRef{玛 25:43}我做客，你们没有收留我；我赤身露体，你们没有给我穿的；我生病和坐监，56你们没有看顾我。\BibleRef{玛 25:44}那时他们也要回答说：主啊，我们什么时候见你饿了，或渴了，或赤身露体，或做客，或生病，或坐监，57而没有服事你呢？\BibleRef{玛 25:45}那时他要回答他们说：我实在告诉你们，你们没有给这些\Emph{最小兄弟}中的一个做\Emph{这些事}，你们也就是没有58给我做\Emph{这些事}。\BibleRef{玛 25:46}这些人要进入永罚，而那些义人却要进入永生。\par

\SourceRecord{Translated by Hope W. Hogg.；Ante-Nicene Fathers；Vol. 9.；Edited by Allan Menzies.；Buffalo, NY: Christian Literature Publishing Co.,；1896.；<http://www.newadvent.org/fathers/100243.htm>.}

% structure: p0001 source=h1 target=section reason=page-title

\section{四福音合参（第四十四节）}

1, 2 \BibleRef{玛 26:1} 耶稣讲完了这一切话，就对他的门徒说：\BibleRef{玛 26:2} 你们知道，两天后便是逾越节，人子要被交出去，被钉在十字架上。3 \BibleRef{玛 26:3} 那时，大司祭、经师和民间长老聚集在大司祭的庭院里，这人名叫盖法；\BibleRef{玛 26:4} 他们一起商议怎样用诡计捉拿耶稣，杀害他。\BibleRef{玛 26:5} 但他们说：不可在庆节期间，免得在民间引起骚动；\BibleRef{路 22:2} 因为他们惧怕百姓。\par

6 \BibleRef{路 22:3} 撒殚进入了那称为依斯加略的犹大，他是十二位之一。7 他就去，与大司祭们、经师们和圣殿的掌权者商议，对他们说：8 你们愿意给我什么（阿拉伯文第167页），我把他交给你们？他们听了\Emph{这话}，很是欢喜，就预备了三十\Emph{块}银钱给他。9 \BibleRef{路 22:6} 他便答应了，从此寻找机会，要把耶稣在没有群众的时候交给他们。\par

10 \BibleRef{谷 14:12} 无酵节的第一天，门徒们来到耶稣跟前，对他说：你愿意我们到哪里去为你预备，好让你吃逾越节晚餐？\par

11 \BibleRef{若 13:1} 逾越节以前，耶稣知道他离开这个世界归回圣父的时辰到了；他既然爱了世上属自己的人，就爱他们到底。12 \BibleRef{若 13:2} 在节期时，撒殚就入了那西满的儿子依斯加略的心，要出卖他。13 \BibleRef{若 13:3} 耶稣因为知道圣父已将一切交在他手里，且知道自己是从圣父出来的，又往天主那里去，14 \BibleRef{若 13:4} 就从席间起来，脱下外衣，15 \BibleRef{若 13:5} 拿了一条毛巾束在腰间，把水倒在盆里，开始洗门徒的脚，并用束在腰间的毛巾擦干。16 \BibleRef{若 13:6} 及至来到西满·刻法跟前，西满对他说：主，你为我洗脚吗？17 \BibleRef{若 13:7} 耶稣回答他说：我所做的，你现在不明白，但后来你会明白。18 \BibleRef{若 13:8} 西满对他说：你绝不可洗我的脚！耶稣对他说：我若不洗你，你就与我无分。19 \BibleRef{若 13:9} 西满·刻法对他说：主，不但我的脚，连手和头也洗吧！20 \BibleRef{若 13:10} 耶稣对他说：洗过澡的人，除了脚以外，全身都是干净的；你们也是干净的，但不都是。21 \BibleRef{若 13:11} 原来耶稣知道谁要出卖他，因此说：你们不都是干净的。\par

22 （阿拉伯文第168页）\BibleRef{若 13:12} 耶稣洗完了他们的脚，就穿上外衣，又坐下，对他们说：你们明白我给你们做的吗？23 \BibleRef{若 13:13} 你们称我为师傅和主，说得对，我确实是。24 \BibleRef{若 13:14} 那么，我身为师傅和主，尚且给你们洗脚，你们也该彼此洗脚。25 \BibleRef{若 13:15} 我给你们立了榜样，叫你们照着我给你们做的去做。26 \BibleRef{若 13:16} 实实在在告诉你们：仆役不能大过主人，宗徒也不能大过派遣他的人。27 \BibleRef{若 13:17} 你们既然知道了这些事，若是实行，便是有福的。28 \BibleRef{若 13:18} 我说这话不是指你们全体；我知道我所拣选的是谁；但为应验圣经上的话：同我一起吃饭的人，向我举起了脚跟。29 \BibleRef{若 13:19} 现在事情发生以前，我告诉你们，好使事情发生时，你们相信我就是那一位。30 \BibleRef{若 13:20} 实实在在告诉你们：谁若接待我所派遣的，就是接待我；接待我的，就是接待那派遣我来的。\par

31 \BibleRef{路 22:27} 是谁为大？是坐席的还是服侍的？岂不是坐席的吗？32 \BibleRef{路 22:28} 我却在你们中间，如同服侍的。\BibleRef{路 22:29} 但是你们是那些同我一起经历试炼的人；33 \BibleRef{路 22:30} 我应许你们，如同我圣父应许我一样，将王国赐给你们，使你们在我的王国里吃喝我筵席上的。\par

34 \BibleRef{路 22:7} 无酵节日到了，那一天犹太人应当宰杀逾越节羔羊。35 \BibleRef{路 22:8} 耶稣就打发刻法和若望两个门徒，说：你们去为我们预备逾越节晚餐，好让我们吃。36,37 \BibleRef{路 22:9} 他们对他说：你愿意我们在哪里预备？\BibleRef{路 22:10} 耶稣对他们说：你们进城去，在你们进入时，会有一个拿着水罐的人迎面而来；\BibleRef{路 22:11} 你们跟随他，到他进入的屋子，对那家的主人说：\BibleRef{玛 26:18} 我们的师傅说：我的时辰近了，并且（阿拉伯文第169页）我要在\Emph{你的家}里守逾越节。\BibleRef{路 22:11} 那么，我同门徒吃逾越节晚餐的客房在哪里？\BibleRef{路 22:12} 他必指给你们一间摆设整齐的大厅；\BibleRef{谷 14:15} 你们就在那里为我们预备。\BibleRef{谷 14:16} 两个门徒就出去，进了城，所遇见的正如他对他们说的；他们便照他所说的预备了逾越节晚餐。\par

41 \BibleRef{路 22:14} 到了晚上，时辰到了，耶稣就入席，十二位宗徒同他一起。42 \BibleRef{路 22:15} 他对他们说：我渴望渴望在受苦之前同你们吃这逾越节晚餐。43 \BibleRef{路 22:16} 我告诉你们：从今以后，我不再吃它，直到它在天主的王国里得以完成。\par

44 \BibleRef{若 13:21} 耶稣说了这话，并在心神里烦乱，就作证说：\Quote{我实实在在，实实在在地告诉你们：你们当中，那与我一同吃饭的\Emph{他}，要出卖我。}他们就非常忧愁，他们开始一个一个地问他说：\Quote{难道是我吗？主？}\BibleRef{谷 14:20} 他回答他们说：\Quote{是十二人中的一个，那与我一同在盘子里蘸手的，要出卖我。}\BibleRef{路 22:21} \Quote{看，那出卖我之人的手，正与我在桌子上。}\BibleRef{谷 14:21} \Quote{人子固然要照所记载的而去，但那出卖人子的人有祸了！因为那人生不生下来为他还更好。}\BibleRef{若 13:22} 门徒们就彼此对看，不知道他说的是谁；\BibleRef{路 22:23} 他们就开始彼此追问，究竟是谁要做\Emph{这}事。\par

\SourceRecord{Translated by Hope W. Hogg.；Ante-Nicene Fathers；Vol. 9.；Edited by Allan Menzies.；Buffalo, NY: Christian Literature Publishing Co.,；1896.；<http://www.newadvent.org/fathers/100244.htm>.}

% structure: p0001 source=h1 target=section reason=page-title

\section{四福音合参（第45节）}

1、2 [阿拉伯文版，第170页] \BibleRef{若 13:23} 他的门徒中，有一个靠在耶稣怀里，就是耶稣所爱的\Emph{那个}。\BibleRef{若 13:24} 西满·刻法向他示意，叫他问耶稣，他所说的\Emph{是}谁。\BibleRef{若 13:25} 那门徒便贴近耶稣的胸膛，问他说：主，是谁？\BibleRef{若 13:26} 耶稣回答说：我蘸一点饼递给谁，谁就是。耶稣就蘸了饼，递给加略人西满的儿子犹达斯。\BibleRef{若 13:27} 他吃了饼以后，撒殚进入了他。耶稣对他说：你要做的，快去做吧。\BibleRef{若 13:28} 在座的人没有一个明白耶稣为什么对他说这话。\BibleRef{若 13:29} 有人因为犹达斯掌管钱囊，以为耶稣是叫他去买过节所需用的，或是拿什么东西周济穷人。\BibleRef{玛 26:25} 出卖耶稣的犹达斯回答说：辣彼，难道是我吗？耶稣对他说：你说的是。\BibleRef{若 13:30} 犹达斯立刻接了饼，就出去；那时正是夜间。\par

10 \BibleRef{若 13:31} 耶稣说：现在人子受到了光荣，天主也在他身上受到了光荣。11 \BibleRef{若 13:32} 如果天主在他身上受到了光荣，天主也要在自己内光荣他，并且立刻光荣他。\par

12 \BibleRef{谷 14:22} 他们正吃的时候，耶稣拿起饼来，祝福了，擘开；\BibleRef{玛 26:26} 递给门徒，说：你们拿去吃吧！这是我的身体。13 \BibleRef{谷 14:23} 又拿起杯来，祝谢了，祝福了，递给他们；\BibleRef{玛 26:27} 说：你们大家拿去喝吧！14、15 \BibleRef{谷 14:23} 众人都喝了。\BibleRef{谷 14:24} 耶稣对他们说：\BibleRef{玛 26:28} 这是我的血，新约的血，为多人流出来，以赦免罪过。16 \BibleRef{玛 26:29} 我告诉你们：从今以后，我不再喝这葡萄汁，直到我在天主的国里，同你们喝新\Emph{酒}的那日子。17 \BibleRef{路 22:19} 你们应这样做，来纪念我。\BibleRef{路 22:31} 耶稣对西满说：西满，看，撒殚求得了许可，要筛你们像筛麦子一样；18 \BibleRef{路 22:32} 但我已为你祈求，使你不致失去信德；你回头以后，要坚固你的弟兄。\par

19 \BibleRef{若 13:33} 孩子们！我同你们在一起的\Emph{时候}不多了；你们将要寻找我，就如我曾对犹太人说过：我去的地方，你们不能去；现在我也对你们说。20 \BibleRef{若 13:34} 我给你们一条新命令：你们该彼此相爱；如同我爱了你们，你们也该照样彼此相爱。21 \BibleRef{若 13:35} 如果你们彼此相爱，众人因此就可认出你们是我的门徒。22 \BibleRef{若 13:36} 西满·刻法问他说：主！你往哪里去？耶稣回答说：我所去的地方，你现在不能跟我去，但后来你要跟我去。\par

23 \BibleRef{玛 26:31} 那时，耶稣对他们说：今夜你们都要为我的缘故跌倒，因为经上记载：我要打击牧人，羊群就四散。24 \BibleRef{玛 26:32} 但是我复活以后，要在你们以先往加里肋亚去。25 \BibleRef{玛 26:33} 西满·刻法回答他说：主！即使众人都为你跌倒，我决不会跌倒。\BibleRef{路 22:33} 我已经准备好同你一起坐监，一同赴死。\BibleRef{若 13:37} 我愿为你舍掉我的性命。27 [阿拉伯文版，第172页] \BibleRef{若 13:38} 耶稣对他说：你愿为我舍掉你的性命吗？\BibleRef{谷 14:30} 我实在告诉你：就在今天，这一夜，鸡叫两遍以前，你要三次不认我。但刻法更加坚决地说：即便我该同你一起死，我也决不会不认你，主！众门徒也都这样说了。\par

29 \BibleRef{若 14:1} 那时耶稣对他们说：\Quote{你们心里不要烦乱：要信天主，30 也要信我。\BibleRef{若 14:2} 在我父的家里有许多住处，不然，我早就告诉你们了。我\Emph{去}，是为给你们预备地方。\BibleRef{若 14:3} 如果我去为你们预备了地方，我必再来接你们到我那里去，为的是我在哪里，你们也在哪里。32, 33 \BibleRef{若 14:4} 我去的地方，你们知道；那条路，你们也知道。}\BibleRef{若 14:5} 多默对他说：\Quote{主，我们不知道你往哪里去，怎么知道那条路呢？}\BibleRef{若 14:6} 耶稣对他说：\Quote{我就是道路、真理、生命：除非经过我，谁也不能到父那里去。35 \BibleRef{若 14:7} 如果你们认识我，也就认识我的父；从现在起，你们认识他，并且已经看见了他。}\BibleRef{若 14:8} 斐理伯对他说：\Quote{主，把父显示给我们，我们就心满意足了。}\BibleRef{若 14:9} 耶稣对他说：\Quote{斐理伯，我与你们同在这么长久，你还不认识我吗？谁看见了我，就看见了父；你怎么说：\InnerQuote{把父显示给我们}呢？38 \BibleRef{若 14:10} 你不信我在父内，父在我内吗？我所说的话，不是凭我自己说的；而是住在我内的父，做他的事业。39 \BibleRef{若 14:11} 你们要信：我在父内，父在我内；40 ［阿拉伯语抄本，第173页］ 否则，至少该因那些事业而信。\BibleRef{若 14:12} 我实实在在告诉你们：凡信我的，我所做的事业，他也要做，并且还要做比这些更大的事业，因为我往父那里去。41 \BibleRef{若 14:13} 你们因我的名无论求什么，我必要践行，为叫父因子得光荣。42 \BibleRef{若 14:14} 你们若因我的名向我求什么，我必要\Emph{践行}。}43, 44 \BibleRef{若 14:15} \Quote{如果你们爱我，就要遵守我的诫命。\BibleRef{若 14:16} 我要请求父，他必会赐给你们另一位\OriginalTerm{护慰者}{\GreekText{παράκλητος}}，使他永远与你们同在；45 \BibleRef{若 14:17} 他\Emph{就是}真理之神，世界不能领受他，因为它看不见他，也不认识他；你们却认识他，因为他与你们同在，并在你们内。46 \BibleRef{若 14:18} 我必不撇下你们为孤儿；我要回到你们这里来。47 \BibleRef{若 14:19} 再过\Emph{不久}，世界就再看不见我，但你们却能看见我，因为我生活，你们也要生活。48 \BibleRef{若 14:20} 到那一天，你们便知道我在我父内，你们在我内，我也在你们内。}\par

\SourceRecord{Translated by Hope W. Hogg.；Ante-Nicene Fathers；Vol. 9.；Edited by Allan Menzies.；Buffalo, NY: Christian Literature Publishing Co.,；1896.；<http://www.newadvent.org/fathers/100245.htm>.}

% structure: p0001 source=h1 target=section reason=page-title

\section{四福音合参（第四十六节）}

1 \BibleRef{若 14:21} \BibleQuote{凡拥有我的诫命并遵守的，就是爱我的；那爱我的，必要蒙我父爱他，我也要爱他，并将我自己显示给他。} 2 \BibleRef{若 14:22} \BibleQuote{犹达斯（不是依斯加略人）对他说：\InnerQuote{主，祢为何要将自己显示给我们，而不显示给世界呢？}} 3 \BibleRef{若 14:23} \BibleQuote{耶稣回答他说：\InnerQuote{凡爱我的，必遵守我的话；我父也必爱他，并且我们要到他那里去，在他那里作住所。}} 4 \BibleRef{若 14:24} \BibleQuote{不爱我的，就不遵守我的话；你们所听到的话不是我的，而是那派遣我来的父的话。}\par

5、6 \BibleRef{若 14:25} \BibleQuote{我还与你们同在的时候，对你们说了这些事。} \BibleRef{若 14:26} \BibleQuote{但护慰者，就是圣神，即我父因我的名所要派遣的，祂必要教训你们一切，} \EditorialNote{阿拉伯文版第174页} \BibleQuote{并叫你们想起我对你们所说的一切。} 7 \BibleRef{若 14:27} \BibleQuote{我把平安留给你们；我将我的平安赐给你们：我所赐给你们的，不像世界所给的。} 8 \BibleRef{若 14:28} \BibleQuote{你们心里不要烦乱，也不要胆怯。你们听见我对你们说过：我去，还要到你们这里来。你们若爱我，就该喜欢我往父那里去，因为父比我大。} 9 \BibleRef{若 14:29} \BibleQuote{现在事情还没有发生，我就告诉你们，为叫你们当事情发生时，能相信我。} 10 \BibleRef{若 14:30} \BibleQuote{现在我不再同你们多说了，因为\OriginalTerm{世界的宰制者}{Archon}将要来到；他在我身上一无所有；} 11 \BibleRef{若 14:31} \BibleQuote{但为叫世界知道我爱父，并且父怎样命令我，我就照样行。}\par

12 \BibleRef{路 22:35} \BibleQuote{祂又对他们说：\InnerQuote{以前我派遣你们不带钱囊、口袋及鞋，你们缺少过什么吗？}他们回答说：\InnerQuote{什么也没有。}} 13 \BibleRef{路 22:36} \BibleQuote{祂对他们说：\InnerQuote{如今，有钱囊的，应当带着；口袋也当带着；没有刀的，要卖掉外衣去买一把。}} 14 \BibleRef{路 22:37} \BibleQuote{我告诉你们：经上所载：\InnerQuote{他被列于叛逆之中}这一句，必须应验在我身上，因为有关我的事，都即将完成。} 15 \BibleRef{路 22:38} \BibleQuote{门徒们说：\InnerQuote{主，请看，这里有两把刀。}祂回答说：\InnerQuote{够了。}} 16 \BibleRef{路 22:38} \BibleQuote{\InnerQuote{起来，我们走吧。}他们就起来，赞美了，便出去，照常往橄榄山去，祂和门徒们同去。}\par

17 \BibleRef{若 15:1} \BibleQuote{祂对他们说：\InnerQuote{我是真葡萄树，我父是园丁。}} 18 \BibleRef{若 15:2} \BibleQuote{凡在我身上不结果实的枝条，祂就剪去；凡结果实的，祂就清理，使它结更多的果实。} 19 \BibleRef{若 15:3} \BibleQuote{你们因我对你们所说的话，已是洁净的了。} 20 \BibleRef{若 15:4} \BibleQuote{你们住在我内，我也住在你们内。正如枝条若不留在葡萄树上，凭自己不能结果实；你们若不留在\Emph{我}内，也照样不能。} 21 \BibleRef{若 15:5} \BibleQuote{我是葡萄树，你们是枝条。那住在我内，我也住在他内的，就结许多果实；因为离了我，你们什么也不能做。} 22 \BibleRef{若 15:6} \BibleQuote{谁若不住在我内，便像枝条被丟在外面而枯干了；人便把它拾起来，投入火中焚烧。} 23 \BibleRef{若 15:7} \BibleQuote{如果你们住在我内，我的话也住在你们内，凡你们愿意求的，就必给你们\Emph{成就}。} 24 \BibleRef{若 15:8} \BibleQuote{我父因此得光荣，那就是你们结许多果实，而你们\Emph{必}成为我的门徒。} 25 \BibleRef{若 15:9} \BibleQuote{正如父爱了我，我也爱了你们；你们应住在我内。} 26 \BibleRef{若 15:10} \BibleQuote{如果你们遵守我的命令，便住在我的爱内，正如我遵守了我父的命令，而住在祂的爱内一样。} 27 \BibleRef{若 15:11} \BibleQuote{我对你们讲了这些，是为叫我的喜乐存留在你们内，并使你们的喜乐圆满无缺。} 28 \BibleRef{若 15:12} \BibleQuote{这是我的命令：你们该彼此相爱，如同我爱了你们一样。} 29 \BibleRef{若 15:13} \BibleQuote{人若为自己的朋友舍掉性命，再没有比这更大的爱情了。} 30 \BibleRef{若 15:14} \BibleQuote{如果你们实行我所命令你们的，你们就是我的朋友。} 31 \BibleRef{若 15:15} \BibleQuote{我不再称你们为仆人，因为仆人不知道主人所做的事。我称你们为朋友，因为凡由我父听来的一切，我都显示给你们了。} 32 \BibleRef{若 15:16} \BibleQuote{不是你们拣选了我，而是我拣选了你们，并任命你们去结果实，而你们的果实应当常存，好使你们因我的名向父求，祂就赐给你们。} 33 \BibleRef{若 15:17} \BibleQuote{这是我命令你们的：你们该彼此相爱。} 34 \BibleRef{若 15:18} \BibleQuote{如果世界恨你们，你们该知道，在你们以前，它已恨了我。} 35 \BibleRef{若 15:19} \BibleQuote{如果你们属于世界，世界必爱属于自己的；但因你们不属于世界，而是我从世界中拣选了你们，为此世界才恨你们。} \EditorialNote{阿拉伯文版第176页} 36 \BibleRef{若 15:20} \BibleQuote{你们要记得我对你们说过的话：没有仆人大于主人的。如果人们迫害了我，也要迫害你们；如果人们遵守了我的话，也要遵守你们的。} 37 \BibleRef{若 15:21} \BibleQuote{但他们为了我名的缘故，要这样对待你们，因为他们不认识那派遣我来的。} 38 \BibleRef{若 15:22} \BibleQuote{假使我没有来，没有教训他们，他们就没有罪；但现在他们对自己的罪无可推诿。} 39 \BibleRef{若 15:23} \BibleQuote{恨我的，也恨我的父。} 40 \BibleRef{若 15:24} \BibleQuote{假使我在他们中，没有做过别人从未做过的事，他们便没有罪；然而，他们现在看见了，仍恨了我和我的父。} 41 \BibleRef{若 15:25} \BibleQuote{这样，那写在他们法律上的话就得应验：\InnerQuote{他们无故恨我}。} 42 \BibleRef{若 15:26} \BibleQuote{当护慰者，就是我从父那里要派遣给你们的，那从父出发的真理之神来到时，祂必要为我作证；} 43 \BibleRef{若 15:27} \BibleQuote{并且你们也要作证，因为你们从起初就与我在一起。}\par

44, 45 \BibleRef{若 16:1} \BibleQuote{我已将这些事告诉了你们，免得你们跌倒。} \BibleRef{若 16:2} \BibleQuote{人要把你们逐出会堂；并且时候必到，凡杀害你们的 46 还以为是尽恭敬天主的义务。} \BibleRef{若 16:3} \BibleQuote{他们这样做，47 是因为没有认识我，也没有认识我的父。} \BibleRef{若 16:4} \BibleQuote{我已将这一切告诉你们，48 为叫你们在这一切发生时，想起我早已告诉你们了。这些事起初我没有告诉你们，因为我还与你们同在。}但现在我往派遣我者那里去，你们中却没有一个人49问我：你往哪里去？ \BibleRef{若 16:6} \BibleQuote{只因我给你们说了这话，50 你们就满心忧闷。} \BibleRef{若 16:7} \BibleQuote{然而，我将真情告诉你们：我去为你们有益，因为我若不去，\OriginalTerm{护慰者}{Paraclete}便不会到你们这里来；51 \EditorialNote{阿拉伯文版，第177页} 我若去了，就要派遣他到你们这里来。} \BibleRef{若 16:8} \BibleQuote{当他来到时，就要指证世界关于罪恶、正义和审判有罪：} 52、53 关于罪恶， \BibleRef{若 16:9} \BibleQuote{因为他们没有信从我；} \BibleRef{若 16:10} \BibleQuote{关于正义，因为我往父那里去54，而你们再见不到我；} \BibleRef{若 16:11} \BibleQuote{关于审判，因为这\OriginalTerm{世界的元首}{Archon}已55受了审判。} \BibleRef{若 16:12} \BibleQuote{我本来还有许多事要告诉你们，然而你们56现在不能担负。} \BibleRef{若 16:13} \BibleQuote{当那一位真理之神来到时，他要把你们引入一切真理，因为他不凭自己讲论，只把他所听到的57讲出来，并把未来的事传告给你们。} \BibleRef{若 16:14} \BibleQuote{他58要光荣我，因为他要把由我所领受的，传告给你们。} \BibleRef{若 16:15} \BibleQuote{凡父所有的一切，都是我的；为此我说：他要由我领受而传告给你们。}\par

\SourceRecord{Translated by Hope W. Hogg.；Ante-Nicene Fathers；Vol. 9.；Edited by Allan Menzies.；Buffalo, NY: Christian Literature Publishing Co.,；1896.；<http://www.newadvent.org/fathers/100246.htm>.}

% structure: p0001 source=h1 target=section reason=page-title

\section{\WorkTitle{迪亚泰萨隆（第四十七段）}}

1 \BibleRef{若 16:16} \Quote{\Emph{片刻}，你们将不再看见我；再\Emph{片刻}，你们又要看见我，因为我往父那里去。} 2 \BibleRef{若 16:17} 于是他的门徒彼此说：\Quote{他对我们说的\InnerQuote{\Emph{片刻}，你们将不再看见我；再\Emph{片刻}，你们又要看见我；又说：我往父那里去，}这是什么意思？} 3 \BibleRef{若 16:18} 他们又说：\Quote{他说的\Emph{片刻}是什么意思？我们不明白他说什么。} 4 \BibleRef{若 16:19} 耶稣知道他们想要问他，便对他们说：\Quote{你们是为我所说的\Emph{片刻}，你们将不再看见我；再\Emph{片刻}，你们又要看见我，而彼此议论么？} 5 \BibleRef{若 16:20} \Quote{我实实在在告诉你们：你们要痛哭哀号，世界却要欢乐；你们将忧愁，但你们的忧愁要变为喜乐。}\par

6 \BibleRef{若 16:21} \Quote{妇人生产的时候，时辰到了，她就忧愁；但一旦生下儿子，就不再记得那痛苦了，因为欢喜一个人生到世上来。} 7 \BibleRef{若 16:22} \Quote{现在你们也忧愁；但我将再见到你们，你们的心要喜乐，}[阿拉伯文，第178页]\Quote{你们的喜乐没有人能夺去。} 8 \BibleRef{若 16:23} \Quote{到那一天，你们什么也不必问我了。我实实在在告诉你们：你们因我的名无论向父求什么，他必赐给你们。} 9 \BibleRef{若 16:24} \Quote{直到现在，你们没有因我的名求过什么；求罢，你们必会得到，好使你们的喜乐得以圆满。}\par

10 \BibleRef{若 16:25} \Quote{我对你们说了这些事，迄今是用\OriginalTerm{谜语}{ænigmas}似的话；但时刻要来到，那时我不再对你们用谜语说话，而要明明地告诉你们有关父的事，到那一天，你们要因我的名祈求。} 11 \BibleRef{若 16:26} \Quote{我不对你们说，我要为你们求父；} 12 \BibleRef{若 16:27} \Quote{因为父自己爱你们，因为你们爱了我，并且相信我是从父那里来的。} 13 \BibleRef{若 16:28} \Quote{我从父那里来，来到了世上；我又离开世界，往父那里去。} 14 \BibleRef{若 16:29} 他的门徒说：\Quote{看，如今你明明地说话，没有用任何谜语。} 15 \BibleRef{若 16:30} \Quote{现在我们晓得你知道一切，不需要任何人问你；因此我们相信你是从天主来的。} 16 \BibleRef{若 16:31} \Quote{现在你们相信吗？} 17 \BibleRef{若 16:32} \Quote{看，时辰要来，且已来到，你们要各自散开，回到自己的地方，留下我独自一人；但我不是独自一人，因为有父与我同在。} 18 \BibleRef{若 16:33} \Quote{我把这些事告诉了你们，为叫你们在我内有平安。在世界上你们要受苦；但放心，我已战胜了世界。}\par

19 \BibleRef{若17:1} 耶稣说了这些话，就举目向天说：\Quote{我的圣父，时辰20来到了；求你光荣你的圣子，好叫圣子也光荣你：\BibleRef{若17:2} 就如你赐给了他权柄21掌管凡有血肉的人，为叫他将永生赐给你所托付给他的人。\BibleRef{若17:3} 永生就是：认识你，唯一的真天主，和你所派遣的\Emph{那一位}22 [阿拉伯文，第179页]，就是耶稣基督。\BibleRef{若17:4} 我在地上已光荣了你，23并且完成了你所委托我所作的工作。\BibleRef{若17:5} 圣父啊！现在，在你面前光荣我吧！赐给我在世界未有以前，我在你前所有的光荣24。\BibleRef{若17:6} 我已将你的名显示给那些你由世界赐给我的人；他们原属于你，你将他们托付给了我，他们也遵守了你的话25,26。\BibleRef{若17:7} 现在他们知道凡你赐给我的，都是由你而来的：\BibleRef{若17:8} 因为你所授给我的话，我都传给了他们；他们接受了\Emph{它们}，确实知道我是出于你，并且相信是你派遣了我。27 \BibleRef{若17:9} 我为他们祈求，不为世界祈求，只为你赐给我的人祈求，28因为他们是属于你的：\BibleRef{若17:10} 我的一切都是你的，你的一切都是29我的：我因他们已受了光荣。\BibleRef{若17:11} 从今以后，我不在世界上了，但他们仍在世界上，我却到你那里去。我的圣父啊！求你因你的30名，保全那些你所赐给我的人，使他们合而为一，如同我们一样。\BibleRef{若17:12} 当我和他们同在时，我因你的名，保全了你所赐给我的人，护卫了他们，其中除了那丧亡之子外，没有丧亡一个，31这是为应验经上的话。\BibleRef{若17:13} 但如今我到你那里去，我在世上讲这话，32是为叫他们心里充满我的喜乐。\BibleRef{若17:14} 我已将你的话授给了他们，世界33却憎恨他们，因为他们不属于世界，如同我不属于世界一样。\BibleRef{若17:15} 我不求你将他们从世界上撤去，只求你保护他们脱免34,35邪恶者。\BibleRef{若17:16} 他们不属于世界，如同我不属于世界一样。\BibleRef{若17:17} 圣父啊！求你以真理祝圣36他们，你的话就是真理。\BibleRef{若17:18} 就如你派遣我到世界上来，我37 [阿拉伯文，第180页] 也照样派遣他们到世界上去。\BibleRef{若17:19} 我为他们的缘故祝圣我自己，为叫他们38也因真理而被祝圣。\BibleRef{若17:20} 我不但为他们祈求，而且也为那些因他们的话而信从我的人39祈求；\BibleRef{若17:21} 愿众人都合而为一！圣父啊！愿他们在我们内合而为一，就如你在我内，我在你内，为叫世界40相信是你派遣了我。\BibleRef{若17:22} 我将你赐给我的41光荣赐给了他们；\BibleRef{若17:23} 为叫他们合而为一，就如我们原为一体一样；我在他们内，你在我内，使他们完全合而为一，并\Emph{为}叫世界知道42是你派遣了我，并且你爱了他们，如爱了我一样。\BibleRef{若17:24} 圣父啊！你所赐给我的人，我愿我在那里，他们也同我在一起，使他们享见你所赐给我的光荣，因为你在创世以前43就爱了我。\BibleRef{若17:25} 公义的圣父啊！世界没有认识你，44我却认识了你，这些人也知道是你派遣了我；我已将你的名宣示给他们了，还要宣示，\BibleRef{若17:26} 使你\Emph{用以}爱我的爱存在他们内，我也在他们内。}\par

\SourceRecord{Translated by Hope W. Hogg.；Ante-Nicene Fathers；Vol. 9.；Edited by Allan Menzies.；Buffalo, NY: Christian Literature Publishing Co.,；1896.；<http://www.newadvent.org/fathers/100247.htm>.}

% structure: p0001 source=h1 target=section reason=page-title

\section{四福音合参（第四十八部分）}

1 \BibleRef{若 18:1} 耶稣说了这话，就和门徒出去，到了一个名叫\Place{革责玛尼}的地方，\BibleRef{玛 26:36} 在\Place{克德龙}平原那边，那山，那里有一个园子；他就和门徒进去了。\BibleRef{若 18:2} 出卖他的\Person{犹达斯}也知道那地方，因为耶稣屡次同门徒在那里聚集。耶稣到了那地方，就对门徒说：\Quote{你们坐在这里，等我去祈祷；4, 5 ［Arabic, p. 181］你们要祈祷，免得陷于诱惑。} 他就带着\Person{刻法}和\Person{载伯德}的儿子们，就是\Person{雅各伯}和\Person{若望}，同去；他开始忧愁起来，并感到焦虑。\BibleRef{玛 26:38} 便对他们说：\Quote{我的心神忧闷得要死，你们留在这里，同我一起醒寤。} \BibleRef{路 22:41} 他就离开他们，约有扔一块石头那么远，\BibleRef{谷 14:35} 就跪下，俯伏在地，祈求，如果\Emph{可能}，使这时辰\Emph{能够}离开他。\BibleRef{谷 14:36} 他说：\Quote{父啊，你凡事都能；请把这杯撤去，\BibleRef{路 22:42} 但不要照我的\Emph{意愿}，而要照你的\Emph{意愿}。} 然后，他来到门徒那里，见他们睡着了，就对\Person{刻法}说：\Quote{\Person{西满}，你睡觉吗？\BibleRef{玛 26:40} 你们竟不能同我醒寤一个时辰吗？醒寤祈祷吧，免得陷于诱惑；心神固然切愿，肉身却软弱。} \BibleRef{玛 26:42} 他第二次又去祈祷说：\Quote{我父，这杯如果不能离开我，非要我喝不可，就愿你的\Emph{意愿}成就。} \BibleRef{谷 14:40} 他又回来，见门徒睡着了，因为他们的眼因忧愁和焦虑而沉重，他们不知道该回答他什么。\BibleRef{玛 26:44} 他又离开他们，第三次去祈祷，又说了同样的话。\BibleRef{路 22:43} 有一位天使从天上显现给他，鼓励他。\BibleRef{路 22:44} 他在极度恐惧中，祈祷越发恳切；他的汗如同血珠滴在地上。\BibleRef{路 22:45} 他从祈祷中起来，来到门徒那里，见他们因忧愁睡着了，就对他们说：\Quote{你们现在睡觉安歇吧！\BibleRef{谷 14:41} 时候到了，人子被交在罪人手里。起来，我们走吧！看，那出卖我的人来了。}\par

22 \BibleRef{玛 26:47} 他还说话的时候，那十二人之一的\Person{犹达斯}来了，同他一起有许多人，带着灯笼、火把\BibleRef{若 18:3} 和刀棒，从司祭长、经师和民间长老那里来，还有罗马的兵士们。出卖者\Person{犹达斯}给了他们一个暗号说：\Quote{我亲吻谁，谁就是。你们要小心拿住他，\Emph{带走。}}\par

24 \BibleRef{若 18:4} 耶稣知道将要临到自己的一切，就出去迎向他们。25 出卖他的\Person{犹达斯}立刻来到耶稣跟前说：\Quote{我的师傅，愿你平安。} 并亲吻了他。\BibleRef{路 22:48} 耶稣对他说：\Quote{\Person{犹达斯}，你用亲吻来出卖人子吗？\Emph{难道是}为这个而来的吗，我的朋友？} 耶稣又对来的人说：\Quote{你们找谁？} 他们回答说：\Quote{\Place{纳匝肋}人耶稣。} 耶稣说：\Quote{我就是。} 出卖他的\Person{犹达斯}也同他们站在一起。\BibleRef{若 18:6} 耶稣一说\Quote{我就是}，他们就倒退，跌倒在地。\BibleRef{若 18:7} 耶稣又问他们：\Quote{你们找谁？} 他们说：\Quote{\Place{纳匝肋}人耶稣。} \BibleRef{若 18:8} 耶稣说：\Quote{我已经告诉你们，我就是。你们若找我，就让这些人走吧。} \BibleRef{若 18:9} 这是为应验他所说的话：\BibleQuote{你所赐给我的人，我没有失落一个。} 33 ［Arabic, p. 183］\BibleRef{玛 26:50} 那时，同\Person{犹达斯}在一起的人上前，抓住耶稣，把他拿住了。\par

34 \BibleRef{路 22:49} 门徒们看见所发生的事，就说：\Quote{主啊，我们可以用刀砍吗？} 35 \BibleRef{若 18:10} \Person{西满} \Person{刻法}有一把刀，就拔出来，砍了大司祭的仆人，削掉他的右耳。那仆人名叫\Person{马尔曷}。36 \BibleRef{若 18:11} 耶稣对\Person{刻法}说：\Quote{我父给我的杯，我岂能不喝吗？把刀收入鞘内，因为凡动刀的，必死于刀下。\BibleRef{玛 26:53} 你想我不能求我父，他就立刻为我派来十二团以上的天使吗？\BibleRef{玛 26:54} 但经书所载必须如此的话，怎能应验呢？} \BibleRef{路 22:51} 耶稣说：\Quote{你们允许这样吧。} \BibleRef{玛 26:55} 耶稣摸了那被砍的耳朵，治好了他。同时耶稣对群众说：\Quote{你们带着刀棒出来拿我，如同拿强盗一样？我天天坐在圣殿里教导人，你们没有拿我。\BibleRef{路 22:53} 但这是你们的时候，黑暗的权势。} \BibleRef{玛 26:56} 这一切事的发生，是为应验先知们的经书。\par

44 \BibleRef{若 18:12} 于是门徒们都离开他逃跑了。兵士、军官和犹太人的差役拿住耶稣，把他带来了。45 \BibleRef{谷 14:51} 有一个少年人，赤身披着一块麻布，跟随耶稣，46 \BibleRef{谷 14:52} 他们就抓住他；他却丢下麻布，赤身逃走了。47 ［Arabic, p. 184］然后他们拿住耶稣，捆绑他，先带他到\Person{亚纳斯}那里，因为他是当年大司祭\Person{盖法}的岳父。48 \BibleRef{若 18:14} 这\Person{盖法}就是曾向犹太人提出忠告说\Quote{一个人替百姓死是有益的}的那位。\par

49 \BibleRef{若 18:15} 西满·\OriginalTerm{刻法}{Cephas}和另一个门徒跟随耶稣。那大司祭50认识那个门徒，他就同耶稣进了院子；\BibleRef{若 18:16} 西满却站在门外。那另一个门徒，是大司祭认识的，51出去对看门的使女说，她就带西满进来。\BibleRef{若 18:17} 那看门的使女看见西满，就定睛看着他，对他说：\Quote{你不也是这个人的门徒之一吗？我是说纳匝肋人52耶稣}但他否认说：\Quote{女人，我不认识他，也不知道53你说的是什么。}\BibleRef{若 18:18} 仆人和差役起身，在院子中间生了54火取暖，因为天冷。\BibleRef{路 22:55} 当55火着起来，他们就围坐四周。西满也来同他们一起坐下取暖，要看看事情如何结局。\par

\SourceRecord{Translated by Hope W. Hogg.；Ante-Nicene Fathers；Vol. 9.；Edited by Allan Menzies.；Buffalo, NY: Christian Literature Publishing Co.,；1896.；<http://www.newadvent.org/fathers/100248.htm>.}

% structure: p0001 source=h1 target=section reason=page-title

\section{\WorkTitle{四福音合参}（第49段）}

1, 2 \BibleRef{若 18:19} 大司祭问耶稣关于祂的门徒和祂的教训。\BibleRef{若 18:20} 耶稣对他说：\Quote{我向来公开对众人说话；我常在会堂和圣殿里教训人，那里是众犹太人聚集的地方；我在暗地里并没有说什么。3 \EditorialNote{阿拉伯文，第185页} 秘密。}\BibleRef{若 18:21} 你为什么问我？可以问那些听见的人，我\OriginalTerm{对4他们}{对4他们}说了什么；因为他们知道我说的一切。\BibleRef{若 18:22} 耶稣说了这话，站在\Emph{那里}的差役之一打了耶稣一巴掌，说：5 \Quote{你就这样回答大司祭吗？}\BibleRef{若 18:23} 耶稣回答说：6 \Quote{我若说得不对，你可以指证哪里不对；若说得对，你为什么打我呢？}\BibleRef{若 18:24} 亚纳斯就把耶稣捆绑着，解送到大司祭盖法那里。\par

7 \BibleRef{若 18:25} 耶稣出去时，西满刻法站在外院取暖。8 \BibleRef{谷 14:69} 那个使女又看见他，开始对\Emph{旁边}站着的人说：\Quote{这个\Emph{人}也是同纳匝肋人耶稣一起的。}\Emph{旁边}站着的人就上前对刻法说：\Quote{你真是祂的门徒之一。}\BibleRef{玛 26:72} 他再次起誓否认说：\Quote{我不认识那个人。}过了一会儿，大司祭的一个仆人——就是西满砍掉其耳朵那人的亲戚——看见他，就争辩说：\Quote{这个\Emph{人}确实同祂在一起：\BibleRef{路 22:59} 他也是加里肋亚人，他的口音也像。}他对西满说：\Quote{我不是在园子里看见你同祂在一起吗？}\BibleRef{谷 14:71} 西满就开始咒骂起誓说：\Quote{我不认识你们所说的这个\Emph{人。}}\BibleRef{路 22:60} 他正说话的时候，鸡立刻叫了两次。\BibleRef{路 22:61} 那时，耶稣转过身来，祂在外面，定睛看着刻法。西满就想起主对他所说的话：17,18 \Quote{鸡叫两次以前，你要三次否认我。}\BibleRef{路 22:62} 西满就出去，痛哭起来。\par

19 \EditorialNote{阿拉伯文，第186页} \BibleRef{路 22:66} 到了早晨，众大司祭的仆役、经师、民间的长老，以及全体群众都聚集，20,21 共同策划；\BibleRef{玛 27:1} 他们商议要把耶稣处死。\BibleRef{玛 26:59} 他们寻找假见证，来控告祂，好把祂处死，22,23 却找不着；\BibleRef{玛 26:60} 但有许多假见证人前来，\BibleRef{谷 14:59} 然而他们的见证却不一致。24,25 \BibleRef{玛 26:60} 最后，有两个假见证人前来，说：\Quote{我们听见祂说：我要拆毁这座人手所造的圣殿，三天内要另造一座非人手所造的。}26,27 但即便如此，他们的见证仍不一致。耶稣却沉默不语。\BibleRef{谷 14:60} 大司祭就起来，站在中间，问耶稣说：28 \BibleRef{玛 26:62} \Quote{你什么都不回答吗？这些人作证反对你，29,30 你怎么办？}\BibleRef{谷 14:61} 耶稣仍然沉默，一言不答。\BibleRef{路 22:66} 他们就把祂带到他们的公会里，对祂说：\Quote{如果你是默西亚，就告诉我们。}\BibleRef{路 22:67} 祂对他们说：\Quote{如果我告诉你们，你们也不会相信；\BibleRef{路 22:68} 如果我问你们，你们也不会回答我一句话，也不放我走。}33 \BibleRef{玛 26:63} 大司祭回答祂说：\Quote{我指着永生天主命令你，告诉我们，你是不是默西亚，永生天主子？}34,35 \BibleRef{玛 26:64} 耶稣对他说：\Quote{你说的是。}\BibleRef{路 22:70} 他们都说：\Quote{那么，你是天主子了？}耶稣说：\Quote{你们说的是。}36 \BibleRef{玛 26:64} \Quote{我告诉你们：从今以后，你们要看见人子坐在37 \EditorialNote{阿拉伯文，第187页} 大能者的右边，乘着天上的云彩降临。}38 大司祭就撕裂自己的衣服，说：\Quote{祂说了亵渎的话。}\BibleRef{路 22:71} 他们都说：\Quote{现在我们何必再需要见证人呢？我们亲自从祂口中听见了亵渎的话。39,40 你们看怎么办呢？}他们全体回答说：\Quote{祂该死。}于是有些人上前，吐唾沫在祂脸上，打祂，戏弄祂。41 差役也打祂的脸，说：\Quote{给我们说预言，\Emph{你}默西亚，42 打你的是谁？}\BibleRef{路 22:65} 他们还说了许多别的话，诬蔑祂，攻击祂。\par

43 他们的全体议会都起来，把耶稣捆绑着，带到44总督府，交给审判官比拉多。\BibleRef{若 18:28} 但他们没有进入总督府，免得吃逾越节时沾染不洁。\par

45 耶稣站在审判官面前。比拉多出来到他们那里，对他们说：46 \Quote{你们对这个人有什么控告？}\BibleRef{若 18:30} 他们回答说：\Quote{如果祂不是作恶的，我们也不会把祂47交给你。}\BibleRef{路 23:2} \Quote{我们发现这个\Emph{人}迷惑我们的百姓，禁止纳税给凯撒，并自称是君王默西亚。}48 \BibleRef{若 18:31} 比拉多对他们说：\Quote{你们自己把祂带去，按照你们的法律审判祂吧。} \EditorialNote{阿拉伯文，第188页} 犹太人对他说：\Quote{我们没有权力处死人。}49 \BibleRef{若 18:32} 这是为应验耶稣所说的话，当祂指明自己将要怎样死的时候。\par

50 \BibleRef{若 18:33} 比拉多进了总督府，叫耶稣来，对他说：你是 51 犹太人的君王吗？\BibleRef{若 18:34} 耶稣对他说：你自己说了这话，或是 52 别人论及我对你说的？\BibleRef{若 18:35} 比拉多对他说：莫非我是个犹太人？你的民族和司祭长把你交给了我，你做了什么？\BibleRef{若 18:36} 耶稣说：我的国不属这世界；假如我的国属这世界，我的臣仆必要争战，使我不致 53 被交给犹太人；但如今我的国不是从这里来的。\BibleRef{若 18:37} 比拉多对他说：那么，\Emph{你是}一个君王吗？耶稣回答说：你说的是，我是君王。我为此而生，也为此来到世界，为给真理作证；\BibleRef{若 18:38} 凡属于真理的，必听我的声音。比拉多对他说：什么是真理？说了这话，他又出去到犹太人那里。\par

\SourceRecord{Translated by Hope W. Hogg.；Ante-Nicene Fathers；Vol. 9.；Edited by Allan Menzies.；Buffalo, NY: Christian Literature Publishing Co.,；1896.；<http://www.newadvent.org/fathers/100249.htm>.}

% structure: p0001 source=h1 target=section reason=page-title

\section{\WorkTitle{四福音合参}（第50节）}

1 \BibleRef{路 23:4} 比拉多对司祭长和群众说：\Quote{我在这人身上查不出什么罪来。} 2 \BibleRef{路 23:5} 但他们喊叫着说：\Quote{祂用自己的教导在犹太全境煽动百姓，从加里肋亚开始，直到这里。} 3 \BibleRef{路 23:6} 比拉多一听见加里肋亚，就问：\Quote{这人是加里肋亚人吗？} 4 \BibleRef{路 23:7} 他一知道祂属黑落德管辖，就把他送到黑落德那里，因为黑落德那时也在耶路撒冷。\par

5 \BibleRef{路 23:8} 黑落德一见耶稣，就极其欢喜；因为他很久以来就想见祂，因他听见关于祂的许多事，并指望［阿拉伯文，第189页］看见祂行些奇迹。6 \BibleRef{路 23:9} 他用许多话问祂，但7耶稣一言不答。\BibleRef{路 23:10} 经师和司祭长站在一旁，极力控告祂。8 \BibleRef{路 23:11} 黑落德和他的士兵藐视祂，戏弄祂，给祂穿上红袍，然后把祂送回比拉多那里。9 \BibleRef{路 23:12} 那一天，比拉多和黑落德成了朋友，因为他们先前彼此有仇。\par

10,11 \BibleRef{路 23:13} 比拉多召集了司祭长和民间的首领，\BibleRef{路 23:14} 对他们说：\Quote{你们把这人带来，当作煽动你们百姓的。我在你们面前审问了祂，并没有查出你们控告祂的任何罪状；12 \BibleRef{路 23:15} 连黑落德也没有，因为我把他送到他那里，他也没有查出祂有什么该处死的事。13 \BibleRef{路 23:16} 所以，我要惩戒祂，然后释放祂。} 14,15 \BibleRef{路 23:18} 众人一齐喊叫说：\Quote{除掉祂，除掉祂！} \BibleRef{谷 15:3} 16 司祭长和长老控告祂许多事。\BibleRef{玛 27:12} 他们控告祂的时候，祂一言不答。17 \BibleRef{玛 27:13} 比拉多对祂说：\Quote{你没听见他们作证控告你这么多事吗？} 18 \BibleRef{玛 27:14} 耶稣连一句话也不回答他，以致总督大为惊奇。\par

19 \BibleRef{玛 27:19} 当总督坐在审判台上时，他的妻子派人来对他说：\Quote{你别管那\Emph{义人}，因为今天我因祂在梦中受了许多苦。}\par

20 \BibleRef{玛 27:15} 每逢节日，总督常照例给群众释放一个他们所愿的囚犯。21 \BibleRef{玛 27:16} 那时，监里有一个出名的囚犯，22,23 名叫巴拉巴。\BibleRef{玛 27:17} 当他们聚集时，比拉多对他们说：\BibleRef{若 18:39} \Quote{你们有个惯例，我可在逾越节给你们释放一个囚犯：你们愿意我给你们释放犹太人的君王吗？} 24 \BibleRef{若 18:40} 他们却一齐喊叫说：\Quote{不要这人，但要巴拉巴！} 这个巴拉巴是个强盗，25 \BibleRef{路 23:19} 他因在城里煽动叛乱和杀人而被下在监里。26 \BibleRef{谷 15:8} 群众一齐喊叫，开始要求\Emph{他照常办理}。27 比拉多回答他们说：\Quote{你们愿意我给你们释放谁？巴拉巴，还是那称为默西亚的耶稣、犹太人的君王？} 28 \BibleRef{玛 27:18} 原来比拉多知道他们是因嫉妒才把祂解来的。29 \BibleRef{玛 27:20} 司祭长和长老却唆使群众要求释放巴拉巴，而除掉耶稣。30 \BibleRef{玛 27:21} 总督回答他们说：\Quote{这两个人中，你们愿意我给你们释放哪一个？} 他们说：\Quote{巴拉巴。} 31 \BibleRef{玛 27:22} 比拉多对他们说：\Quote{那么，那称为默西亚的耶稣，我该怎么处置呢？} 32 \BibleRef{谷 15:13} 他们都喊叫说：\Quote{钉死祂！} 33 \BibleRef{路 23:20} 比拉多又对他们说话，因为他愿意释放耶稣；34 \BibleRef{路 23:21} 但他们却喊叫说：\Quote{钉死祂，钉死祂！释放巴拉巴给我们！} 35 \BibleRef{路 23:22} 比拉多第三次对他们说：\Quote{这人做了什么恶事？我在祂身上查不出什么该死的罪来。所以，我要惩戒祂，然后释放祂。} 36 \BibleRef{路 23:23} 但他们大声坚持，要求钉死祂。他们的声音和司祭长的声音占了优势。37 于是比拉多释放了那个因煽动叛乱和杀人而被下狱的巴拉巴，就是他们所要求的；\BibleRef{玛 27:26} 并用鞭子鞭打了耶稣。\par

38 \BibleRef{玛 27:27} 总督的兵士抓住耶稣，进入总督府，39 \EditorialNote{阿拉伯语本，第191页} 把全队兵士都聚集到他那里。\BibleRef{玛 27:28} 他们剥去他的衣服，给他披上40 一件朱红袍。\BibleRef{若 19:2} 他们给他穿上紫袍，编了41 一顶荆棘冠，戴在他头上，又把一根苇竿放在他右手里；\BibleRef{玛 27:29} 他们戏弄他，嘲笑他，跪在他面前，俯伏42 拜他，说：\Quote{犹太人的王，万岁！}\BibleRef{玛 27:30} 他们吐唾沫在他脸上，又从他手里夺过苇竿，打他的头，\BibleRef{若 19:3} 并掴他的脸颊。43 \BibleRef{若 19:4} 比拉多又到外面去，对犹太人说：\Quote{我把他带到44 你们面前，好叫你们知道，我查不出他有什么罪。}\BibleRef{若 19:5} 耶稣就出来，戴着荆棘冠，穿着紫袍。45 \BibleRef{若 19:6} 比拉多对他们说：\Quote{看，这个人！}司祭长和兵士看见他，就喊叫说：\Quote{钉死他，钉死他！}比拉多对他们说：\Quote{你们自己把他带去钉十字架吧！因为我查不出46 他有什么罪状。}\BibleRef{若 19:7} 犹太人对他说：\Quote{我们有法律，按法律他应47 当死，因为他自命为天主子。}\BibleRef{若 19:8} 比拉多听了这话，48 越发害怕；\BibleRef{若 19:9} 又进了府邸，对耶稣说：\Quote{你49 是哪里来的？}\BibleRef{若 19:10} 耶稣却不回答。比拉多对他说：\Quote{你不对我说话吗？你不知道我有权柄释放你，也有50 权柄钉死你吗？}\BibleRef{若 19:11} 耶稣回答说：\Quote{若不是从上头赐给你的\Emph{它}，你对我就没有任何权柄；因此，把我交给你的人，他的罪51 比你的罪更大。}\BibleRef{若 19:12} 为此，比拉多想要释放他；但犹太人喊叫说：\Quote{你若释放这人，你就不是凯撒的朋友；因为凡自命为王的，就是反对凯撒。}\par

\SourceRecord{Translated by Hope W. Hogg.；Ante-Nicene Fathers；Vol. 9.；Edited by Allan Menzies.；Buffalo, NY: Christian Literature Publishing Co.,；1896.；<http://www.newadvent.org/fathers/100250.htm>.}

% structure: p0001 source=h1 target=section reason=page-title

\section{\WorkTitle{迪亚特萨隆（第51节）}}

1［阿拉伯文版，第192页］\BibleRef{若 19:13} 当\Person{比拉多}听见这话，就把\Person{耶稣}带出来，坐在审判台上，那地方称为石铺地，希伯来语2称为\OriginalTerm{加巴塔}{Gabbatha}。\BibleRef{若 19:14} 那天是逾越节预备日，约莫3第六时辰。\BibleRef{若 19:15} 他对犹太人说：\Quote{看，你们的君王！}他们就喊叫：\Quote{带走他，带走他，钉死他，钉死他！}\Person{比拉多}对他们说：\Quote{我该钉死你们的君王吗？}司祭长对他说：\Quote{除了凯撒4，我们没有君王。}\BibleRef{玛 27:24} \Person{比拉多}看见\Emph{此}，一无所获，而骚动愈烈，就拿水在众人面前洗手，说：\Quote{我流这无辜\Emph{之人}的血，与我无干；你们要知道5。}\BibleRef{玛 27:25} 全体百姓6回答说：\Quote{他的血归在我们和我们的子孙身上！}\BibleRef{若 19:16} 于是\Person{比拉多}命令准他们所求；遂将\Person{耶稣}交出，任凭他们钉死。\par

7\BibleRef{玛 27:3} 那时，出卖耶稣的\Person{犹达斯}看见耶稣受冤，就去把三十块8银钱归还给司祭长和长老，\BibleRef{玛 27:4} 说：\Quote{我出卖了无辜之血，犯了罪。}他们对他说：\Quote{那我们9怎么办？你自己知道。}\BibleRef{玛 27:5} 他把银钱扔在圣殿里，就离去；\Emph{他}10出去吊死了。\BibleRef{玛 27:6} 司祭长拿起银钱说：\Quote{不可放入献仪箱内，因为这是血价11。}\BibleRef{玛 27:7} 他们商议，就用那钱买了陶工的田地12，作为埋葬外乡人之用。\BibleRef{玛 27:8} 因此那块田直到今日13被称为血田。\BibleRef{玛 27:9} 这就应验了先知所说的话：\Quote{我取了三十块银钱，即以色列子民所估定的宝贵14\Emph{之价}，}\BibleRef{玛 27:10} \Quote{我付了陶工田地的价钱，正如主所吩咐我的。}\par

15 犹太人抓住\Person{耶稣}，把他带去钉死。当他背着十字架出去时，他们剥下他身上穿的紫红袍和朱红袍，给他穿上他自己的衣服。他们行进时，遇见一个\Place{基勒乃}人，名叫\Person{西满}，是从乡下来的，是\Person{亚历山大}和\Person{鲁富}的父亲；\BibleRef{玛 27:32} 他们就强迫这个人16背耶稣的十字架。\BibleRef{路 23:26} 他们把十字架放在他身上，叫他背着跟随耶稣；耶稣走，他的十字架在后面。\par

19\BibleRef{路 23:27} 有许多百姓和哭泣哀号20的妇女跟随他。\BibleRef{路 23:28} 耶稣转过身来，对她们说：\Quote{\Place{耶路撒冷}女子，不要为我哭，当为你们自己和你们的21儿女哭。}\BibleRef{路 23:29} \Quote{日子将到，那时人要说：不生育的，未曾怀胎的，未曾哺养过的，是有福的。}\BibleRef{路 23:30} \Quote{那时，人要开始对山说：倒在我们身上！对丘陵说：遮盖我们22！}\BibleRef{路 23:31} \Quote{因为如果他们对青绿的树木这样做，对枯干的将会怎样呢？}\par

24\BibleRef{路 23:32} 他们又带了两个凶犯同耶稣一起处死。\par

25 他们来到一个名为髑髅的地方，希伯来语称作\OriginalTerm{哥耳哥达}{Golgotha}，就在那里钉死他。\BibleRef{路 23:33} 他们把这两个凶犯也钉了，一个在右边，一个在左边26。\BibleRef{谷 15:28} 应验了圣经所说：\Quote{他被列于罪犯之中27。}\BibleRef{谷 15:23} 他们拿搀了没药的酒和醋给他喝；他尝了，却不肯喝，也不接受。\par

28\BibleRef{若 19:23} 兵士钉了耶稣后，拿了他的衣服，分成四份，每个兵一份；他那件长衣是无缝的，从上到下浑然织成29。\BibleRef{若 19:24} 他们彼此说：\Quote{不要撕破，我们抽签，看是谁的。}这就应验了圣经的话：\par

\Quote{他们分了我的外衣，\newline{}又为我的里衣拈阄。}\par

30, 31 \BibleRef{玛 27:36} 兵士便这样做了。他们坐在那里看守他。\BibleRef{若 19:19} \Person{比拉多}写了一个罪状牌，放在十字架上，在他头上。\BibleRef{玛 27:37} 上面写着：\Quote{这是耶稣，\Place{纳匝勒}人，犹太人的32君王。}\BibleRef{若 19:20} 许多犹太人念了这牌子，因为耶稣被钉的地方离城很近，而且是用希伯来、罗马和希腊文写的。33 \BibleRef{若 19:21} 司祭长对\Person{比拉多}说：\Quote{不要写\InnerQuote{犹太人的君王}，而要写\InnerQuote{他说：我是犹太人的君王}}34。\BibleRef{若 19:22} \Person{比拉多}说：\Quote{我所写的，已经写了。}35 百姓站着观望；过路的人36羞辱他，摇摇头说：\Quote{你这拆毁圣殿、三日内重建的37，}\BibleRef{玛 27:40} \Quote{如果你是天主子，救你自己吧，从十字架上下来！}\BibleRef{玛 27:41} 司祭长、经师、长老和\Person{法利塞人}也同样讥诮他，彼此笑着说：38,39 \BibleRef{玛 27:42} \Quote{他救了别人，却不能救自己！如果他真是默西亚，天主的选民，以色列的君王，现在就从十字架上下来，让我们看见而信从他}40。\BibleRef{玛 27:43} \Quote{他倚靠天主——如果天主喜悦他，现在就救他吧41；因为他曾说过：我是天主子。}\BibleRef{路 23:36} 兵士也戏弄他，上前来，\BibleRef{路 23:37} 拿醋给他喝42，对他说：\Quote{如果你是犹太人的君王，救你自己吧！43}\BibleRef{玛 27:44} 同他一起钉的凶犯也都辱骂他。\par

44 \BibleRef{路 23:39} 与祂同钉十字架的两个凶犯中，有一个辱骂祂，说：\Quote{如果你是\OriginalTerm{默西亚}{Messiah}，救你自己，也救我们吧。} 45 \BibleRef{路 23:40} 但他的同伴斥责他，说：\Quote{你同样受这刑罚，难道也不敬畏天主吗？\BibleRef{路 23:41} 我们受罚是公道的，因为我们所受的与我们所行的相称；但这个\Emph{人}没有做过任何不法的事。} 47 \BibleRef{路 23:42} 于是他对耶稣说：\Quote{主啊，当祢进入祢的国度时，求祢记念我。} 48 \BibleRef{路 23:43} 耶稣对他说：\Quote{我实在告诉你，今天你就要同我在\OriginalTerm{乐园}{Paradise}里了。}\par

49 \BibleRef{若 19:25} 在耶稣的十字架旁，站着祂的母亲，祂母亲的姊妹，50 \Person{克罗帕}的玛利亚，和玛利亚玛达肋纳。\BibleRef{若 19:26} 耶稣看见母亲和祂所爱的门徒站\Emph{在近旁}，就对母亲说：\Quote{女人，看，你的儿子！}\BibleRef{若 19:27} 然后对那门徒说：\Quote{看，你的母亲！} 从那时起，那门徒就把她接到\Emph{自己}家里。\par

52 \EditorialNote{阿拉伯文，第196页} 从第六时辰起，黑暗笼罩了全地，直到第九53时辰，太阳也变暗了。到了第九时辰，耶稣大声呼喊说：\Quote{\OriginalTerm{厄里}{Yail, Yaili}，你为什么舍弃了我？}意思是：\Quote{我的天主，我的天主，祢为什么舍弃了我？}\BibleRef{玛 27:47} 站在那里的人中，有几个听见了，就说：\Quote{这\Emph{人}呼叫\OriginalTerm{厄里亚}{Elijah}。}\par

\SourceRecord{Translated by Hope W. Hogg.；Ante-Nicene Fathers；Vol. 9.；Edited by Allan Menzies.；Buffalo, NY: Christian Literature Publishing Co.,；1896.；<http://www.newadvent.org/fathers/100251.htm>.}

% structure: p0001 source=h1 target=section reason=page-title

\section{《迪亚塔撒隆》（第五十二部分）}

1 \BibleRef{若 19:28}此后，耶稣知道一切事都已完成了，为应验经上的话，就说：\BibleQuote{我渴了。}那里放着一个盛满醋的器皿；那时，有一个人赶快跑去，拿了一个海绵，蘸满了那醋，\BibleRef{谷 15:36}绑在芦苇上，送到他口边给他喝。\BibleRef{若 19:30}耶稣尝了那醋后，就说：\BibleQuote{一切都完成了。}5 但其他人却说：\Quote{等着，让我们看看厄里亚是否来救他。}6,7 \BibleRef{路 23:46}耶稣说：\BibleQuote{父啊，宽恕他们吧，因为他们不知道自己在做什么。}耶稣又大声呼喊说：\BibleQuote{父啊，我把我的灵魂交在你手中。}\BibleRef{若 19:30}他说了这话，就低下头，断了气。\par

8 \BibleRef{玛 27:51}立时，圣所的门帘从上到下裂为两半；\BibleRef{玛 27:52}大地震动，磐石崩裂；坟墓开了，许多已安眠的圣人的身体复活了，10 \BibleRef{玛 27:53}在耶稣复活后，他们进入圣城，向许多人显现。\BibleRef{玛 27:54}百夫长和与他一起看守耶稣的人，看见地震和所发生的事，就非常害怕，赞美天主说：\BibleQuote{这个人真是义人；}13 \BibleQuote{这真是天主子。}\BibleRef{路 23:48}所有聚集来看这景象的群众，看见所发生的事，都捶着胸回去了。\par

14 \BibleRef{若 19:31}犹太人因为那天是预备日，就说：不要让这些尸体留在十字架上，因为安息日已到（那安息日是个大日子）；他们请求比拉多打断那些被钉者的腿，好把他们取下来。15 \BibleRef{若 19:32}兵士便来了，打断了第一个人的腿，以及与他同钉的另一人的腿；16 \BibleRef{若 19:33}但来到耶稣那里，见他已死了，就没有打断他的腿；17 \BibleRef{若 19:34}有一个兵士用枪刺透了他的肋旁，立刻流出了血和水。18 \BibleRef{若 19:35}看见这事的人就作了见证，他的见证是真实的，并且他知道自己说的是真话，为叫你们也相信。19 \BibleRef{若 19:36}这事发生，为应验经上的话：\BibleQuote{他的骨头一根也不可折断。}20 \BibleRef{若 19:37}另一处经上也说：\BibleQuote{他们要仰望自己所刺透的人。}\par

21 远处站着耶稣的所有熟人，以及那些从加里肋亚跟随他、服事他的妇女。22 其中有玛利亚玛达肋纳、小雅各伯和若瑟的母亲玛利亚、载伯德儿子的母亲，以及撒罗默，还有许多其他与他一同上耶路撒冷来的妇女。\BibleRef{路 23:49}她们都看见了这些事。\par

24 \BibleRef{谷 15:42}到了预备日的傍晚，因为安息日就要开始，25 \BibleRef{路 23:50}来了一个富人，\BibleRef{玛 27:57}是阿黎玛特雅城的尊贵议员，\BibleRef{路 23:51}那城属犹大地，名叫若瑟，他是个善良正直的人；\BibleRef{若 19:38}他曾是耶稣的门徒，但因害怕犹太人而隐瞒身份。\BibleRef{路 23:51}他并没有附和那些控告者的意愿和行为；\BibleRef{路 23:51}他一直在等候天主的国。\BibleRef{谷 15:43}这人就大胆地进去见比拉多，向他求取耶稣的身体。\BibleRef{谷 15:44}比拉多诧异耶稣已经死了，便叫来百夫长，问他耶稣是否已死。30 \BibleRef{谷 15:45}得知实情后，他就吩咐将耶稣的身体交给若瑟。若瑟买了洁白的细麻布，取下耶稣的身体，用细麻布裹好；他们来取走了遗体。31 尼苛德摩也来了，他从前曾在夜间来见耶稣，带着约一百斤的没药和沉香调和的香料。\BibleRef{若 19:40}他们便取下耶稣的身体，用香料和细麻布裹好，一如犹太人的埋葬习俗。\par

33 \BibleRef{若 19:41}在耶稣被钉十字架的地方，有一个园子；园子里有一座新坟墓，是在岩石中凿成的，里面从来没有安葬过人。34 \BibleRef{若 19:42}他们就把耶稣安放在那里，因为那天是犹太人的预备日，而且坟墓就在附近。35 \BibleRef{玛 27:60}他们滚来一块大石头，堵住墓门，就离开了。36 玛利亚玛达肋纳和若瑟的亲属玛利亚跟在后面，37 \BibleRef{路 23:55}来到坟墓前，对着坟墓坐着，\BibleRef{玛 27:61} 38 看见遗体是怎样被安放进去的。她们回去后，购买了香料和香膏，预备好\Emph{它们}，以便来抹他。39 \BibleRef{路 23:56}在安息日，她们便遵照诫命安息了。\par

40,41 \BibleRef{玛 27:62}司祭长和法利塞人聚集到比拉多那里，对他说：\Quote{\BibleRef{玛 27:63}大人，我们记得那个骗子活着的时候说过：\InnerQuote{三天后 42 我要复活。}\BibleRef{玛 27:64}所以，请你派人看守坟墓直到第三天，免得他的门徒夜间来把他偷走，然后对百姓说：\InnerQuote{他从死人中 43 复活了。}那后来的欺骗就比先前的更坏了。}\BibleRef{玛 27:65}比拉多对他们说：\Quote{你们不是有卫兵吗？去，你们自己尽量防范吧。}44 \BibleRef{玛 27:66}他们就去，在坟墓前布置了\Emph{守卫}，又封了那块石头，并派卫兵看守。\par

安息日的傍晚，即一周第一天的清晨，天还黑的时候，\newline{}\Person{玛利亚玛达肋纳}和另一个\Person{玛利亚}以及其他妇女来看坟墓。她们带着预备好的香料，彼此说：\newline{}\BibleRef{谷 16:3}\BibleQuote{谁为我们从坟墓门口滚开那块石头呢？因为那块石头很大。}\newline{}她们说这话的时候，忽然发生大地震；有天使从天降下来，把石头从门口滚开。她们来到，发现石头已经从坟墓滚开，天使坐在石头上。\newline{}\BibleRef{玛 28:3}\BibleQuote{他的相貌如同闪电，衣服洁白如雪：}\newline{}\BibleRef{玛 28:4}\BibleQuote{看守的人因怕他，吓得浑身乱战，甚至像死\Emph{人}一样。}\newline{}\BibleRef{路 24:3}\BibleQuote{天使离开后，妇女们进了坟墓，不见\Person{耶稣}的身体。}\newline{}\BibleRef{谷 16:5}\BibleQuote{她们看见一个少年人坐在右边，穿着白袍，就甚惊恐。}\newline{}\BibleRef{玛 28:5}\BibleQuote{天使对妇女说：\InnerQuote{不要害怕！我知道你们寻找那被钉十字架的纳匝肋人\Person{耶稣}。他不在这里，照他所说的，已经复活了。}}\newline{}\BibleRef{玛 28:6}\BibleQuote{你们来看安放我们主的地方。}\par

\SourceRecord{Translated by Hope W. Hogg.；Ante-Nicene Fathers；Vol. 9.；Edited by Allan Menzies.；Buffalo, NY: Christian Literature Publishing Co.,；1896.；<http://www.newadvent.org/fathers/100252.htm>.}

% structure: p0001 source=h1 target=section reason=page-title

\section{\WorkTitle{四福音合编（第53节）}}

1 \BibleRef{路 24:4} 当他们对此感到惊奇时，看，有两个男子站在她们旁边，他们的衣裳闪烁发光。2 \BibleRef{路 24:5} 她们害怕，就俯伏在地，他们对她们说：你们为什么在死人中寻找\Emph{生活者}呢？3 \BibleRef{路 24:6} 他不在这里；他已复活了。要记得他还在加里肋亚时对你们所说的话，4 \BibleRef{路 24:7} 并说：人子必须被交在罪人手中，并被钉在十字架上，并在第三日复活。5 \BibleRef{玛 28:7} 但是你们快去，告诉他的门徒和刻法：他已从死人中复活了；而且看，他要在你们之先去加里肋亚；在那里你们要看见他，正如他对你们所说的：看，我已经告诉你们了。6 \BibleRef{谷 16:8} 她们想起了他的话；就急忙从坟墓那里离开，又欢喜又大大地害怕，跑去报告；并且惊惶和恐惧笼罩着她们；她们什么也不告诉人，因为害怕。7 \BibleRef{若 20:2} 玛利亚就跑去见西满刻法和耶稣所爱的那另一位门徒，对他们说：有人把主从坟墓中搬走了，我不知道他们放在哪里。8 \BibleRef{若 20:3} 西满和那另一个门徒就出来，往坟墓那里去。9 \BibleRef{若 20:4} 两个人一起跑，那另一个门徒比西满跑得快，先到了坟墓；10 \BibleRef{若 20:5} 他低头一看，看见殓布放在那里，但没有进去。11 \BibleRef{若 20:6} 西满随后也到了，进了坟墓〔阿拉伯语，第201页〕，看见殓布放在那里；12 \BibleRef{若 20:7} 也看见耶稣头上的那块头巾没有和殓布放在一起，而是卷着放在另一个地方。13 \BibleRef{若 20:8} 那时，先到坟墓的那个门徒也进去了，一看见就相信了。14 \BibleRef{若 20:9} 因为他们还不明白圣经所说的默西亚必须从死人中复活。15 \BibleRef{若 20:10} 然后两个门徒就回家去了。\par

18 \BibleRef{若 20:11} 玛利亚却站在坟墓外面哭泣。她哭泣的时候，低头往坟墓里看，19 \BibleRef{若 20:12} 看见两个天使穿着白衣，坐在安放耶稣身体的地方，一个在头，一个在脚。20 \BibleRef{若 20:13} 天使对她说：女人，你为什么哭？她对他说：有人把我的主搬走了，我不知道他们放在哪里。21 \BibleRef{若 20:14} 她说了这些话，就转过身来，看见耶稣站在那里，却不知道是耶稣。22 \BibleRef{若 20:15} 耶稣对她说：女人，你为什么哭？你找谁？她以为他\Emph{是}园丁，就对他说：先生，如果你把他搬走了，请告诉我你把他放在哪里，我好去取他。23 \BibleRef{若 20:16} 耶稣对她说：玛利亚！她转过身来，用希伯来话对他说：\OriginalTerm{辣步尼}{Rabboni}，意思是师傅。24 \BibleRef{若 20:17} 耶稣对她说：不要摸我，因为我还没有升到父那里去。你到我弟兄那里去，告诉他们：我要升到我的父和你们的父那里，到我的天主和你们的天主那里。\par

25 \BibleRef{谷 16:9} 在一周的第一天，他复活后，首先显现给玛利亚玛达肋纳，他曾从她身上驱逐了七个魔鬼。\par

26 \BibleRef{玛 28:11} 有几个卫兵进城去，将所发生的一切事都报告给司祭长。27 \BibleRef{玛 28:12} 他们就和长老聚集，商议对策。28 \BibleRef{玛 28:13} 给了卫兵不少银钱，说：\Quote{他的门徒夜间来了，趁我们睡觉的时候把他偷走了。}\par

29 \BibleRef{玛 28:14} 如果总督听见了，我们可以说服他，使你们免于责任。30 \BibleRef{玛 28:15} 他们拿了银钱，就照他们所教的去做了。这话就在犹太人中间传开了，直到今日。\par

31 \BibleRef{若 20:18} 玛利亚玛达肋纳就去向门徒报告，说她看见了主，并且他对她说了这些话。\par

\BibleRef{玛 28:8} 当那些妇女前去告知祂的门徒时，\BibleRef{玛 28:9} 耶稣遇见她们，对她们说：\Quote{愿你们平安！}她们就上前抱住祂的脚，朝拜了祂。\BibleRef{玛 28:10} 耶稣对她们说：\Quote{不要害怕！你们去告诉我的弟兄，叫他们往加里肋亚去，在那里他们要看见我。}\BibleRef{路 24:9} 那些妇女回去，把这一切的事告诉了那十一位宗徒和其余的门徒；\BibleRef{谷 16:10} 也告诉了那些曾同祂在一起的人，因为他们正忧愁哭泣。\BibleRef{路 24:10} 她们是玛利亚玛达肋纳、约安纳、雅各伯的母亲玛利亚，以及其余同她们一起的妇女；正是她们把这些事告诉了宗徒们。\BibleRef{谷 16:11} 他们听见她们说祂活着，且向她们显现了，却不相信；\BibleRef{路 24:11} 这些话在他们看来好像是胡言乱语。此后，祂在同一天向其中两个人显现，当时他们正往一个名叫厄玛乌的村庄去，那村庄离耶路撒冷有六十斯塔狄。\BibleRef{路 24:14} 他们彼此谈论所发生的这一切事。\BibleRef{路 24:15} 正当他们彼此谈论和探讨的时候，耶稣走近他们，与他们同行。\BibleRef{路 24:16} 但他们的眼睛被蒙蔽，以致认不出祂来。\BibleRef{路 24:17} 耶稣对他们说：\Quote{你们走路时彼此谈论的是什么事？}\BibleRef{路 24:18} 一个名叫克罗帕的回答说：\Quote{难道你独在耶路撒冷作客，不知道这几天在那里所发生的事吗？}耶稣说：\Quote{什么事？}\BibleRef{路 24:19} 他们说：\Quote{就是关于纳匝肋人耶稣的事。祂是一位先知，在天主和众百姓面前，言语和行动都大有能力。\BibleRef{路 24:20} 司祭长和长老们把祂判了死刑，钉在十字架上。\BibleRef{路 24:21} 我们原指望祂就是要拯救以色列的那一位。不过这些事发生到今天已经是第三天了。\BibleRef{路 24:22} 但是我们中\Emph{某些}妇女也告诉我们，她们曾到过坟墓；\BibleRef{路 24:23} 她们没有找到祂的遗体，就回来告诉我们说，她们见了天使，天使说祂活着。\BibleRef{路 24:24} 我们中也有几个人去过坟墓，所遇见的正如妇女们所说，只是没有看见祂。}\BibleRef{路 24:25} 耶稣对他们说：\Quote{无知的人哪！你们的心信得太迟钝了！\BibleRef{路 24:26} 先知们所说的一切，默西亚不是必须受这些苦难，然后进入祂的光荣吗？}\BibleRef{路 24:27} 于是从梅瑟和众先知起，把全部圣经关于祂的话都给他们讲解明白了。\BibleRef{路 24:28} 他们临近了他们所去的村庄，祂好像还要往前走。\BibleRef{路 24:29} 他们强留祂说：\Quote{请同我们一起住下吧！因为天已经晚了，日头已经平西了。}祂就进去，同他们住下。\BibleRef{路 24:30} 当祂与他们坐席时，就拿起饼来，祝福了，擘开，递给他们。\BibleRef{路 24:31} 他们的眼睛立刻开了，认出了祂；但祂却从他们眼前不见了。\BibleRef{路 24:32} 他们彼此说：\Quote{在路上祂和我们说话，给我们讲解圣经的时候，我们的心岂不是火热的吗？}\par

\BibleRef{路 24:33} 他们就在那时起身，返回耶路撒冷，遇见那十一位宗徒和同他们一起的人聚在一起，\BibleRef{路 24:34} 说：\Quote{主真的复活了，并向西满显现了。}\BibleRef{路 24:35} 他们就述说了路上所发生的事，以及当祂擘饼时他们怎样认出了祂。\BibleRef{谷 16:13} 他们也不相信。\par

\SourceRecord{Translated by Hope W. Hogg.；Ante-Nicene Fathers；Vol. 9.；Edited by Allan Menzies.；Buffalo, NY: Christian Literature Publishing Co.,；1896.；<http://www.newadvent.org/fathers/100253.htm>.}

% structure: p0001 source=h1 target=section reason=page-title

\section{\WorkTitle{四福音合参}（第五十四段）}

1 当他们彼此交谈时，那天傍晚，即一周的第一天来临；门徒因害怕犹太人，就把门关上；\BibleRef{路 24:36}耶稣来站在他们中间，对他们说：\Quote{愿平安与你们同在；\Emph{是}我，不要害怕。}\BibleRef{路 24:37}然而他们惊慌害怕，以为看见了一个鬼魂。\BibleRef{路 24:38}耶稣对他们说：\Quote{你们为什么惊慌？为什么心里起疑念？\BibleRef{路 24:39}看我的手和脚，就知道是我；你们摸摸我，要知道鬼魂是没有骨头的，如同你们看见我有的那样。}\BibleRef{路 24:40}说了这话，就把手和脚和肋旁指给他们看。\BibleRef{路 24:41}他们因为欢喜惊奇，还是不信。\BibleRef{路 24:42}耶稣说：\Quote{你们这里有吃的吗？}他们就给了他一片烤鱼和蜜。\BibleRef{路 24:43}他接过来，在他们面前吃了。\par

8 他对他们说：\Quote{这便是我还与你们同在时，对你们说过的话：凡梅瑟法律、先知书和圣咏上关于我所记载的，都必须应验。}\BibleRef{路 24:44} 9 于是开启他们的心，使他们理解圣经；\BibleRef{路 24:45} 10 又对他们说：\BibleRef{路 24:46}\Quote{经上这样记载：默西亚必须受苦，第三日从死者中复活；11 并且要因他的名，从耶路撒冷开始，向万民宣讲悔改，以得罪赦。\BibleRef{路 24:47} 12 你们就是这些事的见证。\BibleRef{路 24:48} 13 我要将我父所恩许的，遣发到你们身上。}门徒听见这话，都很欢喜。14 耶稣又对他们说：\Quote{愿平安与你们同在：父怎样派遣了我，我也同样派遣你们。}\BibleRef{若 20:21} 15 说了这话，就向他们嘘了一口气，说：\Quote{你们领受圣神吧！}\BibleRef{若 20:22} 16 又对他们说：\Quote{你们赦免谁的罪，谁就得赦免；你们保留谁的罪，谁就被保留。}\BibleRef{若 20:23}\par

17 但十二人之一，称为第杜默的多默，当耶稣来时，没有同门徒在一起。\BibleRef{若 20:24} 18 其他门徒就对他说：\Quote{我们看见了主。}他对他们说：\Quote{我除非看见他手上的钉痕，用我的指头探入钉痕，用我的手探入他的肋旁，我决不信。}\BibleRef{若 20:25}\par

19 八天后，再次的一周的第一天，门徒又在屋里聚集，多默也同他们在一起。门关着，耶稣来了，站在中间，对他们说：\Quote{愿平安与你们同在。}\BibleRef{若 20:26} 20 然后对多默说：\Quote{伸过你的指头来，看我的手；伸出你的手来，探入我的肋旁；不要作不信的人，但要作相信的人。}\BibleRef{若 20:27} 21 多默回答说：\Quote{我的主！我的天主！}\BibleRef{若 20:28} 22 耶稣对他说：\Quote{你因为看见了我才信吗？那些没有看见而相信的，才是有福的。}\BibleRef{若 20:29}\par

23 耶稣在门徒前还行了许多别的神迹，没有记在这本书上。\BibleRef{若 20:30} 24 但这些所记的，是为叫你们信耶稣是默西亚、天主子；并使你们信的人，因他的名而得永生。\BibleRef{若 20:31}\par

25 此后，耶稣又在提庇黎雅海边，再次向门徒显现；他是这样显现的：\BibleRef{若 21:1} 26 当时西满·刻法、称为双胞胎的多默、加里肋亚加纳的纳塔乃耳、载伯德的两个儿子，还有另外两位门徒，都在一起。\BibleRef{若 21:2} 27 西满·刻法对他们说：\Quote{我去打鱼。}他们对他说：\Quote{我们也同你一起去。}他们就出去，上了船；但那夜什么也没捕到。\BibleRef{若 21:3} 28 到了早晨，耶稣站在海岸上，门徒却没有认出他是耶稣。\BibleRef{若 21:4} 29 耶稣对他们说：\Quote{孩子们，你们有吃的吗？}他们回答说：\Quote{没有。}\BibleRef{若 21:5} 30 他对他们说：\Quote{向船右边撒网，就会捕到。}他们便撒下网，竟拉不上来了，因为鱼很多。\BibleRef{若 21:6} 31 耶稣所爱的那门徒对刻法说：\Quote{是主。}西满一听说是主，就束上外衣（因为他赤着身子），跳入海里，向耶稣走去。\BibleRef{若 21:7} 32 其他门徒划着小船（离岸不远，约有二百肘），拖着那网鱼过来了。\BibleRef{若 21:8} 33 他们上了岸，看见那里放着炭火，上面有鱼和饼。\BibleRef{若 21:9} 34 耶稣对他们说：\Quote{把你们刚才捕到的鱼拿几条来。}\BibleRef{若 21:10} 35 西满·刻法便上去，把网拉到岸上，网里装满了大鱼，共一百五十三条；虽然这么多，网却没有破。\BibleRef{若 21:11} 36 耶稣对他们说：\Quote{你们来吃早餐。}门徒中没有人敢问他他是谁，因为知道是主。但主并未以\Emph{自己}本来的形像显现。\BibleRef{若 21:12} 37,38 耶稣来，拿起饼和鱼，递给他们。\BibleRef{若 21:13} 这是耶稣从死者中复活后，第三次向门徒显现。\BibleRef{若 21:14}\par

39 \BibleRef{若 21:15}\BibleQuote{他们吃完早餐后，耶稣对西满·刻法说：\Quote{西满，若纳的儿子，你爱我比这些更深吗？}他对他说：\Quote{是的，主，你知道我爱你。}耶稣对他说：\Quote{你喂养我的羔羊。}} 40 \BibleRef{若 21:16}\BibleQuote{耶稣第二次又对他说：\Quote{西满，若纳的儿子，你爱我吗？}他对他说：\Quote{是的，主，你知道我爱你。}耶稣对他说：\Quote{你牧放我的羊。}} 41 \BibleRef{若 21:17}\BibleQuote{第三次对他说：\Quote{西满，若纳的儿子，你爱我吗？}刻法因耶稣第三次说\Quote{你爱我吗}便忧愁起来，对他说：\Quote{主，你无所不知；你知道我爱你。}耶稣对他说：\Quote{你牧放我的羊。}} 42 \BibleRef{若 21:18}\BibleQuote{耶稣对他说：\Quote{你牧放我的羊。我实实在在告诉你：你年少时，自己束上腰，任意往来；但年老时，你要伸出手来，别人要给你束上腰，带你到你不愿意去的地方。}} 43 \BibleRef{若 21:19}\BibleQuote{他说这话，是指明他要以怎样的死来光荣天主。说了这话，便对他说：\Quote{跟随我吧。}} 44 \BibleRef{若 21:20}\BibleQuote{西满·刻法转过身来，看见耶稣所爱的那个门徒跟着，就是那在晚餐时靠在耶稣胸膛上并说\Quote{主，出卖你的是谁}的那一位。} 45 \BibleRef{若 21:21}\BibleQuote{刻法看见他，就对耶稣说：\Quote{主，这个人将怎样呢？}} 46 \BibleRef{若 21:22}\BibleQuote{耶稣对他说：\Quote{我若要他留到我来，与你何干？你跟随我吧。}} 47 \BibleRef{若 21:23}\BibleQuote{这话就在弟兄们中间传开，说那门徒不会死；其实耶稣并没有说他不死，而是说：\Quote{我若要他留到我来，与你何干？}}\par

48 \BibleRef{若 21:24}\BibleQuote{这就是为这些事作证并写下这些事的门徒；我们知道他的见证是真实的。}\par

\SourceRecord{Translated by Hope W. Hogg.；Ante-Nicene Fathers；Vol. 9.；Edited by Allan Menzies.；Buffalo, NY: Christian Literature Publishing Co.,；1896.；<http://www.newadvent.org/fathers/100254.htm>.}

% structure: p0001 source=h1 target=section reason=page-title

\section{\WorkTitle{四福音合参（第55节）}}

1 \BibleRef{玛 28:16} 但是那十一位门徒去了加里肋亚，到耶稣所指定的那座山上。2 \BibleRef{玛 28:17} 他们看见祂，就朝拜了祂；但是有些人心中疑惑。3 \BibleRef{谷 16:14} 当他们坐在那里时，祂又向他们显现，并责备\Emph{他们}不信和心硬，因为他们看见了祂复活后，仍不相信。\par

4 [阿拉伯语，第209页] \Quote{\BibleRef{玛 28:18} 那时耶稣对他们说：天上地下的一切权柄都已赐给了我；5 就如父派遣了我，我也派遣你们。现在你们往普天下去，向一切受造物宣讲我的福音；\BibleRef{玛 28:19} 使所有民族成为门徒，7 因父、子及圣神之名给他们授洗；\BibleRef{玛 28:20} 教导他们遵守我所吩咐你们的一切。看，我常与你们同在，直到世界的终结。8 \BibleRef{谷 16:16} 信而受洗的，必要得救；但不信的，必被定罪。9 \BibleRef{谷 16:17} 信我的人必有这些记号：因我的名驱逐魔鬼，10 说新语言；\BibleRef{谷 16:18} 并能\Emph{拿起}蛇，若喝了致命的毒药，也决不受害；按手在病人身上，他们就必痊愈。11 \BibleRef{路 24:49} 但是你们，要留在耶路撒冷城里，直到你们领受自高天而来的权能。}\par

12 我们的主耶稣对他们说完话后，就带他们出去，到了伯达尼；13 祂举起手来，祝福了他们。当祂祝福他们时，就离开了他们，被提升天，坐在天主的右边。14, 15 \BibleRef{路 24:52} 他们朝拜了祂，满心欢喜地返回了耶路撒冷；\BibleRef{路 24:53} 时常在圣殿里赞美颂扬天主。阿们。\par

16 \BibleRef{谷 16:20} 他们从那里出发，到处宣讲；主与他们合作，并以所行的神迹证实了他们所传的道。\par

17 \BibleRef{若 21:25} 此外，耶稣所行的事还有许多，若是一一写出来，我想，所写的书，就是世界也容纳不下。\par

\SourceRecord{Translated by Hope W. Hogg.；Ante-Nicene Fathers；Vol. 9.；Edited by Allan Menzies.；Buffalo, NY: Christian Literature Publishing Co.,；1896.；<http://www.newadvent.org/fathers/100255.htm>.}
